% !TEX program = pdflatex
\documentclass[12pt]{article}

\usepackage[utf8]{inputenc}
\usepackage{CJKutf8}
\usepackage{amsmath,amssymb}
\usepackage{mathtools}
\usepackage{amsthm}
\usepackage{geometry}
\geometry{
  a4paper,
  top=25mm,
  bottom=25mm,
  left=20mm,
  right=20mm
}
\usepackage{xcolor}
\usepackage{url}
\usepackage[unicode,colorlinks=true,linkcolor=blue,urlcolor=blue]{hyperref}
\usepackage{xurl}
\Urlmuskip=0mu plus 2mu
\sloppy
\usepackage{tocloft}

\renewcommand{\cftsecleader}{\cftdotfill{\cftdotsep}}
\renewcommand{\refname}{参考文献}

\newtheorem{definition}{定义}[section]
\newtheorem{axiom}{公理}[section]
\newtheorem{assumption}{假设}[section]
\newtheorem{theorem}{定理}[section]
\newtheorem{proposition}{命题}[section]
\newtheorem{corollary}{推论}[section]
\newtheorem{lemma}{引理}[section]
\newtheorem{remark}{备注}[section]
\newtheorem{Certificate}{证书}[section]

\newcommand{\NN}{\mathbb{N}}
\newcommand{\RR}{\mathbb{R}}
\newcommand{\code}[1]{\path{#1}}
\newcommand{\GFramework}{\textsf{G-Framework}}
\newcommand{\Hist}{\operatorname{Hist}}
\newcommand{\Comp}{\operatorname{Comp}}
\newcommand{\diameter}{\operatorname{diam}}
\newcommand{\Collapse}{\operatorname{Collapse}}
\newcommand{\Hole}{\operatorname{Hole}}
\newcommand{\GMS}{\mathbb{G}\mathrm{MS}}
\newcommand{\GSet}{\mathbb{G}\mathrm{Set}}
\newcommand{\id}{\mathrm{id}}

\setlength{\emergencystretch}{6em}
\setlength{\parindent}{0pt}
\setlength{\parskip}{0.6em}

\begin{document}
\begin{CJK*}{UTF8}{gkai}

\title{钱学森弹道、钱学森障碍与观测--推断--障碍修复：\\能力元组、边缘算子与全局截面障碍\\（工作笔记：形式系统版）}
\author{Gao Zheng（高政）}
\date{2026-04-13}
\maketitle

\section*{前言}
\addcontentsline{toc}{section}{前言}

本文的目标不是讨论某条具体轨迹如何被预测，也不是把“钱学森弹道”解释为某个单一工程对象，
而是在 \GFramework{} 的形式系统中，把“由中前段观测历史稳定推出终端判定”的问题改写为
能力元组、集合值终端纤维、全局截面障碍、边缘算子与修复塔的组合问题。本文采用的核心立场是：
真正困难不在于曲线是否复杂，也不在于方程形式是否高阶，而在于给定观测、模型、算力、时延与先验
之后，观测历史能否被稳定压缩为唯一、连续、可验证的终端结论。

本文的第一层立场是观测--推断问题的对象化。传统表述常问“能不能从早期轨迹看出末端结果”，
但这个问法隐藏了太多条件：观测分辨率、模型类、计算资源、刷新时延、先验族、隐藏模式和终端标签
都会改变可推断范围。本文把这些条件合并为能力元组 \(\Lambda_t\)，并由此定义可容许终端纤维
\(\Gamma_t^{\Lambda_t}(b)\)。于是，终端先验推断不再是直观判断，而成为集合值逆映射是否能够局部
单值化、连续化和证书化的问题。

本文的第二层立场是“钱学森障碍”的严格化。在本文中，障碍并不被理解为某条轨迹难以计算，
而是从观测历史 \(b\) 到终端集合 \(\Gamma_t^{\Lambda_t}(b)\) 的压缩失败：同一观测前缀可能对应多个
相互不可区分的可容许终端，或者隐藏模式、模型切换、时延预算和资源限制使局部截面不存在。
因此，钱学森障碍的形式本质是全局截面障碍、局部单值化失败与稳定终端判定失败。

本文的第三层立场是边缘算子的修复作用。能力元组不是静态背景，边缘算子可以改变观测接口、
模型类、算力预算、时延规则和先验族，从而改变有效域、失效域与换模域。边缘算子的职责不是
魔术般消除障碍，而是明确地补入缺失信息、改变纤维结构、压缩剩余歧义或给出失败证书。
由此，障碍修复不是简单增加算力或调整模型，而是对能力元组和可容许终端纤维的结构性重写。

本文的第四层立场是修复塔的层级性。若隐藏模式未入状态，方程本身不闭合；若结构完备性不足，
内部修复无法找回被遗漏的真值；若预算与时延不匹配，局部截面即使存在也可能不可交付。因此，
修复必须按层推进：先补动力学与观测，再补模式与换模，再补资源与时延，再补结构完备性与证书接口。
修复塔的意义在于把“能否推断”拆成可定位的层级义务，而不是把所有失败归结为单一算法不足。

本文的第五层立场是结构完备性优先于性价比。若对象域、模式域或终端标签域本身不完整，
再高的单位预算收益也只会优化错误问题。性价比只有在结构完备性、真值保持和失败证书边界成立之后，
才有排序意义。因此，本文将钱学森弹道/钱学森障碍问题进一步改写为障碍修复--预算匹配问题：
先判断修复是否进入正确对象层，再比较单位预算收益、达阈时间与边际收益饱和。

本文的第六层立场是从轨迹问题转向同伦命中问题。随着后续章节引入极小极大泛函、变分修复、
同伦扭曲命中、轨迹商空间、正规形子空间、连续统崩溃和可数超限离散统，本文逐步把“轨迹推断”
重写为“在修复后的对象空间中命中稳定截面或低残差终端”的问题。这个转化为后续统计力学同伦命中、
PDE 策略簇与求解范式替代提供了几何和控制论前史：命中并不是凭空出现的概率词，而是从截面障碍、
修复塔和边缘算子中自然长出的终端判定机制。

因此，本文在整体 \(\GFramework{}\) 中承担的是“障碍修复母稿”的作用。它向前连接观测历史、
混合系统、集合值分析、逆问题和纤维丛语言，向中间连接能力元组、边缘算子、修复塔和结构完备性，
向后连接同伦扭曲命中、PDE 表象升维、正规形强变形收缩、连续统崩溃和高层 \(L\ge3\) 修复职责。
若后续稿件讨论微观隧穿、统计力学命中或 PDE 求解范式替代，那么本文给出的基础问题是：
什么是障碍，障碍位于哪个纤维，边缘算子如何改变它，修复塔何时给出截面，何时只能签发失败证书。

概括地说，本文把“从观测到终端”的直觉问题改写为“从能力元组到截面障碍”的形式问题，
把“轨迹难预测”改写为“集合值逆映射难单值化”，把“增加资源”改写为“边缘算子改变有效域”，
把“求一个结论”改写为“构造局部截面或失败证书”。它提供的不是具体载体方案，而是一套面向
开放系统推断的障碍定位、修复分层与终端判定形式系统。

\setcounter{tocdepth}{3}
\setcounter{secnumdepth}{3}
\tableofcontents
\clearpage

%%%%%%%%%%%%%%%%%%%%%%%%%%%%%%%%%%%%%%%%%%%%%%%%%%%%%%%%%%%%
\section{形式逻辑：钱学森弹道、钱学森障碍与从观测历史到终端判定的能力元组、边缘算子、全局截面障碍}
\label{sec:tex1-ch1}
%%%%%%%%%%%%%%%%%%%%%%%%%%%%%%%%%%%%%%%%%%%%%%%%%%%%%%%%%%%%

\begin{remark}[本章意义备注：高超音速攻防的抽象语境]
本章的意义，在于把高超音速攻防中“能否由中前段观测稳定判定末段结果”的问题，
从直观轨迹讨论改写成能力元组、可容许终端纤维与截面障碍问题。
对钱学森弹道/钱学森障碍而言，它给出最基础的判别坐标：
真正的难点不是曲线是否复杂，而是观测历史在给定感知、模型、算力与时延条件下能否被稳定单值化。
\end{remark}

\subsection{严格主张、引文定位与章节目标}
\label{subsec:tex1-ch1-claim}

\begin{remark}[核心陈述]
本章把“到达终端之前能否稳定推出唯一终端结论”严格改写为一个集合值逆问题：
\[
  \text{观测历史}
  \xRightarrow{\ \Lambda_t\ }
  \text{可容许终端纤维}
  \xRightarrow{\ \text{截面或失败证书}\ }
  \text{稳定推断边界}.
\]
因此，问题的本质不再被表述为“轨迹长什么样”或“方程写成几阶”，
而被表述为：在给定能力元组 $\Lambda_t$ 与边缘算子 $\Delta_t$ 的条件下，
集合值映射
\[
  b\longmapsto \Gamma_t^{\Lambda_t}(b)
\]
是否在局部上收缩为连续单值截面。
在本章语境下，``钱学森弹道'' 与 ``弹道'' 只作为\textbf{观测侧的术语标签}出现，
指称“由早期观测历史指向终端功能量/终端标签的弹道化表象问题”；
``钱学森障碍'' 则被严格收紧为：在能力元组 $\Lambda_t$ 下，
该集合值逆映射无法稳定压缩为唯一终端结论的结构性失败。
这一组织方式继承仓内 \cite{RepoTex37SemanticGaugeRuntime,RepoTex17EdgeOperator,RepoTex18HomotopyAxiom,
  RepoTex29RepairableObstruction,RepoTex38JacobiBianchiRepair,RepoTex43OpenHeterogeneous}
关于“对象域 + 证书 + 谓词因果链 + 修复塔”的写法；其外部数学背景则以非线性可观测性、逆问题、切换/混合系统、
集合值分析与纤维丛语言为准，可参见
\cite{HermannKrener1977,Isakov2006InversePDE,Liberzon2003Switching,AubinFrankowska2009SetValued,
  Husemoller1994FibreBundles}。
\end{remark}

\begin{remark}[本章的严格目标]
本章只处理抽象控制论、集合值逆问题与开放系统对象层，不绑定任何具体工程载体。
我们要严格保留的正确结论可压缩为四句：
\begin{enumerate}
  \item 终端先验推断是否成立，取决于集合值逆映射 $b\mapsto \Gamma_t^{\Lambda_t}(b)$ 的局部单值化与稳定性；
  \item 边缘算子 $\Delta_t$ 通过改变能力元组而改变有效域、失效域与换模域；
  \item 一旦隐藏模式或模型类切换进入问题，本体表达更接近混合系统或微分包含，而不是单一静态方程；
  \item 修复塔的职责，是按层补齐动力学、观测、模式、换模与资源时延信息，直到得到局部截面或失败证书。
\end{enumerate}
\end{remark}

\begin{remark}[术语边界：允许使用“钱学森弹道/钱学森障碍/弹道”，但仅限抽象口径]
本章不回避 ``钱学森弹道''、``钱学森障碍''、``弹道'' 这些词，
但它们在全文中都只被用作抽象控制论与逆问题语言里的\textbf{命名标签}：
\begin{enumerate}
  \item ``钱学森弹道'' 或 ``弹道'' 指观测者眼中的轨迹/终端表象层；
  \item ``钱学森障碍'' 指在受限能力元组下，从观测历史到唯一终端结论的稳定压缩失败；
  \item 本章不把这些词延伸为任何具体实现、反求、规避或对抗设计问题。
\end{enumerate}
\end{remark}

\subsection{定义链：抽象系统、观测历史与可容许终端纤维}
\label{subsec:tex1-ch1-defs-system}

\begin{definition}[抽象系统（钱学森弹道问题的抽象控制论口径）]
\label{def:tex1-abstract-system}
先记基础系统数据为
\[
  \Sigma_0=(X,U,\Theta,f,h,\psi,T),
\]
其中：
\begin{enumerate}
  \item $X$ 为状态空间，$U$ 为控制/输入历史空间，$\Theta$ 为参数空间；
  \item $f$ 为动力学右端，$h$ 为观测算子，$\psi$ 为终端功能量映射；
  \item $T>0$ 为终端时刻。
\end{enumerate}
并记
\[
  \dot x(t)=f(x(t),u(t),\theta,t),
  \qquad
  z(t)=h(x(t),t)+\eta(t),
  \qquad
  y=\psi(x(T)).
\]
这里：
\begin{enumerate}
  \item $x(t)\in X$，$u(t)\in U$，$\theta\in\Theta$；
  \item $z_{[0,t]}$ 表示到时刻 $t$ 的观测历史；
  \item $y$ 为终端功能量，可解释为终端分类结果、终端标签或终端功能值，不预设具体物理语义。
\end{enumerate}

为容纳隐藏换模、混合系统与集合值动力学，再把基础系统扩充为
\[
  \Sigma=(X,U,\Theta,\mathcal S,f,h,\psi,T,\mathcal A,Y),
\]
其中：
\begin{enumerate}
  \item $\mathcal S$ 为隐藏模式集合，$Y$ 为终端标签或终端功能量所在的度量空间；
  \item $\mathcal A$ 为全部\textbf{可容许演化数据}的集合。每个
  \[
    a=(x_0,u,\theta,\sigma)\in\mathcal A
  \]
  对应一条演化轨迹 $x_a:[0,T]\to X$、一条观测轨迹 $z_a:[0,T]\to Z$ 与一个终端值
  \[
    y_a:=\psi(x_a(T))\in Y.
  \]
\end{enumerate}
这里不预设问题最终必须停留在单一常微分方程口径；它可以是常微分系统、偏微分系统、切换系统，或者更一般的
\[
  \dot x(t)\in F_{\sigma(t)}(x(t),u(t),\theta,t)
\]
形式的微分包含。
\end{definition}

\begin{definition}[观测历史映射]
\label{def:tex1-history-map}
对任意 $t\in[0,T)$，定义历史映射
\[
  \Hist_t:\mathcal A\to \mathcal B_t,
  \qquad
  \Hist_t(a):=z_a|_{[0,t]},
\]
其中 $\mathcal B_t$ 表示所有长度为 $t$ 的可实现观测历史所构成的空间。
记 $d_B$ 为 $\mathcal B_t$ 上的度量，$d_Y$ 为 $Y$ 上的度量。
\end{definition}

\begin{definition}[能力元组与相容谓词]
\label{def:tex1-capability}
对每个 $t<T$，定义观察方的能力元组
\[
  \Lambda_t=(\mathcal O_t,\mathcal M_t,\mathcal C_t,\mathcal L_t,\Pi_t),
\]
其中：
\begin{enumerate}
  \item $\mathcal O_t$ 为观测接口与观测分辨率；
  \item $\mathcal M_t$ 为采用的模型类；
  \item $\mathcal C_t$ 为可用计算资源；
  \item $\mathcal L_t$ 为时延、更新周期与刷新规则；
  \item $\Pi_t$ 为先验族。
\end{enumerate}
定义相容谓词
\[
  \Comp_t^{\Lambda_t}:\mathcal A\to\{\top,\bot\},
\]
其含义是：某个可容许演化数据 $a\in\mathcal A$ 在时刻 $t$ 的能力元组 $\Lambda_t$ 下，
既与已知观测一致，也不违反当前模型类、资源边界、时延规则与先验族。
\end{definition}

\begin{definition}[可容许终端集、终端纤维丛与局部截面]
\label{def:tex1-terminal-fiber}
给定 $t<T$ 与能力元组 $\Lambda_t$，定义相应的可实现历史基空间
\[
  B_t^{\Lambda_t}
  :=
  \bigl\{
    \Hist_t(a)\in \mathcal B_t
    \,\big|\,
    a\in\mathcal A,\ \Comp_t^{\Lambda_t}(a)=\top
  \bigr\}.
\]
对每个 $b\in B_t^{\Lambda_t}$，定义其\textbf{可容许终端集}
\[
  \Gamma_t^{\Lambda_t}(b)
  :=
  \bigl\{
    y_a\in Y
    \,\big|\,
    a\in\mathcal A,\ \Hist_t(a)=b,\ \Comp_t^{\Lambda_t}(a)=\top
  \bigr\}.
\]
进而定义总空间
\[
  E_t^{\Lambda_t}
  :=
  \bigsqcup_{b\in B_t^{\Lambda_t}} \Gamma_t^{\Lambda_t}(b),
\]
以及投影
\[
  p_t^{\Lambda_t}:E_t^{\Lambda_t}\to B_t^{\Lambda_t},
  \qquad
  p_t^{\Lambda_t}(b,y):=b.
\]

若存在开集 $U\subseteq B_t^{\Lambda_t}$ 与连续映射
\[
  s:U\to Y
\]
满足
\[
  s(b)\in\Gamma_t^{\Lambda_t}(b)
  \qquad
  (\forall b\in U),
\]
则称 $s$ 为该纤维丛上的一个\textbf{连续局部截面}；
等价地，$s^\sharp(b):=(b,s(b))$ 满足
\[
  p_t^{\Lambda_t}\circ s^\sharp=\id_U.
\]
\end{definition}

\begin{definition}[钱学森弹道与钱学森障碍的抽象口径]
\label{def:tex1-qxs-terms}
在定义~\ref{def:tex1-abstract-system}--\ref{def:tex1-terminal-fiber} 的对象域下，
作如下术语约定：
\begin{enumerate}
  \item 若某个问题在观测侧被表述为“由历史观测推断终端功能量/终端标签的轨迹化或弹道化问题”，
  则称其\textbf{钱学森弹道表象}由三元组
  \[
    (\Hist_t,\Gamma_t^{\Lambda_t},\psi)
  \]
  给定；
  \item 若对某个时刻 $t<T$ 与某个历史域 $U\subseteq B_t^{\Lambda_t}$，
  不存在连续局部截面
  \[
    s:U\to Y,\qquad s(b)\in\Gamma_t^{\Lambda_t}(b),
  \]
  则称在 $(t,\Lambda_t,U)$ 上存在\textbf{局部钱学森障碍}；
  \item 若在整个基空间 $B_t^{\Lambda_t}$ 上不存在连续全局截面
  \[
    s:B_t^{\Lambda_t}\to Y,\qquad s(b)\in\Gamma_t^{\Lambda_t}(b),
  \]
  则称在 $(t,\Lambda_t)$ 上存在\textbf{全局钱学森障碍}。
\end{enumerate}
\end{definition}

\begin{definition}[局部终端估计映射、正则历史与奇异历史]
\label{def:tex1-local-estimator}
固定 $t<T$、能力元组 $\Lambda_t$ 与开集 $U\subseteq B_t^{\Lambda_t}$。
若存在一个 $C^1$ 映射
\[
  \Phi_t:U\to Y
\]
满足
\[
  \Phi_t(b)\in\Gamma_t^{\Lambda_t}(b)
  \qquad
  (\forall b\in U),
\]
则称 $\Phi_t$ 为一个\textbf{局部终端估计映射}。

若 $Y$ 在 $\Phi_t(b_\ast)$ 的邻域内可由 $\RR^m$ 局部坐标化，并且
\[
  \operatorname{rank}D\Phi_t(b_\ast)=m=\dim Y,
\]
则称 $b_\ast$ 为该估计映射下的\textbf{正则历史}；
若
\[
  \operatorname{rank}D\Phi_t(b_\ast)<m,
\]
则称 $b_\ast$ 为该估计映射下的\textbf{奇异历史}，并称该不等式为\textbf{雅可比退化证书}。
\end{definition}

\begin{remark}[关于“上同调障碍”的边界]
定义~\ref{def:tex1-qxs-terms} 已经足以严格定义 ``钱学森障碍''，
但还\textbf{不足以}自动推出 ``上同调障碍''。
若要把钱学森障碍进一步提升为上同调障碍，
仍需显式构造障碍类、给出相应的复形或纤维化数据，并证明其不可延拓性。
因此在本章里，
\[
  \text{钱学森障碍}
  \quad\text{是已定义对象，}
  \qquad
  \text{上同调障碍}
  \quad\text{仍需额外证书。}
\]
\end{remark}

\begin{definition}[\texorpdfstring{$\Gamma$}{Gamma}-等价表示]
\label{def:tex1-gamma-equivalent}
设 $\Sigma_1,\Sigma_2$ 为两种对同一观测--终端问题的表示。
若对某个固定的 $t<T$、能力元组 $\Lambda_t$ 与开集 $U\subseteq \mathcal B_t$，
恒有
\[
  \Gamma_{t,\Sigma_1}^{\Lambda_t}(b)
  =
  \Gamma_{t,\Sigma_2}^{\Lambda_t}(b)
  \qquad
  (\forall b\in U),
\]
则称 $\Sigma_1,\Sigma_2$ 在 $(t,\Lambda_t,U)$ 上\textbf{\(\Gamma\)-等价}。
\end{definition}

\subsection{定义链：边缘算子、有效域与修复塔}
\label{subsec:tex1-ch1-defs-edge}

\begin{definition}[边缘算子与三类推断域]
\label{def:tex1-edge-operator}
给定阈值 $\varepsilon>0$，定义
\[
  \mathcal E_t^{\Lambda_t}(\varepsilon)
  :=
  \bigl\{
    b\in B_t^{\Lambda_t}
    \,\big|\,
    \diameter_{d_Y}\Gamma_t^{\Lambda_t}(b)\le \varepsilon
  \bigr\},
\]
称其为\textbf{有效推断域}；
\[
  \mathcal F_t^{\Lambda_t}(\varepsilon)
  :=
  B_t^{\Lambda_t}\setminus \mathcal E_t^{\Lambda_t}(\varepsilon),
\]
称其为\textbf{失效域}。

定义\textbf{边缘算子}
\[
  \Delta_t:\Lambda_t\mapsto \Lambda_t',
\]
它允许同时改变观测接口、模型类、资源预算、时延规则与先验族。
由此诱导\textbf{换模域}
\[
  \mathcal S_t^{\Delta_t}(\varepsilon)
  :=
  \mathcal E_t^{\Lambda_t}(\varepsilon)
  \triangle
  \mathcal E_t^{\Lambda_t'}(\varepsilon),
\]
其中 $\triangle$ 表示对称差。
\end{definition}

\begin{definition}[观测--推断修复塔]
\label{def:tex1-repair-tower}
记六个阶段性谓词为
\[
  C_0:\ \text{动力学一致性},
  \qquad
  C_1:\ \text{观测一致性},
  \qquad
  C_2:\ \text{隐藏模式显式化},
\]
\[
  C_3:\ \text{模型类切换一致性},
  \qquad
  C_4:\ \text{资源--时延一致性},
  \qquad
  C_5:\ \text{局部截面或失败证书完备}.
\]
称有限序列
\[
  \mathfrak T_t^{(0)}
  \xrightarrow{R_0}
  \mathfrak T_t^{(1)}
  \xrightarrow{R_1}
  \cdots
  \xrightarrow{R_4}
  \mathfrak T_t^{(5)}
\]
为一个\textbf{观测--推断修复塔}，若对每个 $r=0,\dots,5$，
\[
  \mathfrak T_t^{(r)}
  =
  \bigl(
    \Lambda_t^{(r)},
    B_t^{(r)},
    E_t^{(r)},
    p_t^{(r)}
  \bigr),
\]
并且：
\begin{enumerate}
  \item $\mathfrak T_t^{(r)}$ 满足谓词 $C_0,\dots,C_r$；
  \item 修复算子 $R_r$ 把第 $r$ 层对象送到第 $r+1$ 层对象；
  \item $R_r$ 不破坏已满足的 $C_0,\dots,C_r$，并使新的 $C_{r+1}$ 成立。
\end{enumerate}
\end{definition}

\begin{remark}[关于“通过修复塔加入边缘算子即可突破钱学森障碍”的边界]
定义~\ref{def:tex1-repair-tower} 只把 ``修复'' 严格写成一条分层补齐链，
并不自动推出
\[
  \text{“加入边缘算子并沿修复塔推进，就一定能够突破钱学森障碍”}.
\]
若要把这句话升级为定理，至少还需要三类额外证书：
\begin{enumerate}
  \item 钱学森障碍的严格定义证书；
  \item 修复后连续全局截面或连续局部截面的存在性证书；
  \item 修复代价、资源上界与时延上界的定量证书。
\end{enumerate}
在这些证书未被单独构造完成之前，上述强表述仍应读作\textbf{未证明}，而不是已完成定理。
\end{remark}

\subsection{公理（降级为定理）：修复塔可按层递归构造}
\label{subsec:tex1-ch1-axiom-downgrade}

\begin{theorem}[公理降级：六层观测--推断修复塔可递归构造]
\label{thm:tex1-axiom-downgrade}
设在固定 $t<T$ 下满足：
\begin{enumerate}
  \item 存在初始对象
  \[
    \mathfrak T_t^{(0)}
    =
    \bigl(
      \Lambda_t^{(0)},
      B_t^{(0)},
      E_t^{(0)},
      p_t^{(0)}
    \bigr)
  \]
  且谓词 $C_0$ 成立；
  \item 对每个 $r\in\{0,1,2,3,4\}$，只要 $\mathfrak T_t^{(r)}$ 已构成并满足 $C_0,\dots,C_r$，
  就存在唯一修复算子
  \[
    R_r:\mathfrak T_t^{(r)}\to \mathfrak T_t^{(r+1)}
  \]
  使得 $\mathfrak T_t^{(r+1)}$ 同时满足 $C_0,\dots,C_{r+1}$。
\end{enumerate}
则存在唯一一条六层观测--推断修复塔
\[
  \mathfrak T_t^{(0)}
  \xrightarrow{R_0}
  \mathfrak T_t^{(1)}
  \xrightarrow{R_1}
  \mathfrak T_t^{(2)}
  \xrightarrow{R_2}
  \mathfrak T_t^{(3)}
  \xrightarrow{R_3}
  \mathfrak T_t^{(4)}
  \xrightarrow{R_4}
  \mathfrak T_t^{(5)}.
\]
\end{theorem}

\begin{Certificate}[数学归纳法证书：按层数的修复塔谓词]
\label{cert:tex1-axiom-downgrade}
定义谓词
\[
  P(n):
  \text{“存在唯一长度为 $n$ 的前缀修复塔 }
  \mathfrak T_t^{(0)}\to\cdots\to \mathfrak T_t^{(n)},
  \text{ 且 }C_0,\dots,C_n\text{ 全部成立。”}
\]
则证书登记为
\[
  \pi_{\mathrm{ind}}
  =
  \Bigl(
    P(0),\
    \forall n\in\{0,1,2,3,4\},\ P(n)\Rightarrow P(n+1)
  \Bigr).
\]
\end{Certificate}

\begin{proof}[证明（数学归纳法证书；基于数学符号推演逻辑的谓词因果链）]
\textbf{基步：}由条件(1)，初始对象 $\mathfrak T_t^{(0)}$ 已存在且满足 $C_0$，
因此 $P(0)$ 成立。

\textbf{归纳步：}任取 $n\in\{0,1,2,3,4\}$，设 $P(n)$ 成立。
则已经存在唯一前缀链
\[
  \mathfrak T_t^{(0)}
  \xrightarrow{R_0}
  \cdots
  \xrightarrow{R_{n-1}}
  \mathfrak T_t^{(n)},
\]
并且 $C_0,\dots,C_n$ 全部成立。

由条件(2)，在 $C_0,\dots,C_n$ 已成立的前提下，存在唯一修复算子
\[
  R_n:\mathfrak T_t^{(n)}\to \mathfrak T_t^{(n+1)}
\]
使得新的对象 $\mathfrak T_t^{(n+1)}$ 同时满足
\[
  C_0,\dots,C_n,C_{n+1}.
\]
于是前缀链可唯一延长为
\[
  \mathfrak T_t^{(0)}
  \to
  \cdots
  \to
  \mathfrak T_t^{(n)}
  \xrightarrow{R_n}
  \mathfrak T_t^{(n+1)}.
\]
这就推出 $P(n+1)$。

综上，
\[
  P(0)
  \wedge
  \forall n\in\{0,1,2,3,4\},\ P(n)\Rightarrow P(n+1)
\]
成立。
由数学归纳法，$P(5)$ 成立，即六层观测--推断修复塔存在且唯一。证毕。
\end{proof}

\subsection{定理：钱学森弹道的本质不在高阶偏导本身，而在 \texorpdfstring{$\Gamma$}{Gamma}-纤维与逆问题可适定性}
\label{subsec:tex1-ch1-theorem}

\begin{theorem}[钱学森弹道/钱学森障碍的逆问题本质定理：本质不在高阶偏导本身，而在 \texorpdfstring{$\Gamma$}{Gamma}-纤维与逆问题可适定性]
\label{thm:tex1-gamma-invariant}
固定 $t<T$、能力元组 $\Lambda_t$ 与开集 $U\subseteq \mathcal B_t$。
若两种表示 $\Sigma_1,\Sigma_2$ 在 $(t,\Lambda_t,U)$ 上 \(\Gamma\)-等价，
则下列命题对二者同时真或同时假：
\begin{enumerate}
  \item 在 $U$ 上存在连续局部截面；
  \item 对每个 $b\in U$，纤维 $\Gamma_t^{\Lambda_t}(b)$ 是否为单点；
  \item 对任意阈值 $\varepsilon>0$，$U$ 中各点是否落入有效推断域 $\mathcal E_t^{\Lambda_t}(\varepsilon)$。
\end{enumerate}
因此，到达终端之前的可稳定推断性完全由集合值逆映射
\[
  b\longmapsto \Gamma_t^{\Lambda_t}(b)
\]
决定，而不是由某个表示中出现的微分阶数决定。
特别地，把任意有限阶高阶动力表示写成一阶增广状态表示时，
其 \(\Gamma\)-映射保持不变，故“方程阶数”只是表示特征而不是本质判据。
换言之，若把问题称为 ``钱学森弹道''，则其本质也不在 ``弹道写成几阶方程''，
而在
\[
  b\mapsto \Gamma_t^{\Lambda_t}(b)
\]
是否允许稳定单值化；更一般地，它取决于
\[
  (\text{动力学})+(\text{观测})+(\text{先验})+(\text{计算资源})+(\text{时间窗})
\]
共同定义的逆问题是否可适定。
\end{theorem}

\begin{Certificate}[证书：\texorpdfstring{$\Gamma$}{Gamma}-相同 $\Rightarrow$ 截面/单点性/直径阈值相同]
\label{cert:tex1-gamma-invariant}
证书登记谓词链：
\[
\begin{aligned}
  T_1 &: \forall b\in U,\ \Gamma_{t,\Sigma_1}^{\Lambda_t}(b)=\Gamma_{t,\Sigma_2}^{\Lambda_t}(b),\\
  T_2 &: s:U\to Y\ \text{是 }\Sigma_1\text{ 的连续局部截面}
        \Rightarrow \forall b\in U,\ s(b)\in \Gamma_{t,\Sigma_1}^{\Lambda_t}(b),\\
  T_3 &: T_1\wedge T_2
        \Rightarrow \forall b\in U,\ s(b)\in \Gamma_{t,\Sigma_2}^{\Lambda_t}(b),\\
  T_4 &: T_1
        \Rightarrow
        \Bigl(
          \Gamma_{t,\Sigma_1}^{\Lambda_t}(b)\text{ 为单点}
          \Leftrightarrow
          \Gamma_{t,\Sigma_2}^{\Lambda_t}(b)\text{ 为单点}
        \Bigr),\\
  T_5 &: T_1
        \Rightarrow
        \diameter\Gamma_{t,\Sigma_1}^{\Lambda_t}(b)
        =
        \diameter\Gamma_{t,\Sigma_2}^{\Lambda_t}(b),\\
  T_6 &: \text{高阶表示的一阶增广写法保持同一观测历史与终端映射},\\
  T_7 &: T_3\wedge T_4\wedge T_5\wedge T_6
        \Rightarrow \text{推断性与微分阶数无关，而只与 }\Gamma\text{ 有关}.
\end{aligned}
\]
\end{Certificate}

\begin{proof}[证明（基于数学符号推演逻辑的谓词因果链）]
由 \(T_1\)，两种表示在 $U$ 上对每个观测历史 $b$ 诱导同一个终端可容许集合。

若 $s:U\to Y$ 是 $\Sigma_1$ 的连续局部截面，则由 \(T_2\) 可知
\[
  s(b)\in \Gamma_{t,\Sigma_1}^{\Lambda_t}(b)
  \qquad
  (\forall b\in U).
\]
再由 \(T_1\)，立即得
\[
  s(b)\in \Gamma_{t,\Sigma_2}^{\Lambda_t}(b)
  \qquad
  (\forall b\in U),
\]
这就是 \(T_3\)。
由于连续性只取决于 $s$ 与基空间/值空间的拓扑，
而不取决于把同一纤维集合写成哪一种表示，
故 $s$ 也是 $\Sigma_2$ 的连续局部截面；交换 $\Sigma_1,\Sigma_2$ 的角色可得逆向结论。
于是“连续局部截面的存在性”对 \(\Gamma\)-等价表示不变。

同理，由 \(T_1\) 对每个 $b\in U$ 逐点比较纤维可得 \(T_4\)；
由于集合完全相同，其直径也完全相同，从而 \(T_5\) 成立。
于是有效域、失效域等一切只依赖纤维直径与单点性的对象，均在 \(\Gamma\)-等价下保持不变。

最后，对任何有限阶高阶动力表示，都可通过标准状态增广
\[
  \widetilde x=(x,\dot x,\dots,x^{(k-1)})
\]
把 $k$ 阶表示改写为一阶系统，而观测历史 $\Hist_t$ 与终端投影 $\psi$ 不变；
因此两种写法诱导同一个集合值映射 $b\mapsto \Gamma_t^{\Lambda_t}(b)$，即 \(T_6\)。
把 \(T_3\)、\(T_4\)、\(T_5\)、\(T_6\) 合并，即得 \(T_7\)：
真正决定“到达终端前是否可稳定推断”的是 \(\Gamma\)-映射本身，
而不是某个特定表示中的微分阶数。证毕。
\end{proof}

\begin{corollary}[高阶偏导形成钱学森弹道至多是局部描述]
\label{cor:tex1-high-order-not-essential}
固定 $t<T$、能力元组 $\Lambda_t$、开集 $U\subseteq B_t^{\Lambda_t}$ 与局部终端估计映射
\[
  \Phi_t:U\to Y.
\]
若下列任一证书成立：
\begin{enumerate}
  \item 存在 $b_\ast\in U$ 使得
  \[
    \Gamma_t^{\Lambda_t}(b_\ast)
    \ \text{不是单点};
  \]
  \item 存在 $b_\ast\in U$ 使得
  \[
    \operatorname{rank}D\Phi_t(b_\ast)<\dim Y;
  \]
\end{enumerate}
则在 $b_\ast$ 的邻域上，不存在把早期观测历史稳定压缩为满维唯一终端结论的正则参数化。
因此，
\[
  \text{“高阶偏导形成钱学森弹道/控制轨迹”}
\]
至多只是局部表示描述；若把它提升为问题本质，则在本章对象层下仍属\textbf{未证明强化表达}。
\end{corollary}

\begin{Certificate}[证书：纤维非单点或雅可比退化 $\Rightarrow$ 正则单值化失败]
\label{cert:tex1-high-order-not-essential}
证书登记谓词链：
\[
\begin{aligned}
  H_1 &: \Gamma_t^{\Lambda_t}(b_\ast)\ \text{不是单点},\\
  H_2 &: H_1\Rightarrow \text{在 }b_\ast\text{ 处不存在唯一终端选择},\\
  H_3 &: \operatorname{rank}D\Phi_t(b_\ast)<\dim Y,\\
  H_4 &: H_3\Rightarrow \Phi_t\ \text{在 }b_\ast\text{ 处不是满维正则参数化},\\
  H_5 &: H_2\ \text{或}\ H_4\Rightarrow \text{无法把早期观测历史稳定压缩为满维唯一终端结论},\\
  H_6 &: H_5\Rightarrow \text{“高阶偏导形成钱学森弹道”不能充当问题本质判据}.
\end{aligned}
\]
\end{Certificate}

\begin{proof}[证明（基于数学符号推演逻辑的谓词因果链）]
若 \(H_1\) 成立，则同一历史 $b_\ast$ 至少对应两个不同的可容许终端值，
故由 \(H_2\) 可知，在 $b_\ast$ 处不存在唯一终端选择，
更不可能存在稳定的单值压缩。

若 \(H_3\) 成立，则按定义~\ref{def:tex1-local-estimator}，
$b_\ast$ 是局部终端估计映射 $\Phi_t$ 的奇异历史。
因此 $\Phi_t$ 在 $b_\ast$ 处不能给出满维正则参数化，这就是 \(H_4\)。
于是无论由 \(H_2\) 还是由 \(H_4\) 进入，都得到 \(H_5\)：
在 $b_\ast$ 的邻域上，早期观测历史都不能被稳定压缩为满维唯一终端结论。

因此，问题的关键不在“方程是否写成高阶偏导形式”，
而在相应 \(\Gamma\)-纤维是否单值、局部终端估计是否正则，以及整体逆问题是否可适定。
这就推出 \(H_6\)，即该推论成立。证毕。
\end{proof}

\subsection{引理：前缀不可区分而终端不同，排除唯一稳定截面}
\label{subsec:tex1-ch1-lemma}

\begin{lemma}[钱学森障碍的不可区分前缀引理]
\label{lem:tex1-indistinguishable-prefix}
固定 $t<T$ 与能力元组 $\Lambda_t$。
\begin{enumerate}
  \item 若存在同一历史 $b\in B_t^{\Lambda_t}$ 与两个终端值 $y_1\neq y_2$ 满足
  \[
    y_1,y_2\in \Gamma_t^{\Lambda_t}(b),
  \]
  则在点 $b$ 处不存在唯一终端选择；
  \item 若存在序列
  \[
    b_n',b_n''\in B_t^{\Lambda_t},
    \qquad
    b_n',b_n''\to b_\ast,
  \]
  以及对应的终端值
  \[
    y_n'\in \Gamma_t^{\Lambda_t}(b_n'),
    \qquad
    y_n''\in \Gamma_t^{\Lambda_t}(b_n'')
  \]
  与常数 $\varepsilon_0>0$，使得
  \[
    d_Y(y_n',y_n'')\ge \varepsilon_0
    \qquad
    (\forall n\in\NN),
  \]
  则在 $b_\ast$ 的任意邻域上都不存在连续局部截面。
\end{enumerate}
\end{lemma}

\begin{Certificate}[证书：纤维非单点或近历史分离 $\Rightarrow$ 截面失败]
\label{cert:tex1-indistinguishable-prefix}
证书登记谓词链：
\[
\begin{aligned}
  L_1 &: y_1,y_2\in \Gamma_t^{\Lambda_t}(b)\ \wedge\ y_1\neq y_2,\\
  L_2 &: L_1\Rightarrow \Gamma_t^{\Lambda_t}(b)\ \text{不是单点},\\
  L_3 &: b_n',b_n''\to b_\ast,\ y_n'\in \Gamma_t^{\Lambda_t}(b_n'),\ y_n''\in \Gamma_t^{\Lambda_t}(b_n''),\\
  L_4 &: \forall n,\ d_Y(y_n',y_n'')\ge \varepsilon_0>0,\\
  L_5 &: L_3\wedge L_4\wedge s\ \text{连续}
        \Rightarrow d_Y\bigl(s(b_n'),s(b_n'')\bigr)\to 0,\\
  L_6 &: L_5\ \text{与}\ L_4\ \text{矛盾}
        \Rightarrow \text{不存在连续局部截面}.
\end{aligned}
\]
\end{Certificate}

\begin{proof}[证明（基于数学符号推演逻辑的谓词因果链）]
对第(1)项，由 \(L_1\) 直接得 \(L_2\)：
\[
  \Gamma_t^{\Lambda_t}(b)
  \supseteq
  \{y_1,y_2\},
  \qquad
  y_1\neq y_2.
\]
因此在同一历史 $b$ 上至少存在两个可容许终端结论，故唯一终端选择不成立。

对第(2)项，反设在 $b_\ast$ 的某个邻域 $U$ 上存在连续局部截面
\[
  s:U\to Y,
  \qquad
  s(b)\in\Gamma_t^{\Lambda_t}(b).
\]
由 \(L_3\) 可知，充分大的 $n$ 时 $b_n',b_n''\in U$。
由连续性，
\[
  b_n'\to b_\ast,\ b_n''\to b_\ast
  \Rightarrow
  s(b_n')\to s(b_\ast),\ s(b_n'')\to s(b_\ast),
\]
故
\[
  d_Y\bigl(s(b_n'),s(b_n'')\bigr)\to 0,
\]
这就是 \(L_5\)。
但由于 $s(b_n')\in \Gamma_t^{\Lambda_t}(b_n')$、$s(b_n'')\in \Gamma_t^{\Lambda_t}(b_n'')$，
而 \(L_4\) 要求任意可容许的两支终端值始终保持至少 $\varepsilon_0$ 的分离，
便与
\[
  d_Y\bigl(s(b_n'),s(b_n'')\bigr)\to 0
\]
矛盾。
故不存在这样的连续局部截面，即得 \(L_6\)。证毕。
\end{proof}

\subsection{命题：边缘算子改变有效域、失效域与换模域}
\label{subsec:tex1-ch1-proposition}

\begin{proposition}[钱学森障碍的边缘算子相变命题]
\label{prop:tex1-edge-phase}
固定 $t<T$、阈值 $\varepsilon>0$ 与边缘算子
\[
  \Delta_t:\Lambda_t\mapsto\Lambda_t'.
\]
若存在某个观测历史 $b_\ast$ 满足
\[
  \diameter\Gamma_t^{\Lambda_t}(b_\ast)>\varepsilon
  \qquad\text{且}\qquad
  \diameter\Gamma_t^{\Lambda_t'}(b_\ast)\le \varepsilon,
\]
或者反向不等式成立，
则
\[
  b_\ast\in \mathcal S_t^{\Delta_t}(\varepsilon).
\]
进一步地，若 $\Delta_t$ 改变了隐藏模式集合或其切换规则，
则系统的最小一致表达不应写成单一全局方程
\[
  \dot x=f(x,u,\theta,t),
\]
而应写成切换族或微分包含
\[
  \dot x=f_{\sigma(t)}(x,u,\theta,t),
  \qquad
  \sigma(t)\in\mathcal S_t,
\]
或
\[
  \dot x\in F(x,u,\theta,t).
\]
\end{proposition}

\begin{Certificate}[证书：纤维直径跨阈值 $\Rightarrow$ 域切换；模式变化 $\Rightarrow$ 混合表达]
\label{cert:tex1-edge-phase}
证书登记谓词链：
\[
\begin{aligned}
  P_1 &: \diameter\Gamma_t^{\Lambda_t}(b_\ast)>\varepsilon
        \ \wedge\
        \diameter\Gamma_t^{\Lambda_t'}(b_\ast)\le \varepsilon,\\
  P_2 &: P_1\Rightarrow
        b_\ast\notin \mathcal E_t^{\Lambda_t}(\varepsilon)
        \ \wedge\
        b_\ast\in \mathcal E_t^{\Lambda_t'}(\varepsilon),\\
  P_3 &: P_2\Rightarrow
        b_\ast\in
        \mathcal E_t^{\Lambda_t}(\varepsilon)
        \triangle
        \mathcal E_t^{\Lambda_t'}(\varepsilon)
        =
        \mathcal S_t^{\Delta_t}(\varepsilon),\\
  P_4 &: \Delta_t\ \text{改变模式集合或切换规则},\\
  P_5 &: P_4\Rightarrow
        \text{候选动力学不再由单一 }f\text{ 穷尽，而必须由 } \{f_\sigma\}\text{ 或 }F\text{ 表达}.
\end{aligned}
\]
\end{Certificate}

\begin{proof}[证明（基于数学符号推演逻辑的谓词因果链）]
若 \(P_1\) 成立，则按定义~\ref{def:tex1-edge-operator}，
\[
  \diameter\Gamma_t^{\Lambda_t}(b_\ast)>\varepsilon
  \Rightarrow
  b_\ast\notin \mathcal E_t^{\Lambda_t}(\varepsilon),
\]
且
\[
  \diameter\Gamma_t^{\Lambda_t'}(b_\ast)\le \varepsilon
  \Rightarrow
  b_\ast\in \mathcal E_t^{\Lambda_t'}(\varepsilon),
\]
故 \(P_2\) 成立。
由对称差的定义，立得 \(P_3\)：
\[
  b_\ast\in \mathcal S_t^{\Delta_t}(\varepsilon).
\]
反向不等式的情形完全对称。

若再有 \(P_4\)，则边缘算子已经改变可用模式集合 $\mathcal S_t$ 或模式切换规则。
此时，单一向量场 $f$ 无法同时编码“当前处于哪一模式”与“何时切换模式”两类信息；
因此必须引入模式索引 $\sigma(t)$ 或集合值右端 $F$，
这就得到 \(P_5\)。
于是命题成立。证毕。
\end{proof}

\subsection{推论：未显式建模能力元组与边缘算子时，只能得到局部近似}
\label{subsec:tex1-ch1-corollary}

\begin{corollary}[钱学森弹道/钱学森障碍的全局认知边界推论]
\label{cor:tex1-global-cognition}
设分析者采用的代理能力元组与代理边缘算子分别为
\[
  \widehat\Lambda_t,
  \qquad
  \widehat\Delta_t,
\]
而真实对象分别为
\[
  \Lambda_t,
  \qquad
  \Delta_t.
\]
若存在区间 $I\subset[0,T)$、阈值 $\varepsilon>0$、时刻 $t_\ast\in I$ 及历史 $b_\ast$，
使得
\[
  b_\ast\in
  \mathcal E_{t_\ast}^{\widehat\Lambda_{t_\ast}}(\varepsilon)
  \triangle
  \mathcal E_{t_\ast}^{\Lambda_{t_\ast}}(\varepsilon),
\]
或
\[
  b_\ast\in
  \mathcal S_{t_\ast}^{\widehat\Delta_{t_\ast}}(\varepsilon)
  \triangle
  \mathcal S_{t_\ast}^{\Delta_{t_\ast}}(\varepsilon),
\]
则基于代理对象得到的“有效域/失效域/换模域”判断，不是对真实对象的严格全局结论。
特别地，若在建模中完全省略 $\Lambda_t,\Delta_t$，
则任何单一、静态、全局唯一的方程表述，至多只能被解释为局部近似，而不能被解释为严格全局模型。
\end{corollary}

\begin{Certificate}[证书：代理域与真实域不一致 $\Rightarrow$ 全局判断失真]
\label{cert:tex1-global-cognition}
证书登记谓词链：
\[
\begin{aligned}
  C_1 &: b_\ast\in
        \mathcal E_{t_\ast}^{\widehat\Lambda_{t_\ast}}(\varepsilon)
        \triangle
        \mathcal E_{t_\ast}^{\Lambda_{t_\ast}}(\varepsilon)
        \ \text{或}\
        b_\ast\in
        \mathcal S_{t_\ast}^{\widehat\Delta_{t_\ast}}(\varepsilon)
        \triangle
        \mathcal S_{t_\ast}^{\Delta_{t_\ast}}(\varepsilon),\\
  C_2 &: C_1\Rightarrow
        \text{代理对象与真实对象在 }b_\ast\text{ 处给出不同的域归属},\\
  C_3 &: C_2\Rightarrow
        \text{由代理对象推出的全局域判断不能严格代表真实对象},\\
  C_4 &: \Lambda_t,\Delta_t\ \text{被完全省略}
        \Rightarrow
        \text{模型甚至没有给出域归属所依赖的对象},\\
  C_5 &: C_3\wedge C_4
        \Rightarrow
        \text{单一静态全局方程至多是局部近似}.
\end{aligned}
\]
\end{Certificate}

\begin{proof}[证明（基于数学符号推演逻辑的谓词因果链）]
由 \(C_1\)，在某个具体历史 $b_\ast$ 上，代理对象与真实对象已经给出不同的域归属。
换言之，
\[
  b_\ast
  \text{ 在代理图景中与在真实图景中并不属于同一推断区域。}
\]
这就是 \(C_2\)。

既然存在至少一个 $b_\ast$ 使代理与真实的域判定不同，
则任何声称“代理对象给出的域划分就是全局真域划分”的说法都不成立，
故 \(C_3\) 成立。
若进一步把 $\Lambda_t,\Delta_t$ 从模型中完全省略，
则模型根本没有显式记录观测接口、模型类、资源--时延条件与边缘切换规则，
于是域归属所依赖的对象层本身缺失，这就是 \(C_4\)。

把 \(C_3\) 与 \(C_4\) 合并，便得到 \(C_5\)：
在这种情形下，单一、静态、全局唯一的方程至多只是局部近似，
而不能承担严格全局模型的证明义务。证毕。
\end{proof}

\subsection{备注：最适合保留的压缩表述}
\label{subsec:tex1-ch1-remarks}

\begin{remark}[最短硬表达]
本章最终建议保留的最短硬表达是：
\[
  \boxed{
    \begin{aligned}
      &\text{所谓“钱学森弹道”若只停留在弹道/轨迹表象层，则并未触及问题本质；}\\
      &\text{到达终端之前是否可稳定推出唯一终端结论，}\\
      &\text{由集合值逆映射 }b\mapsto \Gamma_t^{\Lambda_t}(b)\text{ 的局部单值化决定；}\\
      &\text{所谓“钱学森障碍”可严格写成该单值化失败的局部或全局截面障碍；}\\
      &\text{边缘算子 }\Delta_t\text{ 通过改变能力元组而改变有效域、失效域与换模域；}\\
      &\text{修复塔的职责，是按层补齐动力学、观测、模式、换模与资源时延信息，}\\
      &\text{直到得到连续局部截面或失败证书。}
    \end{aligned}
  }.
\]
\end{remark}

\begin{remark}[术语收紧]
因此，以下三种说法应当被严格区分：
\begin{enumerate}
  \item ``钱学森弹道'' 或 ``弹道'' 若仅指轨迹形状、轨迹簇外观或高阶偏导写法，则它只描述某种表示方式；
  \item ``可容许终端纤维''、``局部/全局截面'' 与 ``钱学森障碍'' 描述的是推断问题本身；
  \item ``边缘算子'' 与 ``修复塔'' 描述的是域边界如何变化以及如何被分层补齐。
\end{enumerate}
在正式写作中，应把第二、三类对象置于第一类对象之前。
\end{remark}

\clearpage

%%%%%%%%%%%%%%%%%%%%%%%%%%%%%%%%%%%%%%%%%%%%%%%%%%%%%%%%%%%%
\clearpage
\section{形式逻辑：钱学森弹道/钱学森障碍的混合系统、全局截面障碍、修复塔与预算单调性}
\label{sec:tex1-ch2}
%%%%%%%%%%%%%%%%%%%%%%%%%%%%%%%%%%%%%%%%%%%%%%%%%%%%%%%%%%%%

\begin{remark}[本章意义备注：高超音速攻防的抽象语境]
本章的意义，在于说明高超音速攻防中的困难常来自隐藏模式、换模与预算限制，
而不只是连续动力学本身。
对钱学森障碍而言，这一章解释了为何“看起来可解释的前段轨迹”仍可能对应多个终端模式，
并说明修复必须沿分层对象链推进。
\end{remark}

\subsection{严格主张、引文定位与章节目标}
\label{subsec:tex1-ch2-claim}

\begin{remark}[本章定位]
本章把钱学森弹道/钱学森障碍进一步压成一个完全脱离具体应用语境的对象链：
\[
  \text{混合系统}
  \Longrightarrow
  \text{可容许终端纤维}
  \Longrightarrow
  \text{全局截面障碍}
  \Longrightarrow
  \text{边缘算子}
  \Longrightarrow
  \text{修复塔与预算单调}.
\]
与第~\ref{sec:tex1-ch1} 章相比，本章不再追问“高阶偏导是不是本质”，
而直接给出五类对象的第一版书稿化形式：
\begin{enumerate}
  \item 观测历史空间与可容许终端纤维；
  \item 局部歧义障碍与全局截面障碍；
  \item 隐藏模式未显式入状态时的闭合失败；
  \item 真值保持且单调收缩的修复塔；
  \item 算力/时延预算扩张所诱导的最优剩余歧义单调性。
\end{enumerate}
本章仍然保留 ``钱学森弹道''、``钱学森障碍''、``弹道'' 这些专业应用名词，
但只把它们作为对象命名标签，而不落入任何实现性语境。
仓内接口主要承接 \cite{RepoTex29RepairableObstruction,RepoTex32FiberSectionHole,RepoTex38JacobiBianchiRepair,
  RepoTex17EdgeOperator,RepoTex18HomotopyAxiom}；
外部拓扑/障碍理论背景可参见 \cite{Husemoller1994FibreBundles,Hatcher2002AT}。
\end{remark}

\begin{remark}[本章不设公理]
为避免把尚未构造完成的内容误写成已证定理，
本章不单独设公理，只给出定义、引理、定理、命题、推论与备注。
凡涉及上同调升级、障碍类构造或更强修复结论而本章尚未给出明确对象者，
一律保留为\textbf{未证明升级方向}。
\end{remark}

\subsection{定义组：混合系统、观测历史与可容许终端纤维}
\label{subsec:tex1-ch2-defs-system}

\begin{definition}[混合系统的基础设定]
\label{def:tex1-ch2-hybrid}
设
\[
  X\ \text{为状态空间},\qquad
  U\ \text{为控制空间},\qquad
  \Theta\ \text{为参数空间},
\]
\[
  S\ \text{为有限隐藏模式集},\qquad
  Z\ \text{为观测空间},\qquad
  Y\ \text{为终端标签空间}.
\]
其中 $Y$ 取有限离散拓扑。给定终端时刻 $T>0$。

对任意 $t\in[0,T)$，考虑混合系统
\[
  \dot x(\tau)=f_{\sigma(\tau)}(x(\tau),u(\tau),\theta,\tau),
  \qquad \tau\in[0,T],
\]
\[
  z(\tau)=h_{\sigma(\tau)}(x(\tau),\tau)+\eta(\tau),
\]
\[
  y=\psi(x(T))\in Y.
\]
这里 $\sigma:[0,T]\to S$ 为隐藏模式轨道，$\eta$ 为扰动，$\psi:X\to Y$ 为终端投影。
\end{definition}

\begin{definition}[观测历史空间]
\label{def:tex1-ch2-history}
对任意 $t<T$，定义观测历史空间
\[
  \mathcal B_t:=D([0,t],Z),
\]
即到时刻 $t$ 为止的右连续有左极限观测轨道空间。
若讨论对象不含跳变，也可将其替换为
\[
  C([0,t],Z).
\]
\end{definition}

\begin{definition}[可容许终端纤维]
\label{def:tex1-ch2-gamma}
对任意 $b\in\mathcal B_t$，定义
\[
  \Gamma_t(b):=
  \Bigl\{
    y\in Y
    \,\Big|\,
    \exists(x_0,u,\theta,\sigma,\eta)
    \text{ 使系统方程成立，且 }
    z|_{[0,t]}=b,\ \psi(x(T))=y
  \Bigr\}.
\]
它表示：在给定观测历史 $b$ 时，到终端仍然与系统一致的全部终端标签集合。
\end{definition}

\begin{definition}[局部障碍与全局截面障碍]
\label{def:tex1-ch2-section-obstruction}
若对某个 $b\in\mathcal B_t$ 有
\[
  |\Gamma_t(b)|\ge 2,
\]
则称该点存在\textbf{局部歧义障碍}。

定义总空间
\[
  E_t:=\{(b,y)\in\mathcal B_t\times Y:\ y\in\Gamma_t(b)\},
\]
以及投影
\[
  p_t:E_t\to\mathcal B_t,
  \qquad
  p_t(b,y)=b.
\]
若不存在连续映射
\[
  s:\mathcal B_t\to E_t,
  \qquad
  p_t\circ s=\id_{\mathcal B_t},
\]
则称 $p_t:E_t\to\mathcal B_t$ 存在\textbf{全局截面障碍}。
在这一口径下，钱学森障碍可被重新理解为：
\[
  \text{从观测历史到唯一终端标签的连续全局选择是否存在。}
\]
\end{definition}

\begin{definition}[观察方能力元组与边缘算子]
\label{def:tex1-ch2-capability-edge}
定义观察方在时刻 $t$ 的能力元组
\[
  \Lambda_t=(\mathcal O_t,\mathcal M_t,\Pi_t,\mathcal C_t,\mathcal L_t),
\]
其中：
\[
  \mathcal O_t\ \text{为观测算子类},\qquad
  \mathcal M_t\ \text{为模型类},
\]
\[
  \Pi_t\ \text{为先验集合},\qquad
  \mathcal C_t\ \text{为算力预算},\qquad
  \mathcal L_t\ \text{为时延预算}.
\]
于是更精确的可容许终端纤维应记为
\[
  \Gamma_t^{\Lambda}(b).
\]

定义边缘算子为作用在能力元组或系统约束上的变换
\[
  \mathfrak e:\Lambda\mapsto \mathfrak e(\Lambda),
\]
它诱导出可容许终端纤维的变化
\[
  \Gamma_t^{\Lambda}(b)\mapsto \Gamma_t^{\mathfrak e(\Lambda)}(b).
\]
因此，边缘算子的本质不是“改变一条已知曲线的某一段”，而是
\[
  \text{改变逆问题的可适定性边界、可观测性边界与可辨识性边界。}
\]
\end{definition}

\subsection{引理：隐藏模式未入状态时，方程不闭合}
\label{subsec:tex1-ch2-lemma-closure}

\begin{lemma}[隐藏模式未入状态时，方程不闭合]
\label{lem:tex1-ch2-hidden-mode}
若存在
\[
  \sigma_1\neq \sigma_2,
  \qquad
  (x,u,\theta,\tau)
\]
使得
\[
  f_{\sigma_1}(x,u,\theta,\tau)\neq f_{\sigma_2}(x,u,\theta,\tau),
\]
则不存在单值映射
\[
  F:X\times U\times \Theta\times [0,T]\to TX
\]
使得对所有可容许轨道都有
\[
  \dot x(\tau)=F(x(\tau),u(\tau),\theta,\tau).
\]
\end{lemma}

\begin{Certificate}[证书：反证法证书]
\label{cert:tex1-ch2-hidden-mode}
证书登记谓词链：
\[
\begin{aligned}
  L_1 &: \exists \sigma_1\neq \sigma_2,\ (x,u,\theta,\tau),\\
  L_2 &: f_{\sigma_1}(x,u,\theta,\tau)\neq f_{\sigma_2}(x,u,\theta,\tau),\\
  L_3 &: \text{反设存在单值 }F:X\times U\times \Theta\times[0,T]\to TX,\\
  L_4 &: L_3\Rightarrow
        F(x,u,\theta,\tau)=f_{\sigma_1}(x,u,\theta,\tau)
        =f_{\sigma_2}(x,u,\theta,\tau),\\
  L_5 &: L_2\wedge L_4\Rightarrow \bot.
\end{aligned}
\]
\end{Certificate}

\begin{proof}[证明（基于数学符号推演逻辑的谓词因果链）]
由 \(L_1\) 与 \(L_2\)，存在同一组状态/控制/参数/时间数据，
在两个不同隐藏模式下给出不同动力学右端。

反设存在单值映射
\[
  F:X\times U\times\Theta\times[0,T]\to TX
\]
可统一表示全部可容许轨道。
则在同一组 $(x,u,\theta,\tau)$ 上，必同时有
\[
  F(x,u,\theta,\tau)=f_{\sigma_1}(x,u,\theta,\tau),
\]
以及
\[
  F(x,u,\theta,\tau)=f_{\sigma_2}(x,u,\theta,\tau),
\]
从而推出
\[
  f_{\sigma_1}(x,u,\theta,\tau)=f_{\sigma_2}(x,u,\theta,\tau).
\]
这与 \(L_2\) 矛盾，故 \(L_5\) 成立。
因此不存在这样的单值 $F$，结论得证。
\end{proof}

\begin{corollary}[隐藏模式、能力更新未并入扩展状态时，单一方程一般不闭合]
\label{cor:tex1-ch2-hidden-mode}
若隐藏模式、切换时刻、能力更新规则未被并入扩展状态，
则只在 $x$ 上书写的偏微分方程或常微分方程一般不是闭合方程。
\end{corollary}

\begin{Certificate}[证书：引理~\ref{lem:tex1-ch2-hidden-mode} 的直接调用]
\label{cert:tex1-ch2-hidden-mode-cor}
证书登记谓词链：
\[
\begin{aligned}
  C_1 &: \text{隐藏模式或能力更新未进入扩展状态变量},\\
  C_2 &: C_1\Rightarrow \text{存在同一 }(x,u,\theta,\tau)\text{ 下的多种候选右端},\\
  C_3 &: C_2\Rightarrow \text{调用引理~\ref{lem:tex1-ch2-hidden-mode}},\\
  C_4 &: C_3\Rightarrow \text{单一方程一般不闭合}.
\end{aligned}
\]
\end{Certificate}

\begin{proof}[证明（基于数学符号推演逻辑的谓词因果链）]
由 \(C_1\) 可知，模式差异与能力更新没有被纳入状态描述；
因此同一可见状态数据可能对应不同的内部演化右端，即 \(C_2\)。
由 \(C_2\) 调用引理~\ref{lem:tex1-ch2-hidden-mode}，得到 \(C_3\)；
于是立即推出 \(C_4\)，即该推论成立。
\end{proof}

\subsection{引理：连通基底到离散目标的连续映射必局部常值}
\label{subsec:tex1-ch2-lemma-connected}

\begin{lemma}[连通基底到离散目标的连续映射必为常值]
\label{lem:tex1-ch2-connected-discrete}
若 $B$ 连通，$Y$ 为离散拓扑空间，且
\[
  \varphi:B\to Y
\]
连续，则 $\varphi(B)$ 为离散空间中的连通子集，从而必为单点。
\end{lemma}

\begin{Certificate}[证书：拓扑连通性证书]
\label{cert:tex1-ch2-connected-discrete}
证书登记谓词链：
\[
\begin{aligned}
  L'_1 &: B\ \text{连通},\\
  L'_2 &: \varphi:B\to Y\ \text{连续},\\
  L'_3 &: L'_1\wedge L'_2\Rightarrow \varphi(B)\ \text{连通},\\
  L'_4 &: Y\ \text{离散}\Rightarrow \text{含两个不同点的子集都不连通},\\
  L'_5 &: L'_3\wedge L'_4\Rightarrow \varphi(B)\ \text{只能为单点}.
\end{aligned}
\]
\end{Certificate}

\begin{proof}[证明（基于数学符号推演逻辑的谓词因果链）]
由 \(L'_1\) 与 \(L'_2\)，连续映射将连通集映为连通集，故 \(L'_3\) 成立。
又因 $Y$ 为离散空间，任何含两个不同点的子集都可被分离为两个开闭子集，
所以必不连通，这就是 \(L'_4\)。
因此由 \(L'_3\) 与 \(L'_4\) 可得 \(L'_5\)：$\varphi(B)$ 只能为单点。
于是 $\varphi$ 必为常值映射。证毕。
\end{proof}

\subsection{定理：全局截面障碍的充分条件}
\label{subsec:tex1-ch2-theorem-section}

\begin{theorem}[全局截面障碍的充分条件]
\label{thm:tex1-ch2-section-obstruction}
设 $\mathcal B_t$ 连通，$Y$ 离散。
若存在
\[
  b_0,b_1\in\mathcal B_t
\]
满足
\[
  \Gamma_t(b_0)=\{y_0\},
  \qquad
  \Gamma_t(b_1)=\{y_1\},
  \qquad
  y_0\neq y_1,
\]
则投影
\[
  p_t:E_t\to\mathcal B_t
\]
不存在连续全局截面。
\end{theorem}

\begin{Certificate}[证书：引理~\ref{lem:tex1-ch2-connected-discrete} + 反证法证书]
\label{cert:tex1-ch2-section-obstruction}
证书登记谓词链：
\[
\begin{aligned}
  T_1 &: \mathcal B_t\ \text{连通},\quad Y\ \text{离散},\\
  T_2 &: \Gamma_t(b_0)=\{y_0\},\ \Gamma_t(b_1)=\{y_1\},\ y_0\neq y_1,\\
  T_3 &: \text{反设存在连续全局截面 }s:\mathcal B_t\to E_t,\\
  T_4 &: T_3\Rightarrow \exists \bar y:\mathcal B_t\to Y\ \text{连续，且 }s(b)=(b,\bar y(b)),\\
  T_5 &: T_1\wedge T_4\Rightarrow \bar y\ \text{常值（引理~\ref{lem:tex1-ch2-connected-discrete}）},\\
  T_6 &: T_2\wedge T_4\Rightarrow \bar y(b_0)=y_0,\ \bar y(b_1)=y_1,\\
  T_7 &: T_5\wedge T_6\Rightarrow \bot.
\end{aligned}
\]
\end{Certificate}

\begin{proof}[证明（基于数学符号推演逻辑的谓词因果链）]
反设存在连续全局截面
\[
  s:\mathcal B_t\to E_t,
  \qquad
  p_t\circ s=\id_{\mathcal B_t}.
\]
则每个 $b\in\mathcal B_t$ 都可唯一写成
\[
  s(b)=(b,\bar y(b)),
\]
其中
\[
  \bar y:\mathcal B_t\to Y
\]
连续，且满足
\[
  \bar y(b)\in\Gamma_t(b).
\]
这就是 \(T_4\)。

由 \(T_1\) 与引理~\ref{lem:tex1-ch2-connected-discrete}，$\bar y$ 必为常值，即 \(T_5\)。
另一方面，由 \(T_2\) 可知，任意截面都必须满足
\[
  \bar y(b_0)=y_0,
  \qquad
  \bar y(b_1)=y_1.
\]
这就是 \(T_6\)。
但 \(y_0\neq y_1\)，所以 \(\bar y\) 不可能同时常值且在 $b_0,b_1$ 取不同值，
从而得到 \(T_7\) 的矛盾。
故连续全局截面不存在，定理得证。
\end{proof}

\begin{corollary}[连通观测区域内的双唯一标签强制导致全局选择失败]
\label{cor:tex1-ch2-section-obstruction}
只要在同一连通观测区域内，存在两个历史点被系统强制指向两个不同的唯一终端标签，
则“从历史到唯一终端标签”的全局连续选择必然失败。
这已经构成严格的全局截面障碍。
\end{corollary}

\begin{Certificate}[证书：定理~\ref{thm:tex1-ch2-section-obstruction} 的直接调用]
\label{cert:tex1-ch2-section-obstruction-cor}
证书登记谓词链：
\[
\begin{aligned}
  C'_1 &: \exists b_0,b_1\in\mathcal B_t,\ \Gamma_t(b_0)=\{y_0\},\ \Gamma_t(b_1)=\{y_1\},\ y_0\neq y_1,\\
  C'_2 &: \mathcal B_t\ \text{连通},\\
  C'_3 &: C'_1\wedge C'_2\Rightarrow \text{调用定理~\ref{thm:tex1-ch2-section-obstruction}},\\
  C'_4 &: C'_3\Rightarrow \text{全局连续选择失败}.
\end{aligned}
\]
\end{Certificate}

\begin{proof}[证明（基于数学符号推演逻辑的谓词因果链）]
由 \(C'_1\) 与 \(C'_2\) 正好满足定理~\ref{thm:tex1-ch2-section-obstruction} 的全部前提，
于是由 \(C'_3\) 直接推出 \(C'_4\)。
故该推论成立。证毕。
\end{proof}

\subsection{定义：修复塔}
\label{subsec:tex1-ch2-def-repair}

\begin{definition}[修复塔]
\label{def:tex1-ch2-repair-tower}
固定一个指定真值选择
\[
  y_\star:\mathcal B_t\to Y,
  \qquad
  y_\star(b)\in \Gamma_t^{(0)}(b),
\]
并定义一列集合值映射
\[
  \Gamma_t^{(0)}:=\Gamma_t,
  \qquad
  \Gamma_t^{(n+1)}:=R_n\bigl(\Gamma_t^{(n)}\bigr),
  \quad n\in\NN,
\]
其中每个 $R_n$ 满足：
\[
  \textnormal{(R1) 真值保持： }
  y_\star(b)\in\Gamma_t^{(n)}(b)
  \Rightarrow
  y_\star(b)\in\Gamma_t^{(n+1)}(b),
\]
\[
  \textnormal{(R2) 单调收缩： }
  \Gamma_t^{(n+1)}(b)\subseteq \Gamma_t^{(n)}(b),
\]
\[
  \textnormal{(R3) 严格降歧义： }
  |\Gamma_t^{(n)}(b)|>1
  \Rightarrow
  |\Gamma_t^{(n+1)}(b)|<|\Gamma_t^{(n)}(b)|.
\]
这就是抽象意义下的修复塔。
\end{definition}

\subsection{引理：修复塔嵌套性}
\label{subsec:tex1-ch2-lemma-repair}

\begin{lemma}[修复塔嵌套性]
\label{lem:tex1-ch2-repair-nested}
对任意 $n\in\NN$ 与任意 $b\in\mathcal B_t$，有
\[
  \Gamma_t^{(n)}(b)\subseteq \Gamma_t^{(0)}(b).
\]
更强地，若 $m\le n$，则
\[
  \Gamma_t^{(n)}(b)\subseteq \Gamma_t^{(m)}(b).
\]
\end{lemma}

\begin{Certificate}[证书：关于 \(n\) 的数学归纳法证书]
\label{cert:tex1-ch2-repair-nested}
定义谓词
\[
  P(n):
  \forall b\in\mathcal B_t,\ \Gamma_t^{(n)}(b)\subseteq \Gamma_t^{(0)}(b).
\]
证书登记为
\[
  \bigl(P(0),\ \forall n\in\NN,\ P(n)\Rightarrow P(n+1)\bigr),
\]
并结合 \textnormal{(R2)} 得到任意 $m\le n$ 时的层间嵌套。
\end{Certificate}

\begin{proof}[证明（基于数学符号推演逻辑的谓词因果链）]
\textbf{基步：}当 $n=0$ 时，
\[
  \Gamma_t^{(0)}(b)\subseteq \Gamma_t^{(0)}(b)
\]
显然成立，故 $P(0)$ 成立。

\textbf{归纳步：}设 $P(n)$ 成立，即对任意 $b$，
\[
  \Gamma_t^{(n)}(b)\subseteq \Gamma_t^{(0)}(b).
\]
由修复塔定义中的 \textnormal{(R2)}，
\[
  \Gamma_t^{(n+1)}(b)\subseteq \Gamma_t^{(n)}(b).
\]
与归纳假设联立，得到
\[
  \Gamma_t^{(n+1)}(b)\subseteq \Gamma_t^{(n)}(b)\subseteq \Gamma_t^{(0)}(b).
\]
故 $P(n+1)$ 成立。

于是由数学归纳法，$\forall n\in\NN,\ P(n)$。
更强结论也由 \textnormal{(R2)} 反复作用直接给出：
若 $m\le n$，则
\[
  \Gamma_t^{(n)}(b)\subseteq \Gamma_t^{(n-1)}(b)\subseteq\cdots\subseteq \Gamma_t^{(m)}(b).
\]
引理得证。
\end{proof}

\subsection{定理：有限标签下修复塔有限收敛}
\label{subsec:tex1-ch2-theorem-repair}

\begin{theorem}[有限标签下修复塔有限收敛]
\label{thm:tex1-ch2-repair-finite}
若 $Y$ 有限，且修复塔满足 \textnormal{(R1)(R2)(R3)}，
则对任意 $b\in\mathcal B_t$，序列
\[
  |\Gamma_t^{(n)}(b)|
\]
在至多
\[
  |Y|-1
\]
步后收敛到 $1$。
即存在
\[
  N(b)\le |Y|-1,
\]
使得
\[
  |\Gamma_t^{(N(b))}(b)|=1.
\]
\end{theorem}

\begin{Certificate}[证书：自然数严格下降证书]
\label{cert:tex1-ch2-repair-finite}
证书登记谓词链：
\[
\begin{aligned}
  T''_1 &: a_n:=|\Gamma_t^{(n)}(b)|,\qquad a_0\le |Y|,\\
  T''_2 &: \textnormal{(R1)}\Rightarrow a_n\ge 1,\\
  T''_3 &: a_n>1\ \wedge\ \textnormal{(R3)}\Rightarrow a_{n+1}<a_n,\\
  T''_4 &: \{a_n\}\subseteq \NN_{>0}\ \wedge\ \text{严格下降不可能无限持续},\\
  T''_5 &: T''_1\wedge T''_2\wedge T''_3\wedge T''_4
          \Rightarrow \exists N(b)\le |Y|-1,\ a_{N(b)}=1.
\end{aligned}
\]
\end{Certificate}

\begin{proof}[证明（基于数学符号推演逻辑的谓词因果链）]
固定 $b\in\mathcal B_t$，记
\[
  a_n:=|\Gamma_t^{(n)}(b)|.
\]
由 \(T''_1\) 可知
\[
  a_0=|\Gamma_t^{(0)}(b)|\le |Y|.
\]
由真值保持 \textnormal{(R1)}，指定真值 $y_\star(b)$ 永远不被删除，
故对任意 $n$ 都有
\[
  a_n\ge 1,
\]
这就是 \(T''_2\)。

只要 $a_n>1$，由严格降歧义 \textnormal{(R3)} 可知
\[
  a_{n+1}<a_n,
\]
即 \(T''_3\)。
于是 $\{a_n\}$ 成为一个取值于正整数的严格下降序列。
但正整数中的严格下降序列不可能无限持续，这就是 \(T''_4\)。
因此，序列必须在有限步后达到其最小可能值 $1$。

又因为初值至多为 $|Y|$，每次严格下降至少减少 $1$，
所以最多只能下降 $|Y|-1$ 次。
故存在
\[
  N(b)\le |Y|-1
\]
使得
\[
  a_{N(b)}=1.
\]
这正是 \(T''_5\)，定理得证。
\end{proof}

\begin{corollary}[有限标签下的修复塔不是修辞，而是有限收敛机制]
\label{cor:tex1-ch2-repair-finite}
在有限标签空间内，只要每一级修复都保持真值并严格减少歧义，
则修复塔不是修辞，而是可严格证明的有限收敛机制。
\end{corollary}

\begin{Certificate}[证书：定理~\ref{thm:tex1-ch2-repair-finite} 的直接推论]
\label{cert:tex1-ch2-repair-finite-cor}
证书登记谓词链：
\[
\begin{aligned}
  C''_1 &: Y\ \text{有限},\\
  C''_2 &: \Gamma_t^{(n)}\ \text{满足 \textnormal{(R1)(R2)(R3)}},\\
  C''_3 &: C''_1\wedge C''_2\Rightarrow \text{调用定理~\ref{thm:tex1-ch2-repair-finite}},\\
  C''_4 &: C''_3\Rightarrow \text{修复塔在每个 }b\text{ 上有限步收敛到单点}.
\end{aligned}
\]
\end{Certificate}

\begin{proof}[证明（基于数学符号推演逻辑的谓词因果链）]
由 \(C''_1\) 与 \(C''_2\) 直接满足定理~\ref{thm:tex1-ch2-repair-finite} 的全部假设，
故由 \(C''_3\) 立得 \(C''_4\)。
因此该推论成立。证毕。
\end{proof}

\subsection{定理：同一基底上的内部形变不能凭空消除截面障碍}
\label{subsec:tex1-ch2-theorem-homotopy}

\begin{theorem}[仅在同一基底上做内部形变，不能凭空消除截面障碍]
\label{thm:tex1-ch2-no-free-section}
设
\[
  q:E'\to B,
  \qquad
  p:E\to B
\]
为同一基底 $B$ 上的两个投影。
若存在连续映射
\[
  g:E'\to E,
  \qquad
  p\circ g=q,
\]
且 $q$ 有连续截面
\[
  s':B\to E',
  \qquad
  q\circ s'=\id_B,
\]
则 $p$ 也有连续截面。
\end{theorem}

\begin{Certificate}[证书：直接构造证书]
\label{cert:tex1-ch2-no-free-section}
证书登记谓词链：
\[
\begin{aligned}
  T'''_1 &: g:E'\to E\ \text{连续},\quad p\circ g=q,\\
  T'''_2 &: s':B\to E'\ \text{连续},\quad q\circ s'=\id_B,\\
  T'''_3 &: \text{定义 }s:=g\circ s':B\to E,\\
  T'''_4 &: T'''_1\wedge T'''_2\wedge T'''_3
           \Rightarrow p\circ s=p\circ g\circ s'=q\circ s'=\id_B,\\
  T'''_5 &: T'''_4\Rightarrow s\ \text{是 }p\text{ 的连续截面}.
\end{aligned}
\]
\end{Certificate}

\begin{proof}[证明（基于数学符号推演逻辑的谓词因果链）]
按 \(T'''_3\) 定义
\[
  s:=g\circ s':B\to E.
\]
由于 $g$ 与 $s'$ 都连续，故 $s$ 连续。
再由 \(T'''_1\) 与 \(T'''_2\)，有
\[
  p\circ s
  =
  p\circ g\circ s'
  =
  q\circ s'
  =
  \id_B.
\]
这正是 \(T'''_4\)。
因此 $s$ 是 $p$ 的连续截面，即 \(T'''_5\)。
定理得证。
\end{proof}

\begin{corollary}[真正的障碍修复必须改变至少一个结构层]
\label{cor:tex1-ch2-no-free-section}
若原系统 $p:E\to B$ 已经不存在连续截面，
那么任何仅通过“同一基底上的内部形变”得到、且仍可经基底保持映射回 $E$ 的结构，
都不可能凭空产生连续截面。

换言之，真正的障碍修复必须改变至少一个结构层：
\[
  B,\qquad E,\qquad p,\qquad \Gamma,\qquad \Lambda
\]
中的至少一个，而不是只在原障碍结构内部做形式变形。
\end{corollary}

\begin{Certificate}[证书：逆否命题证书]
\label{cert:tex1-ch2-no-free-section-cor}
证书登记谓词链：
\[
\begin{aligned}
  C'''_1 &: p:E\to B\ \text{无连续截面},\\
  C'''_2 &: \text{若同一基底内部形变后的 }q:E'\to B\ \text{有连续截面},\\
  C'''_3 &: C'''_2\Rightarrow \text{由定理~\ref{thm:tex1-ch2-no-free-section} 可推出 }p\text{ 也有截面},\\
  C'''_4 &: C'''_1\wedge C'''_3\Rightarrow \bot,\\
  C'''_5 &: C'''_4\Rightarrow \text{真正修复必须改变至少一个结构层}.
\end{aligned}
\]
\end{Certificate}

\begin{proof}[证明（基于数学符号推演逻辑的谓词因果链）]
由 \(C'''_1\) 设定原系统 $p$ 不存在连续截面。
若反设某个同一基底上的内部形变后系统 $q$ 已存在连续截面，
则由 \(C'''_3\) 调用定理~\ref{thm:tex1-ch2-no-free-section}，
可反推出原系统 $p$ 也有连续截面。
这与 \(C'''_1\) 矛盾，于是得到 \(C'''_4\)。
因此只能有 \(C'''_5\)：若要真正消除截面障碍，必须改变至少一个结构层。证毕。
\end{proof}

\subsection{命题：不完整认知导致伪确定性}
\label{subsec:tex1-ch2-prop-pseudo}

\begin{proposition}[不完整认知导致伪确定性]
\label{prop:tex1-ch2-pseudo-determinacy}
设真实能力元组为 $\Lambda$，简化能力元组为 $\widehat\Lambda$。
若对某个 $b\in\mathcal B_t$ 有
\[
  \Gamma_t^{\widehat\Lambda}(b)\subsetneq \Gamma_t^{\Lambda}(b),
\]
且
\[
  |\Gamma_t^{\widehat\Lambda}(b)|=1,
  \qquad
  |\Gamma_t^{\Lambda}(b)|\ge 2,
\]
则简化模型在点 $b$ 上给出了\textbf{伪确定性}。
\end{proposition}

\begin{Certificate}[证书：集合严格包含证书]
\label{cert:tex1-ch2-pseudo-determinacy}
证书登记谓词链：
\[
\begin{aligned}
  P^\star_1 &: \Gamma_t^{\widehat\Lambda}(b)\subsetneq \Gamma_t^{\Lambda}(b),\\
  P^\star_2 &: |\Gamma_t^{\widehat\Lambda}(b)|=1,\\
  P^\star_3 &: |\Gamma_t^{\Lambda}(b)|\ge 2,\\
  P^\star_4 &: P^\star_2\wedge P^\star_3\Rightarrow
               \text{简化模型给出唯一结论，而真实系统仍保留多值歧义},\\
  P^\star_5 &: P^\star_1\wedge P^\star_4\Rightarrow \text{该唯一性是模型删减后的假唯一性}.
\end{aligned}
\]
\end{Certificate}

\begin{proof}[证明（基于数学符号推演逻辑的谓词因果链）]
由 \(P^\star_1\)，简化模型只保留了真实可容许终端纤维的一部分。
由 \(P^\star_2\)，该被删减后的纤维恰好收缩为单点；
而由 \(P^\star_3\)，真实系统的纤维仍至少含有两个可容许终端。
故 \(P^\star_4\) 成立：简化模型给出的唯一结论，并不等于真实系统的唯一结论。

再把 \(P^\star_1\) 与 \(P^\star_4\) 合并，便得到 \(P^\star_5\)：
这种“唯一性”只是删减能力元组之后得到的假唯一性，
即所谓伪确定性。命题得证。
\end{proof}

\begin{corollary}[未显式建模隐藏模式、边缘插入、算力扩张与时延边界时，判断易落入伪确定性]
\label{cor:tex1-ch2-pseudo-determinacy}
若未把隐藏模式、边缘算子插入时刻、算力扩张规则、时延边界写入系统，
则对生效条件、失效条件、最省算力策略与动态扩张策略的判断，都可能落入伪确定性。
\end{corollary}

\begin{Certificate}[证书：模型删减证书]
\label{cert:tex1-ch2-pseudo-determinacy-cor}
证书登记谓词链：
\[
\begin{aligned}
  C^\star_1 &: \text{隐藏模式/边缘插入/算力扩张/时延边界未被显式建模},\\
  C^\star_2 &: C^\star_1\Rightarrow \widehat\Lambda\ \text{是}\ \Lambda\ \text{的删减版本},\\
  C^\star_3 &: C^\star_2\Rightarrow \exists b,\ \Gamma_t^{\widehat\Lambda}(b)\subsetneq \Gamma_t^\Lambda(b)\ \text{是可能的},
    \\
  C^\star_4 &: C^\star_3\Rightarrow \text{调用命题~\ref{prop:tex1-ch2-pseudo-determinacy}},\\
  C^\star_5 &: C^\star_4\Rightarrow \text{相关判断可能落入伪确定性}.
\end{aligned}
\]
\end{Certificate}

\begin{proof}[证明（基于数学符号推演逻辑的谓词因果链）]
由 \(C^\star_1\) 可知，模型未显式写入一部分真实边界条件，
因此只是在真实能力元组基础上的删减近似，即 \(C^\star_2\)。
一旦发生删减，便可能出现
\[
  \Gamma_t^{\widehat\Lambda}(b)\subsetneq \Gamma_t^\Lambda(b),
\]
这就是 \(C^\star_3\)。
此时调用命题~\ref{prop:tex1-ch2-pseudo-determinacy}，得到 \(C^\star_4\)；
于是 \(C^\star_5\) 成立，即相关判断可能落入伪确定性。证毕。
\end{proof}

\subsection{命题：算力与时延预算的单调性}
\label{subsec:tex1-ch2-prop-budget}

\begin{proposition}[算力与时延预算的单调性]
\label{prop:tex1-ch2-budget}
定义预算可行修复族
\[
  \mathfrak R(C,L)
  :=
  \{R:\ \operatorname{cost}(R)\le C,\ \operatorname{latency}(R)\le L\}.
\]
给定歧义度量 $\mu$，例如
\[
  \mu(A)=|A|
  \qquad\text{或}\qquad
  \mu(A)=\log |A|.
\]
定义最优剩余歧义
\[
  A_{C,L}(b):=
  \inf_{R\in\mathfrak R(C,L)}
  \mu\bigl(R(\Gamma_t(b))\bigr).
\]
若
\[
  C_1\le C_2,
  \qquad
  L_1\le L_2,
\]
则有
\[
  A_{C_2,L_2}(b)\le A_{C_1,L_1}(b).
\]
\end{proposition}

\begin{Certificate}[证书：可行集包含证书]
\label{cert:tex1-ch2-budget}
证书登记谓词链：
\[
\begin{aligned}
  P^\dagger_1 &: C_1\le C_2,\quad L_1\le L_2,\\
  P^\dagger_2 &: P^\dagger_1\Rightarrow \mathfrak R(C_1,L_1)\subseteq \mathfrak R(C_2,L_2),\\
  P^\dagger_3 &: \text{对更大的可行集取 }\inf\ \text{不会得到更大的值},\\
  P^\dagger_4 &: P^\dagger_2\wedge P^\dagger_3\Rightarrow A_{C_2,L_2}(b)\le A_{C_1,L_1}(b).
\end{aligned}
\]
\end{Certificate}

\begin{proof}[证明（基于数学符号推演逻辑的谓词因果链）]
由 \(P^\dagger_1\) 立得
\[
  \mathfrak R(C_1,L_1)\subseteq \mathfrak R(C_2,L_2),
\]
这就是 \(P^\dagger_2\)。
由于 $A_{C,L}(b)$ 定义为在预算可行集上对剩余歧义所取的下确界，
而对更大的可行集取下确界不可能比对更小的可行集取下确界更大，
故 \(P^\dagger_3\) 成立。
将二者合并，即得
\[
  A_{C_2,L_2}(b)\le A_{C_1,L_1}(b),
\]
即 \(P^\dagger_4\)。命题得证。
\end{proof}

\begin{corollary}[算力扩张与时延放宽可被严格重写为最优剩余歧义单调下降]
\label{cor:tex1-ch2-budget}
“节约算力”“扩张算力”“动态适应”都可以被严格重写为：
预算可行修复族的单调扩展，以及由此诱导的最优剩余歧义单调下降。
\end{corollary}

\begin{Certificate}[证书：命题~\ref{prop:tex1-ch2-budget} 的语义改写]
\label{cert:tex1-ch2-budget-cor}
证书登记谓词链：
\[
\begin{aligned}
  C^\dagger_1 &: \text{算力扩张/时延放宽} \Rightarrow (C,L)\ \text{单调增大},\\
  C^\dagger_2 &: C^\dagger_1\Rightarrow \text{调用命题~\ref{prop:tex1-ch2-budget}},\\
  C^\dagger_3 &: C^\dagger_2\Rightarrow A_{C,L}(b)\ \text{单调不增},\\
  C^\dagger_4 &: C^\dagger_3\Rightarrow \text{最优剩余歧义单调下降的严格改写成立}.
\end{aligned}
\]
\end{Certificate}

\begin{proof}[证明（基于数学符号推演逻辑的谓词因果链）]
由 \(C^\dagger_1\) 可知，所谓算力扩张或时延放宽，
在形式系统中只是预算参数 $(C,L)$ 的单调增大。
于是由 \(C^\dagger_2\) 调用命题~\ref{prop:tex1-ch2-budget}，
得到 \(A_{C,L}(b)\) 的单调不增，即 \(C^\dagger_3\)。
这正是“预算可行修复族扩张 $\Rightarrow$ 最优剩余歧义下降”的含义，
因此 \(C^\dagger_4\) 成立。证毕。
\end{proof}

\subsection{备注：边界、上同调升级条件与最短压缩版}
\label{subsec:tex1-ch2-remarks}

\begin{remark}[本章已严格建立的四个内核]
到这里为止，已经严格建立的是
\[
  \text{闭合性障碍}
  +
  \text{全局截面障碍}
  +
  \text{修复塔有限收敛条件}
  +
  \text{预算单调性}.
\]
这四部分都不依赖具体应用语境。
\end{remark}

\begin{remark}[上同调升级条件仍然未证明]
“上同调障碍”目前在本章仍是\textbf{未证明升级}。
若要把这里的全局截面障碍升级为上同调障碍，至少还需要补充以下对象：
\[
  \text{层 }\mathcal F,
  \qquad
  \check C^\bullet(B,\mathcal F),
  \qquad
  \delta,
  \qquad
  [o]\in \check H^k(B,\mathcal F).
\]
在这些对象没有被单独构造完成之前，本章只能严格称为
\[
  \text{全局截面障碍},
\]
不能直接升级称为
\[
  \text{上同调障碍}.
\]
\end{remark}

\begin{remark}[本章最短压缩版]
本章可以压缩为一句话：
\begin{center}
  \fbox{\parbox{0.88\linewidth}{
    \centering
    所谓“钱学森弹道问题”，在严格对象层上可重写为“受限观测与受限计算下的全局截面障碍问题”。
  }}
\end{center}
\begin{center}
  \fbox{\parbox{0.88\linewidth}{
    \centering
    而边缘算子、修复塔与预算扩张，则分别对应结构边界变化、歧义收缩机制与最优剩余歧义的单调下降机制。
  }}
\end{center}
\end{remark}

%%%%%%%%%%%%%%%%%%%%%%%%%%%%%%%%%%%%%%%%%%%%%%%%%%%%%%%%%%%%
\clearpage
\section{形式逻辑：钱学森弹道/钱学森障碍的结构完备性约束、障碍修复、预算匹配与单位预算收益}
\label{sec:tex1-ch3}
%%%%%%%%%%%%%%%%%%%%%%%%%%%%%%%%%%%%%%%%%%%%%%%%%%%%%%%%%%%%

\begin{remark}[本章意义备注：高超音速攻防的抽象语境]
本章的意义，在于明确结构完备性先于预算效率；
如果关键状态、模式或约束未被纳入对象域，追加算力也只会加快错误结论。
对钱学森弹道/钱学森障碍的抽象研究而言，它把焦点从“算得快”改成“先把该看的结构看全”。
\end{remark}

\subsection{严格主张、引文定位与章节目标}
\label{subsec:tex1-ch3-claim}

\begin{remark}[核心陈述]
本章在第~\ref{sec:tex1-ch1} 章与第~\ref{sec:tex1-ch2} 章已建立的
能力元组、全局截面障碍、伪确定性与预算单调性链条之上，
再加一层严格化：
\[
  \text{钱学森弹道/钱学森障碍问题，不是单纯的“算力大小问题”，}
\]
\[
  \text{而是“结构完备性约束下的障碍修复--预算匹配问题”。}
\]
也就是说，只有当模型闭合、真值保持、边缘算子已被并入能力元组之后，
“性价比”这个说法才获得严格意义；
否则它只是在错误对象层内部比较压缩速度，而不是比较对真实系统的有效性。
本章沿用仓内
\cite{RepoTex17EdgeOperator,RepoTex29RepairableObstruction,RepoTex32FiberSectionHole,RepoTex38JacobiBianchiRepair,
  RepoTex43OpenHeterogeneous}
关于“对象层 + 证书 + 谓词因果链 + 修复塔”的写法，
并以集合值分析、逆问题与混合系统对象层为外部背景，
参见
\cite{Isakov2006InversePDE,Liberzon2003Switching,AubinFrankowska2009SetValued,Husemoller1994FibreBundles}。
\end{remark}

\begin{remark}[本章的严格目标]
本章只讨论一般对抗推断系统的抽象对象层，
不对应任何现实攻击、规避、反求或工程实现方案。
我们要保留的结论有四个：
\begin{enumerate}
  \item 比较“谁更强”之前，必须先比较该方对象层是否结构完备；
  \item 剩余歧义 $V_i$、单位预算收益 $\kappa_i$ 与达阈时间 $\tau_i$ 是严格比较量；
  \item 预算扩张只保证“最优可达歧义不恶化”，不保证“真实有效性自动增强”；
  \item 在有限标签空间里，预算收益存在阈值--饱和结构，而不是无限线性增长。
\end{enumerate}
\end{remark}

\subsection{定义组：剩余歧义、结构完备性、真值保持、单位预算收益与达阈时间}
\label{subsec:tex1-ch3-defs}

\begin{definition}[分离型歧义度量与最优剩余歧义]
\label{def:tex1-ch3-ambiguity}
设 $Y$ 为有限标签空间。
称映射
\[
  \mu:2^Y\to \RR_{\ge 0}
\]
为\textbf{分离型歧义度量}，若满足：
\begin{enumerate}
  \item 单调性：$A\subseteq B\Rightarrow \mu(A)\le \mu(B)$；
  \item 单点零歧义：$\mu(A)=0\Longleftrightarrow |A|=1$。
\end{enumerate}
典型例子包括
\[
  \mu(A)=|A|-1,
  \qquad
  \mu(A)=\log |A|.
\]

对任意一方 $i$、时刻 $t<T$、观测历史 $b\in\mathcal B_t$ 与能力元组 $\Lambda_i(t)$，
记其可容许终端纤维为
\[
  \Gamma_t^{\Lambda_i}(b)\subseteq Y.
\]
给定预算对 $(C_i,L_i)$，定义第 $i$ 方的预算可行修复族为
\[
  \mathfrak R_i(C_i,L_i),
\]
其中每个
\[
  R\in \mathfrak R_i(C_i,L_i)
\]
都是映射 $R:2^Y\to 2^Y$，并满足
\[
  R(A)\subseteq A,
  \qquad
  \id\in \mathfrak R_i(C_i,L_i),
  \qquad
  \id(A)=A
\]
只要预算对 $(C_i,L_i)$ 是可行预算。
若 $\mathfrak R_i(C_i,L_i)=\varnothing$，则约定
\[
  V_i(t,b;C_i,L_i):=+\infty.
\]
若 $\mathfrak R_i(C_i,L_i)\neq\varnothing$，则定义第 $i$ 方在 $(t,b)$ 处的\textbf{最优剩余歧义}为
\[
  V_i(t,b;C_i,L_i)
  :=
  \inf_{R\in \mathfrak R_i(C_i,L_i)}
  \mu\bigl(R(\Gamma_t^{\Lambda_i}(b))\bigr).
\]
\end{definition}

\begin{definition}[结构完备性与真值保持]
\label{def:tex1-ch3-structure}
记真实终端标签为
\[
  y_\star(b)\in Y.
\]
称第 $i$ 方在点 $(t,b)$ \textbf{结构完备}，若
\[
  y_\star(b)\in \Gamma_t^{\Lambda_i}(b).
\]
这表示：构成 $\Gamma_t^{\Lambda_i}(b)$ 的对象层已经把与该点相关的隐藏模式、
边缘算子插入、预算更新规则及其他真实约束纳入能力元组之中。

进一步地，称预算可行修复族 $\mathfrak R_i(C_i,L_i)$ 满足\textbf{真值保持}，若对任意
\[
  R\in \mathfrak R_i(C_i,L_i)
\]
都有
\[
  y_\star(b)\in \Gamma_t^{\Lambda_i}(b)
  \Rightarrow
  y_\star(b)\in R(\Gamma_t^{\Lambda_i}(b)).
\]
\end{definition}

\begin{definition}[单位预算收益]
\label{def:tex1-ch3-kappa}
给定权重
\[
  \alpha,\beta>0,
\]
定义加权预算代价为
\[
  \operatorname{Cost}_i(C_i,L_i):=\alpha C_i+\beta L_i.
\]
当
\[
  \operatorname{Cost}_i(C_i,L_i)>0
  \qquad\text{且}\qquad
  V_i(t,b;C_i,L_i)<+\infty
\]
时，定义第 $i$ 方在 $(t,b)$ 处的\textbf{单位预算收益}为
\[
  \kappa_i(t,b;C_i,L_i)
  :=
  \frac{
    \mu(\Gamma_t^{\Lambda_i}(b))
    -
    V_i(t,b;C_i,L_i)
  }{
    \operatorname{Cost}_i(C_i,L_i)
  }.
\]
它刻画的是：单位预算可压低多少歧义。
\end{definition}

\begin{definition}[达阈时间与预算匹配]
\label{def:tex1-ch3-threshold}
给定判定阈值 $\varepsilon\ge 0$，
并允许预算随时间变化为 $C_i(t),L_i(t)$。
定义第 $i$ 方在观测历史 $b$ 上的\textbf{达阈时间}为
\[
  \tau_i(\varepsilon,b)
  :=
  \inf
  \Bigl\{
    t<T:\ V_i\bigl(t,b;C_i(t),L_i(t)\bigr)\le \varepsilon
  \Bigr\}.
\]
若该集合为空，则记
\[
  \tau_i(\varepsilon,b)=+\infty.
\]
若
\[
  \tau_i(\varepsilon,b)<T,
\]
则称第 $i$ 方在阈值 $\varepsilon$ 下与该点 $(t,b)$ 的判定任务\textbf{预算匹配}。
\end{definition}

\subsection{公理（降级为定理）：结构完备性优先于性价比的层级链}
\label{subsec:tex1-ch3-axiom}

\begin{theorem}[公理（降级为定理）：结构完备性优先于性价比的层级链]
\label{thm:tex1-ch3-priority}
在本章形式系统中，关于钱学森弹道/钱学森障碍的有效比较，
必须按如下四层对象递归登记：
\[
  \mathfrak L_0:\ \text{结构完备性},
\]
\[
  \mathfrak L_1:\ \text{真值保持},
\]
\[
  \mathfrak L_2:\ \text{预算匹配},
\]
\[
  \mathfrak L_3:\ \text{单位预算收益}.
\]
更具体地说，若某个更高层对象具有严格语义，
则其全部前缀层也必须已经成立。
因此有逻辑优先序
\[
  \text{结构完备性}
  \succ
  \text{真值保持}
  \succ
  \text{预算匹配}
  \succ
  \text{单位预算收益}.
\]
\end{theorem}

\begin{Certificate}[数学归纳法证书：四层优先链的前缀登记]
\label{cert:tex1-ch3-priority}
定义谓词
\[
  P(n):
  \exists!\,\text{有效前缀链 }\mathfrak L_0\to\cdots\to \mathfrak L_n,
  \quad
  \text{其长度为 }n+1.
\]
并登记递归兼容条件
\[
  K_0:\ \mathfrak L_0\Rightarrow \mathfrak L_1\ \text{可被检验},
\]
\[
  K_1:\ \mathfrak L_1\Rightarrow \mathfrak L_2\ \text{可被赋予真实有效性含义},
\]
\[
  K_2:\ \mathfrak L_2\Rightarrow \mathfrak L_3\ \text{可被赋予单位预算收益含义}.
\]
则归纳证书登记为
\[
  \Pi_{\mathrm{prio}}
  =
  \Bigl(
    P(0),\
    K_0,\
    K_1,\
    K_2,\
    \forall n\in\{0,1,2\},\ P(n)\Rightarrow P(n+1)
  \Bigr).
\]
\end{Certificate}

\begin{proof}[证明（数学归纳法证书；基于数学符号推演逻辑的谓词因果链）]
先看基步 $n=0$。
由定义~\ref{def:tex1-ch3-structure}，结构完备性给出了
\[
  y_\star(b)\in \Gamma_t^{\Lambda_i}(b),
\]
也即为整个判定问题提供了“真实标签是否被对象层保留”的锚点。
于是长度为 $1$ 的前缀链 $\mathfrak L_0$ 唯一可登记，这就是 \(P(0)\)。

下面做归纳步。
若已知 \(P(0)\) 成立，则 \(K_0\) 说明：
只有在真实标签已被对象层保留时，
“真值保持”才不是悬空表述，而可以被检验为
\[
  y_\star(b)\in \Gamma_t^{\Lambda_i}(b)
  \Rightarrow
  y_\star(b)\in R(\Gamma_t^{\Lambda_i}(b)).
\]
故可把 $\mathfrak L_1$ 追加到前缀链中，从而得到 \(P(1)\)。

若已知 \(P(1)\) 成立，则 \(K_1\) 说明：
只有在结构完备且真值保持都已登记之后，
\[
  V_i(t,b;C_i,L_i),\qquad \tau_i(\varepsilon,b)
\]
才表示“对真实系统有效的歧义压缩与达阈能力”，
而不只是错误纤维内部的形式压缩。
因此可得到 \(P(2)\)。

若已知 \(P(2)\) 成立，则 \(K_2\) 说明：
只有在预算匹配意义已经成立之后，
\[
  \kappa_i(t,b;C_i,L_i)
\]
才可被解释为“单位预算带来的真实有效收益”，
而不是对伪确定性的内部收益。
于是可得到 \(P(3)\)。

综上，由有限归纳法得到长度为 $4$ 的唯一有效前缀链
\[
  \mathfrak L_0\to \mathfrak L_1\to \mathfrak L_2\to \mathfrak L_3.
\]
因此逻辑优先序
\[
  \text{结构完备性}
  \succ
  \text{真值保持}
  \succ
  \text{预算匹配}
  \succ
  \text{单位预算收益}
\]
成立。定理得证。
\end{proof}

\begin{corollary}[“性价比”只有在结构完备且真值保持之后才有严格意义]
\label{cor:tex1-ch3-priority}
若未先建立结构完备性与真值保持，
则单位预算收益 $\kappa_i$ 不能被解释为真实有效性的收益，
只能被解释为某个代理对象内部的压缩速度。
\end{corollary}

\begin{Certificate}[证书：定理~\ref{thm:tex1-ch3-priority} 的直接推论]
\label{cert:tex1-ch3-priority-cor}
证书登记谓词链：
\[
\begin{aligned}
  C^{(3)}_1 &: \text{调用定理~\ref{thm:tex1-ch3-priority} 得到四层优先序},\\
  C^{(3)}_2 &: \kappa_i\ \text{位于最末层 }\mathfrak L_3,\\
  C^{(3)}_3 &: C^{(3)}_1\wedge C^{(3)}_2
              \Rightarrow \mathfrak L_0,\mathfrak L_1\ \text{必须先成立},\\
  C^{(3)}_4 &: C^{(3)}_3\Rightarrow \kappa_i\ \text{只能在结构完备与真值保持后获得严格意义}.
\end{aligned}
\]
\end{Certificate}

\begin{proof}[证明（基于数学符号推演逻辑的谓词因果链）]
由 \(C^{(3)}_1\) 可知，单位预算收益位于优先链末层；
由 \(C^{(3)}_2\) 可知，讨论 $\kappa_i$ 时必然已经调用该末层对象。
于是由 \(C^{(3)}_3\) 可知，所有前缀层必须先被建立。
特别地，结构完备性与真值保持必须先成立，
因此 \(C^{(3)}_4\) 成立。推论得证。
\end{proof}

\subsection{引理：结构不完备时，内部修复无法找回被遗漏的真值}
\label{subsec:tex1-ch3-lemma}

\begin{lemma}[结构不完备时，内部修复无法找回被遗漏的真值]
\label{lem:tex1-ch3-no-recover}
若对某个 $(t,b)$ 有
\[
  y_\star(b)\notin \Gamma_t^{\Lambda_i}(b),
\]
且每个预算可行修复都满足
\[
  R(A)\subseteq A,
\]
则对任意
\[
  R\in \mathfrak R_i(C_i,L_i)
\]
都有
\[
  y_\star(b)\notin R(\Gamma_t^{\Lambda_i}(b)).
\]
\end{lemma}

\begin{Certificate}[证书：集合包含与补集继承证书]
\label{cert:tex1-ch3-no-recover}
证书登记谓词链：
\[
\begin{aligned}
  L^{(3)}_1 &: y_\star(b)\notin \Gamma_t^{\Lambda_i}(b),\\
  L^{(3)}_2 &: \forall R\in\mathfrak R_i(C_i,L_i),\ R(\Gamma_t^{\Lambda_i}(b))\subseteq \Gamma_t^{\Lambda_i}(b),\\
  L^{(3)}_3 &: L^{(3)}_1\wedge L^{(3)}_2
              \Rightarrow y_\star(b)\notin R(\Gamma_t^{\Lambda_i}(b)).
\end{aligned}
\]
\end{Certificate}

\begin{proof}[证明（基于数学符号推演逻辑的谓词因果链）]
由 \(L^{(3)}_1\) 可知，真实标签不在原始可容许终端纤维中。
由 \(L^{(3)}_2\) 可知，任意预算可行修复都只能把该纤维收缩到其内部子集。
因此，任何
\[
  R(\Gamma_t^{\Lambda_i}(b))
\]
都仍是
\[
  \Gamma_t^{\Lambda_i}(b)
\]
的子集，故不可能突然包含原先已经缺失的元素 $y_\star(b)$。
这就是 \(L^{(3)}_3\)，引理得证。
\end{proof}

\begin{corollary}[结构不完备时，零剩余歧义只可能是错误纤维内部的零歧义]
\label{cor:tex1-ch3-no-recover}
若结构完备性失败，则即便某个预算可行修复把歧义压到零，
该零值也只能发生在错误纤维内部，而不能恢复真实标签。
\end{corollary}

\begin{Certificate}[证书：引理~\ref{lem:tex1-ch3-no-recover} 与单点零歧义公理化条件]
\label{cert:tex1-ch3-no-recover-cor}
证书登记谓词链：
\[
\begin{aligned}
  \widehat C^{(3)}_1 &: \text{结构完备性失败，即 }y_\star(b)\notin \Gamma_t^{\Lambda_i}(b),\\
  \widehat C^{(3)}_2 &: \text{调用引理~\ref{lem:tex1-ch3-no-recover}},\\
  \widehat C^{(3)}_3 &: \mu(A)=0\Longleftrightarrow |A|=1,\\
  \widehat C^{(3)}_4 &: \widehat C^{(3)}_1\wedge \widehat C^{(3)}_2\wedge \widehat C^{(3)}_3
                       \Rightarrow \text{零歧义只能对应错误纤维内部的单点化}.
\end{aligned}
\]
\end{Certificate}

\begin{proof}[证明（基于数学符号推演逻辑的谓词因果链）]
由 \(\widehat C^{(3)}_1\) 与 \(\widehat C^{(3)}_2\)，
任何预算可行修复后的纤维都不包含真实标签。
若某个修复把歧义压到零，则由 \(\widehat C^{(3)}_3\) 可知，
该修复后的纤维必为某个单点，但这个单点不可能等于 $y_\star(b)$。
故 \(\widehat C^{(3)}_4\) 成立。推论得证。
\end{proof}

\subsection{定理：钱学森弹道/钱学森障碍问题可重写为障碍修复--预算匹配问题}
\label{subsec:tex1-ch3-theorem-redef}

\begin{theorem}[钱学森弹道/钱学森障碍问题向障碍修复--预算匹配问题的重定义]
\label{thm:tex1-ch3-redef}
若一般对抗推断系统的胜负判据是：
\[
  \text{在终端时刻 }T\text{ 之前，把剩余歧义压到阈值 }\varepsilon\text{ 以下，且满足真值保持},
\]
则该问题等价于比较各方的
\[
  V_i(t,b;C_i,L_i),
  \qquad
  \kappa_i(t,b;C_i,L_i),
  \qquad
  \tau_i(\varepsilon,b).
\]
也就是说，钱学森弹道/钱学森障碍问题在严格对象层上，
可被重写为
\[
  \text{结构完备性约束}
  +
  \text{障碍修复能力}
  +
  \text{预算匹配能力}
  +
  \text{单位预算收益比较}.
\]
\end{theorem}

\begin{Certificate}[证书：定义展开证书]
\label{cert:tex1-ch3-redef}
证书登记谓词链：
\[
\begin{aligned}
  T^{(3)}_1 &: V_i(t,b;C_i,L_i)\le \varepsilon
              \Longleftrightarrow \text{在该时刻已把剩余歧义压到阈值以下},\\
  T^{(3)}_2 &: \tau_i(\varepsilon,b)
              =\inf\{t<T:\ V_i(t,b;C_i(t),L_i(t))\le \varepsilon\},\\
  T^{(3)}_3 &: \kappa_i(t,b;C_i,L_i)
              \text{刻画单位预算带来的歧义压缩收益},\\
  T^{(3)}_4 &: T^{(3)}_1\wedge T^{(3)}_2\wedge T^{(3)}_3
              \Rightarrow \text{原问题完全可由 }(V_i,\kappa_i,\tau_i)\text{ 描述}.
\end{aligned}
\]
\end{Certificate}

\begin{proof}[证明（基于数学符号推演逻辑的谓词因果链）]
由 \(T^{(3)}_1\) 可知，
所谓“是否成功”首先只取决于
\[
  V_i(t,b;C_i,L_i)\le \varepsilon
\]
是否成立。
也就是说，只要知道最优剩余歧义，就知道该方是否已经跨过判定阈值。

由 \(T^{(3)}_2\) 可知，
若把“谁更早达到阈值”作为比较内容，
那么该比较完全由达阈时间
\[
  \tau_i(\varepsilon,b)
\]
承担。

由 \(T^{(3)}_3\) 可知，
若要比较“谁以更省的预算压低更多歧义”，
那么该比较由单位预算收益
\[
  \kappa_i(t,b;C_i,L_i)
\]
承担。

于是把 \(T^{(3)}_1\)、\(T^{(3)}_2\) 与 \(T^{(3)}_3\) 合并，
便得到 \(T^{(3)}_4\)：
原问题不再需要以轨迹外观来表述，
而可以被等价改写为
\[
  (V_i,\kappa_i,\tau_i)
\]
的比较问题。
再结合定理~\ref{thm:tex1-ch3-priority} 中“结构完备性优先”的前缀约束，
便得到本定理的完整表述。证毕。
\end{proof}

\begin{corollary}[核心不再是几何外观，而是剩余歧义的有效压缩效率]
\label{cor:tex1-ch3-redef}
在结构完备且真值保持的对象层内，
钱学森弹道/钱学森障碍问题的核心比较量，
不再是几何表示本身，而是剩余歧义的有效压缩效率。
\end{corollary}

\begin{Certificate}[证书：定理~\ref{thm:tex1-ch3-redef} 的直接推论]
\label{cert:tex1-ch3-redef-cor}
证书登记谓词链：
\[
\begin{aligned}
  C^{(3)}_5 &: \text{调用定理~\ref{thm:tex1-ch3-redef}},\\
  C^{(3)}_6 &: C^{(3)}_5\Rightarrow \text{核心比较量为 }(V_i,\kappa_i,\tau_i),\\
  C^{(3)}_7 &: C^{(3)}_6\Rightarrow \text{几何外观退回为表示层，而非核心判据}.
\end{aligned}
\]
\end{Certificate}

\begin{proof}[证明（基于数学符号推演逻辑的谓词因果链）]
由 \(C^{(3)}_5\) 可知，原问题已被重写为
\[
  (V_i,\kappa_i,\tau_i)
\]
的比较。
因此 \(C^{(3)}_6\) 成立。
而一旦核心判据已经完全由这些量承担，
“轨迹外观是否漂亮”就不再决定真值判定，
于是 \(C^{(3)}_7\) 成立。推论得证。
\end{proof}

\subsection{定理：性价比不是第一性对象；结构完备性先于性价比}
\label{subsec:tex1-ch3-theorem-priority}

\begin{theorem}[性价比不是第一性对象；结构完备性先于性价比]
\label{thm:tex1-ch3-not-first}
若第 $i$ 方在点 $(t,b)$ 结构不完备，即
\[
  y_\star(b)\notin \Gamma_t^{\Lambda_i}(b),
\]
则命题
\[
  V_i(t,b;C_i,L_i)=0
  \Rightarrow
  \text{第 }i\text{ 方已经正确判定}
\]
不成立。
更强地说：若零值由某个预算可行修复达到，
则对应单点必不等于真实标签 $y_\star(b)$。
\end{theorem}

\begin{Certificate}[证书：反例式证书]
\label{cert:tex1-ch3-not-first}
证书登记谓词链：
\[
\begin{aligned}
  T^{(3)}_5 &: y_\star(b)\notin \Gamma_t^{\Lambda_i}(b),\\
  T^{(3)}_6 &: \forall R\in\mathfrak R_i(C_i,L_i),\ R(\Gamma_t^{\Lambda_i}(b))\subseteq \Gamma_t^{\Lambda_i}(b),\\
  T^{(3)}_7 &: T^{(3)}_5\wedge T^{(3)}_6\Rightarrow
              \text{调用引理~\ref{lem:tex1-ch3-no-recover}},\\
  T^{(3)}_8 &: \mu(A)=0\Longleftrightarrow |A|=1,\\
  T^{(3)}_9 &: T^{(3)}_7\wedge T^{(3)}_8
              \Rightarrow \text{任何零歧义单点若存在，必不是 }y_\star(b),\\
  T^{(3)}_{10} &: T^{(3)}_9\Rightarrow
                  \bigl(V_i=0\Rightarrow \text{正确判定}\bigr)\ \text{不成立}.
\end{aligned}
\]
\end{Certificate}

\begin{proof}[证明（基于数学符号推演逻辑的谓词因果链）]
由 \(T^{(3)}_5\) 与 \(T^{(3)}_6\)，
调用引理~\ref{lem:tex1-ch3-no-recover}，
可知任意预算可行修复之后的纤维都仍然不含真实标签。
这就是 \(T^{(3)}_7\)。

再由 \(T^{(3)}_8\) 可知，
若某个预算可行修复把歧义压到零，
则其输出必为单点。
但该单点既然仍位于错误纤维内部，
便不可能等于 $y_\star(b)$。
于是 \(T^{(3)}_9\) 成立。

因此，即便出现
\[
  V_i(t,b;C_i,L_i)=0,
\]
它也至多说明“该方在自己的错误对象层内部完成了单点化”，
而不说明“该方已经对真实系统做出正确判定”。
这正是 \(T^{(3)}_{10}\)。
故定理成立。证毕。
\end{proof}

\begin{corollary}[更准确的表述是：先有结构完备性问题，再有性价比问题]
\label{cor:tex1-ch3-not-first}
在本章形式系统中，更准确的优先序不是
\[
  \text{只剩性价比问题},
\]
而是
\[
  \text{先有结构完备性问题，再有性价比问题}.
\]
也即
\[
  \text{结构完备性}
  \succ
  \text{真值保持}
  \succ
  \text{预算匹配}
  \succ
  \text{单位预算收益}.
\]
\end{corollary}

\begin{Certificate}[证书：定理~\ref{thm:tex1-ch3-priority} 与定理~\ref{thm:tex1-ch3-not-first} 的合成]
\label{cert:tex1-ch3-not-first-cor}
证书登记谓词链：
\[
\begin{aligned}
  C^{(3)}_8 &: \text{调用定理~\ref{thm:tex1-ch3-priority}},\\
  C^{(3)}_9 &: \text{调用定理~\ref{thm:tex1-ch3-not-first}},\\
  C^{(3)}_{10} &: C^{(3)}_8\wedge C^{(3)}_9
                  \Rightarrow \text{结构完备性是性价比比较的先决层},\\
  C^{(3)}_{11} &: C^{(3)}_{10}\Rightarrow
                  \text{“先有结构完备性问题，再有性价比问题”}.
\end{aligned}
\]
\end{Certificate}

\begin{proof}[证明（基于数学符号推演逻辑的谓词因果链）]
由 \(C^{(3)}_8\) 得到四层优先链；
由 \(C^{(3)}_9\) 得到“结构不完备时，零歧义不推出正确判定”。
二者合并即得 \(C^{(3)}_{10}\)：
性价比只有在结构完备性已建立之后才有语义意义。
于是 \(C^{(3)}_{11}\) 成立。推论得证。
\end{proof}

\subsection{定理：预算单调性与达阈能力的单调改进}
\label{subsec:tex1-ch3-theorem-budget}

\begin{theorem}[预算单调性]
\label{thm:tex1-ch3-budget}
若
\[
  C_i^{(1)}\le C_i^{(2)},
  \qquad
  L_i^{(1)}\le L_i^{(2)},
\]
并且预算扩张意味着可行修复族扩张：
\[
  \mathfrak R_i(C_i^{(1)},L_i^{(1)})
  \subseteq
  \mathfrak R_i(C_i^{(2)},L_i^{(2)}),
\]
则有
\[
  V_i\bigl(t,b;C_i^{(2)},L_i^{(2)}\bigr)
  \le
  V_i\bigl(t,b;C_i^{(1)},L_i^{(1)}\bigr).
\]
\end{theorem}

\begin{Certificate}[证书：可行集包含证书]
\label{cert:tex1-ch3-budget}
证书登记谓词链：
\[
\begin{aligned}
  T^{(3)}_{11} &: C_i^{(1)}\le C_i^{(2)},\quad L_i^{(1)}\le L_i^{(2)},\\
  T^{(3)}_{12} &: T^{(3)}_{11}\Rightarrow
                  \mathfrak R_i(C_i^{(1)},L_i^{(1)})
                  \subseteq
                  \mathfrak R_i(C_i^{(2)},L_i^{(2)}),\\
  T^{(3)}_{13} &: \text{对更大的可行集取 }\inf\ \text{不会得到更大的值},\\
  T^{(3)}_{14} &: T^{(3)}_{12}\wedge T^{(3)}_{13}
                  \Rightarrow
                  V_i(t,b;C_i^{(2)},L_i^{(2)})
                  \le
                  V_i(t,b;C_i^{(1)},L_i^{(1)}).
\end{aligned}
\]
\end{Certificate}

\begin{proof}[证明（基于数学符号推演逻辑的谓词因果链）]
由 \(T^{(3)}_{11}\) 得到预算对按分量扩张。
再由预算扩张假设，立即得到 \(T^{(3)}_{12}\)。
由于对更大的可行集取下确界，不会比对更小的可行集取下确界更大，
于是 \(T^{(3)}_{13}\) 成立。
将 \(T^{(3)}_{12}\) 与 \(T^{(3)}_{13}\) 合并，
便得到 \(T^{(3)}_{14}\)，即本定理。
\end{proof}

\begin{corollary}[算力扩张与时延放宽不会恶化最优剩余歧义]
\label{cor:tex1-ch3-budget}
预算扩张只说明
\[
  \text{预算上界}\uparrow
  \Rightarrow
  \text{最优可达歧义}\downarrow\ \text{或保持不变},
\]
并不自动说明“真实有效性一定增强”；
真实有效性仍受定理~\ref{thm:tex1-ch3-not-first} 的结构完备性约束。
\end{corollary}

\begin{Certificate}[证书：定理~\ref{thm:tex1-ch3-budget} 与定理~\ref{thm:tex1-ch3-not-first} 的合成]
\label{cert:tex1-ch3-budget-cor}
证书登记谓词链：
\[
\begin{aligned}
  C^{(3)}_{12} &: \text{调用定理~\ref{thm:tex1-ch3-budget}},\\
  C^{(3)}_{13} &: \text{调用定理~\ref{thm:tex1-ch3-not-first}},\\
  C^{(3)}_{14} &: C^{(3)}_{12}\Rightarrow \text{预算扩张不会恶化 }V_i,\\
  C^{(3)}_{15} &: C^{(3)}_{13}\Rightarrow \text{结构不完备时 }V_i\text{ 的改进不保证真实有效性},\\
  C^{(3)}_{16} &: C^{(3)}_{14}\wedge C^{(3)}_{15}
                  \Rightarrow \text{预算扩张只是必要变量之一，而非第一性对象}.
\end{aligned}
\]
\end{Certificate}

\begin{proof}[证明（基于数学符号推演逻辑的谓词因果链）]
由 \(C^{(3)}_{12}\) 得到预算扩张只能使最优剩余歧义下降或不变；
由 \(C^{(3)}_{13}\) 得到结构不完备时，
即便零歧义也不推出正确判定。
因此，把二者合并便得到 \(C^{(3)}_{16}\)：
预算扩张固然重要，但必须被放在结构完备性之后理解。推论得证。
\end{proof}

\subsection{定理：有限标签下的边际收益饱和}
\label{subsec:tex1-ch3-theorem-saturation}

\begin{theorem}[有限标签下的边际收益饱和]
\label{thm:tex1-ch3-saturation}
设 $Y$ 有限，并设对固定点 $(t,b)$ 存在一条真值保持且严格降歧义的修复序列
\[
  R_0,\dots,R_{N(b)-1},
\]
其中每一级修复都具有有限代价与有限时延：
\[
  \operatorname{cost}(R_n)<+\infty,
  \qquad
  \operatorname{latency}(R_n)<+\infty.
\]
若该序列在 $N(b)$ 步后把纤维压到单点，
则存在有限预算阈值
\[
  (C_i^\star,L_i^\star),
\]
使得只要
\[
  C_i\ge C_i^\star,
  \qquad
  L_i\ge L_i^\star,
\]
就有
\[
  V_i(t,b;C_i,L_i)=0.
\]
并且一旦达到该零值，
进一步增加预算也不会继续改善该点的判定结果。
\end{theorem}

\begin{Certificate}[证书：有限收敛 + 可行预算求和证书]
\label{cert:tex1-ch3-saturation}
证书登记谓词链：
\[
\begin{aligned}
  T^{(3)}_{15} &: \text{调用定理~\ref{thm:tex1-ch2-repair-finite}，得到有限步收敛},\\
  T^{(3)}_{16} &: \text{存在 }R_0,\dots,R_{N(b)-1}\ \text{使得 }
                  R_{N(b)-1}\circ\cdots\circ R_0(\Gamma_t^{\Lambda_i}(b))
                  \text{ 为单点},\\
  T^{(3)}_{17} &: C_i^\star:=\sum_{n=0}^{N(b)-1}\operatorname{cost}(R_n),\quad
                  L_i^\star:=\sum_{n=0}^{N(b)-1}\operatorname{latency}(R_n),\\
  T^{(3)}_{18} &: C_i\ge C_i^\star,\ L_i\ge L_i^\star
                  \Rightarrow R_{N(b)-1}\circ\cdots\circ R_0\in\mathfrak R_i(C_i,L_i),\\
  T^{(3)}_{19} &: T^{(3)}_{16}\wedge T^{(3)}_{18}\wedge \mu(A)=0\Longleftrightarrow |A|=1
                  \Rightarrow V_i(t,b;C_i,L_i)=0,\\
  T^{(3)}_{20} &: \text{调用定理~\ref{thm:tex1-ch3-budget}}
                  \Rightarrow \text{更大预算下仍有 }V_i=0.
\end{aligned}
\]
\end{Certificate}

\begin{proof}[证明（基于数学符号推演逻辑的谓词因果链）]
由定理~\ref{thm:tex1-ch2-repair-finite}，
有限标签空间中的真值保持且严格降歧义修复塔在每个点上都只能有限步下降。
因此可取某个有限长度 $N(b)$ 的修复序列
\[
  R_0,\dots,R_{N(b)-1},
\]
使其在该点把纤维压到单点，这就是 \(T^{(3)}_{15}\) 与 \(T^{(3)}_{16}\)。

由于每一级修复的代价与时延都有限，
可定义
\[
  C_i^\star:=\sum_{n=0}^{N(b)-1}\operatorname{cost}(R_n),
  \qquad
  L_i^\star:=\sum_{n=0}^{N(b)-1}\operatorname{latency}(R_n),
\]
这就是 \(T^{(3)}_{17}\)。
于是当
\[
  C_i\ge C_i^\star,\qquad L_i\ge L_i^\star
\]
时，整条复合修复序列都成为预算可行修复，即 \(T^{(3)}_{18}\)。

又因为该复合修复输出单点，
且分离型歧义度量满足
\[
  \mu(A)=0\Longleftrightarrow |A|=1,
\]
故 \(T^{(3)}_{19}\) 成立，即
\[
  V_i(t,b;C_i,L_i)=0.
\]

最后，对任何更大预算对，
由定理~\ref{thm:tex1-ch3-budget} 得
\[
  V_i(\text{更大预算})\le V_i(t,b;C_i,L_i)=0.
\]
但 $V_i\ge 0$，
故更大预算下仍只能有
\[
  V_i=0.
\]
这就是 \(T^{(3)}_{20}\)。
因此，一旦跨过阈值，进一步增加预算不会继续改善该点判定结果。定理得证。
\end{proof}

\begin{corollary}[单位预算收益存在阈值--饱和结构]
\label{cor:tex1-ch3-saturation}
在有限标签空间内，单位预算收益不是线性无限增长的，
而是存在如下阈值--饱和结构：
\[
  \text{前期预算用于跨越关键阈值，后期预算的边际收益趋于 }0.
\]
因此真正重要的不是无上限堆积预算，
而是
\[
  \text{在正确的结构层上，以足够但不过量的预算跨过关键阈值}.
\]
\end{corollary}

\begin{Certificate}[证书：零值饱和证书]
\label{cert:tex1-ch3-saturation-cor}
证书登记谓词链：
\[
\begin{aligned}
  C^{(3)}_{17} &: \text{调用定理~\ref{thm:tex1-ch3-saturation}},\\
  C^{(3)}_{18} &: C^{(3)}_{17}\Rightarrow \exists(C_i^\star,L_i^\star),\ \forall C_i\ge C_i^\star,\ \forall L_i\ge
    L_i^\star,\ V_i=0,\\
  C^{(3)}_{19} &: V_i=0\Rightarrow \mu(\Gamma_t^{\Lambda_i}(b))-V_i
                  \text{已达到该点可实现的最大压缩值},\\
  C^{(3)}_{20} &: C^{(3)}_{18}\wedge C^{(3)}_{19}
                  \Rightarrow \text{更大预算的边际收益为 }0.
\end{aligned}
\]
\end{Certificate}

\begin{proof}[证明（基于数学符号推演逻辑的谓词因果链）]
由 \(C^{(3)}_{17}\) 得到零剩余歧义在有限预算阈值之后保持不变。
于是 \(C^{(3)}_{18}\) 成立。
一旦 $V_i=0$，说明该点已经被压到单点，
故可压缩歧义的空间已被耗尽，这就是 \(C^{(3)}_{19}\)。
于是再增加预算也不会带来新的歧义压缩，
从而 \(C^{(3)}_{20}\) 成立。推论得证。
\end{proof}

\subsection{命题：结构完备性约束下的障碍修复--预算匹配最严压缩式}
\label{subsec:tex1-ch3-prop}

\begin{proposition}[结构完备性约束下的障碍修复--预算匹配最严压缩式]
\label{prop:tex1-ch3-compressed}
在模型闭合、真值保持且以歧义阈值为判据的一般对抗推断系统中，
钱学森弹道/钱学森障碍问题
可严格重写为预算约束下的障碍修复最优化问题；
其核心比较量不是轨迹外观，而是
\[
  V_i(t,b;C_i,L_i),
  \qquad
  \kappa_i(t,b;C_i,L_i),
  \qquad
  \tau_i(\varepsilon,b).
\]
\end{proposition}

\begin{Certificate}[证书：四定理合成证书]
\label{cert:tex1-ch3-compressed}
证书登记谓词链：
\[
\begin{aligned}
  P^{(3)}_1 &: \text{调用定理~\ref{thm:tex1-ch3-priority}},\\
  P^{(3)}_2 &: \text{调用定理~\ref{thm:tex1-ch3-redef}},\\
  P^{(3)}_3 &: \text{调用定理~\ref{thm:tex1-ch3-not-first}},\\
  P^{(3)}_4 &: \text{调用定理~\ref{thm:tex1-ch3-budget} 与定理~\ref{thm:tex1-ch3-saturation}},\\
  P^{(3)}_5 &: P^{(3)}_1\wedge P^{(3)}_2\wedge P^{(3)}_3\wedge P^{(3)}_4
              \Rightarrow \text{命题成立}.
\end{aligned}
\]
\end{Certificate}

\begin{proof}[证明（基于数学符号推演逻辑的谓词因果链）]
由 \(P^{(3)}_1\) 得到结构完备性先于性价比的优先链；
由 \(P^{(3)}_2\) 得到原问题可由
\[
  (V_i,\kappa_i,\tau_i)
\]
完全描述；
由 \(P^{(3)}_3\) 得到“若结构不完备，则性价比不能充当第一性对象”；
由 \(P^{(3)}_4\) 得到预算扩张与收益饱和的边界性质。
将四者合并，即得 \(P^{(3)}_5\)：
在结构完备性约束之下，
钱学森弹道/钱学森障碍问题可以严格改写为预算约束下的障碍修复最优化问题。
命题得证。
\end{proof}

\subsection{备注：前提边界、完整压缩句与下一层极小极大泛函}
\label{subsec:tex1-ch3-remarks}

\begin{remark}[前提不是空的]
因此，“障碍修复和算力匹配的性价比问题”这个说法在抽象层上成立，
但它的前提不是空的。
这些前提至少包括：
\[
  \text{模型闭合},
  \qquad
  \text{真值保持},
  \qquad
  \text{边缘算子已并入能力元组}.
\]
少了这些前提，
“性价比”就会退化为错误模型内部的压缩效率，
而不是真实系统上的有效性比较。
\end{remark}

\begin{remark}[本章的一句完整压缩句]
本章可以进一步压缩为：
\begin{center}
  \fbox{\parbox{0.9\linewidth}{
    \centering
    钱学森弹道/钱学森障碍问题在严格对象层上，
    可重写为“结构完备性约束下的障碍修复--预算匹配--阈值跨越问题”；
    其中比较量不是轨迹外观，而是剩余歧义、单位预算收益与达阈时间。
  }}
\end{center}
\end{remark}

\begin{remark}[下一层自然对象：极小极大目标泛函]
若再向前推进一步，
下一层自然对象不再是继续讨论“轨迹表象”，
而是直接写一个极小极大目标泛函，例如
\[
  \inf_{(\Lambda_1,\mathfrak R_1,C_1,L_1)}
  \sup_{(\Lambda_2,\mathfrak e_2,C_2,L_2)}
  J\Bigl(
    V_1,V_2,\tau_1,\tau_2,\operatorname{Cost}_1,\operatorname{Cost}_2,\mathbf 1_{\mathrm{err}}
  \Bigr),
\]
其中 $J$ 由剩余歧义、时延、预算与误判罚项共同组成。
一旦进入这一层，
整个表述就会从“直觉判断”完全转入“可证明优化问题”。
\end{remark}

%%%%%%%%%%%%%%%%%%%%%%%%%%%%%%%%%%%%%%%%%%%%%%%%%%%%%%%%%%%%
\clearpage
\section{形式逻辑：钱学森弹道/钱学森障碍的极小极大泛函、稳健值函数与单位预算收益}
\label{sec:tex1-ch4}
%%%%%%%%%%%%%%%%%%%%%%%%%%%%%%%%%%%%%%%%%%%%%%%%%%%%%%%%%%%%

\begin{remark}[本章意义备注：高超音速攻防的抽象语境]
本章的意义，在于把抽象攻防关系整理为稳健极小极大问题：
先控制最坏情形下的剩余歧义，再比较单位预算收益。
对钱学森障碍而言，这说明任何有效判定都应先满足稳健可判，再谈效率优化。
\end{remark}

\subsection{严格主张、引文定位与章节目标}
\label{subsec:tex1-ch4-claim}

\begin{remark}[核心陈述]
本章把“障碍修复--算力匹配--单位预算收益”进一步收紧为一个两层优化结构：
\[
  \text{一级目标：稳健极小极大值 }\mathcal V,
\]
\[
  \text{二级目标：近最优集上的单位预算收益 }\mathrm{CE}.
\]
这样做的原因是：若直接把“性价比比值”放在第一层，
它往往会压过真值保持、结构完备性与最坏情形稳健性。
因此，严格的逻辑次序必须先是
\[
  \min\max J,
\]
再是在一级近最优解集内比较单位预算收益。
为避免与第~\ref{sec:tex1-ch3} 章中的点态剩余歧义 $V_i$ 混淆，
本章把极小极大值函数记作
\[
  \mathcal V(t,b;\bar K_A,\bar K_B).
\]
本章继续沿用仓内
\cite{RepoTex17EdgeOperator,RepoTex18HomotopyAxiom,RepoTex29RepairableObstruction,RepoTex38JacobiBianchiRepair,
  RepoTex43OpenHeterogeneous}
关于“对象层 + 证书 + 谓词因果链 + 修复塔”的写法；
其外部数学背景则由集合值分析、混合系统、非线性优化与极小极大理论支撑，可参见
\cite{AubinFrankowska2009SetValued,Liberzon2003Switching,Sion1958Minimax,Bertsekas2016NonlinearProgramming}。
\end{remark}

\begin{remark}[本章的严格目标]
本章仍只处理抽象对象层，不进入任何现实攻击方案、规避策略或工程反求。
我们要建立的结论有五个：
\begin{enumerate}
  \item 钱学森弹道/钱学森障碍问题可被写成预算约束下的稳健极小极大问题；
  \item 非真值保持或把真实标签删去的策略，应在一级目标层被直接排除；
  \item 甲方预算扩张使稳健值不增，乙方预算扩张使稳健值不减；
  \item 在紧凸与半连续条件下，极小极大值具有鞍点存在的充分条件；
  \item 单位预算收益不是第一性对象，而是一级近最优集上的二级排序对象。
\end{enumerate}
\end{remark}

\subsection{定义组：两方策略、有效可容许终端集、一级目标泛函与二级收益指标}
\label{subsec:tex1-ch4-defs}

\begin{definition}[两方策略与预算标量]
\label{def:tex1-ch4-strategy}
固定时刻 $t<T$ 与当前观测历史 $b\in\mathcal B_t$。

甲方为修复--判定方，其纯策略记为
\[
  \pi_A=(R,d,C_A,L_A),
\]
其中：
\begin{enumerate}
  \item $R$ 为修复算子链；
  \item $d:\mathcal P_\ast(Y)\to Y$ 为终端选择规则，
        且
        \[
          \mathcal P_\ast(Y):=2^Y\setminus\{\varnothing\};
        \]
  \item $C_A,L_A\ge 0$ 分别为甲方算力预算与时间窗口预算。
\end{enumerate}

乙方为结构变化方，其纯策略记为
\[
  \pi_B=(\mathfrak e,C_B,L_B),
\]
其中：
\begin{enumerate}
  \item $\mathfrak e$ 为结构变化链，它作用于能力元组、约束边界与边缘算子插入规则；
  \item $C_B,L_B\ge 0$ 分别为乙方预算。
\end{enumerate}

再定义两方的加权预算标量
\[
  K_A:=\alpha_A C_A+\beta_A L_A,
  \qquad
  K_B:=\alpha_B C_B+\beta_B L_B,
\]
其中
\[
  \alpha_A,\beta_A,\alpha_B,\beta_B>0.
\]
\end{definition}

\begin{definition}[有效可容许终端集与不可接受策略]
\label{def:tex1-ch4-gamma}
在乙方结构变化 $\mathfrak e$ 作用后，
原可容许终端纤维变为
\[
  \Gamma_t^{\mathfrak e}(b).
\]
再经甲方修复算子链 $R$ 后，
得到有效可容许终端集
\[
  \Gamma_t^{\pi_A,\pi_B}(b)
  :=
  R\bigl(\Gamma_t^{\mathfrak e}(b)\bigr).
\]

若对某个策略对 $(\pi_A,\pi_B)$ 出现以下任一情形：
\begin{enumerate}
  \item 真值保持失败；
  \item $\Gamma_t^{\pi_A,\pi_B}(b)=\varnothing$；
\end{enumerate}
则称该策略对在点 $(t,b)$ \textbf{不可接受}。
后文约定：不可接受策略对在一级目标中被赋值为
\[
  J_{t,b}(\pi_A,\pi_B):=+\infty,
\]
并在二级指标中被赋值为
\[
  \mathrm{CE}_{t,b}(\pi_A,\pi_B):=-\infty.
\]
\end{definition}

\begin{definition}[四个基本代价项]
\label{def:tex1-ch4-cost}
设
\[
  \mu:2^Y\to\RR_{\ge 0}
\]
为分离型歧义度量。
对任一可接受策略对 $(\pi_A,\pi_B)$，定义四个基本代价项如下。

第一项，剩余歧义：
\[
  \mathsf A_{t,b}(\pi_A,\pi_B)
  :=
  \mu\bigl(\Gamma_t^{\pi_A,\pi_B}(b)\bigr).
\]

第二项，达阈时延。
给定阈值 $\varepsilon\ge 0$，
记由策略对 $(\pi_A,\pi_B)$ 诱导的延长观测历史为
\[
  b^{\pi_A,\pi_B}_{[0,t+s]}.
\]
定义
\[
  \mathsf D_{t,b}(\pi_A,\pi_B)
  :=
  \inf\Bigl\{
    s\in[0,T-t]:
    \mu\bigl(
      \Gamma_{t+s}^{\pi_A,\pi_B}(b^{\pi_A,\pi_B}_{[0,t+s]})
    \bigr)\le \varepsilon
  \Bigr\}.
\]
若该集合为空，则记
\[
  \mathsf D_{t,b}(\pi_A,\pi_B):=T-t.
\]

第三项，误判罚项。
设真实终端标签为 $y_\star(b)\in Y$，则
\[
  \mathsf E_{t,b}(\pi_A,\pi_B)
  :=
  \mathbf 1\!\left\{
    d\bigl(\Gamma_t^{\pi_A,\pi_B}(b)\bigr)\neq y_\star(b)
  \right\}.
\]

第四项，预算项为
\[
  K_A,\qquad K_B.
\]
\end{definition}

\begin{definition}[一级目标泛函、极小极大值函数与二级收益指标]
\label{def:tex1-ch4-objective}
取权重
\[
  \lambda_1,\lambda_2,\lambda_3,\lambda_4,\lambda_5>0.
\]
对任一可接受策略对 $(\pi_A,\pi_B)$，定义从甲方视角出发的一级目标泛函
\[
  J_{t,b}(\pi_A,\pi_B)
  :=
  \lambda_1 \mathsf A_{t,b}(\pi_A,\pi_B)
  +
  \lambda_2 \mathsf D_{t,b}(\pi_A,\pi_B)
  +
  \lambda_3 K_A
  +
  \lambda_4 \mathsf E_{t,b}(\pi_A,\pi_B)
  -
  \lambda_5 K_B.
\]
其中负号 $-\lambda_5 K_B$ 的含义是：
乙方在增大甲方损失时，也必须支付自己的结构变化预算。

给定预算上界
\[
  \bar K_A,\bar K_B>0,
\]
定义可行策略集
\[
  \Pi_A(\bar K_A)
  :=
  \{\pi_A:K_A\le \bar K_A\},
  \qquad
  \Pi_B(\bar K_B)
  :=
  \{\pi_B:K_B\le \bar K_B\}.
\]
再定义一级稳健值函数
\[
  \mathcal V(t,b;\bar K_A,\bar K_B)
  :=
  \inf_{\pi_A\in\Pi_A(\bar K_A)}
  \sup_{\pi_B\in\Pi_B(\bar K_B)}
  J_{t,b}(\pi_A,\pi_B).
\]

最后，定义甲方在给定对局下的歧义压缩量
\[
  \Delta\mu_{t,b}(\pi_A,\pi_B)
  :=
  \mu\bigl(\Gamma_t^{\mathfrak e}(b)\bigr)
  -
  \mu\bigl(\Gamma_t^{\pi_A,\pi_B}(b)\bigr),
\]
并取任意固定的 $\delta>0$，定义单位预算收益
\[
  \mathrm{CE}_{t,b}(\pi_A,\pi_B)
  :=
  \frac{\Delta\mu_{t,b}(\pi_A,\pi_B)}{K_A+\delta}.
\]
再定义甲方的最坏情形单位预算收益
\[
  \underline{\mathrm{CE}}(t,b;\bar K_A,\bar K_B)
  :=
  \sup_{\pi_A\in\Pi_A(\bar K_A)}
  \inf_{\pi_B\in\Pi_B(\bar K_B)}
  \mathrm{CE}_{t,b}(\pi_A,\pi_B).
\]
\end{definition}

\subsection{公理（降级为定理）：先稳健后收益的两级目标链}
\label{subsec:tex1-ch4-axiom}

\begin{theorem}[公理（降级为定理）：先稳健后收益的两级目标链]
\label{thm:tex1-ch4-priority}
设
\[
  \mathfrak G_0:\ \text{可接受策略/真值保持},
\]
\[
  \mathfrak G_1:\ \text{一级目标泛函 }J_{t,b},
\]
\[
  \mathfrak G_2:\ \text{稳健值函数 }\mathcal V,
\]
\[
  \mathfrak G_3:\ \text{近最优集上的单位预算收益 } \mathrm{CE}.
\]
则关于钱学森弹道/钱学森障碍的严格策略选择，
必须按如下前缀链递归登记：
\[
  \mathfrak G_0
  \to
  \mathfrak G_1
  \to
  \mathfrak G_2
  \to
  \mathfrak G_3.
\]
也就是说，只有当前一层对象已经获得严格语义时，
后一层对象才允许被调用。
\end{theorem}

\begin{Certificate}[数学归纳法证书：两级目标链的四层前缀登记]
\label{cert:tex1-ch4-priority}
定义谓词
\[
  P(n):
  \exists!\,\text{有效前缀链 }\mathfrak G_0\to\cdots\to \mathfrak G_n,
  \quad
  \text{其长度为 }n+1.
\]
并登记递归兼容条件
\[
  H_0:\ \mathfrak G_0\Rightarrow \mathfrak G_1\ \text{可被定义},
\]
\[
  H_1:\ \mathfrak G_1\Rightarrow \mathfrak G_2\ \text{可被定义},
\]
\[
  H_2:\ \mathfrak G_2\Rightarrow \mathfrak G_3\ \text{可被定义}.
\]
则归纳证书登记为
\[
  \Pi_{\mathrm{lex}}
  =
  \Bigl(
    P(0),\
    H_0,\
    H_1,\
    H_2,\
    \forall n\in\{0,1,2\},\ P(n)\Rightarrow P(n+1)
  \Bigr).
\]
\end{Certificate}

\begin{proof}[证明（数学归纳法证书；基于数学符号推演逻辑的谓词因果链）]
基步 $n=0$。
由定义~\ref{def:tex1-ch4-gamma}，
若策略对不可接受，则其一级目标被直接赋值为 $+\infty$，
二级收益被直接赋值为 $-\infty$。
因此，任何后续比较都必须先建立
\[
  \mathfrak G_0:\ \text{可接受策略/真值保持}
\]
这一层。
故长度为 $1$ 的前缀链 $\mathfrak G_0$ 唯一可登记，这就是 \(P(0)\)。

若已知 \(P(0)\) 成立，则由 \(H_0\) 知，
只有在可接受策略已经被筛出之后，
\[
  J_{t,b}(\pi_A,\pi_B)
\]
才不是混合着 $+\infty$ 罚项的悬空对象，而是具有可比性的一级总代价。
故可从 \(\mathfrak G_0\) 延拓到 \(\mathfrak G_1\)，从而得到 \(P(1)\)。

若已知 \(P(1)\) 成立，则由 \(H_1\) 知，
只有一级目标泛函已经被明确定义之后，
\[
  \mathcal V(t,b;\bar K_A,\bar K_B)
  =
  \inf_{\pi_A}\sup_{\pi_B}J_{t,b}(\pi_A,\pi_B)
\]
才有严格语义。
故可从 \(\mathfrak G_1\) 延拓到 \(\mathfrak G_2\)，得到 \(P(2)\)。

若已知 \(P(2)\) 成立，则由 \(H_2\) 知，
只有稳健值函数已经首先给出最坏情形稳健边界之后，
单位预算收益才允许作为第二层排序对象被调用。
否则 \(\mathrm{CE}\) 可能压过真值保持与最坏情形控制。
故可从 \(\mathfrak G_2\) 延拓到 \(\mathfrak G_3\)，得到 \(P(3)\)。

综上，由有限归纳法得到唯一有效前缀链
\[
  \mathfrak G_0
  \to
  \mathfrak G_1
  \to
  \mathfrak G_2
  \to
  \mathfrak G_3.
\]
定理得证。
\end{proof}

\begin{corollary}[不应直接把单位预算收益放在第一层目标]
\label{cor:tex1-ch4-priority}
若跳过一级稳健值函数 $\mathcal V$ 而直接最大化 $\mathrm{CE}$，
则该选择规则并不保证真值保持与最坏情形稳健性仍被优先保留。
\end{corollary}

\begin{Certificate}[证书：定理~\ref{thm:tex1-ch4-priority} 的直接推论]
\label{cert:tex1-ch4-priority-cor}
证书登记谓词链：
\[
\begin{aligned}
  C^{(4)}_1 &: \text{调用定理~\ref{thm:tex1-ch4-priority}},\\
  C^{(4)}_2 &: \mathrm{CE}\ \text{位于链末层 }\mathfrak G_3,\\
  C^{(4)}_3 &: C^{(4)}_1\wedge C^{(4)}_2
              \Rightarrow \mathfrak G_0,\mathfrak G_1,\mathfrak G_2\ \text{必须先成立},\\
  C^{(4)}_4 &: C^{(4)}_3\Rightarrow \mathrm{CE}\ \text{不能直接充当第一层目标}.
\end{aligned}
\]
\end{Certificate}

\begin{proof}[证明（基于数学符号推演逻辑的谓词因果链）]
由 \(C^{(4)}_1\) 可知，单位预算收益位于四层链的末层。
由 \(C^{(4)}_2\) 可知，任何对 \(\mathrm{CE}\) 的调用都必须被看作对末层对象的调用。
因此，由 \(C^{(4)}_3\) 可知，前缀层 \(\mathfrak G_0,\mathfrak G_1,\mathfrak G_2\) 都必须先被建立。
这正是 \(C^{(4)}_4\)：不能跳过一级稳健层而直接以 \(\mathrm{CE}\) 作为第一目标。推论得证。
\end{proof}

\subsection{引理：无穷罚项约定自动排除不可接受策略}
\label{subsec:tex1-ch4-lemma}

\begin{lemma}[无穷罚项约定自动排除不可接受策略]
\label{lem:tex1-ch4-infty-penalty}
若某个甲方策略 $\pi_A$ 存在乙方策略 $\pi_B\in\Pi_B(\bar K_B)$，
使得策略对 $(\pi_A,\pi_B)$ 不可接受，
则有
\[
  \sup_{\pi_B\in\Pi_B(\bar K_B)}J_{t,b}(\pi_A,\pi_B)=+\infty.
\]
从而，该甲方策略不可能成为任何有限稳健值问题中的最优策略。
\end{lemma}

\begin{Certificate}[证书：上确界吸收无穷证书]
\label{cert:tex1-ch4-infty-penalty}
证书登记谓词链：
\[
\begin{aligned}
  L^{(4)}_1 &: \exists \pi_B^\times\in\Pi_B(\bar K_B),\ (\pi_A,\pi_B^\times)\ \text{不可接受},\\
  L^{(4)}_2 &: L^{(4)}_1\Rightarrow J_{t,b}(\pi_A,\pi_B^\times)=+\infty,\\
  L^{(4)}_3 &: L^{(4)}_2\Rightarrow
              \sup_{\pi_B\in\Pi_B(\bar K_B)}J_{t,b}(\pi_A,\pi_B)\ge +\infty,\\
  L^{(4)}_4 &: L^{(4)}_3\Rightarrow
              \sup_{\pi_B\in\Pi_B(\bar K_B)}J_{t,b}(\pi_A,\pi_B)=+\infty.
\end{aligned}
\]
\end{Certificate}

\begin{proof}[证明（基于数学符号推演逻辑的谓词因果链）]
由 \(L^{(4)}_1\) 可知，存在某个乙方策略 \(\pi_B^\times\)
会把甲方策略 \(\pi_A\) 拉入不可接受区域。
于是由定义~\ref{def:tex1-ch4-gamma} 立得
\[
  J_{t,b}(\pi_A,\pi_B^\times)=+\infty,
\]
即 \(L^{(4)}_2\)。

而上确界至少不小于其定义域内任一取值，
故
\[
  \sup_{\pi_B\in\Pi_B(\bar K_B)}J_{t,b}(\pi_A,\pi_B)
  \ge
  J_{t,b}(\pi_A,\pi_B^\times)
  =
  +\infty.
\]
于是 \(L^{(4)}_3\) 与 \(L^{(4)}_4\) 同时成立。
因此该甲方策略的最坏情形总代价为 \(+\infty\)，
它不可能在任何有限稳健值问题中成为最优策略。引理得证。
\end{proof}

\subsection{定理：极小极大值的语义}
\label{subsec:tex1-ch4-theorem-semantics}

\begin{theorem}[极小极大值的语义：精确版与近似版]
\label{thm:tex1-ch4-semantics}
给定任意阈值 $\theta\in\RR$。

\textnormal{(1)} 若下确界在甲方可行集上被某个策略 $\pi_A^\sharp\in\Pi_A(\bar K_A)$ 取到，
则
\[
  \mathcal V(t,b;\bar K_A,\bar K_B)\le \theta
\]
当且仅当存在某个
\[
  \pi_A^\sharp\in\Pi_A(\bar K_A)
\]
使得对所有
\[
  \pi_B\in\Pi_B(\bar K_B)
\]
都有
\[
  J_{t,b}(\pi_A^\sharp,\pi_B)\le \theta.
\]

\textnormal{(2)} 在不假设下确界可取到的一般情形下，
\[
  \mathcal V(t,b;\bar K_A,\bar K_B)\le \theta
\]
当且仅当对任意 $\delta>0$，
都存在某个
\[
  \pi_A^\delta\in\Pi_A(\bar K_A)
\]
使得对所有
\[
  \pi_B\in\Pi_B(\bar K_B)
\]
都有
\[
  J_{t,b}(\pi_A^\delta,\pi_B)\le \theta+\delta.
\]
\end{theorem}

\begin{Certificate}[证书：极小极大定义展开证书]
\label{cert:tex1-ch4-semantics}
证书登记谓词链：
\[
\begin{aligned}
  T^{(4)}_1 &: \mathcal V=\inf_{\pi_A\in\Pi_A(\bar K_A)}\sup_{\pi_B\in\Pi_B(\bar K_B)}J_{t,b}(\pi_A,\pi_B),\\
  T^{(4)}_2 &: \text{若下确界可取到，则 } \exists \pi_A^\sharp,\
               \mathcal V=\sup_{\pi_B}J_{t,b}(\pi_A^\sharp,\pi_B),\\
  T^{(4)}_3 &: T^{(4)}_2\Rightarrow
               \bigl(\mathcal V\le \theta \Longleftrightarrow
               \forall \pi_B,\ J_{t,b}(\pi_A^\sharp,\pi_B)\le \theta\bigr),\\
  T^{(4)}_4 &: \text{一般情形下，}\mathcal V\le\theta
               \Longleftrightarrow
               \forall \delta>0,\ \exists \pi_A^\delta,\
               \sup_{\pi_B}J_{t,b}(\pi_A^\delta,\pi_B)\le \theta+\delta.
\end{aligned}
\]
\end{Certificate}

\begin{proof}[证明（基于数学符号推演逻辑的谓词因果链）]
由 \(T^{(4)}_1\) 直接展开定义，
\[
  \mathcal V(t,b;\bar K_A,\bar K_B)
  =
  \inf_{\pi_A\in\Pi_A(\bar K_A)}
  \sup_{\pi_B\in\Pi_B(\bar K_B)}
  J_{t,b}(\pi_A,\pi_B).
\]

先证精确版。
若下确界由某个 \(\pi_A^\sharp\) 取到，
则由 \(T^{(4)}_2\) 有
\[
  \mathcal V
  =
  \sup_{\pi_B\in\Pi_B(\bar K_B)}J_{t,b}(\pi_A^\sharp,\pi_B).
\]
于是
\[
  \mathcal V\le \theta
  \Longleftrightarrow
  \sup_{\pi_B}J_{t,b}(\pi_A^\sharp,\pi_B)\le \theta
  \Longleftrightarrow
  \forall \pi_B,\ J_{t,b}(\pi_A^\sharp,\pi_B)\le \theta,
\]
这就是 \(T^{(4)}_3\)。

再证近似版。
若 \(\mathcal V\le \theta\)，则按下确界的定义，
对任意 \(\delta>0\)，存在某个 \(\pi_A^\delta\) 使得
\[
  \sup_{\pi_B\in\Pi_B(\bar K_B)}J_{t,b}(\pi_A^\delta,\pi_B)
  \le
  \mathcal V+\delta
  \le
  \theta+\delta.
\]
反之，若对任意 \(\delta>0\) 都有此性质，
则对所有 \(\delta>0\) 成立
\[
  \mathcal V
  \le
  \sup_{\pi_B}J_{t,b}(\pi_A^\delta,\pi_B)
  \le
  \theta+\delta.
\]
令 \(\delta\downarrow 0\)，便得 \(\mathcal V\le \theta\)。
这就是 \(T^{(4)}_4\)。
定理得证。
\end{proof}

\begin{corollary}[钱学森弹道/钱学森障碍问题被改写为预算上界内压住最坏总代价]
\label{cor:tex1-ch4-semantics}
因此，钱学森弹道/钱学森障碍问题可以被严格改写为：
\[
  \text{甲方能否在预算上界内，压住所有乙方结构变化下的最坏总代价。}
\]
\end{corollary}

\begin{Certificate}[证书：定理~\ref{thm:tex1-ch4-semantics} 的直接推论]
\label{cert:tex1-ch4-semantics-cor}
证书登记谓词链：
\[
\begin{aligned}
  C^{(4)}_5 &: \text{调用定理~\ref{thm:tex1-ch4-semantics}},\\
  C^{(4)}_6 &: C^{(4)}_5\Rightarrow \mathcal V\ \text{表示预算上界内的最坏总代价边界},\\
  C^{(4)}_7 &: C^{(4)}_6\Rightarrow \text{原问题可被改写为“压住最坏总代价”的问题}.
\end{aligned}
\]
\end{Certificate}

\begin{proof}[证明（基于数学符号推演逻辑的谓词因果链）]
由 \(C^{(4)}_5\) 知，\(\mathcal V\) 的定义与语义已经被完全展开。
因此，\(\mathcal V\) 就是“甲方先选、乙方后扰动”这一零和近似模型下的最坏总代价边界，
这正是 \(C^{(4)}_6\)。
故 \(C^{(4)}_7\) 成立。推论得证。
\end{proof}

\subsection{定理：预算单调性}
\label{subsec:tex1-ch4-theorem-budget}

\begin{theorem}[极小极大值函数的预算单调性]
\label{thm:tex1-ch4-budget}
若
\[
  \bar K_A^{(1)}\le \bar K_A^{(2)},
\]
则有
\[
  \mathcal V(t,b;\bar K_A^{(2)},\bar K_B)
  \le
  \mathcal V(t,b;\bar K_A^{(1)},\bar K_B).
\]
若
\[
  \bar K_B^{(1)}\le \bar K_B^{(2)},
\]
则有
\[
  \mathcal V(t,b;\bar K_A,\bar K_B^{(1)})
  \le
  \mathcal V(t,b;\bar K_A,\bar K_B^{(2)}).
\]
\end{theorem}

\begin{Certificate}[证书：可行集包含关系证书]
\label{cert:tex1-ch4-budget}
证书登记谓词链：
\[
\begin{aligned}
  T^{(4)}_5 &: \bar K_A^{(1)}\le \bar K_A^{(2)}
               \Rightarrow
               \Pi_A(\bar K_A^{(1)})\subseteq \Pi_A(\bar K_A^{(2)}),\\
  T^{(4)}_6 &: \bar K_B^{(1)}\le \bar K_B^{(2)}
               \Rightarrow
               \Pi_B(\bar K_B^{(1)})\subseteq \Pi_B(\bar K_B^{(2)}),\\
  T^{(4)}_7 &: \text{对更大的甲方可行集取 }\inf\ \text{不会更差},\\
  T^{(4)}_8 &: \text{对更大的乙方可行集取 }\sup\ \text{不会更小},\\
  T^{(4)}_9 &: T^{(4)}_5\wedge T^{(4)}_7
               \Rightarrow
               \mathcal V(t,b;\bar K_A^{(2)},\bar K_B)
               \le
               \mathcal V(t,b;\bar K_A^{(1)},\bar K_B),\\
  T^{(4)}_{10} &: T^{(4)}_6\wedge T^{(4)}_8
                  \Rightarrow
                  \mathcal V(t,b;\bar K_A,\bar K_B^{(1)})
                  \le
                  \mathcal V(t,b;\bar K_A,\bar K_B^{(2)}).
\end{aligned}
\]
\end{Certificate}

\begin{proof}[证明（基于数学符号推演逻辑的谓词因果链）]
由 \(\bar K_A^{(1)}\le \bar K_A^{(2)}\) 得
\[
  \Pi_A(\bar K_A^{(1)})\subseteq \Pi_A(\bar K_A^{(2)}),
\]
即 \(T^{(4)}_5\)。
对更大的甲方可行集取下确界，不会比对更小的可行集取下确界更差，
故 \(T^{(4)}_7\) 成立。
于是得到 \(T^{(4)}_9\)，即第一条不等式。

同理，由 \(\bar K_B^{(1)}\le \bar K_B^{(2)}\) 得
\[
  \Pi_B(\bar K_B^{(1)})\subseteq \Pi_B(\bar K_B^{(2)}),
\]
即 \(T^{(4)}_6\)。
对更大的乙方可行集取上确界，不会更小，故 \(T^{(4)}_8\) 成立，
从而得到 \(T^{(4)}_{10}\)，即第二条不等式。定理得证。
\end{proof}

\begin{corollary}[算力匹配被写成严格不等式]
\label{cor:tex1-ch4-budget}
于是，“算力匹配”可以被严格写成：
\[
  \text{甲方预算扩张}
  \Rightarrow
  \mathcal V\downarrow,
  \qquad
  \text{乙方预算扩张}
  \Rightarrow
  \mathcal V\uparrow.
\]
\end{corollary}

\begin{Certificate}[证书：定理~\ref{thm:tex1-ch4-budget} 的压缩式推论]
\label{cert:tex1-ch4-budget-cor}
证书登记谓词链：
\[
\begin{aligned}
  C^{(4)}_8 &: \text{调用定理~\ref{thm:tex1-ch4-budget}},\\
  C^{(4)}_9 &: C^{(4)}_8\Rightarrow \bar K_A\uparrow\ \text{时 }\mathcal V\ \text{不增},\\
  C^{(4)}_{10} &: C^{(4)}_8\Rightarrow \bar K_B\uparrow\ \text{时 }\mathcal V\ \text{不减},\\
  C^{(4)}_{11} &: C^{(4)}_9\wedge C^{(4)}_{10}\Rightarrow \text{推论成立}.
\end{aligned}
\]
\end{Certificate}

\begin{proof}[证明（基于数学符号推演逻辑的谓词因果链）]
由 \(C^{(4)}_8\) 得到甲乙双方预算扩张下的两条单调不等式；
这两条不等式分别正是 \(C^{(4)}_9\) 与 \(C^{(4)}_{10}\)。
将二者合并便得到 \(C^{(4)}_{11}\)，推论得证。
\end{proof}

\subsection{定理：鞍点存在的充分条件}
\label{subsec:tex1-ch4-theorem-saddle}

\begin{theorem}[鞍点存在的充分条件]
\label{thm:tex1-ch4-saddle}
若满足以下条件：
\begin{enumerate}
  \item \(\Pi_A(\bar K_A)\) 与 \(\Pi_B(\bar K_B)\) 为紧凸集；
  \item \(J_{t,b}(\pi_A,\pi_B)\) 关于 \(\pi_A\) 下半连续且凸，关于 \(\pi_B\) 上半连续且凹；
  \item 若纯策略集本身不凸，则转而在其概率单纯形上考虑混合策略。
\end{enumerate}
则有
\[
  \inf_{\pi_A\in\Pi_A(\bar K_A)}
  \sup_{\pi_B\in\Pi_B(\bar K_B)}
  J_{t,b}(\pi_A,\pi_B)
  =
  \sup_{\pi_B\in\Pi_B(\bar K_B)}
  \inf_{\pi_A\in\Pi_A(\bar K_A)}
  J_{t,b}(\pi_A,\pi_B),
\]
并且存在鞍点
\[
  (\pi_A^\ast,\pi_B^\ast).
\]
\end{theorem}

\begin{Certificate}[证书：标准极小极大定理证书]
\label{cert:tex1-ch4-saddle}
证书登记谓词链：
\[
\begin{aligned}
  T^{(4)}_{11} &: \Pi_A(\bar K_A),\Pi_B(\bar K_B)\ \text{紧且凸},\\
  T^{(4)}_{12} &: J_{t,b}\ \text{对 }\pi_A\ \text{下半连续且凸，对 }\pi_B\ \text{上半连续且凹},\\
  T^{(4)}_{13} &: \text{若纯策略不凸，则转入混合策略概率单纯形},\\
  T^{(4)}_{14} &: T^{(4)}_{11}\wedge T^{(4)}_{12}\wedge T^{(4)}_{13}
                  \Rightarrow \text{可调用 Sion 型极小极大定理},\\
  T^{(4)}_{15} &: T^{(4)}_{14}\Rightarrow \inf\sup J=\sup\inf J\ \text{且鞍点存在}.
\end{aligned}
\]
\end{Certificate}

\begin{proof}[证明（基于数学符号推演逻辑的谓词因果链）]
在 \(T^{(4)}_{11}\)、\(T^{(4)}_{12}\) 与 \(T^{(4)}_{13}\) 的条件下，
策略域已满足紧凸性要求，目标泛函也满足标准极小极大定理所需的凸凹与半连续条件。
因此可调用 Sion 型极小极大定理
\cite{Sion1958Minimax,Bertsekas2016NonlinearProgramming}，
得到
\[
  \inf_{\pi_A}\sup_{\pi_B}J_{t,b}(\pi_A,\pi_B)
  =
  \sup_{\pi_B}\inf_{\pi_A}J_{t,b}(\pi_A,\pi_B),
\]
并且存在鞍点 \((\pi_A^\ast,\pi_B^\ast)\)。
这正是 \(T^{(4)}_{15}\)。定理得证。
\end{proof}

\begin{corollary}[直觉上的对抗被提升为有值的可证明优化问题]
\label{cor:tex1-ch4-saddle}
在上述紧凸与半连续条件下，
钱学森弹道/钱学森障碍问题中的“直觉对抗”被提升为一个有值的可证明优化问题。
\end{corollary}

\begin{Certificate}[证书：定理~\ref{thm:tex1-ch4-saddle} 的解释性推论]
\label{cert:tex1-ch4-saddle-cor}
证书登记谓词链：
\[
\begin{aligned}
  C^{(4)}_{12} &: \text{调用定理~\ref{thm:tex1-ch4-saddle}},\\
  C^{(4)}_{13} &: C^{(4)}_{12}\Rightarrow \inf\sup J=\sup\inf J,\\
  C^{(4)}_{14} &: C^{(4)}_{13}\Rightarrow \mathcal V\ \text{是有值对象而非仅靠直觉的比较量},\\
  C^{(4)}_{15} &: C^{(4)}_{14}\Rightarrow \text{推论成立}.
\end{aligned}
\]
\end{Certificate}

\begin{proof}[证明（基于数学符号推演逻辑的谓词因果链）]
由 \(C^{(4)}_{12}\) 得到极小极大相等与鞍点存在。
于是 \(\mathcal V\) 不再只是口头上的最坏情形代价，
而是一个具有严格存在性的数值对象，这就是 \(C^{(4)}_{14}\)。
故 \(C^{(4)}_{15}\) 成立。推论得证。
\end{proof}

\subsection{命题：一级目标优先，二级指标排序}
\label{subsec:tex1-ch4-prop-lexicographic}

\begin{proposition}[一级目标优先，二级指标排序]
\label{prop:tex1-ch4-lexicographic}
给定任意 \(\eta\ge 0\)，定义甲方的近最优集
\[
  \Pi_A^\eta(t,b)
  :=
  \Bigl\{
    \pi_A\in\Pi_A(\bar K_A):
    \sup_{\pi_B\in\Pi_B(\bar K_B)}J_{t,b}(\pi_A,\pi_B)
    \le
    \mathcal V(t,b;\bar K_A,\bar K_B)+\eta
  \Bigr\}.
\]
则更稳妥的选择规则不是直接最大化
\[
  \mathrm{CE}_{t,b},
\]
而是先限制在 \(\Pi_A^\eta(t,b)\) 中，
再求
\[
  \pi_A^{\mathrm{ce}}
  \in
  \operatorname*{arg\,max}_{\pi_A\in\Pi_A^\eta(t,b)}
  \inf_{\pi_B\in\Pi_B(\bar K_B)}
  \mathrm{CE}_{t,b}(\pi_A,\pi_B).
\]
\end{proposition}

\begin{Certificate}[证书：词典序优化证书]
\label{cert:tex1-ch4-lexicographic}
证书登记谓词链：
\[
\begin{aligned}
  P^{(4)}_1 &: \Pi_A^\eta(t,b)\ \text{把一级总代价控制在 }\mathcal V+\eta\ \text{以内},\\
  P^{(4)}_2 &: \text{若直接最大化 }\mathrm{CE},\ \text{则可能选出一级总代价过大的策略},\\
  P^{(4)}_3 &: P^{(4)}_1\Rightarrow \text{先锁定稳健层，再比较收益层},\\
  P^{(4)}_4 &: P^{(4)}_3\Rightarrow
               \operatorname*{arg\,max}_{\pi_A\in\Pi_A^\eta(t,b)}
               \inf_{\pi_B}\mathrm{CE}_{t,b}(\pi_A,\pi_B)
               \text{是更稳妥的二级排序规则}.
\end{aligned}
\]
\end{Certificate}

\begin{proof}[证明（基于数学符号推演逻辑的谓词因果链）]
由 \(P^{(4)}_1\) 可知，
一旦把策略限制在近最优集 \(\Pi_A^\eta(t,b)\) 中，
一级最坏情形总代价就已经被锁定在
\[
  \mathcal V(t,b;\bar K_A,\bar K_B)+\eta
\]
以内。

若直接在整个 \(\Pi_A(\bar K_A)\) 上最大化 \(\mathrm{CE}\)，
则可能出现“预算效率看似很高、但一级最坏总代价很差”的策略，
这就是 \(P^{(4)}_2\)。

因此，必须先利用 \(P^{(4)}_1\) 锁定一级稳健层，
再在该层中最大化二级收益，
即得到 \(P^{(4)}_3\)。
于是 \(P^{(4)}_4\) 成立，
命题得证。
\end{proof}

\begin{corollary}[严格表达不是只看性价比，而是先稳健后性价比]
\label{cor:tex1-ch4-lexicographic}
因此，严格表达不是
\[
  \text{只看性价比},
\]
而是
\[
  \text{先求最坏情形稳健最优，再在近最优集里比较单位预算收益}.
\]
\end{corollary}

\begin{Certificate}[证书：命题~\ref{prop:tex1-ch4-lexicographic} 的压缩推论]
\label{cert:tex1-ch4-lexicographic-cor}
证书登记谓词链：
\[
\begin{aligned}
  C^{(4)}_{16} &: \text{调用命题~\ref{prop:tex1-ch4-lexicographic}},\\
  C^{(4)}_{17} &: C^{(4)}_{16}\Rightarrow \text{直接最大化 }\mathrm{CE}\ \text{被排除},\\
  C^{(4)}_{18} &: C^{(4)}_{17}\Rightarrow \text{必须先稳健后收益}.
\end{aligned}
\]
\end{Certificate}

\begin{proof}[证明（基于数学符号推演逻辑的谓词因果链）]
由 \(C^{(4)}_{16}\) 可知，二级收益比较只能在近最优集内部进行。
于是 \(C^{(4)}_{17}\) 成立：直接最大化 \(\mathrm{CE}\) 的规则被排除。
因此 \(C^{(4)}_{18}\) 成立。推论得证。
\end{proof}

\subsection{命题：误判罚项的主导条件}
\label{subsec:tex1-ch4-prop-penalty}

\begin{proposition}[误判罚项的主导条件]
\label{prop:tex1-ch4-penalty}
设存在上界
\[
  0\le \mathsf A_{t,b}\le M_A^{\mathrm{amb}},
  \qquad
  0\le \mathsf D_{t,b}\le M_A^{\mathrm{delay}},
\]
\[
  0\le K_A\le \bar K_A,
  \qquad
  0\le K_B\le \bar K_B.
\]
若选择系数满足
\[
  \lambda_4
  >
  \lambda_1 M_A^{\mathrm{amb}}
  +
  \lambda_2 M_A^{\mathrm{delay}}
  +
  \lambda_3 \bar K_A
  +
  \lambda_5 \bar K_B,
\]
则误判一次所带来的总代价增加，
严格大于其余各项在全可行域内所能带来的最大改善。
\end{proposition}

\begin{Certificate}[证书：有界差值估计证书]
\label{cert:tex1-ch4-penalty}
证书登记谓词链：
\[
\begin{aligned}
  P^{(4)}_5 &: 0\le \mathsf A_{t,b}\le M_A^{\mathrm{amb}},\quad 0\le \mathsf D_{t,b}\le M_A^{\mathrm{delay}},\\
  P^{(4)}_6 &: 0\le K_A\le \bar K_A,\quad 0\le K_B\le \bar K_B,\\
  P^{(4)}_7 &: P^{(4)}_5\wedge P^{(4)}_6
               \Rightarrow
               \text{除误判项外，其余各项总改变量}
               \le
               \lambda_1 M_A^{\mathrm{amb}}
               +
               \lambda_2 M_A^{\mathrm{delay}}
               +
               \lambda_3 \bar K_A
               +
               \lambda_5 \bar K_B,\\
  P^{(4)}_8 &: \lambda_4
               >
               \lambda_1 M_A^{\mathrm{amb}}
               +
               \lambda_2 M_A^{\mathrm{delay}}
               +
               \lambda_3 \bar K_A
               +
               \lambda_5 \bar K_B,\\
  P^{(4)}_9 &: P^{(4)}_7\wedge P^{(4)}_8
               \Rightarrow
               \text{一次误判的代价增加不可能被其余项的改善抵消}.
\end{aligned}
\]
\end{Certificate}

\begin{proof}[证明（基于数学符号推演逻辑的谓词因果链）]
由 \(P^{(4)}_5\) 与 \(P^{(4)}_6\) 可知，
除误判项之外，其余项在全可行域内的总改变量至多为
\[
  \lambda_1 M_A^{\mathrm{amb}}
  +
  \lambda_2 M_A^{\mathrm{delay}}
  +
  \lambda_3 \bar K_A
  +
  \lambda_5 \bar K_B,
\]
这就是 \(P^{(4)}_7\)。

若再满足 \(P^{(4)}_8\)，
则从
\[
  \mathsf E_{t,b}=0
\]
跳到
\[
  \mathsf E_{t,b}=1
\]
所增加的代价 \(\lambda_4\)
严格大于其余各项能够带来的任何最大改善。
因此，一次误判所造成的总代价增加不可能被其余项抵消，
即 \(P^{(4)}_9\)。
命题得证。
\end{proof}

\begin{corollary}[若把正确性置于首位，就让 \texorpdfstring{$\lambda_4$}{lambda4} 成为主导系数]
\label{cor:tex1-ch4-penalty}
因此，若希望把正确性放在首位，
则应让误判罚项系数 \(\lambda_4\) 满足命题~\ref{prop:tex1-ch4-penalty} 的主导条件。
\end{corollary}

\begin{Certificate}[证书：命题~\ref{prop:tex1-ch4-penalty} 的直接应用]
\label{cert:tex1-ch4-penalty-cor}
证书登记谓词链：
\[
\begin{aligned}
  C^{(4)}_{19} &: \text{调用命题~\ref{prop:tex1-ch4-penalty}},\\
  C^{(4)}_{20} &: C^{(4)}_{19}\Rightarrow \lambda_4\ \text{主导时，误判代价不可被其余项抵消},\\
  C^{(4)}_{21} &: C^{(4)}_{20}\Rightarrow \text{若把正确性置于首位，应取主导 }\lambda_4.
\end{aligned}
\]
\end{Certificate}

\begin{proof}[证明（基于数学符号推演逻辑的谓词因果链）]
由 \(C^{(4)}_{19}\) 得，主导条件一旦满足，
一次误判就不可能被其余项改善所抵消。
这正是 \(C^{(4)}_{20}\)。
故若把正确性置于首位，
便应按 \(C^{(4)}_{21}\) 选择主导的 \(\lambda_4\)。推论得证。
\end{proof}

\subsection{备注：层级关系、完整压缩句与下一步离散化方向}
\label{subsec:tex1-ch4-remarks}

\begin{remark}[完整层级关系]
到这里为止，本章把对象层次进一步收紧为
\[
  \text{闭合性}
  \succ
  \text{真值保持}
  \succ
  \text{最坏情形总代价 }\mathcal V
  \succ
  \text{单位预算收益 }\mathrm{CE}.
\]
因此，“性价比”只能作为稳健层之后的二级比较量，
不能越过结构完备性与真值保持而直接充当第一性对象。
\end{remark}

\begin{remark}[本章的一句完整压缩句]
本章可以压缩为：
\begin{center}
  \fbox{\parbox{0.9\linewidth}{
    \centering
    钱学森弹道/钱学森障碍问题在严格对象层上，
    先被写成预算约束下的稳健极小极大问题 $\mathcal V$，
    再在近最优解集内用 $\mathrm{CE}$ 比较单位预算收益；
    因而“障碍修复和算力匹配的性价比”只能作为第二层排序，而不是第一层判据。
  }}
\end{center}
\end{remark}

\begin{remark}[下一步自然方向：离散时间 Bellman 递推或马尔可夫博弈化]
若继续推进，
下一步最自然的工作不再是继续停留在口头直觉，
而是把
\[
  R,\qquad \mathfrak e,\qquad \mu,\qquad d,\qquad y_\star
\]
逐个具体化，
做成离散时间版本
\[
  t=0,1,\dots,N
\]
上的动态规划或马尔可夫博弈模型。
一旦进入这一层，
便可以继续推出 Bellman 型递推关系或鞍点迭代公式。
\end{remark}

%%%%%%%%%%%%%%%%%%%%%%%%%%%%%%%%%%%%%%%%%%%%%%%%%%%%%%%%%%%%
\clearpage
\section{形式逻辑：钱学森弹道/钱学森障碍的障碍修复变分问题、同伦扭曲命中与条件性求解定理}
\label{sec:tex1-ch5}
%%%%%%%%%%%%%%%%%%%%%%%%%%%%%%%%%%%%%%%%%%%%%%%%%%%%%%%%%%%%

\begin{remark}[本章意义备注：高超音速攻防的抽象语境]
本章的意义，在于给出障碍修复的条件性求解框架，
把“能否逼近可判区”写成变分问题，并把关键命中步嵌入递推链。
对高超音速攻防的抽象理解而言，这一章说明从不可判到可判必须依赖可验证的修复路径与条件证书。
\end{remark}

\subsection{章节目标与仓内证据链}
\label{subsec:tex1-ch5-goal}

\begin{remark}[本章要压成的严格对象]
在钱学森弹道/钱学森障碍的抽象对象层上，
本章把问题进一步收紧为两个顺序明确的对象：
\[
  \text{一级对象：障碍修复的变分求解问题},
\]
\[
  \text{二级对象：同伦扭曲命中作为变分可行域中的关键支配步}.
\]
因此，真正的求解对象不再是原始表象弹道本身，
而是“从局部最优到终端全局最优正规形的合法提升链极小化”。
这一重写同时调用五段仓内证据链：
\cite{RepoTex36HomotopyTwistEquivalence,RepoTex38JacobiBianchiRepair,RepoTex39OpenRepairFamily,
  RepoTex41HomotopyRepairBundle,RepoTex53PDEStrategyCluster}。
其中，\cite{RepoTex36HomotopyTwistEquivalence} 负责“高阶障碍可消去 + 存在性与最优性双重收敛”，
\cite{RepoTex41HomotopyRepairBundle} 负责“可行域上的变分极小元存在”，
\cite{RepoTex38JacobiBianchiRepair} 负责“有限终止、终端唯一最优正规形、唯一最优提升路径与命中步必要性”，
而 \cite{RepoTex39OpenRepairFamily} 负责把“命中”收紧为带结构判据的合法分支命中。
\cite{RepoTex53PDEStrategyCluster} 进一步把“同伦扭曲命中”落到统计力学参数格点、有限候选评分与策略簇自由能变分解释上。
统计力学采样、Gibbs 分布、最大熵与退火式优化背景可参照
\cite{Gibbs1902Elementary,Jaynes1957Information,Kirkpatrick1983Optimization,GemanGeman1984StochasticRelaxation}。
\end{remark}

\begin{remark}[本章的严格边界]
本章仍只讨论抽象修复塔、变分路径族、Bellman 递推与同伦扭曲命中之间的对象层关系；
``钱学森弹道''、``钱学森障碍'' 与 ``弹道'' 只作为专业应用名词被完整保留，
不承载任何具体实现、反求或对抗设计信息。
\end{remark}

\subsection{定义组：终端导向的障碍修复变分问题与同伦扭曲命中}
\label{subsec:tex1-ch5-defs}

\begin{definition}[终端导向的障碍修复变分问题]
\label{def:tex1-ch5-variational}
设终端唯一全局最优正规形为
\[
  \tau^\ast\in \mathcal T_{N_\ast}(s,t),
\]
并记从各层局部最优对象提升到终端点 \(\tau^\ast\) 的合法提升路径族为
\[
  \mathfrak L^{\mathrm{loc}\to\tau^\ast}.
\]
对任一路径
\[
  \ell
  =
  (\tau_{n_0},\tau_{n_0+1},\dots,\tau_{N_\ast})
  \in
  \mathfrak L^{\mathrm{loc}\to\tau^\ast},
\]
定义路径作用量
\[
  \mathbb A(\ell)
  :=
  \Phi_{n_0}(\tau_{n_0})
  +
  \sum_{k=n_0}^{N_\ast-1} c_k(\tau_k,\tau_{k+1}),
\]
并定义严格化路径作用量
\[
  \Psi(\ell)
  :=
  \mathbb A(\ell)
  +
  \delta\,\Lambda_{\mathrm{path}}(\ell),
  \qquad \delta>0.
\]
相应的终端导向障碍修复变分问题定义为
\[
  \ell^\ast
  =
  \operatorname*{arg\,min}_{\ell\in\mathfrak L^{\mathrm{loc}\to\tau^\ast}}
  \Psi(\ell).
\]
因此，在钱学森弹道/钱学森障碍的抽象对象层上，
求解对象已不再是原始表象轨迹，
而是“从局部最优到终端最优正规形的合法提升链的极小化问题”。
\end{definition}

\begin{definition}[同伦扭曲命中]
\label{def:tex1-ch5-hit}
固定层级 \(n\) 与局部对象 \(\tau\in\mathcal T_n\)。
记
\[
  \mathfrak B_n(\tau)\subseteq \mathcal T_n
\]
为 \(\tau\) 所属旧局部盆的同层可见像，
\[
  \iota_n^{\mathcal T}:\mathcal T_n\hookrightarrow \mathcal T_{n+1}
\]
为层间嵌入，
并记
\[
  \nu_n(\tau)\in \NN
\]
为 \(\tau\) 处的活跃障碍计数。
若某一步合法提升
\[
  \tau \rightsquigarrow \eta,
  \qquad
  \eta\in \mathrm{Lift}_n(\tau),
\]
同时满足
\[
  \eta\notin \iota_n^{\mathcal T}\bigl(\mathfrak B_n(\tau)\bigr),
\]
以及
\[
  \nu_{n+1}(\eta)\le \nu_n(\tau)-1,
\]
则称该步为一次\textbf{同伦扭曲命中}，
亦即一次带严格判据的“微观隧穿命中”。
换言之，所谓“命中”在这里被严格写成
\[
  \text{离开旧局部盆的像}
  \quad+\quad
  \text{活跃障碍计数严格下降}.
\]
\end{definition}

\subsection{形式逻辑：从 PDE 技术求解到同伦扭曲微观隧穿命中}
\label{subsec:tex1-ch5-pde-to-tunnel-hit}

\begin{definition}[PDE 表象求解与结构性命中问题]
\label{def:tex1-ch5-pde-hit-problem}
固定一个观测--终端任务 \(q\)，令
\[
  \mathsf{PDE}(q)
\]
表示把 \(q\) 表述为偏微分方程、连续轨迹追踪或外推弹道的技术性求解问题。
令
\[
  \mathsf{Obs}(q)
  :=
  \bigl(
    B_t^{\Lambda_t},
    \Gamma_t^{\Lambda_t},
    \mathcal F_t^{\Lambda_t}(\varepsilon),
    \nu_n,
    \mathfrak B_n
  \bigr)
\]
表示由第~\ref{sec:tex1-ch1} 节和本节定义给出的障碍对象包。
其中 \(B_t^{\Lambda_t}\) 是观测历史基底，
\(\Gamma_t^{\Lambda_t}\) 是可容许终端纤维，
\(\mathcal F_t^{\Lambda_t}(\varepsilon)\) 是失效域，
\(\nu_n\) 是活跃障碍计数，
\(\mathfrak B_n\) 是旧局部盆像。

称 \(q\) 的\textbf{同伦扭曲微观隧穿命中问题}为如下存在性问题：
\[
  \mathsf{TunHit}(q)
  :
  \Longleftrightarrow
  \exists n,\tau,\eta,K
  \quad
  \eta\in\mathcal C_n(\tau;K)
  \wedge
  \eta\notin\iota_n^{\mathcal T}\bigl(\mathfrak B_n(\tau)\bigr)
  \wedge
  \nu_{n+1}(\eta)\le\nu_n(\tau)-1
  \wedge
  \operatorname{end}(\eta)\sim\tau^\ast.
\]
这里 \(\mathcal C_n(\tau;K)\) 是在算力--时间预算 \(K\) 下可枚举或可采样的局部候选族，
\(\operatorname{end}(\eta)\sim\tau^\ast\) 表示该候选指向与终端全局最优正规形同类的终端模式。
因此，结构性命中问题不是更快求解原 PDE，
而是把原问题改写为“在同伦扭曲后的障碍空间中寻找满足降障碍、离盆像与终端模式命中的微观隧穿路径”。
\end{definition}

\begin{theorem}[公理（降级为定理）：PDE 表象可按前缀链升维为微观隧穿命中问题]
\label{thm:tex1-ch5-pde-uplift-induction}
设存在四段对象变换
\[
  \mathcal X_0
  \xrightarrow{U_0}
  \mathcal X_1
  \xrightarrow{U_1}
  \mathcal X_2
  \xrightarrow{U_2}
  \mathcal X_3
  \xrightarrow{U_3}
  \mathcal X_4,
\]
其中
\[
\begin{aligned}
  \mathcal X_0&:=\mathsf{PDE}(q),\\
  \mathcal X_1&:=\bigl(B_t^{\Lambda_t},\Gamma_t^{\Lambda_t},\mathcal F_t^{\Lambda_t}(\varepsilon)\bigr),\\
  \mathcal X_2&:=\bigl(\mathcal X_1,\mathfrak B_n,\nu_n,\mathrm{Lift}_n\bigr),\\
  \mathcal X_3&:=\bigl(\mathcal C_n(\tau;K),E_n,\beta\bigr),\\
  \mathcal X_4&:=\mathsf{TunHit}(q).
\end{aligned}
\]
并设每个 \(U_i\) 都保持终端可判谓词
\[
  \mathsf{Term}(q)
\]
的真值方向，即
\[
  P(i):\quad
  \mathcal X_i\text{ 已定义且 }
  \mathcal X_i\Rightarrow \mathsf{Term}(q)
\]
满足
\[
  P(i)\Rightarrow P(i+1)
  \qquad (i=0,1,2,3).
\]
则
\[
  \mathsf{PDE}(q)
  \Longrightarrow
  \mathsf{TunHit}(q)
\]
作为结构升维链成立。
\end{theorem}

\begin{Certificate}[数学归纳法证书：四段升维前缀]
\label{cert:tex1-ch5-pde-uplift-induction}
归纳谓词为
\[
  P(i):
  \mathcal X_i\text{ 已定义且 }
  \mathcal X_i\Rightarrow\mathsf{Term}(q).
\]
证书登记为
\[
  \pi_{\mathrm{PDE}\to\mathrm{hit}}
  =
  \Bigl(
    P(0),\
    \forall i\in\{0,1,2,3\},\ P(i)\Rightarrow P(i+1)
  \Bigr).
\]
四个后继箭头分别表示：
\[
\begin{aligned}
  U_0 &: \text{PDE 表象}\Rightarrow \Gamma\text{-纤维障碍空间},\\
  U_1 &: \Gamma\text{-纤维障碍空间}\Rightarrow \text{同伦扭曲候选空间},\\
  U_2 &: \text{同伦扭曲候选空间}\Rightarrow \text{统计力学采样空间},\\
  U_3 &: \text{统计力学采样空间}\Rightarrow \text{微观隧穿命中谓词}.
\end{aligned}
\]
\end{Certificate}

\begin{proof}[证明（数学归纳法证书；基于数学符号推演逻辑的谓词因果链）]
对 \(i\in\{0,1,2,3,4\}\) 作有限数学归纳。

\textbf{基步：}
\(\mathcal X_0=\mathsf{PDE}(q)\) 是原始技术性表象，
其终端目标正是判断 \(\mathsf{Term}(q)\)，故 \(P(0)\) 成立。

\textbf{归纳步：}
设 \(P(i)\) 成立。
若 \(i=0\)，由定义~\ref{def:tex1-terminal-fiber} 与定义~\ref{def:tex1-qxs-terms}，
PDE 表象中的轨迹求解被送入
\[
  b\longmapsto\Gamma_t^{\Lambda_t}(b)
\]
的集合值逆问题，故 \(P(1)\) 成立。
若 \(i=1\)，由定义~\ref{def:tex1-ch5-hit}，
障碍纤维进一步被写成旧盆像、活跃障碍计数与合法提升候选的命中判据，故 \(P(2)\) 成立。
若 \(i=2\)，由预算 \(K\) 截断出有限或可采样候选族
\[
  \mathcal C_n(\tau;K)
\]
并赋予能量 \(E_n\)，候选比较进入统计力学采样空间，故 \(P(3)\) 成立。
若 \(i=3\)，由离开旧盆像、严格降障碍与终端模式命中三个谓词的合取，
得到 \(\mathsf{TunHit}(q)\)，故 \(P(4)\) 成立。

于是
\[
  P(0)\wedge
  \bigwedge_{i=0}^{3}(P(i)\Rightarrow P(i+1))
  \Longrightarrow
  P(4).
\]
由有限数学归纳法，PDE 表象可按该前缀链升维为微观隧穿命中问题。证毕。
\end{proof}

\begin{lemma}[统计力学采样把候选比较压成命中概率]
\label{lem:tex1-ch5-stat-hit}
固定 \(n,\tau,K\)，设候选族
\[
  \mathcal C_n(\tau;K)
\]
有限，能量函数为
\[
  E_n:\mathcal C_n(\tau;K)\to\RR,
\]
并定义命中候选集
\[
  \mathcal H_n(\tau;K)
  :=
  \bigl\{
    \eta\in\mathcal C_n(\tau;K)
    \mid
    \eta\notin\iota_n^{\mathcal T}(\mathfrak B_n(\tau)),
    \ \nu_{n+1}(\eta)\le\nu_n(\tau)-1
  \bigr\}.
\]
对 \(\beta>0\)，定义 Gibbs 采样分布
\[
  \mathbb P_\beta(\eta)
  :=
  \frac{\exp(-\beta E_n(\eta))}
       {\sum_{\xi\in\mathcal C_n(\tau;K)}\exp(-\beta E_n(\xi))}.
\]
若存在严格能隙
\[
  \min_{\eta\in\mathcal H_n(\tau;K)}E_n(\eta)
  <
  \min_{\xi\in\mathcal C_n(\tau;K)\setminus\mathcal H_n(\tau;K)}E_n(\xi),
\]
则
\[
  \lim_{\beta\to\infty}
  \mathbb P_\beta\bigl(\mathcal H_n(\tau;K)\bigr)=1.
\]
\end{lemma}

\begin{Certificate}[证书：有限 Gibbs 分布 + 严格能隙 + 零温极限]
\label{cert:tex1-ch5-stat-hit}
证书由三段组成：
\[
\begin{aligned}
  S_1 &: \mathcal C_n(\tau;K)\text{ 有限}
        \Rightarrow Z_\beta<\infty;\\
  S_2 &: \min E_n|_{\mathcal H}
        <
        \min E_n|_{\mathcal C\setminus\mathcal H}
        \Rightarrow
        \text{命中候选具有严格能量优势};\\
  S_3 &: S_1\wedge S_2
        \Rightarrow
        \beta\to\infty\text{ 时 Gibbs 质量集中到命中极小元}.
\end{aligned}
\]
该证书的统计力学背景来自 Gibbs 分布、最大熵和退火优化口径
\cite{Gibbs1902Elementary,Jaynes1957Information,Kirkpatrick1983Optimization,GemanGeman1984StochasticRelaxation}。
\end{Certificate}

\begin{proof}[证明（基于数学符号推演逻辑的谓词因果链）]
记
\[
  e_H:=\min_{\eta\in\mathcal H_n(\tau;K)}E_n(\eta),
  \qquad
  e_N:=\min_{\xi\in\mathcal C_n(\tau;K)\setminus\mathcal H_n(\tau;K)}E_n(\xi).
\]
由假设 \(e_H<e_N\)，存在 \(\Delta=e_N-e_H>0\)。
于是非命中候选的总权重相对命中极小能级满足
\[
  \sum_{\xi\notin\mathcal H_n}
  \exp(-\beta E_n(\xi))
  \le
  |\mathcal C_n|\exp(-\beta(e_H+\Delta)).
\]
而命中候选中至少有一个 \(\eta_H\) 满足 \(E_n(\eta_H)=e_H\)，故
\[
  \sum_{\eta\in\mathcal H_n}
  \exp(-\beta E_n(\eta))
  \ge
  \exp(-\beta e_H).
\]
因此
\[
  \mathbb P_\beta(\mathcal C_n\setminus\mathcal H_n)
  \le
  |\mathcal C_n|\exp(-\beta\Delta)
  \longrightarrow 0.
\]
于是
\[
  \mathbb P_\beta(\mathcal H_n)
  =
  1-\mathbb P_\beta(\mathcal C_n\setminus\mathcal H_n)
  \longrightarrow 1.
\]
引理得证。
\end{proof}

\begin{theorem}[结构性命中定理：PDE 技术求解问题可改写为同伦扭曲障碍空间中的命中问题]
\label{thm:tex1-ch5-structural-hit}
在固定任务 \(q\) 上，若满足以下条件：
\begin{enumerate}
  \item \(\mathsf{PDE}(q)\) 与 \(\mathsf{Obs}(q)\) 在终端纤维上 \(\Gamma\)-等价；
  \item 对每个非零障碍对象 \(\tau\)，存在预算可枚举或可采样候选族 \(\mathcal C_n(\tau;K)\)；
  \item 命中候选集 \(\mathcal H_n(\tau;K)\) 非空；
  \item 能量 \(E_n\) 对命中候选具有严格优势；
  \item 命中候选的终端模式与 \(\tau^\ast\) 相容。
\end{enumerate}
则 \(q\) 的求解核心可严格改写为
\[
  \mathsf{PDE}(q)
  \Longrightarrow
  \mathsf{Obs}(q)
  \Longrightarrow
  \text{统计力学采样}
  \Longrightarrow
  \mathsf{TunHit}(q).
\]
也即，原问题的本质不再是单纯求解 PDE、连续轨迹外推或低维曲线预测，
而是在同伦扭曲后的障碍空间中，
借助算力预算和统计力学式采样，
寻找能够穿透旧盆像、降低活跃障碍计数并命中真实终端模式的微观隧穿路径。
\end{theorem}

\begin{Certificate}[证书：\(\Gamma\)-等价 + 候选有限化 + Gibbs 命中 + 终端相容]
\label{cert:tex1-ch5-structural-hit}
证书登记为
\[
\begin{aligned}
  T_1 &: \mathsf{PDE}(q)\sim_\Gamma\mathsf{Obs}(q),\\
  T_2 &: \nu_n(\tau)>0\Rightarrow \mathcal C_n(\tau;K)\neq\varnothing,\\
  T_3 &: \mathcal H_n(\tau;K)\neq\varnothing,\\
  T_4 &: \min E_n|_{\mathcal H}<\min E_n|_{\mathcal C\setminus\mathcal H},\\
  T_5 &: \operatorname{end}(\eta)\sim\tau^\ast\quad(\eta\in\mathcal H).
\end{aligned}
\]
其中 \(T_1\) 回挂定义~\ref{def:tex1-gamma-equivalent} 与定理~\ref{thm:tex1-gamma-invariant}，
\(T_2,T_3,T_5\) 回挂定义~\ref{def:tex1-ch5-hit} 与 \cite{RepoTex38JacobiBianchiRepair,RepoTex39OpenRepairFamily}，
\(T_4\) 回挂引理~\ref{lem:tex1-ch5-stat-hit} 与 \cite{RepoTex53PDEStrategyCluster}。
\end{Certificate}

\begin{proof}[证明（基于数学符号推演逻辑的谓词因果链）]
由 \(T_1\)，PDE 表象与障碍对象包在终端纤维上 \(\Gamma\)-等价。
因此，由定理~\ref{thm:tex1-gamma-invariant}，终端可判性不取决于 PDE 表示本身，
而取决于集合值逆映射和可容许终端纤维。

由 \(T_2\)，非零障碍阶段存在预算内候选族。
由 \(T_3\)，其中包含离开旧盆像并严格降低活跃障碍计数的命中候选。
由 \(T_4\) 与引理~\ref{lem:tex1-ch5-stat-hit}，
统计力学采样在零温极限下把概率质量集中到命中候选上。
由 \(T_5\)，该命中候选的终端模式与 \(\tau^\ast\) 相容。
于是
\[
  \exists\eta\in\mathcal C_n(\tau;K)
  \quad
  \eta\in\mathcal H_n(\tau;K)
  \wedge
  \operatorname{end}(\eta)\sim\tau^\ast,
\]
即 \(\mathsf{TunHit}(q)\) 成立。

故技术性 PDE 求解问题已经被改写为同伦扭曲障碍空间中的结构性命中问题。定理得证。
\end{proof}

\begin{proposition}[微观隧穿命中是障碍修复的局部执行判据]
\label{prop:tex1-ch5-tunnel-hit-local}
在定理~\ref{thm:tex1-ch5-structural-hit} 的条件下，若某个候选
\[
  \eta^\star\in\mathcal C_n(\tau;K)
\]
同时满足
\[
  \eta^\star\notin\iota_n^{\mathcal T}\bigl(\mathfrak B_n(\tau)\bigr),
  \qquad
  \nu_{n+1}(\eta^\star)\le\nu_n(\tau)-1,
\]
\[
  E_n(\eta^\star)
  =
  \min_{\eta\in\mathcal C_n(\tau;K)}E_n(\eta),
  \qquad
  \operatorname{end}(\eta^\star)\sim\tau^\ast,
\]
则 \(\eta^\star\) 可签发为一次微观隧穿命中。
\end{proposition}

\begin{Certificate}[证书：离盆像 + 降障碍 + 能量极小 + 终端命中]
\label{cert:tex1-ch5-tunnel-hit-local}
证书链为
\[
\begin{aligned}
  M_1 &: \eta^\star\notin\iota_n^{\mathcal T}(\mathfrak B_n(\tau));\\
  M_2 &: \nu_{n+1}(\eta^\star)\le\nu_n(\tau)-1;\\
  M_3 &: E_n(\eta^\star)=\min_{\eta\in\mathcal C_n(\tau;K)}E_n(\eta);\\
  M_4 &: \operatorname{end}(\eta^\star)\sim\tau^\ast;\\
  M_5 &: M_1\wedge M_2\wedge M_3\wedge M_4
        \Rightarrow \mathsf{TunHit}(q).
\end{aligned}
\]
该证书与 tex53\_pub 中“统计力学参数格点--同伦扭曲命中--有限候选评分--策略簇重编码”的证书口径一致
\cite{RepoTex53PDEStrategyCluster}。
\end{Certificate}

\begin{proof}[证明（基于数学符号推演逻辑的谓词因果链）]
由 \(M_1\)，候选已经离开旧局部盆像。
由 \(M_2\)，候选严格降低或至少按定义降低活跃障碍计数。
二者合取正是定义~\ref{def:tex1-ch5-hit} 的同伦扭曲命中核心判据。
由 \(M_3\)，该候选又是预算候选族中的变分极小元。
由 \(M_4\)，其终端模式与唯一全局最优正规形相容。
因此该候选同时满足局部命中、预算内极小与终端相容三项条件，
可签发为 \(\mathsf{TunHit}(q)\) 的局部执行实例。
命题得证。
\end{proof}

\begin{corollary}[PDE 求解到结构性命中的压缩式]
\label{cor:tex1-ch5-pde-to-hit}
在上述对象域与证书条件下，正确的压缩链为
\[
  \boxed{
  \text{PDE 求解}
  \Longrightarrow
  \text{同伦扭曲障碍空间}
  \Longrightarrow
  \text{统计力学采样}
  \Longrightarrow
  \text{微观隧穿命中}.}
\]
因此，本章所说的求解升级不是“更快算原题”，而是
\[
  \boxed{
  \text{把技术性求解问题提升为结构性命中问题}.}
\]
\end{corollary}

\begin{Certificate}[证书：四段压缩链]
\label{cert:tex1-ch5-pde-to-hit-cor}
证书登记为
\[
\begin{aligned}
  D_1 &: \text{PDE 表象由 \(\Gamma\)-纤维等价重写为障碍空间};\\
  D_2 &: \text{障碍空间由同伦扭曲候选与命中判据局部化};\\
  D_3 &: \text{候选比较由 Gibbs/变分采样转为命中概率集中};\\
  D_4 &: \text{命中候选由离盆像、降障碍、终端相容签发为微观隧穿命中}.
\end{aligned}
\]
\end{Certificate}

\begin{proof}[证明（基于数学符号推演逻辑的谓词因果链）]
\(D_1\) 由定义~\ref{def:tex1-gamma-equivalent} 与定理~\ref{thm:tex1-gamma-invariant} 给出；
\(D_2\) 由定义~\ref{def:tex1-ch5-hit} 给出；
\(D_3\) 由引理~\ref{lem:tex1-ch5-stat-hit} 给出；
\(D_4\) 由命题~\ref{prop:tex1-ch5-tunnel-hit-local} 给出。
四段合取即得压缩链，推论成立。
\end{proof}

\begin{remark}[边界：结构命中不是无条件 PDE 解算器]
这里的结论只签发对象层改写：
\[
  \text{技术性 PDE 表象}
  \quad\leadsto\quad
  \text{同伦扭曲障碍空间中的结构性命中问题}.
\]
它不声称任意偏微分方程都可被无条件求解，
也不声称统计力学采样自动给出全局正确答案。
严格可签发的是：当 \(\Gamma\)-等价、预算候选族、命中候选非空、严格能量优势与终端相容证书同时闭合时，
原问题的可判核心可以由微观隧穿命中路径承担。
\end{remark}

\subsection{公理（降级为定理）：先变分收敛，后命中支配的两层求解链}
\label{subsec:tex1-ch5-axiom-downgraded}

\begin{theorem}[公理（降级为定理）：先变分收敛，后命中支配的两层求解链]
\label{thm:tex1-ch5-axiom-downgraded}
设对每个层级 \(n\le N_\ast\) 与每个可终端到达对象 \(\tau\in\mathcal T_n\)，
记从 \(\tau\) 出发到 \(\tau^\ast\) 的合法后缀路径族为
\[
  \mathfrak L_n(\tau,\tau^\ast)
  \subseteq
  \mathfrak L^{\mathrm{loc}\to\tau^\ast},
\]
并定义 Bellman 型尾值函数
\[
  V_{N_\ast}(\tau)
  :=
  \begin{cases}
    0, & \tau=\tau^\ast,\\
    +\infty, & \tau\neq \tau^\ast,
  \end{cases}
\]
\[
  V_n(\tau)
  :=
  \min_{\eta\in \mathrm{Lift}_n(\tau)}
  \bigl(c_n(\tau,\eta)+V_{n+1}(\eta)\bigr).
\]
若满足：
\begin{enumerate}
  \item 合法提升族有限且终端可达；
  \item 终端唯一全局最优正规形 \(\tau^\ast\) 存在；
  \item 盆内刚性成立；
  \item 命中步存在性成立；
  \item 严格隧穿优势成立；
\end{enumerate}
则对任意 \(n\le N_\ast\) 与任意可终端到达的 \(\tau\in\mathcal T_n\)，同时有：
\begin{enumerate}
  \item \(\mathfrak L_n(\tau,\tau^\ast)\neq\varnothing\)；
  \item 存在唯一尾路径
  \[
    \ell_n^\ast(\tau)\in\mathfrak L_n(\tau,\tau^\ast)
  \]
  使
  \[
    \Psi\bigl(\ell_n^\ast(\tau)\bigr)
    =
    \min_{\ell\in\mathfrak L_n(\tau,\tau^\ast)}
    \Psi(\ell) ;
  \]
  \item 只要 \(\nu_n(\tau)>0\)，则 \(\ell_n^\ast(\tau)\) 的下一步必须取为同伦扭曲命中步。
\end{enumerate}
因此，在抽象对象层上，
\[
  \text{先变分收敛到合法最优后缀}
  \Longrightarrow
  \text{再由命中步支配局部更新}
\]
不再只是口号，而是条件闭合后的定理链。
\end{theorem}

\begin{Certificate}[数学归纳法证书：定理~\ref{thm:tex1-ch5-axiom-downgraded}]
\label{cert:tex1-ch5-axiom-downgraded}
对 \(n\) 作从 \(N_\ast\) 到起始层的反向归纳。
定义三个谓词：
\[
\begin{aligned}
  P_n(\tau) &: \mathfrak L_n(\tau,\tau^\ast)\neq\varnothing,\\
  Q_n(\tau) &: \exists!\,\ell_n^\ast(\tau)\in\mathfrak L_n(\tau,\tau^\ast)
              \ \text{使}\
              \Psi(\ell_n^\ast(\tau))
              =
              \min_{\ell\in\mathfrak L_n(\tau,\tau^\ast)}\Psi(\ell),\\
  R_n(\tau) &: \nu_n(\tau)>0
              \Rightarrow
              \text{最优尾路径的首步是同伦扭曲命中步}.
\end{aligned}
\]
归纳证书分三段：
\begin{enumerate}
  \item \textbf{基步：} 在终端层 \(n=N_\ast\) 上，\(\tau^\ast\) 的唯一后缀路径是单点路径 \((\tau^\ast)\)，故 \(P_{N_\ast},Q_{N_\ast}\) 成立；
  \item \textbf{归纳步：} 若对所有下一层对象 \(\eta\in\mathcal T_{n+1}\) 都有 \(P_{n+1}(\eta)\wedge Q_{n+1}(\eta)\wedge R_{n+1}(\eta)\)，

  则由 Bellman 递推与有限可达性推出 \(P_n(\tau)\wedge Q_n(\tau)\)；
  \item \textbf{命中支配：} 当 \(\nu_n(\tau)>0\) 时，由命中步存在性 + 严格隧穿优势，
  非命中步不可能成为 Bellman 极小元的首步，故 \(R_n(\tau)\) 成立。
\end{enumerate}
\end{Certificate}

\begin{proof}[证明（基于数学符号推演逻辑的谓词因果链）]
先作基步。对 \(n=N_\ast\)，若 \(\tau=\tau^\ast\)，
则唯一合法后缀路径就是
\[
  (\tau^\ast)\in \mathfrak L_{N_\ast}(\tau^\ast,\tau^\ast),
\]
且其严格化路径作用量唯一确定，
故 \(P_{N_\ast}(\tau^\ast)\) 与 \(Q_{N_\ast}(\tau^\ast)\) 同时成立。
此时已无更深层提升，\(R_{N_\ast}\) 真值空置。

再作归纳步。固定某个 \(n<N_\ast\)，并假设对全部下一层对象 \(\eta\) 都已有
\[
  P_{n+1}(\eta)\wedge Q_{n+1}(\eta)\wedge R_{n+1}(\eta).
\]
由于合法提升族有限且终端可达，
集合
\[
  \mathrm{Lift}_n(\tau)
\]
非空且可比较。
于是 Bellman 递推
\[
  V_n(\tau)
  =
  \min_{\eta\in \mathrm{Lift}_n(\tau)}
  \bigl(c_n(\tau,\eta)+V_{n+1}(\eta)\bigr)
\]
给出至少一个后继极小元 \(\eta^\ast\)。
由归纳假设中的 \(Q_{n+1}(\eta^\ast)\)，
从 \(\eta^\ast\) 到 \(\tau^\ast\) 存在唯一最优后缀路径；
把首步 \(\tau\rightsquigarrow\eta^\ast\) 接到该唯一后缀路径前面，
就得到 \(\mathfrak L_n(\tau,\tau^\ast)\) 的一个最优元，
故 \(P_n(\tau)\) 与 \(Q_n(\tau)\) 成立。

最后处理命中支配。
若 \(\nu_n(\tau)>0\)，
则由命中步存在性，存在至少一个合法后继 \(\eta_{\mathrm{hit}}\in \mathrm{Lift}_n(\tau)\)
满足定义~\ref{def:tex1-ch5-hit}；
又由严格隧穿优势，任意不离开旧盆像或不严格降低活跃障碍计数的后继 \(\eta_{\mathrm{stay}}\)
都满足
\[
  c_n(\tau,\eta_{\mathrm{hit}})+V_{n+1}(\eta_{\mathrm{hit}})
  <
  c_n(\tau,\eta_{\mathrm{stay}})+V_{n+1}(\eta_{\mathrm{stay}}).
\]
因此 Bellman 极小元的首步不可能是非命中步，
只能取为同伦扭曲命中步。
故 \(R_n(\tau)\) 成立。

由反向归纳法，\(P_n,Q_n,R_n\) 对一切 \(n\le N_\ast\) 同时成立。
于是定理得证。
\end{proof}

\begin{corollary}[两层求解链的严格压缩句]
\label{cor:tex1-ch5-axiom-downgraded}
在定理~\ref{thm:tex1-ch5-axiom-downgraded} 的条件下，
钱学森弹道/钱学森障碍问题在对象层上的求解链可压缩为
\[
  \begin{aligned}
    &\text{合法提升族上的变分极小化}\\
    \Longrightarrow\ &
    \text{Bellman 递推锁定唯一最优后缀}\\
    \Longrightarrow\ &
    \text{非零障碍阶段由命中步支配局部更新}.
  \end{aligned}
\]
\end{corollary}

\begin{Certificate}[证书：推论~\ref{cor:tex1-ch5-axiom-downgraded} 的直接压缩]
\label{cert:tex1-ch5-axiom-downgraded-cor}
证书登记谓词链：
\[
\begin{aligned}
  C^{(5)}_1 &: \text{调用定理~\ref{thm:tex1-ch5-axiom-downgraded} 的结论 (1)--(3)},\\
  C^{(5)}_2 &: C^{(5)}_1
              \Rightarrow
              \text{极小化对象、唯一后缀与命中支配已连成闭链},\\
  C^{(5)}_3 &: C^{(5)}_2
              \Rightarrow
              \text{得到三段式严格压缩句}.
\end{aligned}
\]
\end{Certificate}

\begin{proof}[证明（基于数学符号推演逻辑的谓词因果链）]
由 \(C^{(5)}_1\) 可知，
变分极小化对象、唯一最优后缀与命中步必要性已经在同一层级链上闭合。
这正是 \(C^{(5)}_2\)。
把三段结论直接合并，即得 \(C^{(5)}_3\) 所表述的压缩句。推论得证。
\end{proof}

\subsection{引理与主定理：开放修复问题向命中约束变分求解的压缩}
\label{subsec:tex1-ch5-main}

\begin{lemma}[高阶可消去与双重收敛把开放修复问题压成有限合法提升族]
\label{lem:tex1-ch5-finite-family}
若存在某个高阶层 \(m>n\) 使
\[
  \Pi_n^m(\mathsf{Ob}_n)=0,
\]
并且存在投影曲线序列
\[
  \bar{\gamma}_{n,r}\to \gamma_n^\ast,
\qquad
  \mathcal S_n(\bar{\gamma}_{n,r})\to \inf \mathcal S_n,
\qquad
  \mathcal L_n(\bar{\gamma}_{n,r})\to \inf \mathcal L_n,
\]
再结合极限层可行路径集非空、作用量下半连续且强制有界下，
以及终端层有限终止与终端唯一全局最优正规形成立，
则原来的开放障碍修复问题可在对象层上压缩为一个非空有限路径族
\[
  \mathfrak L^{\mathrm{loc}\to\tau^\ast},
\]
并且严格化路径作用量 \(\Psi\) 在该有限集上存在极小元。
\end{lemma}

\begin{Certificate}[证书：高阶修复信道 + 双重收敛 + 直接法极小化 + 有限终止]
\label{cert:tex1-ch5-finite-family}
证书分三段：
\begin{enumerate}
  \item 由 \cite{RepoTex36HomotopyTwistEquivalence}，
  \[
    \Pi_n^m(\mathsf{Ob}_n)=0
    \Rightarrow
    \text{低阶障碍在高阶层成为可吸收边界项};
  \]
  \item 由 \cite{RepoTex36HomotopyTwistEquivalence,RepoTex41HomotopyRepairBundle}，
  投影曲线双重收敛 + 下半连续 + 强制有界下
  \[
    \Rightarrow
    \text{可行变分极小元存在};
  \]
  \item 由 \cite{RepoTex38JacobiBianchiRepair}，
  有限终止 + 终端唯一全局最优正规形
  \[
    \Rightarrow
    \mathfrak L^{\mathrm{loc}\to\tau^\ast}
    \text{ 非空且有限}.
  \]
\end{enumerate}
\end{Certificate}

\begin{proof}[证明（基于数学符号推演逻辑的谓词因果链）]
先由第一段，
\[
  \Pi_n^m(\mathsf{Ob}_n)=0
\]
说明原始低层障碍不是绝对断裂，
而是可被高阶边界项吸收的可修复障碍。
于是开放修复问题不必再在原层对象上作无结构硬求极值，
而可以被提升到“可修复的扩展候选族”中比较。

再由第二段，
\[
  \bar{\gamma}_{n,r}\to \gamma_n^\ast,
\qquad
  \mathcal S_n(\bar{\gamma}_{n,r})\to \inf \mathcal S_n,
\qquad
  \mathcal L_n(\bar{\gamma}_{n,r})\to \inf \mathcal L_n
\]
把“高阶修复后可投影回低层”的存在性与最优性同时固定下来。
结合极限层可行域非空、作用量下半连续和强制有界下，
可由直接法得到可行变分极小元。

最后由第三段，终端层有限终止意味着候选对象域最终收敛到有限终端对象域，
而终端唯一全局最优正规形 \(\tau^\ast\) 的存在，
把原先开放式的修复问题压缩为指向 \(\tau^\ast\) 的有限合法提升族
\[
  \mathfrak L^{\mathrm{loc}\to\tau^\ast}.
\]
有限集上的实值函数 \(\Psi\) 必有极小元。
故引理成立。
\end{proof}

\begin{theorem}[条件性定理：障碍修复可收敛为命中约束下的变分求解]
\label{thm:tex1-ch5-conditional}
若满足以下两组条件：
\begin{enumerate}
  \item \textbf{高阶修复与双重收敛条件：}
  存在层间传递
  \[
    \Pi_n^m(\mathsf{Ob}_n)=0,
  \]
  并且投影曲线满足
  \[
    \bar{\gamma}_{n,r}\to \gamma_n^\ast,
  \]
  \[
    \mathcal S_n(\bar{\gamma}_{n,r})\to \inf \mathcal S_n,
    \qquad
    \mathcal L_n(\bar{\gamma}_{n,r})\to \inf \mathcal L_n,
  \]
  同时极限层可行域非空、作用量下半连续且强制有界下；
  \item \textbf{离散修复塔路径条件：}
  A. 有限性与可达性；
  B. 终端唯一全局最优正规形 \(\tau^\ast\)；
  C. 盆内刚性；
  D. 命中步存在性；
  E. 严格隧穿优势。
\end{enumerate}
则钱学森弹道/钱学森障碍问题在对象层上可被严格改写为
\[
  \min_{\ell\in \mathfrak L^{\mathrm{loc}\to\tau^\ast}} \Psi(\ell),
\]
并且存在唯一极小路径
\[
  \ell^\ast\in \mathfrak L^{\mathrm{loc}\to\tau^\ast}
\]
使
\[
  \Psi(\ell^\ast)
  =
  \min_{\ell\in \mathfrak L^{\mathrm{loc}\to\tau^\ast}}
  \Psi(\ell),
  \qquad
  \operatorname{end}(\ell^\ast)=\tau^\ast.
\]
只要沿途仍存在非零提升障碍，
则最优路径 \(\ell^\ast\) 的下一步必须取为同伦扭曲命中步。
换言之，
\[
  \text{障碍修复问题}
  \Longrightarrow
  \text{命中约束下的变分极小化问题}
\]
在上述条件全部闭合时成立。
\end{theorem}

\begin{Certificate}[证书：条件性定理~\ref{thm:tex1-ch5-conditional} 的两段式证书]
\label{cert:tex1-ch5-conditional}
证书分两段：
\begin{enumerate}
  \item 由 \cite{RepoTex36HomotopyTwistEquivalence,RepoTex41HomotopyRepairBundle}，
  \[
    \Pi_n^m(\mathsf{Ob}_n)=0
    \Rightarrow
    \text{高阶修复信道存在},
  \]
  且
  \[
    \text{双重收敛}
    +
    \text{下半连续}
    +
    \text{强制有界下}
    \Rightarrow
    \text{变分极小元存在};
  \]
  \item 由 \cite{RepoTex38JacobiBianchiRepair,RepoTex39OpenRepairFamily}，
  \[
    \begin{aligned}
      &\text{有限终止}
      +
      \tau^\ast
      +
      \text{Bellman 递推}\\
      &\qquad
      +
      \text{命中步存在性}
      +
      \text{严格隧穿优势}\\
      &\Rightarrow
      \text{唯一最优提升路径与命中步必要性}.
    \end{aligned}
  \]
\end{enumerate}
\end{Certificate}

\begin{proof}[证明（基于数学符号推演逻辑的谓词因果链）]
由证书第一段和引理~\ref{lem:tex1-ch5-finite-family}，
原始开放修复问题已被压缩到一个非空有限候选族
\[
  \mathfrak L^{\mathrm{loc}\to\tau^\ast},
\]
并且 \(\Psi\) 在该族上存在极小元。
这一步把“开放修复”收紧为“有限合法提升族上的变分问题”。

再由证书第二段与定理~\ref{thm:tex1-ch5-axiom-downgraded}，
对该有限合法提升族应用 Bellman 型递推，
便得到唯一的最优后缀拼接规则；
终端唯一全局最优正规形 \(\tau^\ast\) 使最优路径的终点唯一固定，
故存在唯一极小路径
\[
  \ell^\ast
  =
  \operatorname*{arg\,min}_{\ell\in \mathfrak L^{\mathrm{loc}\to\tau^\ast}}
  \Psi(\ell),
  \qquad
  \operatorname{end}(\ell^\ast)=\tau^\ast.
\]

最后，若 \(\ell^\ast\) 的某一阶段仍满足非零提升障碍，
则由命中步存在性与严格隧穿优势，
该阶段的 Bellman 极小首步只能是同伦扭曲命中步；
否则就会落入“留在旧盆像内或不严格降障碍”的次优更新，与最优性矛盾。
因此，最优路径在所有非零障碍阶段都由命中步支配局部更新。

故本章所要的最严格改写成立：
问题先收敛为障碍修复的变分求解问题，
再由同伦扭曲命中给出变分可行域中的关键支配步。
定理得证。
\end{proof}

\begin{corollary}[最严格压缩式]
\label{cor:tex1-ch5-conditional}
\begin{center}
  \fbox{\parbox{0.92\linewidth}{
    \centering
    问题可收敛为“从局部最优到终端全局最优正规形”的障碍修复变分问题；\\
    同伦扭曲命中不是附会描述，而是变分可行域中的关键支配步；\\
    终端全局最优正规形 $\tau^\ast$ 由唯一最优提升路径 $\ell^\ast$ 实现。
  }}
\end{center}
\end{corollary}

\begin{Certificate}[证书：推论~\ref{cor:tex1-ch5-conditional} 的主定理压缩]
\label{cert:tex1-ch5-conditional-cor}
证书登记谓词链：
\[
\begin{aligned}
  C^{(5)}_4 &: \text{调用定理~\ref{thm:tex1-ch5-conditional}},\\
  C^{(5)}_5 &: C^{(5)}_4
              \Rightarrow
              \text{变分问题、命中支配步与唯一最优路径同时成立},\\
  C^{(5)}_6 &: C^{(5)}_5
              \Rightarrow
              \text{得到三句式严格压缩表达}.
\end{aligned}
\]
\end{Certificate}

\begin{proof}[证明（基于数学符号推演逻辑的谓词因果链）]
由 \(C^{(5)}_4\) 得，
主定理已经同时给出三件事：
变分求解对象、命中步的局部必要性，以及唯一最优提升路径 \(\ell^\ast\)。
这正是 \(C^{(5)}_5\)。
把三项结论按对象层次排成一句，即得 \(C^{(5)}_6\) 所述压缩式。推论得证。
\end{proof}

\subsection{命题与推论：命中作为 Bellman 递推中的局部支配步}
\label{subsec:tex1-ch5-proposition}

\begin{proposition}[同伦扭曲命中不是额外装饰，而是 Bellman 递推中的局部支配步]
\label{prop:tex1-ch5-hit}
在定理~\ref{thm:tex1-ch5-conditional} 的全部条件下，
对任意阶段 \(n<N_\ast\) 与任意满足 \(\nu_n(\tau)>0\) 的对象 \(\tau\in\mathcal T_n\)，
定义局部 Bellman 更新集
\[
  \mathcal U_n(\tau)
  :=
  \operatorname*{arg\,min}_{\eta\in\mathrm{Lift}_n(\tau)}
  \bigl(c_n(\tau,\eta)+V_{n+1}(\eta)\bigr).
\]
则任意
\[
  \eta^\ast\in \mathcal U_n(\tau)
\]
都满足定义~\ref{def:tex1-ch5-hit}，
即
\[
  \eta^\ast\notin \iota_n^{\mathcal T}\bigl(\mathfrak B_n(\tau)\bigr),
  \qquad
  \nu_{n+1}(\eta^\ast)\le \nu_n(\tau)-1.
\]
因此，
\[
  \text{命中}
  =
  \text{局部变分更新规则的一部分},
\]
而不是附加在路径积分语句外部的一层修辞。
\end{proposition}

\begin{Certificate}[证书：Bellman 极小元的命中判据]
\label{cert:tex1-ch5-hit}
证书分三段：
\begin{enumerate}
  \item 定理~\ref{thm:tex1-ch5-conditional} 已给出：非零障碍阶段的最优路径首步必须是命中步；
  \item \(\mathcal U_n(\tau)\) 正是该阶段 Bellman 极小首步的集合；
  \item 因而 \(\mathcal U_n(\tau)\) 中的任一元素都必须满足命中判据。
\end{enumerate}
\end{Certificate}

\begin{proof}[证明（基于数学符号推演逻辑的谓词因果链）]
由证书第一段，
当 \(\nu_n(\tau)>0\) 时，
从 \(\tau\) 出发的任何最优路径都必须以命中步开头。
而由证书第二段，
\[
  \mathcal U_n(\tau)
\]
恰好枚举了该阶段全部 Bellman 极小首步。
若其中存在某个
\[
  \eta^\sharp\in\mathcal U_n(\tau)
\]
不是命中步，
则以 \(\eta^\sharp\) 为首步的最优后缀将违背主定理的“非零障碍阶段必须命中”结论。
故证书第三段成立，
\(\mathcal U_n(\tau)\) 的全部元素都必须满足命中判据。
命题得证。
\end{proof}

\begin{corollary}[命中不是额外装饰，而是局部求解机制]
\label{cor:tex1-ch5-hit}
因此，“基于同伦扭曲的命中进行求解”的精确含义不是
\[
  \text{随意加一个命中修辞},
\]
而是
\[
  \begin{aligned}
    &\text{在每个仍有非零障碍的层级，}\\
    &\text{把离开旧盆像且严格降低活跃障碍计数的提升步，}\\
    &\text{作为 Bellman 递推中的局部支配步。}
  \end{aligned}
\]
\end{corollary}

\begin{Certificate}[证书：命题~\ref{prop:tex1-ch5-hit} 的直接语义化]
\label{cert:tex1-ch5-hit-cor}
证书登记谓词链：
\[
\begin{aligned}
  C^{(5)}_7 &: \text{调用命题~\ref{prop:tex1-ch5-hit}},\\
  C^{(5)}_8 &: C^{(5)}_7
              \Rightarrow
              \text{命中步已进入局部 Bellman 更新规则},\\
  C^{(5)}_9 &: C^{(5)}_8
              \Rightarrow
              \text{命中是局部求解机制，而非附加修辞}.
\end{aligned}
\]
\end{Certificate}

\begin{proof}[证明（基于数学符号推演逻辑的谓词因果链）]
由 \(C^{(5)}_7\) 得，
命中步已经被证明为 Bellman 递推的局部支配步。
于是 \(C^{(5)}_8\) 成立。
一旦命中进入局部更新规则，
它就不再是路径外部的附着描述，
而是求解过程的组成部分。
这正是 \(C^{(5)}_9\)。推论得证。
\end{proof}

\subsection{备注：强弱关系、条件边界与仓内分工}
\label{subsec:tex1-ch5-remarks}

\begin{remark}[这一步已经比“只说路径积分选最优路径”更强]
若只说
\[
  \text{路径积分会选最优路径},
\]
那么我们只得到了终端比较原则，
却没有得到“中间层到底如何走”的局部机制。
而本章把结构收紧为
\[
  \text{高层障碍消去}
  \Rightarrow
  \text{终端可达}
  \Rightarrow
  \text{路径作用量最小化}
  \Rightarrow
  \text{命中步主导局部更新},
\]
因此已经从“终端选优”推进到“层间求解机制”。
\end{remark}

\begin{remark}[严格边界：这里只能写成条件性定理]
本章最需要守住的边界是：
\[
  \text{“可基于命中进行求解”}
\]
只能写成条件性定理，而不能直接写成无条件定理。
因为其中最强的接口恰好是
\[
  \text{命中步存在性},
  \qquad
  \text{严格隧穿优势}.
\]
若这两条没有单独闭合，
则当前最多只能严格说到
\[
  \text{问题可被改写为障碍修复的变分框架},
\]
而不能进一步断言
\[
  \text{命中机制必然给出求解}.
\]
\end{remark}

\begin{remark}[与仓内文本的对应关系]
若把仓内文本的角色压成一句话，
则 \cite{RepoTex36HomotopyTwistEquivalence,RepoTex41HomotopyRepairBundle}
负责证明
\[
  \text{障碍可修复}
  \quad+\quad
  \text{存在性与最优性双重收敛},
\]
而 \cite{RepoTex38JacobiBianchiRepair,RepoTex39OpenRepairFamily}
负责证明
\[
  \text{有限终止}
  \quad+\quad
  \text{终端唯一最优正规形}
  \quad+\quad
  \text{唯一最优提升路径}
  \quad+\quad
  \text{命中步的局部必要性}.
\]
两者拼接之后，
才得到本章关于“障碍修复变分求解 + 同伦扭曲命中支配步”的严格版本。
\end{remark}

%%%%%%%%%%%%%%%%%%%%%%%%%%%%%%%%%%%%%%%%%%%%%%%%%%%%%%%%%%%%
\clearpage
\section{形式逻辑：轨迹商空间、编码拓扑升级与严格经典同伦等价}
\label{sec:tex1-ch7}

\begin{remark}[本章意义备注：高超音速攻防的抽象语境]
本章的意义，在于把对轨迹外形的关注提升到轨迹等价类与编码拓扑层。
对钱学森弹道研究而言，它提示真正稳定的对象不是单条曲线，
而是经观测与编码后仍保持不变的拓扑/同伦类。
\end{remark}

\subsection{章节目标与证据来源}
\label{subsec:tex1-ch7-goal}

\begin{remark}[本章目标]
本章把“严格同伦等价”分成两层来建立。
第一层，是把仓内 tex36/tex38 已经给出的
\[
  \text{路径积分支撑类}
  \Longleftrightarrow
  \text{同伦扭曲正规形}
  \Longleftrightarrow
  \text{稳定修复序列}
\]
之编码双射链，升级为经典拓扑意义下的同胚链与同伦等价链。
第二层，是把该结论接到
\[
  \text{钱学森弹道（偏导控制解空间）}
  \Longrightarrow
  \text{钱学森障碍（拓扑约束/修复空间）}
\]
这一对象重写上，并明确说明：严格成立的对象不是原始 PDE 解空间本身，
而是其按路径积分支撑类压缩后的商空间。
\end{remark}

\begin{remark}[仓内与外部依据]
本章的直接仓内依据是
\cite{RepoTex36HomotopyTwistEquivalence,RepoTex38JacobiBianchiRepair}。
其中，
\cite{RepoTex36HomotopyTwistEquivalence}
已给出
\[
  \mathcal F_A \cong_{\mathrm{set}} \mathcal{PI}_A \cong_{\mathrm{set}} \mathcal T_A
\]
的显式双射链，
并把 \(\simeq_{\mathrm{enc\text{-}h}}\) 明确界定为“存在互逆编码映射”的编码同伦等价，
同时说明：若再赋予相容拓扑，则可升级为经典拓扑同伦等价。
\cite{RepoTex36HomotopyTwistEquivalence}
还给出
\[
  \mathcal T_A \cong_{\mathrm{set}} \mathcal E_A^{\mathrm{rep}}
\]
的稳定修复序列双射；
\cite{RepoTex38JacobiBianchiRepair}
则给出终端零缺陷正规形、唯一最优路径支撑类与终端代表的严格证书。
外部拓扑背景可参见 \cite{Hatcher2002AT,Husemoller1994FibreBundles}。
\end{remark}

\subsection{定义：对象域、轨迹商空间与相容拓扑}
\label{subsec:tex1-ch7-defs}

\begin{definition}[对象域]
\label{def:tex1-ch7-objects}
固定一个层或一个终端层级 \(A\)。
定义以下对象：
\[
  \mathcal S_A^{\mathrm{adm}}
\]
为可容许控制/PDE 解轨道空间；
\[
  \mathcal{PI}_A
\]
为路径积分支撑类空间；
\[
  \mathcal T_A
\]
为同伦扭曲正规形空间；
\[
  \mathcal F_A
\]
为纤维终点编码空间；
\[
  \mathcal E_A^{\mathrm{rep}}
\]
为稳定修复序列空间。

在仓内 tex36 的口径下，已存在显式双射或逆映射限制
\[
  \overline{\mathrm{PT}}_A^{p_{\mathrm{src}}}:\mathcal{PI}_A\to \mathcal F_A,
  \qquad
  \operatorname{NF}_A:\mathcal{PI}_A\to \mathcal T_A,
  \qquad
  \operatorname{Stab}_A:\mathcal E_A^{\mathrm{rep}}\to \mathcal T_A,
\]
以及
\[
  q_A|_{\mathcal T_A}:\mathcal T_A\to\mathcal{PI}_A
\]
作为 \(\operatorname{NF}_A\) 的互逆映射。
\end{definition}

\begin{definition}[轨迹侧的等价关系与商空间]
\label{def:tex1-ch7-quotient}
定义商映射
\[
  Q_A:\mathcal S_A^{\mathrm{adm}}\to \mathcal{PI}_A,
\]
它把每条可容许解轨道送到其路径积分支撑类。

定义等价关系
\[
  \gamma_1\sim_A\gamma_2
  \quad:\Longleftrightarrow\quad
  Q_A(\gamma_1)=Q_A(\gamma_2).
\]

定义商空间
\[
  \widetilde{\mathcal S}_A
  :=
  \mathcal S_A^{\mathrm{adm}}/\sim_A,
\]
并记商映射为
\[
  \pi_A:\mathcal S_A^{\mathrm{adm}}\to \widetilde{\mathcal S}_A.
\]

在 \(\widetilde{\mathcal S}_A\) 上取由 \(\pi_A\) 诱导的商拓扑。
由商空间的泛性质，存在唯一映射
\[
  \overline Q_A:\widetilde{\mathcal S}_A\to \mathcal{PI}_A
\]
使得
\[
  Q_A=\overline Q_A\circ \pi_A.
\]
\end{definition}

\begin{definition}[相容拓扑]
\label{def:tex1-ch7-topology}
为了把 tex36 的“编码同伦等价”升级为经典拓扑同伦等价，定义如下相容拓扑。

在 \(\mathcal{PI}_A\) 上取由 \(Q_A\) 诱导的商拓扑，即
\[
  U\subseteq \mathcal{PI}_A \text{ 开}
  \quad:\Longleftrightarrow\quad
  Q_A^{-1}(U)\subseteq \mathcal S_A^{\mathrm{adm}} \text{ 开}.
\]

在 \(\mathcal F_A\) 上取由 \(\overline{\mathrm{PT}}_A^{p_{\mathrm{src}}}\) 传递的拓扑，即
\[
  V\subseteq \mathcal F_A \text{ 开}
  \quad:\Longleftrightarrow\quad
  \bigl(\overline{\mathrm{PT}}_A^{p_{\mathrm{src}}}\bigr)^{-1}(V)\subseteq \mathcal{PI}_A \text{ 开}.
\]

在 \(\mathcal T_A\) 上取由 \(\operatorname{NF}_A\) 传递的拓扑，即
\[
  W\subseteq \mathcal T_A \text{ 开}
  \quad:\Longleftrightarrow\quad
  \operatorname{NF}_A^{-1}(W)\subseteq \mathcal{PI}_A \text{ 开}.
\]

在 \(\mathcal E_A^{\mathrm{rep}}\) 上取由 \(\operatorname{Stab}_A\) 传递的拓扑，即
\[
  R\subseteq \mathcal E_A^{\mathrm{rep}} \text{ 开}
  \quad:\Longleftrightarrow\quad
  \operatorname{Stab}_A(R)\subseteq \mathcal T_A \text{ 开}.
\]
\end{definition}

\subsection{公理（降级为定理）：拓扑升级链可按对象层前缀递归闭合}
\label{subsec:tex1-ch7-axiom}

\begin{theorem}[公理（降级为定理）：五段拓扑升级链可按前缀递归闭合]
\label{thm:tex1-ch7-axiom-downgraded}
定义五段升级见证：
\[
\begin{aligned}
  E_0 &: \mathcal S_A^{\mathrm{adm}}
         \Longrightarrow
         \widetilde{\mathcal S}_A
         \xrightarrow{\overline Q_A}
         \mathcal{PI}_A,\\
  E_1 &: \widetilde{\mathcal S}_A\cong_{\mathrm{Top}}\mathcal{PI}_A,\\
  E_2 &: \mathcal{PI}_A\cong_{\mathrm{Top}}\mathcal T_A,\\
  E_3 &: \mathcal T_A\cong_{\mathrm{Top}}\mathcal E_A^{\mathrm{rep}},\\
  E_4 &: \mathcal{PI}_A\cong_{\mathrm{Top}}\mathcal F_A
         \Longrightarrow
         \widetilde{\mathcal S}_A,\ \mathcal F_A,\ \mathcal{PI}_A,\ \mathcal T_A,\ \mathcal E_A^{\mathrm{rep}}
         \text{ 落在同一经典同伦型里。}
\end{aligned}
\]
对每个 \(k\in\{0,1,2,3,4\}\)，定义谓词
\[
  P(k)
  :\Longleftrightarrow
  \text{前缀见证 }E_0,\dots,E_k\text{ 已被文内结果闭合。}
\]
则对全部 \(k\in\{0,1,2,3,4\}\)，命题 \(P(k)\) 成立。
因此，从轨迹商空间到路径积分支撑类、正规形、稳定修复序列与纤维终点编码的拓扑升级链，
可以按对象层前缀递归闭合。
\end{theorem}

\begin{Certificate}[数学归纳法证书：按前缀长度 \(k\) 的闭合]
\label{cert:tex1-ch7-axiom-downgraded}
归纳谓词为
\[
  P(k):
  \text{前缀 }E_0,\dots,E_k\text{ 已闭合。}
\]
证书分两步：
\begin{enumerate}
  \item \textbf{基础步 \(P(0)\)：}
  定义~\ref{def:tex1-ch7-quotient} 与定义~\ref{def:tex1-ch7-topology}
  已经把 \(\widetilde{\mathcal S}_A\)、\(\mathcal{PI}_A\) 及 \(\overline Q_A\) 的对象层接口固定完毕；
  \item \textbf{归纳步 \(P(k)\Rightarrow P(k+1)\)：}
  对每个 \(k<4\)，后继见证分别由引理~\ref{lem:tex1-ch7-qbar}、
  定理~\ref{thm:tex1-ch7-classical} 与推论~\ref{cor:tex1-ch7-fiber}
  提供。
\end{enumerate}
\end{Certificate}

\begin{proof}[证明（数学归纳法证书；基于数学符号推演逻辑的谓词因果链）]
对前缀长度 \(k\) 作数学归纳。

\textbf{基础步：}
当 \(k=0\) 时，
定义~\ref{def:tex1-ch7-quotient} 已给出
\[
  \mathcal S_A^{\mathrm{adm}}
  \Longrightarrow
  \widetilde{\mathcal S}_A
  \xrightarrow{\overline Q_A}
  \mathcal{PI}_A,
\]
同时定义~\ref{def:tex1-ch7-topology} 已固定相关拓扑。
故 \(P(0)\) 成立。

\textbf{归纳步：}
假设 \(P(k)\) 成立。
若 \(k=0\)，则由引理~\ref{lem:tex1-ch7-qbar}，
\[
  \widetilde{\mathcal S}_A\cong_{\mathrm{Top}}\mathcal{PI}_A,
\]
故 \(E_1\) 成立，从而 \(P(1)\) 成立。

若 \(k=1\)，则由 tex36 中 \(\operatorname{NF}_A\) 与 \(q_A|_{\mathcal T_A}\) 的互逆双射链，
再结合定义~\ref{def:tex1-ch7-topology} 的传递拓扑，
可得
\[
  \mathcal{PI}_A\cong_{\mathrm{Top}}\mathcal T_A,
\]
故 \(E_2\) 成立，从而 \(P(2)\) 成立。

若 \(k=2\)，则由 tex36 中
\[
  \operatorname{Stab}_A:\mathcal E_A^{\mathrm{rep}}\to \mathcal T_A
\]
的双射及定义~\ref{def:tex1-ch7-topology} 的传递拓扑，
可得
\[
  \mathcal T_A\cong_{\mathrm{Top}}\mathcal E_A^{\mathrm{rep}},
\]
故 \(E_3\) 成立，从而 \(P(3)\) 成立。

若 \(k=3\)，则由 tex36 中
\[
  \overline{\mathrm{PT}}_A^{p_{\mathrm{src}}}:\mathcal{PI}_A\to\mathcal F_A
\]
的双射链与定义~\ref{def:tex1-ch7-topology} 的传递拓扑，
可得
\[
  \mathcal{PI}_A\cong_{\mathrm{Top}}\mathcal F_A.
\]
再由前面已闭合的同胚链，五个对象都落入同一经典同伦型。
故 \(E_4\) 成立，从而 \(P(4)\) 成立。

因此对全部 \(k\in\{0,1,2,3,4\}\)，都有 \(P(k)\) 成立。
定理得证。
\end{proof}

\subsection{引理：轨迹商空间与路径积分支撑类空间同胚}
\label{subsec:tex1-ch7-lemma}

\begin{lemma}[商空间与路径积分支撑类空间同胚]
\label{lem:tex1-ch7-qbar}
\[
  \overline Q_A:\widetilde{\mathcal S}_A\to \mathcal{PI}_A
\]
是同胚。
\end{lemma}

\begin{Certificate}[证书：商空间纤维识别 + 商拓扑泛性质]
\label{cert:tex1-ch7-qbar}
证书只用两件事：
\begin{enumerate}
  \item
  \[
    \gamma_1\sim_A\gamma_2
    \iff
    Q_A(\gamma_1)=Q_A(\gamma_2)
  \]
  使得 \(\widetilde{\mathcal S}_A\) 的点正好就是 \(Q_A\) 的纤维；
  \item \(\widetilde{\mathcal S}_A\) 与 \(\mathcal{PI}_A\) 分别按 \(\pi_A\) 与 \(Q_A\) 的商拓扑定义。
\end{enumerate}
\end{Certificate}

\begin{proof}[证明（基于数学符号推演逻辑的谓词因果链）]
由 \(\sim_A\) 的定义，
每个商类 \([\gamma]\in\widetilde{\mathcal S}_A\) 唯一对应一个值
\[
  Q_A(\gamma)\in \mathcal{PI}_A.
\]
因此 \(\overline Q_A\) 是良定义且为双射。

又因为
\[
  Q_A=\overline Q_A\circ \pi_A,
\]
并且 \(\mathcal{PI}_A\) 的拓扑正是由 \(Q_A\) 诱导的商拓扑，
故按商拓扑的泛性质，\(\overline Q_A\) 连续。

反过来，由 \(\widetilde{\mathcal S}_A\) 也是由同一纤维识别所得到的商空间，
\(\overline Q_A^{-1}\) 同样满足对应的泛性质，因此也连续。
于是 \(\overline Q_A\) 为同胚。
引理得证。
\end{proof}

\subsection{定理：严格经典同伦等价链}
\label{subsec:tex1-ch7-theorem}

\begin{theorem}[严格经典同伦等价：轨迹商空间、路径积分支撑类、正规形与稳定修复序列]
\label{thm:tex1-ch7-classical}
在定义~\ref{def:tex1-ch7-objects}--定义~\ref{def:tex1-ch7-topology} 的口径下，有同胚链
\[
  \widetilde{\mathcal S}_A
  \cong_{\mathrm{Top}}
  \mathcal{PI}_A
  \cong_{\mathrm{Top}}
  \mathcal T_A
  \cong_{\mathrm{Top}}
  \mathcal E_A^{\mathrm{rep}}.
\]
因此有经典同伦等价链
\[
  \widetilde{\mathcal S}_A
  \simeq_h
  \mathcal{PI}_A
  \simeq_h
  \mathcal T_A
  \simeq_h
  \mathcal E_A^{\mathrm{rep}}.
\]
并且交换关系可写成
\[
\begin{array}{ccccc}
  \mathcal S_A^{\mathrm{adm}} & \xrightarrow{Q_A} & \mathcal{PI}_A & \xrightarrow{\operatorname{NF}_A} & \mathcal T_A \\
  \pi_A\downarrow & & \Vert & & \uparrow q_A|_{\mathcal T_A} \\
  \widetilde{\mathcal S}_A & \xrightarrow{\overline Q_A} & \mathcal{PI}_A & \xleftarrow{\operatorname{Stab}_A} &
    \mathcal E_A^{\mathrm{rep}}
\end{array}
\]
其中 tex36 已给出这些映射的显式双射链与互逆关系。
\end{theorem}

\begin{Certificate}[证书：四步同胚证书]
\label{cert:tex1-ch7-classical}
证书分四步：
\begin{enumerate}
  \item 引理~\ref{lem:tex1-ch7-qbar} 已证
  \[
    \widetilde{\mathcal S}_A
    \cong_{\mathrm{Top}}
    \mathcal{PI}_A;
  \]
  \item tex36 已证
  \[
    \operatorname{NF}_A:\mathcal{PI}_A\to \mathcal T_A
  \]
  与
  \[
    q_A|_{\mathcal T_A}:\mathcal T_A\to\mathcal{PI}_A
  \]
  互为逆映射；由 \(\mathcal T_A\) 的拓扑定义，它们自动构成同胚；
  \item tex36 已证
  \[
    \operatorname{Stab}_A:\mathcal E_A^{\mathrm{rep}}\to\mathcal T_A
  \]
  为双射；由 \(\mathcal E_A^{\mathrm{rep}}\) 的拓扑定义，它自动构成同胚；
  \item 任意同胚都诱导经典同伦等价。
\end{enumerate}
\end{Certificate}

\begin{proof}[证明（基于数学符号推演逻辑的谓词因果链）]
由引理~\ref{lem:tex1-ch7-qbar}，
\[
  \widetilde{\mathcal S}_A
  \cong_{\mathrm{Top}}
  \mathcal{PI}_A.
\]

再由 tex36 的双射链，
\(\operatorname{NF}_A\) 与 \(q_A|_{\mathcal T_A}\) 互为逆映射。
由于 \(\mathcal T_A\) 的拓扑按 \(\operatorname{NF}_A\) 传递定义，
这二者自动互为连续逆映射，从而
\[
  \mathcal{PI}_A
  \cong_{\mathrm{Top}}
  \mathcal T_A.
\]

同理，由 tex36 的稳定修复序列双射，
\(\operatorname{Stab}_A\) 为双射；
由于 \(\mathcal E_A^{\mathrm{rep}}\) 的拓扑按 \(\operatorname{Stab}_A\) 传递定义，
故
\[
  \mathcal T_A
  \cong_{\mathrm{Top}}
  \mathcal E_A^{\mathrm{rep}}.
\]

于是同胚可复合，得到
\[
  \widetilde{\mathcal S}_A
  \cong_{\mathrm{Top}}
  \mathcal{PI}_A
  \cong_{\mathrm{Top}}
  \mathcal T_A
  \cong_{\mathrm{Top}}
  \mathcal E_A^{\mathrm{rep}}.
\]

而任意同胚 \(f:X\to Y\) 都给出经典同伦等价：
取 \(g=f^{-1}\)，并令
\[
  H_X(x,t)=x,
  \qquad
  H_Y(y,t)=y,
\]
则
\[
  g\circ f=\mathrm{id}_X,
  \qquad
  f\circ g=\mathrm{id}_Y,
\]
恒等映射当然与自身同伦。
因此上述四个空间两两经典同伦等价。
定理得证。
\end{proof}

\subsection{推论：钱学森弹道到钱学森障碍的严格落点}
\label{subsec:tex1-ch7-corollaries}

\begin{corollary}[“钱学森弹道 \(\to\) 钱学森障碍”的严格版]
\label{cor:tex1-ch7-qxs}
若把
\[
  \mathcal S_A^{\mathrm{adm}}
\]
理解为偏导控制轨迹的可容许解空间，
把
\[
  \mathcal E_A^{\mathrm{rep}}
\]
理解为稳定障碍修复序列空间，
那么严格成立的命题不是
\[
  \mathcal S_A^{\mathrm{adm}}\simeq_h \mathcal E_A^{\mathrm{rep}},
\]
而是
\[
  \boxed{
    \mathcal S_A^{\mathrm{adm}}/\sim_A
    \simeq_h
    \mathcal E_A^{\mathrm{rep}}.
  }
\]
也就是：
\[
  \boxed{
    \text{原始轨迹问题按支撑类压缩后的商空间，与障碍修复序列空间严格同伦等价。}
  }.
\]
\end{corollary}

\begin{Certificate}[证书：轨迹商压缩 + 四空间同伦链]
\label{cert:tex1-ch7-qxs}
证书分两步：
\begin{enumerate}
  \item 定义~\ref{def:tex1-ch7-quotient} 已把原始轨迹空间压成支撑类商空间 \(\widetilde{\mathcal S}_A\)；
  \item 定理~\ref{thm:tex1-ch7-classical} 已给出
  \[
    \widetilde{\mathcal S}_A
    \simeq_h
    \mathcal E_A^{\mathrm{rep}}.
  \]
\end{enumerate}
\end{Certificate}

\begin{proof}[证明（基于数学符号推演逻辑的谓词因果链）]
由定义~\ref{def:tex1-ch7-quotient}，
\[
  \widetilde{\mathcal S}_A
  =
  \mathcal S_A^{\mathrm{adm}}/\sim_A
\]
正是把原始轨迹空间按路径积分支撑类压缩后的商空间。
再由定理~\ref{thm:tex1-ch7-classical}，
\[
  \widetilde{\mathcal S}_A
  \simeq_h
  \mathcal E_A^{\mathrm{rep}}.
\]
将两式合并，即得结论。
推论得证。
\end{proof}

\begin{corollary}[纤维终点编码也在同一经典同伦型内]
\label{cor:tex1-ch7-fiber}
由于 tex36 还给出
\[
  \mathcal F_A\simeq_{\mathrm{enc\text{-}h}}\mathcal{PI}_A
  \simeq_{\mathrm{enc\text{-}h}}\mathcal T_A,
\]
而上面已经用相容拓扑把它升级为经典拓扑版本，
所以进一步得到
\[
  \widetilde{\mathcal S}_A
  \simeq_h
  \mathcal F_A
  \simeq_h
  \mathcal{PI}_A
  \simeq_h
  \mathcal T_A
  \simeq_h
  \mathcal E_A^{\mathrm{rep}}.
\]
因此，
\[
  \text{纤维终点编码},\quad
  \text{路径支撑类编码},\quad
  \text{正规形编码},\quad
  \text{稳定修复序列编码}
\]
都落在同一个经典同伦型里。
\end{corollary}

\begin{Certificate}[证书：纤维终点编码的拓扑升级证书]
\label{cert:tex1-ch7-fiber}
证书分两段：
\begin{enumerate}
  \item tex36 已给出 \(\mathcal F_A\)、\(\mathcal{PI}_A\)、\(\mathcal T_A\) 的编码双射链；
  \item 定义~\ref{def:tex1-ch7-topology} 已把 \(\mathcal F_A\) 的拓扑按 \(\overline{\mathrm{PT}}_A^{p_{\mathrm{src}}}\) 传递定义，
    因此该双射链可升级为同胚链。
\end{enumerate}
\end{Certificate}

\begin{proof}[证明（基于数学符号推演逻辑的谓词因果链）]
由 tex36，
\[
  \overline{\mathrm{PT}}_A^{p_{\mathrm{src}}}:\mathcal{PI}_A\to \mathcal F_A
\]
与其逆映射在编码层是互逆双射。
又由于 \(\mathcal F_A\) 的拓扑按 \(\overline{\mathrm{PT}}_A^{p_{\mathrm{src}}}\) 传递定义，
故该双射自动升级为同胚。
于是
\[
  \mathcal F_A\cong_{\mathrm{Top}}\mathcal{PI}_A.
\]
再与定理~\ref{thm:tex1-ch7-classical} 复合，即得结论。
推论得证。
\end{proof}

\subsection{命题：原始 PDE 解空间本身的强版仍需额外证书}
\label{subsec:tex1-ch7-proposition}

\begin{proposition}[若要求原始 PDE 解空间本身与障碍空间同伦等价]
\label{prop:tex1-ch7-strong}
若要证明更强命题
\[
  \mathcal S_A^{\mathrm{adm}}
  \simeq_h
  \mathcal T_A,
\]
则还需要额外证书：
存在连续截面
\[
  s_A:\mathcal{PI}_A\to \mathcal S_A^{\mathrm{adm}}
\]
与强变形收缩
\[
  H_A:\mathcal S_A^{\mathrm{adm}}\times I\to \mathcal S_A^{\mathrm{adm}}
\]
满足
\[
  Q_A\circ s_A=\mathrm{id}_{\mathcal{PI}_A},
\]
以及
\[
  H_A(-,0)=\mathrm{id}_{\mathcal S_A^{\mathrm{adm}}},
  \qquad
  H_A(-,1)=s_A\circ Q_A.
\]
一旦这两条成立，就有
\[
  \mathcal S_A^{\mathrm{adm}}
  \simeq_h
  \mathcal{PI}_A
  \simeq_h
  \mathcal T_A.
\]
\end{proposition}

\begin{Certificate}[证书：连续代表选取 + 强变形收缩]
\label{cert:tex1-ch7-strong}
证书分两条：
\begin{enumerate}
  \item \(Q_A\circ s_A=\mathrm{id}_{\mathcal{PI}_A}\) 说明 \(s_A\) 是 \(Q_A\) 的连续右逆；
  \item \(H_A\) 给出
  \[
    \mathrm{id}_{\mathcal S_A^{\mathrm{adm}}}
    \simeq
    s_A\circ Q_A
  \]
  的强变形收缩证书。
\end{enumerate}
\end{Certificate}

\begin{proof}[证明（基于数学符号推演逻辑的谓词因果链）]
由
\[
  Q_A\circ s_A=\mathrm{id}_{\mathcal{PI}_A},
\]
可知 \(Q_A\) 与 \(s_A\) 至少在一侧互逆。
再由
\[
  H_A(-,0)=\mathrm{id}_{\mathcal S_A^{\mathrm{adm}}},
  \qquad
  H_A(-,1)=s_A\circ Q_A,
\]
可知
\[
  \mathrm{id}_{\mathcal S_A^{\mathrm{adm}}}
  \simeq
  s_A\circ Q_A.
\]
因此 \(Q_A\) 与 \(s_A\) 互为同伦逆，故
\[
  \mathcal S_A^{\mathrm{adm}}
  \simeq_h
  \mathcal{PI}_A.
\]
再与定理~\ref{thm:tex1-ch7-classical} 复合，即得
\[
  \mathcal S_A^{\mathrm{adm}}
  \simeq_h
  \mathcal T_A.
\]
命题得证。
\end{proof}

\begin{remark}[本章已经严格完成的层级]
本章已经严格完成的结论是
\[
  \boxed{
    \mathcal S_A^{\mathrm{adm}}/\sim_A
    \cong_{\mathrm{Top}}
    \mathcal{PI}_A
    \cong_{\mathrm{Top}}
    \mathcal T_A
    \cong_{\mathrm{Top}}
    \mathcal E_A^{\mathrm{rep}}
  }
\]
从而
\[
  \boxed{
    \mathcal S_A^{\mathrm{adm}}/\sim_A
    \simeq_h
    \mathcal T_A
    \simeq_h
    \mathcal E_A^{\mathrm{rep}}.
  }
\]
这就是“钱学森弹道（偏导控制解空间）升级为钱学森障碍（拓扑约束/修复空间）”的严格经典同伦等价落点。
\end{remark}

\begin{remark}[强版仍需额外补证]
更强的结论
\[
  \mathcal S_A^{\mathrm{adm}}
  \simeq_h
  \mathcal T_A
\]
目前并不是仓内 tex36/tex38 已直接给出的结论。
仓内当前直接给出的，是编码双射链，以及在补充相容拓扑后可升级为经典同伦等价的那一层；
而“原始 PDE 解空间本身的连续代表选取 \(s_A\) 与强变形收缩 \(H_A\)”仍需要在具体应用中单独补证。
\end{remark}

\clearpage
\section{形式逻辑：原始 admissible 空间到正规形子空间的强变形收缩强版}
\label{sec:tex1-ch8}

\begin{remark}[本章意义备注：高超音速攻防的抽象语境]
本章的意义，在于把庞杂的 admissible 场景压缩到正规形子空间，
从而筛出与障碍形成和修复真正相关的结构自由度。
对钱学森障碍的抽象研究而言，它帮助我们把注意力从冗余表象收回到决定可判性的核心变量。
\end{remark}

\subsection{章节目标与证据来源}
\label{subsec:tex1-ch8-goal}

\begin{remark}[本章目标]
上一章已严格证明
\[
  \mathcal S_A^{\mathrm{adm}}/\sim_A
  \simeq_h
  \mathcal T_A.
\]
本章继续处理更强命题：
\[
  \boxed{
    \mathcal S_A^{\mathrm{adm}}
    \simeq_h
    \mathcal T_A
  }.
\]
更严格地说，我们将证明：在连续截面 \(s_A\) 与类内连续收缩 \(K_A\) 两张证书成立时，
\[
  \boxed{
    \mathcal T_A\subseteq \mathcal S_A^{\mathrm{adm}}
    \text{ 是一个强变形收缩像。}
  }
\]
这就是“钱学森弹道”的原始 admissible 轨迹空间，直接强变形收缩到“钱学森障碍”的正规形子空间之强版落点。
\end{remark}

\begin{remark}[仓内与外部依据]
本章直接承接上一章的商空间同伦链，并进一步调用
\cite{RepoTex36HomotopyTwistEquivalence,RepoTex38JacobiBianchiRepair}。
其中，
\cite{RepoTex36HomotopyTwistEquivalence}
已经给出对象骨架
\[
  q_A:\mathcal S_A^{\mathrm{adm}}\to \mathcal{PI}_A,
  \qquad
  \operatorname{NF}_A:\mathcal{PI}_A\to \mathcal S_A^{\mathrm{adm}},
  \qquad
  q_A\circ \operatorname{NF}_A=\id_{\mathcal{PI}_A},
\]
并明确写出
\[
  \mathcal T_A=\operatorname{im}(\operatorname{NF}_A)
\]
这一正规形子空间；
同文后半又把
\[
  \mathcal E_A^{\mathrm{rep}}
  \cong_{\mathrm{set}}
  \mathcal T_A
\]
做成显式双射。
\cite{RepoTex38JacobiBianchiRepair}
则继续给出终端零缺陷正规形与唯一最优路径支撑类。
外部拓扑上，强变形收缩与同伦逆的标准背景可参见 \cite{Hatcher2002AT}。
\end{remark}

\subsection{定义：原始 admissible 空间、正规形子空间与两张强版证书}
\label{subsec:tex1-ch8-defs}

\begin{definition}[原始 admissible 空间与正规形子空间]
\label{def:tex1-ch8-objects}
令
\[
  \mathcal S_A^{\mathrm{adm}}
  :=
  \operatorname{Path}_A(b_{\mathrm{src}},B_{\mathrm{tar}})
\]
为原始 admissible 路径空间。

令
\[
  q_A:\mathcal S_A^{\mathrm{adm}}\to \mathcal{PI}_A
\]
为自然商映射。

令
\[
  s_A:=\operatorname{NF}_A:\mathcal{PI}_A\to \mathcal S_A^{\mathrm{adm}},
  \qquad
  \mathcal T_A:=\operatorname{im}(s_A)\subseteq \mathcal S_A^{\mathrm{adm}}.
\]
由 tex36 的对象链，
\[
  q_A\circ s_A=\id_{\mathcal{PI}_A}.
\]
因此 \(s_A\) 是 \(q_A\) 的一个截面。
\end{definition}

\begin{definition}[证书 S：连续代表选取]
\label{def:tex1-ch8-cert-s}
称 \(s_A\) 满足\textbf{连续代表选取证书 S}，当且仅当
\[
  s_A:\mathcal{PI}_A\to \mathcal S_A^{\mathrm{adm}}
\]
连续。
\end{definition}

\begin{definition}[证书 R：类内连续收缩]
\label{def:tex1-ch8-cert-r}
称 \(K_A\) 满足\textbf{类内连续收缩证书 R}，当且仅当存在连续映射
\[
  K_A:\mathcal S_A^{\mathrm{adm}}\times I\to \mathcal S_A^{\mathrm{adm}}
\]
满足对任意 \(\gamma\in \mathcal S_A^{\mathrm{adm}}\)、\(t\in I\)：
\[
  q_A\bigl(K_A(\gamma,t)\bigr)=q_A(\gamma),
\]
\[
  K_A(\gamma,0)=\gamma,
  \qquad
  K_A(\gamma,1)=s_A\bigl(q_A(\gamma)\bigr),
\]
\[
  \gamma\in \mathcal T_A
  \Longrightarrow
  K_A(\gamma,t)=\gamma.
\]
也就是说，\(K_A\) 在每个商类纤维内，把任意代表连续缩到该类的正规形代表，并且对正规形点逐点固定。
\end{definition}

\subsection{公理（降级为定理）：强版证书链可按前缀递归闭合}
\label{subsec:tex1-ch8-axiom}

\begin{theorem}[公理（降级为定理）：五段强版证书链可按前缀递归闭合]
\label{thm:tex1-ch8-axiom-downgraded}
定义五段强版见证：
\[
\begin{aligned}
  E_0 &: q_A,\ s_A,\ \mathcal T_A=\operatorname{im}(s_A),\\
  E_1 &: q_A\circ s_A=\id_{\mathcal{PI}_A},\\
  E_2 &: K_A\ \text{满足证书 R，且逐点固定}\ \mathcal T_A,\\
  E_3 &: r_A:=s_A\circ q_A,\ i_A:\mathcal T_A\hookrightarrow \mathcal S_A^{\mathrm{adm}},
         \quad
         r_A\circ i_A=\id_{\mathcal T_A},\\
  E_4 &: \mathcal T_A\subseteq \mathcal S_A^{\mathrm{adm}}
         \text{ 是强变形收缩像，故 }
         \mathcal S_A^{\mathrm{adm}}\simeq_h\mathcal T_A.
\end{aligned}
\]
对每个 \(k\in\{0,1,2,3,4\}\)，定义谓词
\[
  P(k)
  :\Longleftrightarrow
  \text{前缀见证 }E_0,\dots,E_k\text{ 已被文内结果闭合。}
\]
则对全部 \(k\in\{0,1,2,3,4\}\)，命题 \(P(k)\) 成立。
因此，强版同伦等价不是额外口头加强，而是建立在一条可按前缀递归核验的有限证书链之上。
\end{theorem}

\begin{Certificate}[数学归纳法证书：按前缀长度 \(k\) 的闭合]
\label{cert:tex1-ch8-axiom-downgraded}
归纳谓词为
\[
  P(k):
  \text{前缀 }E_0,\dots,E_k\text{ 已闭合。}
\]
证书分两步：
\begin{enumerate}
  \item \textbf{基础步 \(P(0)\)：}
  定义~\ref{def:tex1-ch8-objects} 已固定对象域；
  \item \textbf{归纳步 \(P(k)\Rightarrow P(k+1)\)：}
  对每个 \(k<4\)，后继见证分别由定义~\ref{def:tex1-ch8-objects}、
  定理~\ref{thm:tex1-ch8-strong} 与命题~\ref{prop:tex1-ch8-constructive}
  提供。
\end{enumerate}
\end{Certificate}

\begin{proof}[证明（数学归纳法证书；基于数学符号推演逻辑的谓词因果链）]
对前缀长度 \(k\) 作数学归纳。

\textbf{基础步：}
当 \(k=0\) 时，定义~\ref{def:tex1-ch8-objects} 已给出 \(q_A\)、\(s_A\) 与 \(\mathcal T_A=\operatorname{im}(s_A)\)。
故 \(P(0)\) 成立。

\textbf{归纳步：}
假设 \(P(k)\) 成立。
若 \(k=0\)，则由 tex36 的对象链
\[
  q_A\circ s_A=\id_{\mathcal{PI}_A},
\]
故 \(E_1\) 成立，从而 \(P(1)\) 成立。

若 \(k=1\)，则证书 R 正是定义~\ref{def:tex1-ch8-cert-r} 所固定的条件接口。
一旦该证书被加入，便得到类内连续收缩 \(K_A\) 及其对 \(\mathcal T_A\) 的逐点固定性，故 \(E_2\) 成立，从而 \(P(2)\) 成立。

若 \(k=2\)，则在包含映射 \(i_A\) 与回缩映射 \(r_A:=s_A\circ q_A\) 的记号下，
由 \(q_A\circ s_A=\id_{\mathcal{PI}_A}\) 可推出
\[
  r_A\circ i_A=\id_{\mathcal T_A},
\]
故 \(E_3\) 成立，从而 \(P(3)\) 成立。

若 \(k=3\)，则由定理~\ref{thm:tex1-ch8-strong}，
\(K_A\) 正给出从 \(\id_{\mathcal S_A^{\mathrm{adm}}}\) 到 \(i_A\circ r_A\) 的强变形收缩同伦，
故 \(E_4\) 成立，从而 \(P(4)\) 成立。

因此对全部 \(k\in\{0,1,2,3,4\}\)，都有 \(P(k)\) 成立。
定理得证。
\end{proof}

\subsection{引理：回缩映射在正规形子空间上满足左逆恒等式}
\label{subsec:tex1-ch8-lemma}

\begin{lemma}[回缩映射在正规形子空间上满足左逆恒等式]
\label{lem:tex1-ch8-retraction}
定义包含映射
\[
  i_A:\mathcal T_A\hookrightarrow \mathcal S_A^{\mathrm{adm}}
\]
与回缩映射
\[
  r_A:=s_A\circ q_A:\mathcal S_A^{\mathrm{adm}}\to \mathcal T_A.
\]
则有
\[
  r_A\circ i_A=\id_{\mathcal T_A}.
\]
\end{lemma}

\begin{Certificate}[证书：像空间表示 + 截面恒等式]
\label{cert:tex1-ch8-retraction}
证书分两段：
\begin{enumerate}
  \item 任意 \(\tau\in \mathcal T_A\) 都可写成 \(\tau=s_A(x)\)；
  \item \(q_A\circ s_A=\id_{\mathcal{PI}_A}\)。
\end{enumerate}
\end{Certificate}

\begin{proof}[证明（基于数学符号推演逻辑的谓词因果链）]
任取 \(\tau\in \mathcal T_A\)。
由 \(\mathcal T_A=\operatorname{im}(s_A)\)，存在某个 \(x\in \mathcal{PI}_A\) 使
\[
  \tau=s_A(x).
\]
于是
\[
  r_A(i_A(\tau))
  =
  s_A(q_A(\tau))
  =
  s_A(q_A(s_A(x)))
  =
  s_A(x)
  =
  \tau.
\]
故
\[
  r_A\circ i_A=\id_{\mathcal T_A}.
\]
引理得证。
\end{proof}

\subsection{定理：原始 admissible 空间强变形收缩到正规形子空间}
\label{subsec:tex1-ch8-theorem}

\begin{theorem}[更强版主定理：原始 admissible 空间强变形收缩到正规形子空间]
\label{thm:tex1-ch8-strong}
在证书 S 与证书 R 成立时，
\[
  \mathcal T_A
\]
是
\[
  \mathcal S_A^{\mathrm{adm}}
\]
的强变形收缩像。于是
\[
  \boxed{
    \mathcal S_A^{\mathrm{adm}}
    \simeq_h
    \mathcal T_A.
  }
\]
\end{theorem}

\begin{Certificate}[证书：包含/回缩/强同伦三件套]
\label{cert:tex1-ch8-strong}
只需证明三件事：
\begin{enumerate}
  \item
  \[
    r_A\circ i_A=\id_{\mathcal T_A};
  \]
  \item 取
  \[
    H_A:=K_A,
  \]
  则
  \[
    H_A(\gamma,0)=\gamma,
    \qquad
    H_A(\gamma,1)=i_A\circ r_A(\gamma);
  \]
  \item 对任意 \(\tau\in\mathcal T_A\)，有
  \[
    H_A(\tau,t)=\tau.
  \]
\end{enumerate}
\end{Certificate}

\begin{proof}[证明（基于数学符号推演逻辑的谓词因果链）]
首先，由引理~\ref{lem:tex1-ch8-retraction}，
\[
  r_A\circ i_A=\id_{\mathcal T_A}.
\]

再取
\[
  H_A:=K_A.
\]
由证书 R，
\[
  H_A(\gamma,0)=K_A(\gamma,0)=\gamma.
\]
同时
\[
  H_A(\gamma,1)
  =
  K_A(\gamma,1)
  =
  s_A(q_A(\gamma))
  =
  r_A(\gamma).
\]
由于 \(r_A(\gamma)\in \mathcal T_A\subseteq \mathcal S_A^{\mathrm{adm}}\)，
写成包含映射形式即
\[
  H_A(\gamma,1)=i_A\circ r_A(\gamma).
\]

最后，若 \(\tau\in \mathcal T_A\)，
则由证书 R 的固定性，
\[
  H_A(\tau,t)=K_A(\tau,t)=\tau.
\]

因此 \(H_A\) 是从
\[
  \id_{\mathcal S_A^{\mathrm{adm}}}
\]
到
\[
  i_A\circ r_A
\]
的一个强变形收缩同伦，且逐点固定 \(\mathcal T_A\)。
故
\[
  \mathcal T_A
\]
是
\[
  \mathcal S_A^{\mathrm{adm}}
\]
的强变形收缩像，从而
\[
  \mathcal S_A^{\mathrm{adm}} \simeq_h \mathcal T_A.
\]
定理得证。
\end{proof}

\subsection{推论：与稳定修复序列空间的更强版等价}
\label{subsec:tex1-ch8-corollaries}

\begin{corollary}[与稳定修复序列空间的更强版等价]
\label{cor:tex1-ch8-rep}
tex36 已经把稳定修复序列空间
\[
  \mathcal E_A^{\mathrm{rep}}
\]
与
\[
  \mathcal T_A
\]
做成显式双射 \(\operatorname{Stab}_A\)。
若给 \(\mathcal E_A^{\mathrm{rep}}\) 赋予由 \(\operatorname{Stab}_A\) 传递的拓扑，
则
\[
  \operatorname{Stab}_A:\mathcal E_A^{\mathrm{rep}}\to \mathcal T_A
\]
是同胚。于是
\[
  \boxed{
    \mathcal S_A^{\mathrm{adm}}
    \simeq_h
    \mathcal T_A
    \cong_{\mathrm{Top}}
    \mathcal E_A^{\mathrm{rep}}.
  }
\]
\end{corollary}

\begin{Certificate}[证书：强变形收缩 + 稳定序列同胚]
\label{cert:tex1-ch8-rep}
证书分两步：
\begin{enumerate}
  \item 定理~\ref{thm:tex1-ch8-strong} 已给出
  \[
    \mathcal S_A^{\mathrm{adm}}\simeq_h \mathcal T_A;
  \]
  \item tex36 已给出 \(\operatorname{Stab}_A\) 的双射，而上一章的传递拓扑口径说明它可升级为同胚。
\end{enumerate}
\end{Certificate}

\begin{proof}[证明（基于数学符号推演逻辑的谓词因果链）]
由定理~\ref{thm:tex1-ch8-strong}，
\[
  \mathcal S_A^{\mathrm{adm}}\simeq_h \mathcal T_A.
\]
又由 tex36 的稳定修复序列双射与上一章的拓扑升级口径，
\[
  \operatorname{Stab}_A:\mathcal E_A^{\mathrm{rep}}\to \mathcal T_A
\]
可视为同胚。
故
\[
  \mathcal S_A^{\mathrm{adm}}
  \simeq_h
  \mathcal T_A
  \cong_{\mathrm{Top}}
  \mathcal E_A^{\mathrm{rep}}.
\]
推论得证。
\end{proof}

\begin{proposition}[构造性充分条件：如何显式造出 \(K_A\)]
\label{prop:tex1-ch8-constructive}
设对每个
\[
  x\in \mathcal{PI}_A
\]
都存在一个 Banach 路径图
\[
  \chi_x:q_A^{-1}(x)\to C_x\subset E_x,
\]
满足：
\begin{enumerate}
  \item \(C_x\) 是 \(E_x\) 中的闭凸集；
  \item \(s_A(x)\in q_A^{-1}(x)\) 是该类的规范中心；
  \item 图族 \(\chi_x\) 与中心 \(\chi_x(s_A(x))\) 随 \(x\) 连续变化。
\end{enumerate}
定义
\[
  K_{A,x}(\gamma,t)
  :=
  \chi_x^{-1}
  \Bigl(
    (1-t)\chi_x(\gamma)+t\,\chi_x(s_A(x))
  \Bigr),
\]
再定义全局映射
\[
  K_A(\gamma,t)
  :=
  K_{A,q_A(\gamma)}(\gamma,t).
\]
则 \(K_A\) 满足证书 R，于是定理~\ref{thm:tex1-ch8-strong} 适用。
\end{proposition}

\begin{Certificate}[证书：凸组合封闭 + 类内保持 + 连续依赖]
\label{cert:tex1-ch8-constructive}
证书分三段：
\begin{enumerate}
  \item \(C_x\) 的凸性保证线性插值仍留在 \(C_x\) 内；
  \item \(\chi_x^{-1}\) 把该凸组合送回同一个纤维 \(q_A^{-1}(x)\)；
  \item 图族与中心随 \(x\) 的连续变化给出 \(K_A\) 的全局连续性。
\end{enumerate}
\end{Certificate}

\begin{proof}[证明（基于数学符号推演逻辑的谓词因果链）]
固定一个类
\[
  x\in \mathcal{PI}_A.
\]
由于
\[
  \chi_x(\gamma)\in C_x,
  \qquad
  \chi_x(s_A(x))\in C_x,
\]
而 \(C_x\) 是凸集，所以对任意 \(t\in I\)，
\[
  (1-t)\chi_x(\gamma)+t\,\chi_x(s_A(x))\in C_x.
\]
故 \(K_{A,x}(\gamma,t)\) 良定义，且始终落在同一个纤维
\[
  q_A^{-1}(x)
\]
里。于是
\[
  q_A(K_A(\gamma,t))=q_A(\gamma).
\]

再由定义，
\[
  K_{A,x}(\gamma,0)=\chi_x^{-1}(\chi_x(\gamma))=\gamma,
\]
\[
  K_{A,x}(\gamma,1)=\chi_x^{-1}(\chi_x(s_A(x)))=s_A(x).
\]

若 \(\gamma=s_A(x)\)，则
\[
  K_{A,x}(s_A(x),t)
  =
  \chi_x^{-1}
  \Bigl(
    (1-t)\chi_x(s_A(x))+t\,\chi_x(s_A(x))
  \Bigr)
  =
  s_A(x).
\]
故固定性成立。

最后，由图族与中心对 \(x\) 的连续依赖，
\(K_A\) 全局连续。
因此证书 R 成立。
再由定理~\ref{thm:tex1-ch8-strong}，
\[
  \mathcal S_A^{\mathrm{adm}}\simeq_h \mathcal T_A.
\]
命题得证。
\end{proof}

\begin{corollary}[这就是更强版严格证明的压缩表达]
\label{cor:tex1-ch8-strong}
因此，更强版的严格结论可压成：
\[
  \boxed{
    \text{只要正规化截面 }s_A=\operatorname{NF}_A\text{ 连续，且每个商类纤维可连续收缩到该正规形代表，}
  }
\]
\[
  \boxed{
    \mathcal S_A^{\mathrm{adm}}\text{ 本身就强变形收缩到 } \mathcal T_A,
    \text{而不是只在商空间层面同伦等价。}
  }
\]
换句话说，弱版是
\[
  \mathcal S_A^{\mathrm{adm}}/\sim_A \simeq_h \mathcal T_A,
\]
强版是
\[
  \mathcal S_A^{\mathrm{adm}} \simeq_h \mathcal T_A,
\]
而强版真正新增的内容，正是
\[
  \boxed{
    \text{连续截面 }s_A
    \quad+\quad
    \text{类内强变形收缩 }K_A.
  }
\]
\end{corollary}

\begin{Certificate}[证书：弱版到强版的增量证书]
\label{cert:tex1-ch8-upgrade}
证书可压成两条：
\begin{enumerate}
  \item 上一章已给出弱版
  \[
    \mathcal S_A^{\mathrm{adm}}/\sim_A \simeq_h \mathcal T_A;
  \]
  \item 本章增加的正是定义~\ref{def:tex1-ch8-cert-s} 与定义~\ref{def:tex1-ch8-cert-r} 这两张证书。
\end{enumerate}
\end{Certificate}

\begin{proof}[证明（基于数学符号推演逻辑的谓词因果链）]
由上一章，弱版同伦等价已经在商空间层面成立。
本章新增的并不是对象，
而是把 \(s_A\) 的连续性与 \(K_A\) 的类内连续收缩性单独证书化。
一旦这两张证书成立，
定理~\ref{thm:tex1-ch8-strong} 就把商空间层面的同伦等价提升为原始 admissible 空间本身的强变形收缩结论。
推论得证。
\end{proof}

\begin{remark}[已证层级与条件性边界]
在当前材料下，下列对象骨架已经由 tex36 直接给出：
\[
  q_A,\qquad
  \operatorname{NF}_A,\qquad
  q_A\circ \operatorname{NF}_A=\id,\qquad
  \mathcal T_A=\operatorname{im}(\operatorname{NF}_A),
\]
以及
\[
  \mathcal E_A^{\mathrm{rep}}\cong \mathcal T_A.
\]
这些是已证对象层。

但若要把
\[
  \mathcal S_A^{\mathrm{adm}}\simeq_h \mathcal T_A
\]
写成无条件主定理，
则仍需要在具体模型中闭合下面两件事：
\[
  \text{(i) } s_A=\operatorname{NF}_A \text{ 的连续性，}
\]
\[
  \text{(ii) } q_A^{-1}(x) \text{ 的类内连续收缩 }K_{A,x}.
\]
若这两件事没有被单独立住，则更强命题仍然是条件性定理；
无条件已证的，仍是上一章的商空间版本。
\end{remark}

\begin{remark}[最压缩的一句话]
最压缩的更强版证明可以写成：
\[
  \boxed{
    r_A:=s_A\circ q_A,\qquad
    H_A(\gamma,t):=K_A(\gamma,t).
  }
\]
若满足
\[
  q_A\circ s_A=\id,\qquad
  H_A(-,0)=\id,\qquad
  H_A(-,1)=i_A\circ r_A,\qquad
  H_A|_{\mathcal T_A\times I}=\id,
\]
则
\[
  \boxed{
    \mathcal T_A \text{ 是 } \mathcal S_A^{\mathrm{adm}} \text{ 的强变形收缩像。}
  }
\]
这就是完整的更强版。
\end{remark}

\clearpage
\section{形式逻辑：连续统崩溃、语义收缩边缘算子与组合同伦强版}
\label{sec:tex1-ch9}

\begin{remark}[本章意义备注：高超音速攻防的抽象语境]
本章的意义，在于说明连续观测空间可以在保持关键结构的前提下收缩为离散有效骨架。
对高超音速攻防的抽象判定而言，
这为实时识别、分层告警与符号化证书提供了理论支撑，但仍停留在建模层而非工程流程层。
\end{remark}

\subsection{严格主张、对象边界与引文定位}
\label{subsec:tex1-ch9-claim}

\begin{remark}[核心陈述]
本节把“钱学森弹道”的连续统 admissible 对象侧，与“钱学森障碍”的离散修复/正规形侧，严格接成下面四步：
\[
  \text{连续统崩溃}
  \Longrightarrow
  \text{语义收缩边缘算子}
  \Longrightarrow
  \text{组合同伦强收缩}
  \Longrightarrow
  \text{几何实现后的经典强变形收缩}.
\]
其逻辑含义是：连续统并不是被原样保留，而是先经语义收缩压到有限离散有效骨架；
再由类保持、严格降缺陷的边缘算子在每个商类纤维内做有限步修复；
最后通过组合同伦与几何实现，把离散统上的有限步修复提升为经典拓扑意义下的强变形收缩。
\end{remark}

\begin{remark}[严格边界]
本节不采用“纯离散拓扑 + 经典 \(I=[0,1]\)-同伦”这一会导致非平凡同伦退化的写法，
而采用
\[
  \boxed{
    \text{离散统对象}
    +\text{组合同伦}
    +\text{几何实现}
    +\text{经典强变形收缩}
  }.
\]
在这个口径下，``钱学森弹道''、``钱学森障碍''、``弹道'' 仍只作为抽象命名标签使用：
前者指连续对象空间与其观测/路径表象侧，后者指语义收缩后的障碍修复、正规化与终端代表侧。
\end{remark}

\begin{remark}[仓内文稿与外部背景]
本节的对象骨架与正规形选择口径承接 \cite{RepoTex36HomotopyTwistEquivalence}，
其中“路径积分支撑类 = 同伦扭曲正规形 = 稳定修复序列”的集合级双射已被写成正式链条；
缺陷严格下降、有限终止与终端零缺陷正规形，则承接
\cite{RepoTex38JacobiBianchiRepair}。
外部拓扑上，有限骨架的几何实现、强变形收缩与同伦等价的标准背景，
可参见 \cite{Hatcher2002AT,Husemoller1994FibreBundles}。
\end{remark}

\subsection{定义链：连续统崩溃、离散有效骨架与语义收缩边缘算子}
\label{subsec:tex1-ch9-defs}

\begin{definition}[连续统崩溃的语义收缩]
\label{def:tex1-ch9-collapse}
设
\[
  X_A
\]
是原始 admissible 连续对象空间。
定义一个连续统崩溃的语义收缩映射
\[
  \mathfrak C_A:X_A\to D_A,
\]
其中
\[
  D_A=(V_A,E_A)
\]
是有限有向图，称为\emph{有限离散有效骨架}。

再定义：
\begin{enumerate}
  \item 商映射
  \[
    q_A:V_A\to \Pi_A,
  \]
  其中 \(\Pi_A\) 表示路径积分支撑类集合；
  \item 缺陷计数
  \[
    \mu_A:V_A\to \NN;
  \]
  \item 正规形子集
  \[
    T_A:=\{v\in V_A:\mu_A(v)=0\}.
  \]
\end{enumerate}
在本节语境下，\(V_A\) 记录“连续统崩溃后的有效离散代表”，
\(E_A\) 记录“一步允许的边缘算子修复步”，
\(T_A\) 记录“零缺陷终端正规形代表”。
\end{definition}

\begin{definition}[语义收缩边缘算子]
\label{def:tex1-ch9-operator}
对每个 \(v\in V_A\)，定义语义收缩边缘算子
\[
  \partial^-_{A,\Sigma}:V_A\to V_A
\]
如下：
\begin{enumerate}
  \item 若 \(\mu_A(v)=0\)，则定义
  \[
    \partial^-_{A,\Sigma}(v):=v;
  \]
  \item 若 \(\mu_A(v)>0\)，则从 \(v\) 的允许出边中，按语义规则 \(\Sigma\) 选取一个规范后继 \(w_v\)，并定义
  \[
    \partial^-_{A,\Sigma}(v):=w_v.
  \]
\end{enumerate}
本节稍后的公理（降级为定理）将证明：在有限离散有效骨架上，若存在类保持、严格降缺陷与终端唯一正规形接口，
则 \(\partial^-_{A,\Sigma}\) 自动满足正规化与类内收束的四条闭合律。
\end{definition}

\begin{definition}[类内离散收缩、全局离散同伦与几何实现]
\label{def:tex1-ch9-contraction}
对任意 \(x\in \Pi_A\)，定义其纤维
\[
  q_A^{-1}(x)\subseteq V_A,
\]
并定义
\[
  M_x:=\max\{\mu_A(v):v\in q_A^{-1}(x)\},
  \qquad
  M:=\max_{v\in V_A}\mu_A(v).
\]

对任意 \(x\in \Pi_A\)，定义类内离散收缩
\[
  K_{A,x}:q_A^{-1}(x)\times \{0,1,\dots,M_x\}\to q_A^{-1}(x),
\]
为
\[
  K_{A,x}(v,k)
  :=
  \bigl(\partial^-_{A,\Sigma}\bigr)^{\min(k,\mu_A(v))}(v).
\]

再定义全局离散同伦
\[
  H_A:V_A\times \{0,1,\dots,M\}\to V_A,
\]
为
\[
  H_A(v,k)
  :=
  \bigl(\partial^-_{A,\Sigma}\bigr)^{\min(k,\mu_A(v))}(v).
\]

最后，记 \(G_A^\partial\) 为由顶点集 \(V_A\) 与边集
\[
  E_A^\partial
  :=
  \bigl\{
    \{v,\partial^-_{A,\Sigma}(v)\}
    :
    v\in V_A\setminus T_A
  \bigr\}
\]
构成的无向图。
记
\[
  |V_A|
\]
为 \(G_A^\partial\) 的几何实现，
记
\[
  |T_A|
\]
为其在顶点集 \(T_A\) 上诱导的全子图的几何实现。
\end{definition}

\subsection{公理（降级为定理）：语义收缩边缘算子的四条闭合律}
\label{subsec:tex1-ch9-axiom}

\begin{theorem}[公理（降级为定理）：语义收缩边缘算子的四条闭合律]
\label{thm:tex1-ch9-axiom-downgraded}
设定义~\ref{def:tex1-ch9-collapse} 与定义~\ref{def:tex1-ch9-operator} 中的对象已经固定，
并满足以下接口：
\begin{enumerate}
  \item \(V_A\) 为有限集；
  \item 对任意 \(v\in V_A\setminus T_A\)，存在一条允许修复边
  \[
    v\to w_v\in E_A
  \]
  使
  \[
    q_A(w_v)=q_A(v),
    \qquad
    \mu_A(w_v)\le \mu_A(v)-1;
  \]
  \item 对任意 \(x\in \Pi_A\)，纤维 \(q_A^{-1}(x)\) 中存在唯一零缺陷顶点
  \[
    \tau_x\in q_A^{-1}(x)\cap T_A;
  \]
  \item \(\partial^-_{A,\Sigma}\) 在每个非正规形点上取上述规范后继 \(w_v\)，并在 \(T_A\) 上固定。
\end{enumerate}
则 \(\partial^-_{A,\Sigma}\) 自动满足以下四条闭合律：
\begin{enumerate}
  \item \textbf{类保持：}
  \[
    q_A\!\left(
      \bigl(\partial^-_{A,\Sigma}\bigr)^k(v)
    \right)=q_A(v)
    \qquad
    (0\le k\le \mu_A(v));
  \]
  \item \textbf{严格降缺陷：}
  \[
    \mu_A\!\left(
      \bigl(\partial^-_{A,\Sigma}\bigr)^k(v)
    \right)
    \le
    \mu_A(v)-k
    \qquad
    (0\le k\le \mu_A(v));
  \]
  \item \textbf{正规形固定：}
  \[
    \mu_A(v)=0
    \Rightarrow
    \partial^-_{A,\Sigma}(v)=v;
  \]
  \item \textbf{类内收束一致性：}
  \[
    \bigl(\partial^-_{A,\Sigma}\bigr)^{\mu_A(v)}(v)=\tau_{q_A(v)},
  \]
  并且若 \(q_A(v_1)=q_A(v_2)\)，则
  \[
    \bigl(\partial^-_{A,\Sigma}\bigr)^{\mu_A(v_1)}(v_1)
    =
    \bigl(\partial^-_{A,\Sigma}\bigr)^{\mu_A(v_2)}(v_2).
  \]
\end{enumerate}
\end{theorem}

\begin{Certificate}[数学归纳法证书：按缺陷计数的有限步收束证书]
\label{cert:tex1-ch9-axiom-downgraded}
定义谓词
\[
  P(m)
  :
  \Longleftrightarrow
  \text{对任意满足 }\mu_A(v)=m\text{ 的 }v\in V_A,
  \text{四条闭合律都成立。}
\]
证书分两步：
\begin{enumerate}
  \item \textbf{基础步 \(P(0)\)：}零缺陷点按定义被 \(\partial^-_{A,\Sigma}\) 固定；
  \item \textbf{归纳步 \(P(m)\Rightarrow P(m+1)\)：}对任意 \(\mu_A(v)=m+1\) 的点，
  由规范后继 \(w_v=\partial^-_{A,\Sigma}(v)\) 满足
  \[
    q_A(w_v)=q_A(v),
    \qquad
    \mu_A(w_v)\le m,
  \]
  再把归纳假设施加到 \(w_v\)。
\end{enumerate}
\end{Certificate}

\begin{proof}[证明（数学归纳法证书；基于数学符号推演逻辑的谓词因果链）]
对 \(m\in\NN\) 做数学归纳。

\textbf{基础步：}若 \(\mu_A(v)=0\)，则 \(v\in T_A\)。
由定义~\ref{def:tex1-ch9-operator}，
\[
  \partial^-_{A,\Sigma}(v)=v.
\]
于是
\[
  q_A\!\left(
    \bigl(\partial^-_{A,\Sigma}\bigr)^0(v)
  \right)=q_A(v),
  \qquad
  \mu_A\!\left(
    \bigl(\partial^-_{A,\Sigma}\bigr)^0(v)
  \right)=0,
\]
并且
\[
  \bigl(\partial^-_{A,\Sigma}\bigr)^{\mu_A(v)}(v)=v=\tau_{q_A(v)}.
\]
故 \(P(0)\) 成立。

\textbf{归纳步：}设 \(P(m)\) 成立，任取 \(\mu_A(v)=m+1\)。
记
\[
  w_v:=\partial^-_{A,\Sigma}(v).
\]
由定理假设，
\[
  q_A(w_v)=q_A(v),
  \qquad
  \mu_A(w_v)\le m.
\]
对任意 \(1\le k\le m+1\)，由归纳假设施加于 \(w_v\)，得到
\[
  q_A\!\left(
    \bigl(\partial^-_{A,\Sigma}\bigr)^k(v)
  \right)
  =
  q_A\!\left(
    \bigl(\partial^-_{A,\Sigma}\bigr)^{k-1}(w_v)
  \right)
  =
  q_A(w_v)
  =
  q_A(v),
\]
以及
\[
  \mu_A\!\left(
    \bigl(\partial^-_{A,\Sigma}\bigr)^k(v)
  \right)
  =
  \mu_A\!\left(
    \bigl(\partial^-_{A,\Sigma}\bigr)^{k-1}(w_v)
  \right)
  \le
  \mu_A(w_v)-(k-1)
  \le
  m-(k-1)
  =
  (m+1)-k.
\]
特别地，当 \(k=m+1=\mu_A(v)\) 时，
\[
  \mu_A\!\left(
    \bigl(\partial^-_{A,\Sigma}\bigr)^{\mu_A(v)}(v)
  \right)
  \le 0,
\]
故
\[
  \bigl(\partial^-_{A,\Sigma}\bigr)^{\mu_A(v)}(v)\in T_A\cap q_A^{-1}(q_A(v)).
\]
由零缺陷代表的唯一性，
\[
  \bigl(\partial^-_{A,\Sigma}\bigr)^{\mu_A(v)}(v)=\tau_{q_A(v)}.
\]
于是 \(P(m+1)\) 成立。

由数学归纳法，\(P(m)\) 对所有 \(m\in\NN\) 成立。
最后，若 \(q_A(v_1)=q_A(v_2)=x\)，则由上式分别得到
\[
  \bigl(\partial^-_{A,\Sigma}\bigr)^{\mu_A(v_1)}(v_1)=\tau_x
  =
  \bigl(\partial^-_{A,\Sigma}\bigr)^{\mu_A(v_2)}(v_2).
\]
故四条闭合律全部成立。
\end{proof}

\subsection{定理链：正规化截面、离散连续性与组合强变形收缩}
\label{subsec:tex1-ch9-theorems}

\begin{theorem}[正规化截面由语义收缩边缘算子唯一决定]
\label{thm:tex1-ch9-section}
定义
\[
  s_A:\Pi_A\to T_A
\]
为
\[
  s_A(x)
  :=
  \bigl(\partial^-_{A,\Sigma}\bigr)^{\mu_A(v)}(v),
  \qquad
  v\in q_A^{-1}(x).
\]
则 \(s_A\) 良定义，并且
\[
  q_A\circ s_A=\id_{\Pi_A},
  \qquad
  s_A=\operatorname{NF}_A.
\]
\end{theorem}

\begin{Certificate}[证书：有限步收束 + 唯一终端代表]
\label{cert:tex1-ch9-section}
证书分三段：
\begin{enumerate}
  \item 由定理~\ref{thm:tex1-ch9-axiom-downgraded}，
  对任意 \(v\in q_A^{-1}(x)\)，
  \[
    \bigl(\partial^-_{A,\Sigma}\bigr)^{\mu_A(v)}(v)=\tau_x;
  \]
  \item 因而 \(s_A(x)\) 与代表元 \(v\) 的选取无关；
  \item tex36 中 \(\operatorname{NF}_A\) 的角色正是“把同一支撑类压到唯一正规形代表”，而 tex38 给出终端零缺陷正规形的唯一性。
\end{enumerate}
\end{Certificate}

\begin{proof}[证明（基于数学符号推演逻辑的谓词因果链）]
任取 \(x\in \Pi_A\) 与 \(v\in q_A^{-1}(x)\)。
由定理~\ref{thm:tex1-ch9-axiom-downgraded}，
\[
  \bigl(\partial^-_{A,\Sigma}\bigr)^{\mu_A(v)}(v)=\tau_x\in T_A.
\]
若 \(v_1,v_2\in q_A^{-1}(x)\)，则同理由类内收束一致性，
\[
  \bigl(\partial^-_{A,\Sigma}\bigr)^{\mu_A(v_1)}(v_1)
  =
  \tau_x
  =
  \bigl(\partial^-_{A,\Sigma}\bigr)^{\mu_A(v_2)}(v_2).
\]
故 \(s_A\) 良定义。

再由 \(\tau_x\in q_A^{-1}(x)\)，得到
\[
  q_A(s_A(x))=x,
\]
即
\[
  q_A\circ s_A=\id_{\Pi_A}.
\]

最后，由 \cite{RepoTex36HomotopyTwistEquivalence}，
\(\operatorname{NF}_A\) 的定义角色是：对每个路径积分支撑类选取唯一正规形代表；
由 \cite{RepoTex38JacobiBianchiRepair}，
终端零缺陷正规形代表在每个类内唯一。
而 \(s_A\) 正是把每个 \(x\in\Pi_A\) 送到其唯一零缺陷正规形 \(\tau_x\)。
因此对任意 \(x\in\Pi_A\)，有
\[
  s_A(x)=\operatorname{NF}_A(x),
\]
从而
\[
  s_A=\operatorname{NF}_A.
\]
定理得证。
\end{proof}

\begin{theorem}[离散统里的连续性自动闭合]
\label{thm:tex1-ch9-continuity}
给 \(\Pi_A\)、\(T_A\)、每个纤维 \(q_A^{-1}(x)\) 与时间集
\[
  \{0,1,\dots,M_x\},
  \qquad
  \{0,1,\dots,M\}
\]
都赋离散拓扑，则映射
\[
  s_A:\Pi_A\to T_A,
\]
\[
  K_{A,x}:q_A^{-1}(x)\times\{0,1,\dots,M_x\}\to q_A^{-1}(x),
\]
\[
  H_A:V_A\times\{0,1,\dots,M\}\to V_A
\]
全部连续。
\end{theorem}

\begin{Certificate}[证书：离散拓扑上的任意映射连续]
\label{cert:tex1-ch9-continuity}
若定义域的每个子集都开，则任意映射的任意开集原像都开。
\end{Certificate}

\begin{proof}[证明（基于数学符号推演逻辑的谓词因果链）]
在离散拓扑下，定义域中的任意子集都是开集。
因此对任意映射 \(f:X\to Y\) 与任意开集 \(U\subseteq Y\)，
\[
  f^{-1}(U)\subseteq X
\]
自动为开集。
把这一事实分别施加于 \(s_A\)、\(K_{A,x}\) 与 \(H_A\)，结论即得。
\end{proof}

\begin{remark}[这一连续性并不自动给出非平凡经典同伦]
定理~\ref{thm:tex1-ch9-continuity} 只解决了“截面连续性”和“类内收缩连续性”的闭合问题。
若直接把所有对象都当作离散拓扑空间并使用经典 \(I\)-同伦，
则非平凡同伦会退化。
因此真正的经典强变形收缩，还必须经过后续的组合同伦与几何实现步骤。
\end{remark}

\begin{lemma}[类内离散收缩的端点、类保持与固定性]
\label{lem:tex1-ch9-fiber}
对任意 \(x\in\Pi_A\) 与任意 \(v\in q_A^{-1}(x)\)，映射 \(K_{A,x}\) 满足：
\begin{enumerate}
  \item
  \[
    K_{A,x}(v,0)=v;
  \]
  \item
  \[
    K_{A,x}(v,M_x)=s_A(x);
  \]
  \item 对任意 \(k\in\{0,1,\dots,M_x\}\)，
  \[
    K_{A,x}(v,k)\in q_A^{-1}(x);
  \]
  \item 若 \(v\in T_A\cap q_A^{-1}(x)\)，则对任意 \(k\)，
  \[
    K_{A,x}(v,k)=v.
  \]
\end{enumerate}
\end{lemma}

\begin{Certificate}[证书：类保持 + 有限步收束 + 正规形固定]
\label{cert:tex1-ch9-fiber}
证书只用三条：
\begin{enumerate}
  \item 定理~\ref{thm:tex1-ch9-axiom-downgraded} 的类保持；
  \item 定理~\ref{thm:tex1-ch9-axiom-downgraded} 的类内收束一致性；
  \item 定理~\ref{thm:tex1-ch9-axiom-downgraded} 的正规形固定。
\end{enumerate}
\end{Certificate}

\begin{proof}[证明（基于数学符号推演逻辑的谓词因果链）]
由定义~\ref{def:tex1-ch9-contraction}，
\[
  K_{A,x}(v,0)
  =
  \bigl(\partial^-_{A,\Sigma}\bigr)^0(v)
  =
  v.
\]

再由 \(M_x\ge \mu_A(v)\)，得到
\[
  K_{A,x}(v,M_x)
  =
  \bigl(\partial^-_{A,\Sigma}\bigr)^{\mu_A(v)}(v)
  =
  s_A(x),
\]
其中最后一步由定理~\ref{thm:tex1-ch9-section}。

对任意 \(k\)，
\[
  q_A\bigl(K_{A,x}(v,k)\bigr)
  =
  q_A\!\left(
    \bigl(\partial^-_{A,\Sigma}\bigr)^{\min(k,\mu_A(v))}(v)
  \right)
  =
  q_A(v)
  =
  x,
\]
故 \(K_{A,x}(v,k)\in q_A^{-1}(x)\)。

最后，若 \(v\in T_A\)，则 \(\mu_A(v)=0\)。
因此
\[
  K_{A,x}(v,k)
  =
  \bigl(\partial^-_{A,\Sigma}\bigr)^0(v)
  =
  v.
\]
引理得证。
\end{proof}

\begin{corollary}[离散统下的组合强变形收缩]
\label{cor:tex1-ch9-comb}
全局离散同伦 \(H_A\) 给出一个从 \(V_A\) 到 \(T_A\) 的组合强变形收缩：
\[
  H_A(v,0)=v,
  \qquad
  H_A(v,M)=s_A(q_A(v)),
\]
并且
\[
  v\in T_A
  \Rightarrow
  H_A(v,k)=v.
\]
因此，在组合同伦意义下，
\[
  \boxed{
    V_A \searrow_{\mathrm{comb}} T_A.
  }
\]
\end{corollary}

\begin{Certificate}[证书：逐纤维收缩拼接为全局收缩]
\label{cert:tex1-ch9-comb}
证书分两步：
\begin{enumerate}
  \item 引理~\ref{lem:tex1-ch9-fiber} 已在每个纤维 \(q_A^{-1}(x)\) 上给出离散收缩；
  \item 不同纤维之间互不干扰，因为 \(H_A(v,k)\) 始终保持在 \(q_A(v)\) 所属的同一纤维内。
\end{enumerate}
\end{Certificate}

\begin{proof}[证明（基于数学符号推演逻辑的谓词因果链）]
由定义~\ref{def:tex1-ch9-contraction}，
\[
  H_A(v,0)
  =
  \bigl(\partial^-_{A,\Sigma}\bigr)^0(v)
  =
  v.
\]
又由于 \(M\ge \mu_A(v)\)，
\[
  H_A(v,M)
  =
  \bigl(\partial^-_{A,\Sigma}\bigr)^{\mu_A(v)}(v)
  =
  s_A(q_A(v)).
\]
若 \(v\in T_A\)，则 \(\mu_A(v)=0\)，故
\[
  H_A(v,k)=v.
\]
因此 \(H_A\) 在每个纤维上都把任意点有限步收缩到唯一正规形代表，
并且逐点固定 \(T_A\)。
故它给出组合意义下的强变形收缩。
\end{proof}

\begin{proposition}[几何实现后得到经典强变形收缩]
\label{prop:tex1-ch9-realization}
设 \(i_A:|T_A|\hookrightarrow |V_A|\) 为包含映射，
并把顶点回缩规则
\[
  r_A(v):=s_A(q_A(v))
  \qquad
  (v\in V_A)
\]
定义在 \(V_A\) 上。
记
\[
  \overline r_A:|V_A|\to |T_A|
\]
为沿 \(G_A^\partial\) 各个根树分量向其根顶点连续收缩所得到的回缩。
则组合强变形收缩 \(H_A\) 可以沿 \(G_A^\partial\) 的边做分段线性插值，
得到一个经典同伦
\[
  \widehat H_A:|V_A|\times I\to |V_A|,
\]
满足
\[
  \widehat H_A(-,0)=\id_{|V_A|},
  \qquad
  \widehat H_A(-,1)=i_A\circ \overline r_A,
\]
并且
\[
  \widehat H_A(\tau,t)=\tau
  \qquad
  (\tau\in |T_A|).
\]
从而
\[
  \boxed{
    |V_A| \searrow |T_A|.
  }
\]
\end{proposition}

\begin{Certificate}[证书：收缩轨道沿规范边集分段线性插值]
\label{cert:tex1-ch9-realization}
证书只用两条：
\begin{enumerate}
  \item 对任意顶点 \(v\)，相邻时刻 \(H_A(v,k)\) 与 \(H_A(v,k+1)\) 要么相等，要么由 \(G_A^\partial\) 中的一条边连接；
  \item 有限图的每条边在几何实现中都是一个闭线段，故可逐段线性插值。
\end{enumerate}
\end{Certificate}

\begin{proof}[证明（基于数学符号推演逻辑的谓词因果链）]
先证 \(G_A^\partial\) 的每个连通分量都是一棵以某个 \(\tau_x\in T_A\) 为根的有向树。
事实上，对任意 \(v\in V_A\setminus T_A\)，其唯一出边由
\[
  v\longmapsto \partial^-_{A,\Sigma}(v)
\]
给出，而定理~\ref{thm:tex1-ch9-axiom-downgraded} 保证
\[
  \mu_A(\partial^-_{A,\Sigma}(v))\le \mu_A(v)-1.
\]
故沿有向边前进时，缺陷计数严格下降，因此不存在有向环。
又由同一定理的类内收束一致性，每个分量都在有限步后收束到唯一零缺陷根 \(\tau_x\)。
故 \(G_A^\partial\) 是若干根树分量的不交并。

接着固定一个顶点 \(v\in V_A\) 与一个整数 \(0\le k<M\)。
若 \(k\ge \mu_A(v)\)，则
\[
  H_A(v,k+1)=H_A(v,k)=s_A(q_A(v)).
\]
若 \(k<\mu_A(v)\)，则由定义~\ref{def:tex1-ch9-contraction}，
\[
  H_A(v,k+1)
  =
  \partial^-_{A,\Sigma}\!\bigl(H_A(v,k)\bigr).
\]
因此 \(H_A(v,k)\) 与 \(H_A(v,k+1)\) 要么相等，
要么正好是 \(G_A^\partial\) 中的一对邻接顶点。

于是对每个时间区间
\[
  \left[\frac{k}{M},\frac{k+1}{M}\right],
\]
都可以在对应的边线段上做线性插值。
把所有顶点轨道的线性插值按时间分段拼接，
就得到连续映射
\[
  \widehat H_A:|V_A|\times I\to |V_A|.
\]
其终点映射正是各根树分量向其根的连续回缩 \(\overline r_A\) 与包含 \(i_A\) 的复合。

起点条件
\[
  \widehat H_A(-,0)=\id_{|V_A|}
\]
来自 \(H_A(v,0)=v\)；
终点条件
\[
  \widehat H_A(-,1)=i_A\circ \overline r_A
\]
来自 \(H_A(v,M)=s_A(q_A(v))\)；
而 \(|T_A|\) 上的逐点固定性，则来自 \(v\in T_A\Rightarrow H_A(v,k)=v\)。
因此 \(\widehat H_A\) 是一个经典强变形收缩同伦，
从而
\[
  |V_A|\searrow |T_A|.
\]
\end{proof}

\begin{theorem}[连续统崩溃后的更强版同伦等价]
\label{thm:tex1-ch9-collapse-equivalence}
若存在连续映射
\[
  \mathfrak C_A:X_A\to |V_A|,
  \qquad
  j_A:|V_A|\hookrightarrow X_A
\]
满足
\[
  \mathfrak C_A\circ j_A=\id_{|V_A|},
\]
并且
\[
  j_A\circ \mathfrak C_A \simeq_h \id_{X_A},
\]
则有
\[
  X_A \simeq_h |V_A| \simeq_h |T_A|,
\]
从而
\[
  \boxed{
    X_A \simeq_h |T_A|.
  }
\]
\end{theorem}

\begin{Certificate}[证书：连续统崩溃同伦逆 + 正规形子骨架强收缩]
\label{cert:tex1-ch9-collapse-equivalence}
证书分两段：
\begin{enumerate}
  \item \(\mathfrak C_A\circ j_A=\id_{|V_A|}\) 与 \(j_A\circ \mathfrak C_A \simeq_h \id_{X_A}\) 说明 \(X_A\) 与 \(|V_A|\)
    互为同伦逆；
  \item 命题~\ref{prop:tex1-ch9-realization} 已给出 \(|V_A|\searrow |T_A|\)。
\end{enumerate}
\end{Certificate}

\begin{proof}[证明（基于数学符号推演逻辑的谓词因果链）]
由
\[
  \mathfrak C_A\circ j_A=\id_{|V_A|}
\]
与
\[
  j_A\circ \mathfrak C_A \simeq_h \id_{X_A},
\]
可知 \(X_A\) 与 \(|V_A|\) 同伦等价。
再由命题~\ref{prop:tex1-ch9-realization}，
\[
  |V_A|\searrow |T_A|,
\]
故
\[
  |V_A|\simeq_h |T_A|.
\]
同伦等价可传递，因此
\[
  X_A \simeq_h |T_A|.
\]
定理得证。
\end{proof}

\begin{corollary}[“连续统崩溃的语义收缩边缘算子修复”的严格落点]
\label{cor:tex1-ch9-final}
设钱学森弹道的连续对象侧由 \(X_A\) 表示，
钱学森障碍的离散修复/正规形侧由 \(T_A\) 表示。
则在本节口径下，下面这句话具有严格数学落点：
\[
  \boxed{
    \text{先用连续统崩溃映射 }\mathfrak C_A\text{ 把 }X_A\text{ 压到有限离散有效骨架 }V_A;
  }
\]
\[
  \boxed{
    \text{再用语义收缩边缘算子 }\partial^-_{A,\Sigma}\text{ 在每个商类纤维内做严格降缺陷的有限步修复;}
  }
\]
\[
  \boxed{
    \text{由此得到 }V_A\searrow_{\mathrm{comb}}T_A,\quad |V_A|\searrow |T_A|,\quad X_A\simeq_h |T_A|.
  }
\]
\end{corollary}

\begin{Certificate}[证书：连续统崩溃 + 语义收缩 + 几何实现]
\label{cert:tex1-ch9-final}
证书分三段：
\begin{enumerate}
  \item 定理~\ref{thm:tex1-ch9-section} 与定理~\ref{thm:tex1-ch9-continuity} 给出 \(s_A\) 的内生定义与连续性；
  \item 引理~\ref{lem:tex1-ch9-fiber} 与推论~\ref{cor:tex1-ch9-comb} 给出纤维内有限步修复与组合强变形收缩；
  \item 命题~\ref{prop:tex1-ch9-realization} 与定理~\ref{thm:tex1-ch9-collapse-equivalence} 给出几何实现后的经典强变形收缩与连续统同伦等价。
\end{enumerate}
\end{Certificate}

\begin{proof}[证明（基于数学符号推演逻辑的谓词因果链）]
由定理~\ref{thm:tex1-ch9-axiom-downgraded}，
\(\partial^-_{A,\Sigma}\) 在每个商类纤维内都类保持并严格降缺陷，
且有限步收束到唯一零缺陷代表。
故定理~\ref{thm:tex1-ch9-section} 给出
\[
  s_A=\operatorname{NF}_A,
\]
而定理~\ref{thm:tex1-ch9-continuity} 给出 \(s_A\) 与 \(K_{A,x}\) 的连续性自动闭合。

再由引理~\ref{lem:tex1-ch9-fiber}，
每个纤维 \(q_A^{-1}(x)\) 都在有限步内收缩到唯一正规形代表；
由推论~\ref{cor:tex1-ch9-comb}，
全局上得到
\[
  V_A\searrow_{\mathrm{comb}}T_A.
\]

随后，由命题~\ref{prop:tex1-ch9-realization}，
\[
  |V_A|\searrow |T_A|.
\]
若再加入连续统崩溃映射 \(\mathfrak C_A\) 与嵌入 \(j_A\) 的同伦逆接口，
则定理~\ref{thm:tex1-ch9-collapse-equivalence} 给出
\[
  X_A\simeq_h |T_A|.
\]
故推论成立。
\end{proof}

\begin{remark}[这一步内生化了上一节的两张证书]
在上一节的连续对象空间口径下，
强版结论仍把
\[
  \text{(i) } s_A=\operatorname{NF}_A \text{ 的连续性},
  \qquad
  \text{(ii) } q_A^{-1}(x)\text{ 的类内连续收缩 }K_{A,x}
\]
留作条件接口。
而本节的作用正是：在离散统与几何实现路线下，
把这两件事分别降为定理~\ref{thm:tex1-ch9-continuity} 与引理~\ref{lem:tex1-ch9-fiber} 的内部结论。
\end{remark}

\begin{remark}[最关键的一句话]
真正给出强版的，并不是“离散拓扑本身”，而是下面四步一起成立：
\[
  \boxed{
    \text{连续统崩溃}
    +\text{语义收缩边缘算子}
    +\text{组合同伦}
    +\text{几何实现}
  }.
\]
也正因为如此，
``钱学森弹道'' 从连续 admissible 对象空间 \(X_A\) 升维到
``钱学森障碍'' 的离散修复/正规形骨架 \(|T_A|\)，
才不是一句修辞，而是一个可追踪的对象级同伦链。
\end{remark}

\clearpage
\section{形式逻辑：可数超限离散统与 L大于等于3 纤维丛修复塔强版}
\label{sec:tex1-ch10}

\begin{remark}[本章意义备注：高超音速攻防的抽象语境]
本章的意义，在于把前一章得到的离散骨架继续推入可数超限与高层修复塔，
说明复杂攻防链中的修复可以具有严格的层级推进秩序。
对钱学森障碍而言，这一章刻画的是长期、多阶段不确定性如何被逐层吸收到 \(L_{\ge 3}\) 修复框架中。
\end{remark}

\subsection{严格主张、分层对象与引文定位}
\label{subsec:tex1-ch10-claim}

\begin{remark}[核心陈述]
本节把上一节的“有限离散骨架强版”再向前推进一层：
不是只说明“离散统下可以收束”，而是说明
\[
  \text{为何必须先经历连续统崩溃的语义收缩}
  \Longrightarrow
  \text{再进入 }L_{\ge 3}\text{ 超限递归纤维丛修复塔}.
\]
也就是说，连续统并不是被粗暴替代，而是先由连续统崩溃映射压到可数离散有效骨架，
再由超限递归语义收缩边缘算子在 \(L_{\ge 3}\) 上逐层吸收障碍，
最终把钱学森弹道的连续对象侧 \(X_A\) 与钱学森障碍的正规形侧 \(|T_A|\) 接成同一条同伦链。
\end{remark}

\begin{remark}[仓内文稿与外部背景]
本节直接承接 \cite{RepoTex41HomotopyRepairBundle} 中
\[
  \begin{aligned}
    L_1&:\text{离散（可数无穷）},\qquad
    L_2:\text{连续（不可数无穷）},\\
    L_3&:\omega\text{-递归编码首层},\qquad
    L_{\ge 3}:\text{在 }L_2\text{ 上逐层吸收障碍的离散修复纤维塔}.
  \end{aligned}
\]
这一对象链。
其中 \cite{RepoTex41HomotopyRepairBundle} 已经把“\(L_2\) 的连续几何”与“\(L_{\ge 3}\) 的离散修复语义”
组织为同一个纤维丛正规型；
\cite{RepoTex36HomotopyTwistEquivalence} 给出路径积分支撑类、正规形与稳定修复序列的对象级等价；
\cite{RepoTex38JacobiBianchiRepair} 给出缺陷严格下降、有限终止与终端正规形的接口。
外部拓扑背景仍以 \cite{Hatcher2002AT,Husemoller1994FibreBundles} 为准。
\end{remark}

\begin{remark}[严格边界]
本节采用的不是“纯离散拓扑 + 经典 \(I\)-同伦”，
而是
\[
  \boxed{
    L_1\text{ 可数索引}
    +L_2\text{ 连续基底}
    +L_{\ge 3}\text{ 超限递归修复塔}
    +\text{几何实现后的经典强变形收缩}
  }.
\]
因此，``钱学森弹道'' 仍指连续对象/路径表象侧，
``钱学森障碍'' 则指经过连续统崩溃与超限递归修复后得到的离散修复/正规形侧。
\end{remark}

\subsection{定义链：三层对象、连续统崩溃与超限递归语义收缩}
\label{subsec:tex1-ch10-defs}

\begin{definition}[离散统的三层对象]
\label{def:tex1-ch10-layers}
定义
\[
  L_1:=\NN
\]
为离散可数无穷索引层，
定义
\[
  L_2
\]
为连续本体/连续几何层，
定义
\[
  L_{\ge 3}
\]
为在 \(L_2\) 上逐层吸收孔洞障碍的离散修复纤维塔。

再定义两个语义投影：
\[
  \Pi_{l_2}:\mathsf{Law}_{\mathrm{Cont}}\to \mathsf{Law}_{\mathrm{Cont}}^{(l_2)},
  \qquad
  \Pi_{l_{\ge 3}}:\mathsf{Law}_{\mathrm{Cont}}\to \mathsf{Law}_{\mathrm{Cont}}^{(l_{\ge 3})}.
\]
其中 \(\Pi_{l_2}\) 记录 classical CH 在连续层上的外延阴影，
\(\Pi_{l_{\ge 3}}\) 记录经 \(L_3\) 编码与 \(L_{\ge 3}\) 修复后仍保留的同伦/曲率/语义结构。
\end{definition}

\begin{definition}[连续统崩溃的语义收缩]
\label{def:tex1-ch10-collapse}
设原始连续对象空间为
\[
  X_A.
\]
定义连续统崩溃的语义收缩映射
\[
  \Collapse_{\Sigma}:X_A\to V_A^\omega,
\]
其中
\[
  V_A^\omega=\{v_0,v_1,v_2,\dots\}
\]
是一个可数离散有效骨架。

并进一步定义：
\begin{enumerate}
  \item 商映射
  \[
    q_A:V_A^\omega\to \Pi_A;
  \]
  \item 缺陷高度函数
  \[
    \mu_A:V_A^\omega\to \omega_1;
  \]
  \item 零障碍正规形子集
  \[
    T_A:=\{v\in V_A^\omega:\mu_A(v)=0\}.
  \]
\end{enumerate}
在本节口径下，“连续由离散 breakdown/rewrite 事件链治理”的严格含义就是：
\(X_A\) 的有效修复并不在连续层上自治展开，而是经 \(\Collapse_{\Sigma}\) 压到
\(V_A^\omega\) 后，由可数离散事件链逐步执行。
\end{definition}

\begin{definition}[超限递归语义收缩边缘算子]
\label{def:tex1-ch10-operator}
定义边缘算子
\[
  \partial^-_{A,\Sigma}:V_A^\omega\to V_A^\omega
\]
并对任意 \(v\in V_A^\omega\) 定义超限递归序列
\[
  v^{(0)}:=v,
  \qquad
  v^{(\beta+1)}:=\partial^-_{A,\Sigma}\bigl(v^{(\beta)}\bigr),
  \qquad
  v^{(\lambda)}:=\operatorname*{colim}_{\beta<\lambda}v^{(\beta)}
\]
（其中 \(\lambda\) 为极限序数）。

要求它满足五条条件：
\begin{enumerate}
  \item \textbf{(E1) 类保持：}
  \[
    q_A(\partial^-_{A,\Sigma}(v))=q_A(v);
  \]
  \item \textbf{(E2) 后继步严格降障碍：}
  \[
    \mu_A(v)>0
    \Rightarrow
    \mu_A(\partial^-_{A,\Sigma}(v))<\mu_A(v);
  \]
  \item \textbf{(E3) 零障碍固定：}
  \[
    \mu_A(v)=0
    \Rightarrow
    \partial^-_{A,\Sigma}(v)=v;
  \]
  \item \textbf{(E4) 极限层相容：}
  \[
    q_A(v^{(\lambda)})=q_A(v),
    \qquad
    \mu_A(v^{(\lambda)})\le \inf_{\beta<\lambda}\mu_A(v^{(\beta)});
  \]
  \item \textbf{(E5) 同类稳定终点一致：}
  若 \(q_A(v_1)=q_A(v_2)\)，则在某个稳定序数阶段 \(\alpha_\ast\) 上有
  \[
    v_1^{(\alpha_\ast)}=v_2^{(\alpha_\ast)}.
  \]
\end{enumerate}
\end{definition}

\subsection{公理（降级为定理）：可数超限稳定阶段与共尾枚举存在}
\label{subsec:tex1-ch10-axiom}

\begin{theorem}[公理（降级为定理）：可数超限离散统的统一稳定序数与共尾枚举]
\label{thm:tex1-ch10-axiom-downgraded}
在定义~\ref{def:tex1-ch10-layers}--定义~\ref{def:tex1-ch10-operator} 的对象下，
若进一步满足：
\begin{enumerate}
  \item \(\Collapse_{\Sigma}\) 的有效 breakdown/rewrite 事件链是可数的；
  \item 对每个 \(v\in V_A^\omega\)，由 \((E1)\)--\((E4)\) 所定义的超限递归序列都良定义；
  \item 对每个 \(v\in V_A^\omega\)，存在某个序数 \(\beta_v<\omega_1\) 使
  \[
    v^{(\beta_v)}\in T_A;
  \]
  \item 同一商类中的稳定终点由 \((E5)\) 唯一确定。
\end{enumerate}
则存在一个统一稳定序数
\[
  \alpha_\ast:=\sup_{v\in V_A^\omega}\beta_v<\omega_1,
\]
并存在一个严格递增的可数共尾序列
\[
  0=\alpha_0<\alpha_1<\alpha_2<\cdots<\alpha_\ast,
  \qquad
  \sup_{m\in\NN}\alpha_m=\alpha_\ast,
\]
使得对任意 \(v\in V_A^\omega\)，都有
\[
  v^{(\alpha_\ast)}\in T_A,
  \qquad
  q_A(v^{(\alpha_\ast)})=q_A(v).
\]
\end{theorem}

\begin{Certificate}[数学归纳法证书：按共尾前缀长度的稳定前缀证书]
\label{cert:tex1-ch10-axiom-downgraded}
固定任意 \(v\in V_A^\omega\)，定义谓词
\[
  P_v(m)
  :
  \Longleftrightarrow
  \text{前缀 }
  v^{(\alpha_0)},v^{(\alpha_1)},\dots,v^{(\alpha_m)}
  \text{ 全部良定义，且始终保持在 }q_A(v)\text{ 对应的商类内。}
\]
证书分两步：
\begin{enumerate}
  \item \textbf{基础步 \(P_v(0)\)：}
  \[
    v^{(\alpha_0)}=v
  \]
  显然良定义；
  \item \textbf{归纳步 \(P_v(m)\Rightarrow P_v(m+1)\)：}
  由 \((E1)\)--\((E4)\)，从 \(v^{(\alpha_m)}\) 到 \(v^{(\alpha_{m+1})}\) 的后继/极限推进都类保持且良定义。
\end{enumerate}
再配合序数良基性与 \(V_A^\omega\) 的可数性，就得到统一稳定序数 \(\alpha_\ast\) 与可数共尾枚举。
\end{Certificate}

\begin{proof}[证明（数学归纳法证书；基于数学符号推演逻辑的谓词因果链）]
先固定任意 \(v\in V_A^\omega\)。
由假设 2 与定义~\ref{def:tex1-ch10-operator}，
\[
  v^{(0)}=v,\qquad
  v^{(\beta+1)}=\partial^-_{A,\Sigma}(v^{(\beta)}),\qquad
  v^{(\lambda)}=\operatorname*{colim}_{\beta<\lambda}v^{(\beta)}
\]
在所有 \(\beta<\omega_1\) 上都良定义。

对共尾序列的前缀长度做数学归纳。

\textbf{基础步：}当 \(m=0\) 时，
\[
  v^{(\alpha_0)}=v
\]
良定义，且显然属于商类 \(q_A(v)\)。

\textbf{归纳步：}若 \(P_v(m)\) 成立，则 \(v^{(\alpha_m)}\) 已良定义且满足
\[
  q_A(v^{(\alpha_m)})=q_A(v).
\]
若 \(\alpha_{m+1}=\alpha_m+1\)，则由 \((E1)\) 有
\[
  q_A(v^{(\alpha_{m+1})})
  =
  q_A(\partial^-_{A,\Sigma}(v^{(\alpha_m)}))
  =
  q_A(v^{(\alpha_m)})
  =
  q_A(v).
\]
若 \(\alpha_{m+1}\) 为极限序数，则由 \((E4)\) 有
\[
  q_A(v^{(\alpha_{m+1})})=q_A(v).
\]
故 \(P_v(m+1)\) 成立。
由数学归纳法，\(P_v(m)\) 对所有 \(m\in\NN\) 成立。

再证每个 \(v\) 都在某个 \(\beta_v<\omega_1\) 上稳定到零障碍。
对序数 \(\mu_A(v)\) 做良基归纳：
若 \(\mu_A(v)=0\)，取 \(\beta_v=0\)；
若 \(\mu_A(v)>0\)，则由 \((E2)\)
\[
  \mu_A(\partial^-_{A,\Sigma}(v))<\mu_A(v),
\]
故对更小序数应用归纳假设，存在 \(\gamma<\omega_1\) 使
\[
  \bigl(\partial^-_{A,\Sigma}(v)\bigr)^{(\gamma)}\in T_A.
\]
于是取 \(\beta_v:=\gamma+1\) 即得
\[
  v^{(\beta_v)}\in T_A.
\]

由于 \(V_A^\omega\) 可数，集合
\[
  \{\beta_v:v\in V_A^\omega\}
\]
是可数个可数序数的集合，因此
\[
  \alpha_\ast:=\sup_{v\in V_A^\omega}\beta_v<\omega_1.
\]
任一可数序数都存在一个 \(\NN\)-索引的共尾递增枚举，
故存在
\[
  0=\alpha_0<\alpha_1<\alpha_2<\cdots<\alpha_\ast,
  \qquad
  \sup_{m\in\NN}\alpha_m=\alpha_\ast.
\]
最后，由 \(q_A\) 的类保持性与 \((E5)\)，得到对任意 \(v\in V_A^\omega\)，
\[
  v^{(\alpha_\ast)}\in T_A,
  \qquad
  q_A(v^{(\alpha_\ast)})=q_A(v).
\]
定理得证。
\end{proof}

\subsection{定理链：正规化截面、有限改变引理与进入 L大于等于3修复塔的前置性}
\label{subsec:tex1-ch10-theorems}

\begin{theorem}[超限正规化截面存在定理]
\label{thm:tex1-ch10-section}
在定理~\ref{thm:tex1-ch10-axiom-downgraded} 的条件下，定义
\[
  s_A:\Pi_A\to T_A,
  \qquad
  s_A([v]):=v^{(\alpha_\ast)}.
\]
则 \(s_A\) 良定义，并且
\[
  q_A\circ s_A=\id_{\Pi_A},
  \qquad
  s_A=\operatorname{NF}_A.
\]
\end{theorem}

\begin{Certificate}[证书：统一稳定阶段 + 同类终点一致]
\label{cert:tex1-ch10-section}
证书分三条：
\begin{enumerate}
  \item 定理~\ref{thm:tex1-ch10-axiom-downgraded} 已给出 \(v^{(\alpha_\ast)}\in T_A\)；
  \item \((E5)\) 保证同一商类上的终点一致；
  \item \cite{RepoTex36HomotopyTwistEquivalence} 中 \(\operatorname{NF}_A\) 的角色正是“同类唯一正规形代表”。
\end{enumerate}
\end{Certificate}

\begin{proof}[证明（基于数学符号推演逻辑的谓词因果链）]
任取 \(x=[v]\in\Pi_A\)。
由定理~\ref{thm:tex1-ch10-axiom-downgraded}，
\[
  v^{(\alpha_\ast)}\in T_A.
\]
若 \(v_1,v_2\) 满足 \(q_A(v_1)=q_A(v_2)\)，则由 \((E5)\)
\[
  v_1^{(\alpha_\ast)}=v_2^{(\alpha_\ast)}.
\]
故 \(s_A\) 与代表元无关，良定义。

又由类保持性，
\[
  q_A(s_A([v]))
  =
  q_A(v^{(\alpha_\ast)})
  =
  q_A(v)
  =
  [v].
\]
因此
\[
  q_A\circ s_A=\id_{\Pi_A}.
\]

最后，由 \cite{RepoTex36HomotopyTwistEquivalence}，
\(\operatorname{NF}_A\) 的定义角色是对每个路径积分支撑类选取唯一正规形代表；
而 \(s_A\) 把每个类 \([v]\) 送到其唯一稳定终点 \(v^{(\alpha_\ast)}\)。
因此
\[
  s_A=\operatorname{NF}_A.
\]
定理得证。
\end{proof}

\begin{lemma}[有限改变引理]
\label{lem:tex1-ch10-finite-change}
对任意固定的 \(v\in V_A^\omega\)，存在某个有限整数 \(m(v)\in\NN\)，使得对所有 \(m\ge m(v)\)，都有
\[
  v^{(\alpha_m)}=v^{(\alpha_\ast)}.
\]
换言之，虽然系统层面的稳定阶段是可数超限的，
但每个具体顶点沿共尾事件链只会发生有限次真实改变。
\end{lemma}

\begin{Certificate}[证书：序数不存在 \(\omega\)-长严格下降链]
\label{cert:tex1-ch10-finite-change}
证书只用两条：
\begin{enumerate}
  \item 若 \(\mu_A(v^{(\alpha_m)})>0\)，则真实改变一步就会严格降低序数高度；
  \item 序数不存在以 \(\NN\) 为指标的无限严格下降链。
\end{enumerate}
\end{Certificate}

\begin{proof}[证明（基于数学符号推演逻辑的谓词因果链）]
反设不存在这样的 \(m(v)\)。
则对任意 \(m\in\NN\)，都存在 \(n\ge m\) 使
\[
  v^{(\alpha_n)}\neq v^{(\alpha_\ast)}.
\]
特别地，可以抽出一个无限子序列
\[
  v^{(\alpha_{m_0})},v^{(\alpha_{m_1})},v^{(\alpha_{m_2})},\dots
\]
使其中每一步都是真实改变。
由 \((E2)\)，每次真实改变都会使缺陷高度严格下降：
\[
  \mu_A(v^{(\alpha_{m_0})})
  >
  \mu_A(v^{(\alpha_{m_1})})
  >
  \mu_A(v^{(\alpha_{m_2})})
  >
  \cdots.
\]
这给出了一条以 \(\NN\) 为指标的无限严格下降序数链，
与序数良基性矛盾。
故存在有限整数 \(m(v)\) 使对所有 \(m\ge m(v)\)，
\[
  v^{(\alpha_m)}=v^{(\alpha_\ast)}.
\]
引理得证。
\end{proof}

\begin{proposition}[连续统崩溃是进入 \(L_{\ge 3}\) 修复塔的构造前置]
\label{prop:tex1-ch10-prior}
在 \cite{RepoTex41HomotopyRepairBundle} 的对象口径下，
\(L_{\ge 3}\) 修复塔的层对象由
\[
  (x,\sigma_{\le n},r_n)
\]
这类“连续基底点 + 可数编码前缀 + 离散修复标签”的三元组构成。
因此，若不先给出连续统崩溃/编码映射
\[
  \Collapse_{\Sigma}:X_A\to V_A^\omega,
\]
则无法仅在原始连续对象空间 \(X_A\) 上直接定义：
\begin{enumerate}
  \item 商映射 \(q_A:V_A^\omega\to\Pi_A\)；
  \item 超限递归边缘算子 \(\partial^-_{A,\Sigma}:V_A^\omega\to V_A^\omega\)；
  \item 由可数事件链索引的统一共尾修复序列 \((\alpha_m)_{m\in\NN}\)。
\end{enumerate}
也就是说，连续统崩溃不是附加修辞，而是进入 \(L_{\ge 3}\) 修复塔的构造前置。
\end{proposition}

\begin{Certificate}[证书：对象域不匹配]
\label{cert:tex1-ch10-prior}
证书只用两条：
\begin{enumerate}
  \item \cite{RepoTex41HomotopyRepairBundle} 中 \(L_{\ge 3}\) 的层对象需要离散编码前缀 \(\sigma_{\le n}\) 与修复标签 \(r_n\)；
  \item 这些对象不属于原始连续空间 \(X_A\) 的内禀坐标，而属于经编码/崩溃后得到的离散语义骨架。
\end{enumerate}
\end{Certificate}

\begin{proof}[证明（基于数学符号推演逻辑的谓词因果链）]
由 \cite{RepoTex41HomotopyRepairBundle} 的定义链，
\(L_{\ge 3}\) 的第 \(n\) 层对象写成
\[
  (x,\sigma_{\le n},r_n),
\]
其中 \(x\in L_2\) 是连续基底点，
\(\sigma_{\le n}\) 是可数编码前缀，
\(r_n\) 是离散修复标签。
因此，超限递归修复塔并不是构造在“裸连续对象 \(x\)”上，
而是构造在“连续基底点 + 离散编码/修复语义”这类混合对象上。

若不给出 \(\Collapse_{\Sigma}:X_A\to V_A^\omega\)，
则 \(X_A\) 中没有内部定义好的可数编码顶点集，
也就无法在对象层上写出
\[
  q_A,\qquad \partial^-_{A,\Sigma},\qquad (\alpha_m)_{m\in\NN}.
\]
故连续统崩溃是进入 \(L_{\ge 3}\) 修复塔的构造前置。
\end{proof}

\subsection{定理：可数超限离散统版强变形收缩与连续对象同伦等价}
\label{subsec:tex1-ch10-retract}

\begin{theorem}[可数超限离散统版强变形收缩]
\label{thm:tex1-ch10-retract}
取定理~\ref{thm:tex1-ch10-axiom-downgraded} 的可数共尾序列
\[
  0=\alpha_0<\alpha_1<\alpha_2<\cdots<\alpha_\ast,
  \qquad
  \sup_{m\in\NN}\alpha_m=\alpha_\ast,
\]
并定义时间节点
\[
  t_m:=1-2^{-m}\in[0,1).
\]
再定义图 \(G_A^\omega\)：
其顶点集为 \(V_A^\omega\)，
若对某个 \(v\in V_A^\omega\) 与某个 \(m\in\NN\) 有
\[
  v^{(\alpha_m)}\neq v^{(\alpha_{m+1})},
\]
则在 \(v^{(\alpha_m)}\) 与 \(v^{(\alpha_{m+1})}\) 之间连一条边。
记其几何实现为 \(|V_A^\omega|\)，记 \(T_A\) 所诱导的子图几何实现为 \(|T_A|\)。

定义
\[
  r_A:=s_A\circ q_A:V_A^\omega\to T_A.
\]
则存在连续映射
\[
  H_A:|V_A^\omega|\times I\to |V_A^\omega|
\]
满足
\[
  H_A(-,0)=\id_{|V_A^\omega|},
  \qquad
  H_A(-,1)=i_A\circ r_A,
\]
并且
\[
  H_A(\tau,t)=\tau
  \qquad
  (\tau\in |T_A|),
\]
其中 \(i_A:|T_A|\hookrightarrow |V_A^\omega|\) 是包含映射。
因此
\[
  \boxed{
    |V_A^\omega|\searrow |T_A|.
  }
\]
\end{theorem}

\begin{Certificate}[证书：有限改变 + 时间压缩 + 分段线性插值]
\label{cert:tex1-ch10-retract}
证书分三段：
\begin{enumerate}
  \item 由引理~\ref{lem:tex1-ch10-finite-change}，每个顶点沿 \((\alpha_m)\) 只发生有限次真实改变；
  \item 因而顶点轨道在 \(t\to 1\) 时最终常值于 \(v^{(\alpha_\ast)}\)；
  \item 对每个时间区间 \([t_m,t_{m+1}]\)，都可以沿相邻顶点之间的边做线性插值。
\end{enumerate}
\end{Certificate}

\begin{proof}[证明（基于数学符号推演逻辑的谓词因果链）]
先在顶点集上定义
\[
  H_A(v,t_m):=v^{(\alpha_m)},
  \qquad
  H_A(v,1):=v^{(\alpha_\ast)}=s_A(q_A(v)).
\]
对任意固定顶点 \(v\)，由引理~\ref{lem:tex1-ch10-finite-change}，
存在有限 \(m(v)\) 使得对所有 \(m\ge m(v)\)，
\[
  v^{(\alpha_m)}=v^{(\alpha_\ast)}.
\]
故顶点轨道在 \(t\to 1\) 时最终常值，因而在顶点上连续。

再对每个时间区间 \([t_m,t_{m+1}]\)，
若 \(v^{(\alpha_m)}=v^{(\alpha_{m+1})}\)，则令 \(H_A(v,t)\) 在该区间上常值；
若二者不同，则由图 \(G_A^\omega\) 的定义，
它们之间有一条边，可在该边线段上作线性插值。
把这些分段线性插值在所有顶点上拼接，并沿单形线性延拓到几何实现，
就得到连续映射
\[
  H_A:|V_A^\omega|\times I\to |V_A^\omega|.
\]

当 \(t=0\) 时，
\[
  H_A(v,0)=v^{(\alpha_0)}=v,
\]
故
\[
  H_A(-,0)=\id_{|V_A^\omega|}.
\]
当 \(t=1\) 时，
\[
  H_A(v,1)=v^{(\alpha_\ast)}=s_A(q_A(v))=r_A(v),
\]
故
\[
  H_A(-,1)=i_A\circ r_A.
\]
若 \(\tau\in T_A\)，则由 \((E3)\) 对任意 \(\beta\le \alpha_\ast\) 都有
\[
  \tau^{(\beta)}=\tau,
\]
因此
\[
  H_A(\tau,t)=\tau.
\]
所以 \(H_A\) 是从 \(|V_A^\omega|\) 到 \(|T_A|\) 的强变形收缩。
\end{proof}

\begin{theorem}[更强版：从连续对象到正规形的经典同伦等价]
\label{thm:tex1-ch10-collapse-equivalence}
若存在连续映射
\[
  \Collapse_{\Sigma}:X_A\to |V_A^\omega|,
  \qquad
  j_A:|V_A^\omega|\hookrightarrow X_A
\]
满足
\[
  \Collapse_{\Sigma}\circ j_A=\id_{|V_A^\omega|},
  \qquad
  j_A\circ \Collapse_{\Sigma}\simeq_h \id_{X_A},
\]
则
\[
  \boxed{
    X_A \simeq_h |V_A^\omega| \simeq_h |T_A|.
  }
\]
从而
\[
  \boxed{
    X_A \simeq_h |T_A|.
  }
\]
\end{theorem}

\begin{Certificate}[证书：连续统崩溃同伦逆 + 可数超限强收缩]
\label{cert:tex1-ch10-collapse-equivalence}
证书分两段：
\begin{enumerate}
  \item \(\Collapse_{\Sigma}\) 与 \(j_A\) 给出 \(X_A\) 与 \(|V_A^\omega|\) 的同伦逆；
  \item 定理~\ref{thm:tex1-ch10-retract} 已给出
  \[
    |V_A^\omega|\searrow |T_A|.
  \]
\end{enumerate}
\end{Certificate}

\begin{proof}[证明（基于数学符号推演逻辑的谓词因果链）]
由
\[
  \Collapse_{\Sigma}\circ j_A=\id_{|V_A^\omega|}
\]
与
\[
  j_A\circ \Collapse_{\Sigma}\simeq_h \id_{X_A},
\]
可知
\[
  X_A\simeq_h |V_A^\omega|.
\]
再由定理~\ref{thm:tex1-ch10-retract}，
\[
  |V_A^\omega|\searrow |T_A|,
\]
故
\[
  |V_A^\omega|\simeq_h |T_A|.
\]
同伦等价可传递，因此
\[
  X_A\simeq_h |T_A|.
\]
定理得证。
\end{proof}

\begin{corollary}[“离散统下通过连续统崩溃的语义收缩边缘算子修复而收束”的严格落点]
\label{cor:tex1-ch10-final}
设钱学森弹道的连续对象侧由 \(X_A\) 表示，
钱学森障碍的正规形侧由 \(|T_A|\) 表示。
则用户命题
\[
  \text{“在离散统下是可以收束的，即通过连续统崩溃的语义收缩边缘算子修复”}
\]
在本节中的严格落点就是下面这条定理链：
\[
  \begin{aligned}
    L_1\text{ 离散索引}
    &\Longrightarrow
    L_2\text{ 连续阴影}
    \Longrightarrow
    \text{continuous breakdown 的语义收缩}\\
    &\Longrightarrow
    L_{\ge 3}\text{ 超限递归纤维丛修复塔}
    \Longrightarrow
    T_A\text{ 正规形}.
  \end{aligned}
\]
并且最终得到
\[
  \boxed{
    \begin{aligned}
      X_A&\xrightarrow{\ \Collapse_{\Sigma}\ }|V_A^\omega|
      \searrow |T_A|,\\
      X_A&\simeq_h |T_A|.
    \end{aligned}
  }
\]
\end{corollary}

\begin{Certificate}[证书：统一稳定序数 + 超限强收缩 + 连续统同伦逆]
\label{cert:tex1-ch10-final}
证书分三段：
\begin{enumerate}
  \item 定理~\ref{thm:tex1-ch10-axiom-downgraded} 与定理~\ref{thm:tex1-ch10-section} 给出统一稳定序数 \(\alpha_\ast\) 与正规化截面 \(s_A\)；
  \item 定理~\ref{thm:tex1-ch10-retract} 给出 \(|V_A^\omega|\searrow |T_A|\)；
  \item 定理~\ref{thm:tex1-ch10-collapse-equivalence} 给出 \(X_A\simeq_h |T_A|\)。
\end{enumerate}
\end{Certificate}

\begin{proof}[证明（基于数学符号推演逻辑的谓词因果链）]
由定理~\ref{thm:tex1-ch10-axiom-downgraded}，
连续统崩溃后的可数离散骨架 \(V_A^\omega\) 上存在统一稳定序数 \(\alpha_\ast\) 与共尾修复序列 \((\alpha_m)\)。
由定理~\ref{thm:tex1-ch10-section}，
该稳定阶段定义出正规化截面
\[
  s_A=\operatorname{NF}_A.
\]
再由定理~\ref{thm:tex1-ch10-retract}，
\[
  |V_A^\omega|\searrow |T_A|.
\]
若进一步加入 \(\Collapse_{\Sigma}\) 与 \(j_A\) 的同伦逆接口，
则由定理~\ref{thm:tex1-ch10-collapse-equivalence}，
\[
  X_A\simeq_h |T_A|.
\]
故推论成立。
\end{proof}

\begin{remark}[这一版比上一节的有限步语义收缩更强]
上一节的中心链条是
\[
  \text{有限步严格降缺陷}
  \Longrightarrow
  \text{有限离散骨架上的强变形收缩}.
\]
本节则把它升级为
\[
  \begin{aligned}
    \text{可数 breakdown/rewrite 事件链}
    &\Longrightarrow
    \text{统一稳定序数 }\alpha_\ast\\
    &\Longrightarrow
    \text{超限递归修复塔}
    \Longrightarrow
    \text{稳定阶段正规形}.
  \end{aligned}
\]
因此，这不是简单补充，而是把“有限步版本”推进成“可数超限版本”。
\end{remark}

\begin{remark}[最压缩的一句话]
最压缩地说：
\[
  \boxed{
    \text{上一节给出的是“离散强版”的有限轮廓，}
  }
\]
\[
  \boxed{
    \text{本节则把它提升为“离散统 -- 超限递归 -- 纤维丛修复塔”强版。}
  }
\]
而把这条桥梁真正立起来的关键，
正是
\[
  \Collapse_{\Sigma}
  \quad+\quad
  \partial^-_{A,\Sigma}
  \quad+\quad
  (\alpha_m)_{m\in\NN}
  \quad+\quad
  L_{\ge 3}\text{ 纤维丛修复塔}.
\]
\end{remark}

%%%%%%%%%%%%%%%%%%%%%%%%%%%%%%%%%%%%%%%%%%%%%%%%%%%%%%%%%%%%
\clearpage
\section{形式逻辑：L大于等于3 的双层修复职责、完整性修复与完备性修复}
\label{sec:tex1-ch11}
%%%%%%%%%%%%%%%%%%%%%%%%%%%%%%%%%%%%%%%%%%%%%%%%%%%%%%%%%%%%

\begin{remark}[本章意义备注：高超音速攻防的抽象语境]
本章的意义，在于明确完整性修复与完备性修复不是同一件事，且次序不可交换。
对抽象攻防研究而言，
这相当于区分“先把关键来源补齐”与“再把推断覆盖补满”两类职责，从而避免把不同性质的失败混为一谈。
\end{remark}

\subsection{严格主张、引文定位与术语边界}
\label{subsec:tex1-ch11-claim}

\begin{remark}[核心陈述]
本节把上一节关于 \(L_{\ge 3}\) 修复塔的结论再收紧一次：
\[
  \boxed{
    L_{\ge 3}\text{ 的修复不是单一的 }L_2\text{ 孔洞修复，}
  }
\]
\[
  \boxed{
    \text{而是“}(L_1\to L_2)\text{ 的完整性修复”与“}L_2\text{ 自身的完备性修复”同时成立。}
  }
\]
在这一口径下，钱学森弹道仍对应 \(L_2\) 上的连续对象/连续表象侧，
而钱学森障碍则对应“离散来源链未闭合”或“连续基底内部孔洞未被纤维平凡化”这两类失败的并集。
\end{remark}

\begin{remark}[仓内文稿定位]
本节直接承接 \cite{RepoTex41HomotopyRepairBundle} 中关于
\[
  L_1,\qquad L_2,\qquad L_3,\qquad L_{>3}
\]
的分层职责：其中 \(L_3\) 被明确定义为“双可数离散对对 \(L_2\) 的首层 \(\omega\)-递归编码落点”，
而 \(L_{>3}\) 被定义为“以 \(L_2\) 为基底、沿纤维吸收障碍并恢复全局可行截面的高层修复塔”。
与之配套，\cite{RepoTex36HomotopyTwistEquivalence} 提供正规化截面与稳定修复序列的对象级等价接口，
\cite{RepoTex38JacobiBianchiRepair} 提供缺陷严格下降、有限终止与终端正规形的离散证书接口。
\end{remark}

\begin{remark}[严格边界]
本节不把 \(L_3\) 误读为“可有可无的技术前奏”，也不把 \(L_{>3}\) 误读为“单一孔洞修补器”。
严格地说，
\[
  L_3\text{ 负责把 }L_1\text{ 的离散证书桥接到 }L_2\text{ 的连续基底上，}
\]
\[
  L_{>3}\text{ 负责把 }L_2\text{ 基底内部的孔洞障碍沿纤维方向吸收并平凡化。}
\]
因此，\(L_{\ge 3}=L_3\cup L_{>3}\) 的总职责必定是双层修复，而不是单层修复。
\end{remark}

\subsection{定义链：完整性谓词、完备性谓词与双层修复记号}
\label{subsec:tex1-ch11-defs}

\begin{definition}[完整性修复谓词]
\label{def:tex1-ch11-integrity}
对任意 \(x\in L_2\)，定义 \((L_1\to L_2)\) 的完整性谓词为
\[
  \mathbf{Int}_{1\to2}(x)
  \;:\Longleftrightarrow\;
  \exists \sigma\in \mathcal{C}_\omega,\quad \operatorname{Enc}(\sigma)=x.
\]
它表达的不是“\(x\) 作为连续对象存在”这种弱断言，
而是“\(x\) 具有一条可审计的离散来源链，可追溯到 \(L_1\) 与双离散对编码”的强断言。
在本节术语中，这一层对应的是数域完整性。
\end{definition}

\begin{definition}[完备性修复谓词]
\label{def:tex1-ch11-completeness}
对任意开集 \(U\subseteq L_2\)，若存在孔洞障碍类
\[
  [\omega_{\mathrm{obs}}(U)]\neq 0\in H^{m+1}(U;A),
\]
则称
\[
  \Hole(U)
  \;:\Longleftrightarrow\;
  [\omega_{\mathrm{obs}}(U)]\neq 0
\]
成立。

进一步定义 \(L_2\) 层的完备性修复谓词
\[
  \mathbf{Cmp}_{2}(U)
  \;:\Longleftrightarrow\;
  \pi_\infty^\ast[\omega_{\mathrm{obs}}(U)]=0
  \;\wedge\;
  \exists s_\infty:U\to \mathcal E_\infty.
\]
它的含义是：原本在 \(L_2\) 基底上不可延拓的孔洞障碍，
在极限修复塔总空间中被拉回平凡化，并且恢复全局可行截面。
在本节术语中，这一层对应的是数域完备性。
\end{definition}

\begin{definition}[双层修复语义]
\label{def:tex1-ch11-repair}
定义三种修复语义记号：
\[
  \operatorname{Rep}^{\mathrm{int}}_{1\to2}
  \;:\Longleftrightarrow\;
  \forall x\in L_2,\ \mathbf{Int}_{1\to2}(x),
\]
\[
  \operatorname{Rep}^{\mathrm{cmp}}_{2}
  \;:\Longleftrightarrow\;
  \forall U\subseteq L_2,\ \Hole(U)\Rightarrow \mathbf{Cmp}_{2}(U),
\]
\[
  \operatorname{Rep}_{\ge3}
  \;:\Longleftrightarrow\;
  \operatorname{Rep}^{\mathrm{int}}_{1\to2}
  \wedge
  \operatorname{Rep}^{\mathrm{cmp}}_{2}.
\]
其中 \(\operatorname{Rep}^{\mathrm{int}}_{1\to2}\) 表示首层编码桥所实现的完整性修复，
\(\operatorname{Rep}^{\mathrm{cmp}}_{2}\) 表示高层纤维吸收所实现的完备性修复，
\(\operatorname{Rep}_{\ge3}\) 则表示二者并取的总职责。
\end{definition}

\subsection{公理（降级为定理）：L大于等于3 双层职责可按层前缀递归钉死}
\label{subsec:tex1-ch11-axiom}

\begin{theorem}[公理（降级为定理）：\(L_{\ge 3}\) 的双层职责可按层前缀递归钉死]
\label{thm:tex1-ch11-axiom-downgraded}
定义层前缀
\[
  \mathfrak P_2:=(L_1,L_2),\qquad
  \mathfrak P_3:=(L_1,L_2,L_3),\qquad
  \mathfrak P_4:=(L_1,L_2,L_3,L_{>3}).
\]
再定义谓词 \(P(n)\)：
\[
  P(n)
  \;:\Longleftrightarrow\;
  \text{前缀 }\mathfrak P_n\text{ 所新增并钉死的最高修复职责已被唯一识别。}
\]
则在 \cite{RepoTex41HomotopyRepairBundle} 的对象链下有：
\begin{enumerate}
  \item \(P(2)\) 成立，且 \(\mathfrak P_2\) 只钉死 \(L_2\) 作为唯一连续基底；
  \item \(P(3)\) 成立，且 \(\mathfrak P_3\) 所新增并钉死的职责正是
  \[
    \operatorname{Rep}^{\mathrm{int}}_{1\to2};
  \]
  \item \(P(4)\) 成立，且 \(\mathfrak P_4\) 所新增并钉死的职责正是
  \[
    \operatorname{Rep}^{\mathrm{cmp}}_{2}.
  \]
\end{enumerate}
因此
\[
  \boxed{
    \operatorname{Rep}_{\ge3}
    =
    \operatorname{Rep}^{\mathrm{int}}_{1\to2}
    \wedge
    \operatorname{Rep}^{\mathrm{cmp}}_{2}.
  }
\]
\end{theorem}

\begin{Certificate}[数学归纳法证书：层前缀职责的递归分离证书]
\label{cert:tex1-ch11-axiom-downgraded}
对前缀长度 \(n\in\{2,3,4\}\) 做数学归纳：
\begin{enumerate}
  \item \textbf{基步 \(P(2)\)：}由 \cite{RepoTex41HomotopyRepairBundle}，\((L_1,L_2)\) 只固定“离散索引层 + 连续基底层”的前缀，尚未引入编码桥与高阶修复；
  \item \textbf{归纳步 \(P(2)\Rightarrow P(3)\)：}由同文的首层 \(\omega\)-递归编码定理，\(L_3\) 新增的恰是 \((L_1\to L_2)\) 的完整性修复；
  \item \textbf{归纳步 \(P(3)\Rightarrow P(4)\)：}由同文的修复塔定理，\(L_{>3}\) 新增的恰是 \(L_2\) 基底内部孔洞的完备性修复。
\end{enumerate}
由此前缀链分离，可得 \(L_{\ge3}\) 的总职责正是两类修复的合取。
\end{Certificate}

\begin{proof}[证明（数学归纳法证书；基于数学符号推演逻辑的谓词因果链）]
\textbf{基步 \(P(2)\)。}
由 \cite{RepoTex41HomotopyRepairBundle}，
\[
  L_1
\]
只承担可数离散索引与证书槽，
\[
  L_2
\]
才承担连续基底、路径与障碍类。
故在前缀
\[
  \mathfrak P_2=(L_1,L_2)
\]
中，被钉死的只是“\(L_2\) 为唯一连续基底”这一对象角色，
而
\[
  \operatorname{Rep}^{\mathrm{int}}_{1\to2},
  \qquad
  \operatorname{Rep}^{\mathrm{cmp}}_{2}
\]
尚未出现，因此 \(P(2)\) 成立。

\textbf{归纳步 \(P(2)\Rightarrow P(3)\)。}
由 \cite{RepoTex41HomotopyRepairBundle} 的首层 \(\omega\)-递归编码定理，
\[
  \operatorname{Enc}:\mathcal C_\omega\to L_2
\]
为满射，且
\[
  L_3
\]
正是双可数离散对对 \(L_2\) 的首层编码落点。
因此，对任意 \(x\in L_2\)，都存在
\[
  \sigma\in\mathcal C_\omega
\]
使
\[
  \operatorname{Enc}(\sigma)=x,
\]
也即
\[
  \forall x\in L_2,\ \mathbf{Int}_{1\to2}(x).
\]
故 \(\mathfrak P_3\) 新增并钉死的职责正是
\[
  \operatorname{Rep}^{\mathrm{int}}_{1\to2}.
\]
于是 \(P(3)\) 成立。

\textbf{归纳步 \(P(3)\Rightarrow P(4)\)。}
再由 \cite{RepoTex41HomotopyRepairBundle} 的修复塔定理，
若某开集 \(U\subseteq L_2\) 满足
\[
  \Hole(U),
\]
则在极限修复塔
\[
  \pi_\infty:\mathcal E_\infty\to U
\]
上有
\[
  \pi_\infty^\ast[\omega_{\mathrm{obs}}(U)]=0,
  \qquad
  \exists s_\infty:U\to\mathcal E_\infty.
\]
也即
\[
  \forall U\subseteq L_2,\ \Hole(U)\Rightarrow \mathbf{Cmp}_2(U).
\]
因此 \(\mathfrak P_4\) 新增并钉死的职责正是
\[
  \operatorname{Rep}^{\mathrm{cmp}}_{2}.
\]
于是 \(P(4)\) 成立。

由基步与两次归纳步，
\[
  \operatorname{Rep}_{\ge3}
  =
  \operatorname{Rep}^{\mathrm{int}}_{1\to2}
  \wedge
  \operatorname{Rep}^{\mathrm{cmp}}_{2}.
\]
定理得证。
\end{proof}

\subsection{定理链：完整性修复、完备性修复与次序不可交换}
\label{subsec:tex1-ch11-theorems}

\begin{lemma}[首层 \(\omega\)-递归编码桥给出完整性修复]
\label{lem:tex1-ch11-integrity}
若
\[
  \operatorname{Enc}:\mathcal C_\omega\to L_2
\]
为满射，且 \(L_3\) 是其首层 \(\omega\)-递归编码落点，
则
\[
  \forall x\in L_2,\ \mathbf{Int}_{1\to2}(x).
\]
换言之，
\[
  \operatorname{Rep}^{\mathrm{int}}_{1\to2}
\]
成立。
\end{lemma}

\begin{Certificate}[证书：满射编码证书]
\label{cert:tex1-ch11-integrity}
证书只有一条：
\[
  \operatorname{Enc}:\mathcal C_\omega\to L_2
\]
为满射。
因此每个 \(x\in L_2\) 都有至少一个可审计离散编码源 \(\sigma\)。
\end{Certificate}

\begin{proof}[证明（基于数学符号推演逻辑的谓词因果链）]
任取 \(x\in L_2\)。
由满射性，存在某个
\[
  \sigma\in\mathcal C_\omega
\]
使
\[
  \operatorname{Enc}(\sigma)=x.
\]
于是由定义~\ref{def:tex1-ch11-integrity}，
\[
  \mathbf{Int}_{1\to2}(x)
\]
成立。
由于 \(x\in L_2\) 任意，
\[
  \forall x\in L_2,\ \mathbf{Int}_{1\to2}(x).
\]
故
\[
  \operatorname{Rep}^{\mathrm{int}}_{1\to2}
\]
成立。引理得证。
\end{proof}

\begin{theorem}[\(L_{\ge3}\) 修复的双重职责定理]
\label{thm:tex1-ch11-dual}
在本文体系内，若同时满足：
\begin{enumerate}
  \item \(L_3\) 是双可数离散对对 \(L_2\) 的首层 \(\omega\)-递归编码落点，且
  \[
    \operatorname{Enc}:\mathcal C_\omega\to L_2
  \]
  为满射；
  \item 对任意开集 \(U\subseteq L_2\)，若 \(\Hole(U)\) 成立，则在极限修复塔上有
  \[
    \pi_\infty^\ast[\omega_{\mathrm{obs}}(U)]=0,
    \qquad
    \exists s_\infty:U\to\mathcal E_\infty;
  \]
\end{enumerate}
则
\[
  \boxed{
    \operatorname{Rep}_{\ge3}
    \Longrightarrow
    \Bigl(\forall x\in L_2,\ \mathbf{Int}_{1\to2}(x)\Bigr)
    \wedge
    \Bigl(\forall U\subseteq L_2,\ \Hole(U)\Rightarrow \mathbf{Cmp}_{2}(U)\Bigr).
  }
\]
亦即：
\[
  \boxed{
    \operatorname{Rep}_{\ge3}
    =
    \operatorname{Rep}^{\mathrm{int}}_{1\to2}
    \wedge
    \operatorname{Rep}^{\mathrm{cmp}}_{2}.
  }
\]
\end{theorem}

\begin{Certificate}[证书：满射编码桥 + 障碍拉回平凡化 + 全局截面恢复]
\label{cert:tex1-ch11-dual}
证书分两段：
\begin{enumerate}
  \item \(\operatorname{Enc}\) 的满射性给出完整性修复证书；
  \item \(\pi_\infty^\ast[\omega_{\mathrm{obs}}(U)]=0\) 与 \(s_\infty\) 的存在给出完备性修复证书。
\end{enumerate}
\end{Certificate}

\begin{proof}[证明（基于数学符号推演逻辑的谓词因果链）]
先证完整性部分。
由引理~\ref{lem:tex1-ch11-integrity}，
\[
  \forall x\in L_2,\ \mathbf{Int}_{1\to2}(x).
\]
因此，\(L_2\) 中每个连续对象都不再是“无来源的连续点”，
而都带有一条可追溯到 \(L_1\) 与双离散编码的来源链。
这说明 \(L_{\ge3}\) 的职责首先包含
\[
  \operatorname{Rep}^{\mathrm{int}}_{1\to2}.
\]

再证完备性部分。
任取开集 \(U\subseteq L_2\)，若 \(\Hole(U)\) 成立，
则按假设 2 有
\[
  \pi_\infty^\ast[\omega_{\mathrm{obs}}(U)]=0,
  \qquad
  \exists s_\infty:U\to\mathcal E_\infty.
\]
于是由定义~\ref{def:tex1-ch11-completeness}，
\[
  \mathbf{Cmp}_2(U)
\]
成立。
由于 \(U\subseteq L_2\) 任意，
\[
  \forall U\subseteq L_2,\ \Hole(U)\Rightarrow \mathbf{Cmp}_2(U).
\]
这说明 \(L_{\ge3}\) 的职责同时包含
\[
  \operatorname{Rep}^{\mathrm{cmp}}_2.
\]

综上，
\[
  \operatorname{Rep}_{\ge3}
  =
  \operatorname{Rep}^{\mathrm{int}}_{1\to2}
  \wedge
  \operatorname{Rep}^{\mathrm{cmp}}_2.
\]
定理得证。
\end{proof}

\begin{proposition}[两类修复的逻辑次序不能交换]
\label{prop:tex1-ch11-order}
在 \(L_{\ge3}=L_3\cup L_{>3}\) 的职责分工下，
\[
  \boxed{
    \text{先做完整性修复，再做完备性修复。}
  }
\]
更精确地说：
\begin{enumerate}
  \item 若只做 \(\operatorname{Rep}^{\mathrm{cmp}}_{2}\) 而不做 \(\operatorname{Rep}^{\mathrm{int}}_{1\to2}\)，
  则孔洞虽可在总空间中被平凡化，但连续对象仍缺少离散来源链，体系依然不完整；
  \item 若只做 \(\operatorname{Rep}^{\mathrm{int}}_{1\to2}\) 而不做 \(\operatorname{Rep}^{\mathrm{cmp}}_{2}\)，
  则连续对象虽可回溯到离散来源，但 \(L_2\) 内部孔洞仍未被修复，体系依然不完备。
\end{enumerate}
因此，完整修复必须取两者合取，并按
\[
  L_3\ \text{的首层编码桥}
  \Longrightarrow
  L_{>3}\ \text{的纤维修复塔}
\]
的顺序组织。
\end{proposition}

\begin{Certificate}[证书：对象桥先行 + 纤维修复后续]
\label{cert:tex1-ch11-order}
证书分两条：
\begin{enumerate}
  \item \(L_{>3}\) 的修复塔以 \(L_2\) 为基底、以编码后对象为修复域，故默认其对象桥已经建立；
  \item 单独完成任一修复都只能消去一种失败，不能同时消去“来源断裂”与“孔洞障碍”两种失败。
\end{enumerate}
\end{Certificate}

\begin{proof}[证明（基于数学符号推演逻辑的谓词因果链）]
先看只做完备性修复的情形。
若
\[
  \operatorname{Rep}^{\mathrm{cmp}}_2
\]
成立而
\[
  \neg \operatorname{Rep}^{\mathrm{int}}_{1\to2},
\]
则存在某个 \(x\in L_2\) 使
\[
  \neg \mathbf{Int}_{1\to2}(x).
\]
这意味着 \(x\) 虽处于连续基底中，但没有可审计的离散来源链。
因此，即使某些开集上的孔洞已被拉回平凡化，体系仍缺失从 \(L_1\) 到 \(L_2\) 的对象桥，故不完整。

再看只做完整性修复的情形。
若
\[
  \operatorname{Rep}^{\mathrm{int}}_{1\to2}
\]
成立而
\[
  \neg \operatorname{Rep}^{\mathrm{cmp}}_2,
\]
则存在某个开集 \(U\subseteq L_2\) 使
\[
  \Hole(U)\wedge \neg \mathbf{Cmp}_2(U).
\]
这意味着虽然 \(U\) 中各对象都有可追溯编码，但孔洞障碍尚未在修复塔总空间中被平凡化，故体系仍不完备。

最后，由 \cite{RepoTex41HomotopyRepairBundle}，
\[
  L_3
\]
被钉死为首层编码桥，
\[
  L_{>3}
\]
被钉死为高层纤维修复塔。
故其逻辑次序只能是“先桥接、后修洞”，不能交换。
命题得证。
\end{proof}

\begin{corollary}[对上一节表述的直接替换]
\label{cor:tex1-ch11-replace}
因此，若把上一节中过于简化的说法
\[
  L_{\ge3}\text{ 修复 }=L_2\text{ 孔洞修复}
\]
替换为严格版本，则应直接写成
\[
  \boxed{
    \begin{aligned}
      L_{\ge3}\text{ 修复}
      &=
      \underbrace{\text{$(L_1\to L_2)$ 的完整性修复}}_{\text{首层 }\omega\text{-递归编码桥}}\\
      &\qquad+
      \underbrace{\text{$L_2$ 孔洞的完备性修复}}_{\text{高层纤维吸收与截面恢复}}.
    \end{aligned}
  }
\]
也就是说，\(L_{\ge3}\) 的完整语义不是单层孔洞修复，而是双层修复。
\end{corollary}

\begin{Certificate}[证书：双层职责定理 + 次序命题]
\label{cert:tex1-ch11-replace}
证书分两段：
\begin{enumerate}
  \item 定理~\ref{thm:tex1-ch11-dual} 已给出
  \[
    \operatorname{Rep}_{\ge3}
    =
    \operatorname{Rep}^{\mathrm{int}}_{1\to2}
    \wedge
    \operatorname{Rep}^{\mathrm{cmp}}_{2};
  \]
  \item 命题~\ref{prop:tex1-ch11-order} 已给出两类修复不能互相替代，且逻辑次序不能交换。
\end{enumerate}
\end{Certificate}

\begin{proof}[证明（基于数学符号推演逻辑的谓词因果链）]
由定理~\ref{thm:tex1-ch11-dual}，
\[
  \operatorname{Rep}_{\ge3}
  =
  \operatorname{Rep}^{\mathrm{int}}_{1\to2}
  \wedge
  \operatorname{Rep}^{\mathrm{cmp}}_{2}.
\]
由命题~\ref{prop:tex1-ch11-order}，
两类修复不能互相替代，且必须按“先完整性、后完备性”的顺序组织。
因此，把 \(L_{\ge3}\) 的总职责压缩为单一的 \(L_2\) 孔洞修复是不准确的；
正确表述必须显式保留这两层语义。推论得证。
\end{proof}

\begin{remark}[这一节的重要性：\(L_3\) 的地位被彻底钉死]
如果不加本节这一层，很容易把
\[
  L_3
\]
误读为“通往 \(L_{>3}\) 的技术前奏”。
但本节已经把它钉死为第一类修复本身：
\[
  L_3\text{ 修的是 }(L_1\to L_2)\text{ 的完整性。}
\]
而
\[
  L_{>3}\text{ 修的是 }L_2\text{ 内部孔洞的完备性。}
\]
于是 \(L_{\ge3}\) 的内部职责分解在本文体系内被完全说清了。
\end{remark}

\begin{remark}[最压缩的一句话]
最压缩地说：
\[
  \boxed{
    \text{完整性修复负责把离散来源链接上连续基底，}
  }
\]
\[
  \boxed{
    \text{完备性修复负责把连续基底内部孔洞在纤维塔中平凡化。}
  }
\]
这两层并取，才是 \(L_{\ge3}\) 修复的完整严格语义。
\end{remark}

%%%%%%%%%%%%%%%%%%%%%%%%%%%%%%%%%%%%%%%%%%%%%%%%%%%%%%%%%%%%
\clearpage
\section{形式逻辑：L3 的证书化逻辑占位显著增重层与高层修复入口}
\label{sec:tex1-ch12}
%%%%%%%%%%%%%%%%%%%%%%%%%%%%%%%%%%%%%%%%%%%%%%%%%%%%%%%%%%%%

\begin{remark}[本章意义备注：高超音速攻防的抽象语境]
本章的意义，在于把 \(L_3\) 固定为证书化逻辑占位显著增重层，
说明高层修复不只是增加对象，而是增加可核验、可传递、可说明的判断负载。
对高超音速攻防的抽象问题来说，它对应的是从“能猜”提升到“能给出可审计判定理由”的层级跃迁。
\end{remark}

\subsection{严格主张、引文定位与术语边界}
\label{subsec:tex1-ch12-claim}

\begin{remark}[核心陈述]
本节直接固定 \(L_3\) 的最终形式定位：
\[
  \boxed{
    L_3\in \GMS
  }
\]
\[
  \boxed{
    \operatorname{Sig}_{\mathcal C}(L_3\mid \mathfrak P_2),
    \qquad
    \mathfrak P_2=(L_1,L_2)
  }
\]
\[
  \boxed{
    L_3\text{ 是相对于 }L_1\to L_2\text{ 止步前缀的、基于对应证书的逻辑占位显著增重层。}
  }
\]
在钱学森弹道连续对象进入钱学森障碍修复体系的主链中，
\[
  L_3=\{(\sigma_{\le 3},x):x\in K_{\sigma_{\le 3}}\}
\]
从定义开始就耦合了离散前缀、连续片、归属关系、递归生成接口与局部证书槽，
因此它本身就是富结构对象；
同时，它附带承担富结构入口、\((L_1\to L_2)\) 完整性修复桥与高层修复入口这三类不可约角色。
\end{remark}

\begin{remark}[仓内文稿定位]
本节继续承接 \cite{RepoTex41HomotopyRepairBundle}：
该文已把 \(L_3\) 钉死为“双可数离散对对 \(L_2\) 的首层 \(\omega\)-递归编码落点”，
并把 \(L_{>3}\) 钉死为“以 \(L_2\) 为基底的高层修复/删边/可行域重构塔”；
\cite{RepoTex36HomotopyTwistEquivalence} 提供正规化截面、稳定修复序列与编码同一性的接口；
\cite{RepoTex38JacobiBianchiRepair} 则提供严格下降、有限终止与终端正规形的离散修复证书；
\cite{RepoMetaGMSAxiom} 给出“广义数学结构 \(\GMS\)”的公理口径，
其中 A1、A3、A4、A7 分别钉死结构可集合化封装、有限异构递归构成、经 GRL 积分仍闭合于 \(\GMS\) 内，
以及“结构--集合--路径积分”的三位一体范式等价。
因此，本节将把仓内已存在的对象定义、修复接口与元结构公理，
收束为关于 \(L_3\) 本体定位、职责分解与证书化逻辑占位的严格定理链。
\end{remark}

\begin{remark}[严格边界]
本节采用如下三层严格分工：
\[
  L_2=\text{唯一连续基底},
\]
\[
  L_3=\text{证书化逻辑占位显著增重层，兼 }(L_1\to L_2)\text{ 完整性修复桥},
\]
\[
  L_{>3}=\text{高阶孔洞完备性修复与删边治理层}.
\]
因此，在
\[
  L_1\to L_2\to L_3\to L_{>3}
\]
这条正规链上，\(L_3\) 与 \(L_{>3}\) 都属于富结构对象，
其中 \(L_3\) 承担显式耦合“离散前缀 + 连续片 + 归属关系 + 完整性修复接口”的首层职责，
而 \(L_{>3}\) 继续承担更高阶的孔洞吸收、删边治理与可行域重构职责。
\end{remark}

\subsection{定义链：逻辑占位权重、职责集与证书化度量}
\label{subsec:tex1-ch12-defs}

\begin{definition}[逻辑性度量与逻辑占位权重]
\label{def:tex1-ch12-logicality}
定义对象 \(X\) 的逻辑性度量为
\[
  \Lambda_{\mathrm{log}}(X)
  :=
  \alpha\,\rho(X)
  +
  \beta\,\eta(X)
  +
  \gamma\,\chi(X)
  +
  \delta\,\iota(X),
  \qquad
  \alpha,\beta,\gamma,\delta>0.
\]
其中
\[
  \rho,\eta,\chi,\iota:\operatorname{Obj}\to \mathbb{N}_0
\]
分别表示：
\begin{enumerate}
  \item \(\rho(X)\)：结构异质度，即 \(X\) 内部耦合了多少种不可互相约简的结构成分；
  \item \(\eta(X)\)：递归生成度，即 \(X\) 是否作为后续层级生成的首层种子；
  \item \(\chi(X)\)：证书承载度，即 \(X\) 是否天然承载可审计的局部证书槽；
  \item \(\iota(X)\)：不可消去度，即删去 \(X\) 后主链中有多少定理会失去逻辑闭合。
\end{enumerate}
在此意义下，“富结构”并不等于“对象很复杂”，
而 \(\Lambda_{\mathrm{log}}(X)\) 衡量的是对象在富结构宇宙内部的逻辑展开度，
并不是对象是否属于富结构宇宙的成员资格判据。
\end{definition}

\begin{definition}[富结构成员资格、平庸结构与富结构分层]
\label{def:tex1-ch12-richness}
定义谓词
\[
  \operatorname{Rich}(X)
  :
  \Longleftrightarrow
  X\in \GMS.
\]
给定阈值
\[
  0<\theta_0<\theta_1,
\]
定义
\[
  \operatorname{Plain}(X)
  :
  \Longleftrightarrow
  \operatorname{Rich}(X)\wedge \Lambda_{\mathrm{log}}(X)\le \theta_0.
\]
更一般地，定义富结构宇宙内部的三段分层：
\[
  \GMS_{\mathrm{low}}(\theta_0)
  :=
  \{X\in \GMS:\Lambda_{\mathrm{log}}(X)\le \theta_0\},
\]
\[
  \GMS_{\mathrm{mid}}(\theta_0,\theta_1)
  :=
  \{X\in \GMS:\theta_0<\Lambda_{\mathrm{log}}(X)\le \theta_1\},
\]
\[
  \GMS_{\mathrm{high}}(\theta_1)
  :=
  \{X\in \GMS:\Lambda_{\mathrm{log}}(X)>\theta_1\}.
\]
于是“平庸结构”不是 \(\GMS\) 之外的另一种类型，
而只是 \(\GMS\) 内部的低逻辑性度量截面。
\end{definition}

\begin{theorem}[平庸结构也是富结构]
\label{thm:tex1-ch12-plain-rich}
对任意对象 \(X\)，都有
\[
  \operatorname{Plain}(X)\Rightarrow \operatorname{Rich}(X).
\]
\end{theorem}

\begin{Certificate}[证书：定义展开证书]
\label{cert:tex1-ch12-plain-rich}
证书直接来自定义~\ref{def:tex1-ch12-richness}：
\[
  \operatorname{Plain}(X)
  :
  \Longleftrightarrow
  \operatorname{Rich}(X)\wedge \Lambda_{\mathrm{log}}(X)\le \theta_0.
\]
\end{Certificate}

\begin{proof}[证明（基于数学符号推演逻辑的谓词因果链）]
由定义~\ref{def:tex1-ch12-richness}，
\[
  \operatorname{Plain}(X)
  \Longleftrightarrow
  \operatorname{Rich}(X)\wedge \Lambda_{\mathrm{log}}(X)\le \theta_0.
\]
因此一旦
\[
  \operatorname{Plain}(X)
\]
成立，就必有
\[
  \operatorname{Rich}(X)
\]
成立。定理得证。
\end{proof}

\begin{corollary}[“平庸/丰富”不是成员资格二分，而是富结构宇宙内部的度量分层]
\label{cor:tex1-ch12-layering}
正确的分层不是
\[
  \text{平庸结构}
  \quad/\quad
  \text{富结构},
\]
而是
\[
  \GMS
  =
  \GMS_{\mathrm{low}}(\theta_0)
  \sqcup
  \GMS_{\mathrm{mid}}(\theta_0,\theta_1)
  \sqcup
  \GMS_{\mathrm{high}}(\theta_1).
\]
其中所谓“平庸结构”只对应
\[
  \GMS_{\mathrm{low}}(\theta_0),
\]
而不是 \(\GMS\) 之外的另一个类型域。
\end{corollary}

\begin{Certificate}[证书：成员资格统一 + 阈值分层]
\label{cert:tex1-ch12-layering}
证书分两条：
\begin{enumerate}
  \item 定理~\ref{thm:tex1-ch12-plain-rich} 已证“平庸结构也是富结构”；
  \item 定义~\ref{def:tex1-ch12-richness} 已把 \(\GMS\) 按 \(\Lambda_{\mathrm{log}}\) 划分为 low/mid/high 三段。
\end{enumerate}
\end{Certificate}

\begin{proof}[证明（基于数学符号推演逻辑的谓词因果链）]
由定理~\ref{thm:tex1-ch12-plain-rich}，
\[
  \operatorname{Plain}(X)\Rightarrow X\in \GMS.
\]
因此“平庸结构”并不离开 \(\GMS\)。
再由定义~\ref{def:tex1-ch12-richness}，
\(\GMS\) 已经按 \(\Lambda_{\mathrm{log}}\) 被分成 low/mid/high 三段。
故“平庸/丰富”的正确关系不是成员资格二分，
而是富结构宇宙内部的度量分层。推论得证。
\end{proof}

\begin{definition}[架构占位权重]
\label{def:tex1-ch12-weight}
定义任一层 \(L_i\) 的架构占位权重为
\[
  W(L_i)
  :=
  \sum_{\rho\in \mathcal R(L_i)}\omega_\rho,
  \qquad
  \omega_\rho>0.
\]
其中 \(\mathcal R(L_i)\) 表示该层承担的\textbf{不可约职责集合}，
即不能被更低层或同层其他对象替代的职责。
该量只是逻辑性度量 \(\Lambda_{\mathrm{log}}\) 在“职责集合”方向上的投影，
并不是对象本体性的全部刻画。
\end{definition}

\begin{definition}[比较基准下的 \(L_3\) 极简职责集与完整职责集]
\label{def:tex1-ch12-role-set}
为刻画 \(L_3\) 架构占位权重的严格增量，
先定义其极简职责集为
\[
  \mathcal R_{\mathrm{min}}(L_3)
  :=
  \{\mathbf E_3\},
\]
其中 \(\mathbf E_3\) 表示“对 \(L_2\) 的首层 \(\omega\)-递归离散编码”。

再定义其完整职责集为
\[
  \mathcal R_{\mathrm{full}}(L_3)
  :=
  \{
    \mathbf E_3,\,
    \mathbf{Int}_{1\to2},\,
    \mathbf{Adr}_3,\,
    \mathbf{Cert}_3,\,
    \mathbf{Gate}_{3\to >3}
  \}.
\]
其中：
\begin{enumerate}
  \item \(\mathbf{Int}_{1\to2}\) 表示 \((L_1\to L_2)\) 的完整性修复；
  \item \(\mathbf{Adr}_3\) 表示对 \(L_2\) 连续对象的首层可寻址化；
  \item \(\mathbf{Cert}_3\) 表示首层证书槽与局部可审计编码位；
  \item \(\mathbf{Gate}_{3\to >3}\) 表示把对象正式送入 \(L_{>3}\) 高阶修复塔的入口职责。
\end{enumerate}
\end{definition}

\begin{definition}[下游不可缺依赖度]
\label{def:tex1-ch12-depend}
定义 \(L_3\) 的下游不可缺依赖度为
\[
  D(L_3)
  :=
  \#\bigl\{
    \varphi:
    \varphi\text{ 是下游命题且 }L_3\text{ 对其不可缺}
  \bigr\}.
\]
它统计的不是“有多少命题提到了 \(L_3\)”这种弱计数，
而是“有多少下游命题在去掉 \(L_3\) 后便无法闭合”的强计数。
\end{definition}

\begin{definition}[证书化逻辑占位度量与显著增重层谓词]
\label{def:tex1-ch12-certlogical}
定义层或层前缀 \(B\) 的对应证书族为
\[
  \mathcal C(B).
\]
定义其证书化逻辑占位度量为
\[
  \Lambda_{\mathrm{log}}^{\mathcal C}(B)
  :=
  \sum_{c\in \mathcal C(B)} w(c),
  \qquad
  w(c)>0.
\]
对任一层 \(L_i\) 与比较基准 \(B\)，定义“显著增重层”谓词
\[
  \operatorname{Sig}_{\mathcal C}(L_i\mid B)
  :
  \Longleftrightarrow
  \Lambda_{\mathrm{log}}^{\mathcal C}(L_i)-\Lambda_{\mathrm{log}}^{\mathcal C}(B)
  \ge
  \varepsilon_{\mathcal C},
  \qquad
  \varepsilon_{\mathcal C}>0.
\]
这里“显著增重”刻画的是证书化逻辑占位度量的上升，
而不是对象成员资格的改变。
\end{definition}

\subsection{公理（降级为定理）：L3 富结构入口职责可按层前缀递归钉死}
\label{subsec:tex1-ch12-axiom}

\begin{theorem}[公理（降级为定理）：\(L_3\) 富结构入口职责的层前缀递归钉死]
\label{thm:tex1-ch12-axiom-downgraded}
定义层前缀
\[
  \mathfrak P_2:=(L_1,L_2),\qquad
  \mathfrak P_3:=(L_1,L_2,L_3),\qquad
  \mathfrak P_4:=(L_1,L_2,L_3,L_{>3}).
\]
再定义谓词
\[
  P(n)
  :
  \Longleftrightarrow
  \text{前缀 }\mathfrak P_n\text{ 已把新增的最高不可约职责唯一钉死。}
\]
则在 \cite{RepoTex41HomotopyRepairBundle} 的对象链下：
\begin{enumerate}
  \item \(P(2)\) 成立，且 \(\mathfrak P_2\) 只钉死 \(L_2\) 为连续基底；
  \item \(P(3)\) 成立，且 \(\mathfrak P_3\) 新增并钉死的职责正是
  \[
    \mathbf E_3,\qquad \mathbf{Int}_{1\to2},\qquad \mathbf{Adr}_3,\qquad \mathbf{Cert}_3;
  \]
  \item \(P(4)\) 成立，且 \(\mathfrak P_4\) 新增并钉死的职责正是
  \[
    \mathbf{Gate}_{3\to >3}
    \quad\text{与}\quad
    L_{>3}\text{ 的高层孔洞完备性修复职责。}
  \]
\end{enumerate}
因此
\[
  \boxed{
    \mathcal R_{\mathrm{min}}(L_3)
    \subsetneq
    \mathcal R_{\mathrm{full}}(L_3).
  }
\]
\end{theorem}

\begin{Certificate}[数学归纳法证书：层前缀职责扩张证书]
\label{cert:tex1-ch12-axiom-downgraded}
对前缀长度 \(n\in\{2,3,4\}\) 做数学归纳：
\begin{enumerate}
  \item \textbf{基步 \(P(2)\)：}由 \cite{RepoTex41HomotopyRepairBundle}，\((L_1,L_2)\) 只钉死离散索引层与连续基底层；
  \item \textbf{归纳步 \(P(2)\Rightarrow P(3)\)：}由首层 \(\omega\)-递归编码满射接口，\(L_3\) 同时引入编码、完整性、可寻址化与证书前缀；
  \item \textbf{归纳步 \(P(3)\Rightarrow P(4)\)：}由高层修复塔定理，\(L_{>3}\) 的动作必须作用在已编码、已可寻址、已有证书槽的对象上，因此入口职责反向钉死在 \(L_3\)。
\end{enumerate}
由此前缀链递归分离，得到 \(L_3\) 的完整职责集相对于极简职责集严格扩张。
\end{Certificate}

\begin{proof}[证明（数学归纳法证书；基于数学符号推演逻辑的谓词因果链）]
\textbf{基步 \(P(2)\)。}
在前缀
\[
  \mathfrak P_2=(L_1,L_2)
\]
中，\(L_1\) 只提供离散索引与证书槽占位，\(L_2\) 提供唯一连续基底。
此时既没有首层编码，也没有高层入口，更没有完备性修复塔。
故 \(P(2)\) 成立。

\textbf{归纳步 \(P(2)\Rightarrow P(3)\)。}
由 \cite{RepoTex41HomotopyRepairBundle}，
\[
  \operatorname{Enc}:\mathcal C_\omega\to L_2
\]
为满射，且 \(L_3\) 是双可数离散对对 \(L_2\) 的首层 \(\omega\)-递归编码落点。
因此：
\begin{enumerate}
  \item \(\mathbf E_3\) 被钉死，因为 \(L_3\) 首先负责把双离散码压到连续基底上；
  \item \(\mathbf{Int}_{1\to2}\) 被钉死，因为每个 \(x\in L_2\) 都必须经某个 \(\sigma\in\mathcal C_\omega\) 进入富结构体系；
  \item \(\mathbf{Adr}_3\) 被钉死，因为每个连续对象第一次获得“可被后续修复塔引用的首层地址”；
  \item \(\mathbf{Cert}_3\) 被钉死，因为首层前缀正是局部可审计证书槽的第一落点。
\end{enumerate}
故 \(P(3)\) 成立。

\textbf{归纳步 \(P(3)\Rightarrow P(4)\)。}
再由 \cite{RepoTex41HomotopyRepairBundle}，
\[
  L_{>3}
\]
沿纤维方向执行障碍吸收、删边与可行域重构。
但要使这些动作可执行，所作用对象必须已经具备：
\[
  \text{首层编码地址}
  \quad+\quad
  \text{局部证书前缀}
  \quad+\quad
  \text{进入修复塔的入口位。}
\]
前两项已由 \(P(3)\) 给出，第三项于是被唯一钉死为
\[
  \mathbf{Gate}_{3\to >3}.
\]
故 \(P(4)\) 成立。

由基步与两次归纳步，
\[
  \mathcal R_{\mathrm{min}}(L_3)=\{\mathbf E_3\}
  \subsetneq
  \{\mathbf E_3,\mathbf{Int}_{1\to2},\mathbf{Adr}_3,\mathbf{Cert}_3,\mathbf{Gate}_{3\to >3}\}
  =
  \mathcal R_{\mathrm{full}}(L_3).
\]
定理得证。
\end{proof}

\subsection{定理链：权重严格增大、下游依赖增长与新定位}
\label{subsec:tex1-ch12-theorems}

\begin{lemma}[首层编码桥给出可寻址化、证书前缀与高层入口的前置条件]
\label{lem:tex1-ch12-entry}
若
\[
  \operatorname{Enc}:\mathcal C_\omega\to L_2
\]
为满射，且 \(L_{>3}\) 的修复动作只作用在已编码对象上，
则
\[
  \mathbf{Adr}_3,\qquad \mathbf{Cert}_3,\qquad \mathbf{Gate}_{3\to >3}
\]
都属于 \(L_3\) 的不可约职责。
\end{lemma}

\begin{Certificate}[证书：对象先编码，后修复]
\label{cert:tex1-ch12-entry}
证书分三条：
\begin{enumerate}
  \item 满射编码使每个 \(x\in L_2\) 首次获得首层地址；
  \item 首层前缀是局部证书槽与可审计编码位的首个落点；
  \item 高层修复塔只对已进入编码体系的对象起作用，因此入口职责不能脱离 \(L_3\)。
\end{enumerate}
\end{Certificate}

\begin{proof}[证明（基于数学符号推演逻辑的谓词因果链）]
任取 \(x\in L_2\)。
由满射性，存在
\[
  \sigma\in\mathcal C_\omega
\]
使
\[
  \operatorname{Enc}(\sigma)=x.
\]
因此，\(x\) 第一次被高层体系引用时，必先带有某个首层编码前缀。
这就给出 \(\mathbf{Adr}_3\)。

又因为该前缀不仅是地址，还是局部可审计编码位与证书槽的第一显式落点，
故 \(\mathbf{Cert}_3\) 也是 \(L_3\) 的不可约职责。

最后，由 \(L_{>3}\) 的高层修复动作只作用于已编码对象，
可知“把对象正式送入高层修复塔”这一动作本身不能由 \(L_{>3}\) 逆向自造，
而必须由 \(L_3\) 提供入口。
因此 \(\mathbf{Gate}_{3\to >3}\) 也属于 \(L_3\) 的不可约职责。
引理得证。
\end{proof}

\begin{theorem}[\(L_3\) 的架构占位权重严格增大]
\label{thm:tex1-ch12-weight}
若接受上一节的双层修复分解
\[
  \operatorname{Rep}_{\ge3}
  =
  \operatorname{Rep}^{\mathrm{int}}_{1\to2}
  \wedge
  \operatorname{Rep}^{\mathrm{cmp}}_{2},
\]
并约定
\[
  \operatorname{Rep}^{\mathrm{int}}_{1\to2}\text{ 主要落在 }L_3,
  \qquad
  \operatorname{Rep}^{\mathrm{cmp}}_{2}\text{ 主要落在 }L_{>3},
\]
则有
\[
  \boxed{
    W_{\mathrm{full}}(L_3)>W_{\mathrm{min}}(L_3).
  }
\]
\end{theorem}

\begin{Certificate}[证书：职责集严格包含 + 正权重求和]
\label{cert:tex1-ch12-weight}
证书分两段：
\begin{enumerate}
  \item 定理~\ref{thm:tex1-ch12-axiom-downgraded} 与引理~\ref{lem:tex1-ch12-entry} 给出
  \[
    \mathcal R_{\mathrm{min}}(L_3)\subsetneq \mathcal R_{\mathrm{full}}(L_3);
  \]
  \item 定义~\ref{def:tex1-ch12-weight} 中每个职责权重都满足 \(\omega_\rho>0\)。
\end{enumerate}
\end{Certificate}

\begin{proof}[证明（基于数学符号推演逻辑的谓词因果链）]
由定义~\ref{def:tex1-ch12-role-set}，
\[
  \mathcal R_{\mathrm{min}}(L_3)=\{\mathbf E_3\}.
\]
由定理~\ref{thm:tex1-ch12-axiom-downgraded} 与引理~\ref{lem:tex1-ch12-entry}，
\[
  \mathcal R_{\mathrm{full}}(L_3)
  =
  \{\mathbf E_3,\mathbf{Int}_{1\to2},\mathbf{Adr}_3,\mathbf{Cert}_3,\mathbf{Gate}_{3\to >3}\},
\]
且
\[
  \mathcal R_{\mathrm{min}}(L_3)\subsetneq \mathcal R_{\mathrm{full}}(L_3).
\]

再由定义~\ref{def:tex1-ch12-weight}，
\[
  W_{\mathrm{min}}(L_3)
  =
  \sum_{\rho\in\mathcal R_{\mathrm{min}}(L_3)}\omega_\rho,
  \qquad
  W_{\mathrm{full}}(L_3)
  =
  \sum_{\rho\in\mathcal R_{\mathrm{full}}(L_3)}\omega_\rho.
\]
由于新增职责的权重都满足
\[
  \omega_\rho>0,
\]
严格包含立即推出
\[
  W_{\mathrm{full}}(L_3)>W_{\mathrm{min}}(L_3).
\]
定理得证。
\end{proof}

\begin{corollary}[\(L_3\) 的架构占位权重的正式分解]
\label{cor:tex1-ch12-weight-decompose}
记
\[
  W_{\mathrm{enc}}(L_3):=\omega_{\mathbf E_3},
  \qquad
  W_{\mathrm{int}}(L_3):=\omega_{\mathbf{Int}_{1\to2}},
\]
\[
  W_{\mathrm{gate}}(L_3)
  :=
  \omega_{\mathbf{Adr}_3}
  +
  \omega_{\mathbf{Cert}_3}
  +
  \omega_{\mathbf{Gate}_{3\to >3}}.
\]
则
\[
  \boxed{
    W_{\mathrm{full}}(L_3)
    =
    W_{\mathrm{enc}}(L_3)
    +
    W_{\mathrm{int}}(L_3)
    +
    W_{\mathrm{gate}}(L_3).
  }
\]
亦即
\[
  \boxed{
    L_3\text{ 的架构占位权重，不再只是首层编码权重，}
  }
\]
\[
  \boxed{
    \text{而是首层编码权重 }+\text{ 完整性修复权重 }+\text{ 高层修复入口权重。}
  }
\]
\end{corollary}

\begin{Certificate}[证书：完整职责集的三块求和分解]
\label{cert:tex1-ch12-weight-decompose}
证书由定理~\ref{thm:tex1-ch12-weight} 与定义~\ref{def:tex1-ch12-role-set} 共同给出：
\[
  \mathcal R_{\mathrm{full}}(L_3)
  =
  \{\mathbf E_3\}
  \sqcup
  \{\mathbf{Int}_{1\to2}\}
  \sqcup
  \{\mathbf{Adr}_3,\mathbf{Cert}_3,\mathbf{Gate}_{3\to >3}\}.
\]
\end{Certificate}

\begin{proof}[证明（基于数学符号推演逻辑的谓词因果链）]
由定义~\ref{def:tex1-ch12-role-set}，
\[
  \mathcal R_{\mathrm{full}}(L_3)
  =
  \{\mathbf E_3,\mathbf{Int}_{1\to2},\mathbf{Adr}_3,\mathbf{Cert}_3,\mathbf{Gate}_{3\to >3}\}.
\]
再由定义~\ref{def:tex1-ch12-weight}，
\[
  W_{\mathrm{full}}(L_3)
  =
  \omega_{\mathbf E_3}
  +
  \omega_{\mathbf{Int}_{1\to2}}
  +
  \omega_{\mathbf{Adr}_3}
  +
  \omega_{\mathbf{Cert}_3}
  +
  \omega_{\mathbf{Gate}_{3\to >3}}.
\]
按
\[
  \{\mathbf E_3\},
  \qquad
  \{\mathbf{Int}_{1\to2}\},
  \qquad
  \{\mathbf{Adr}_3,\mathbf{Cert}_3,\mathbf{Gate}_{3\to >3}\}
\]
三块重新分组，即得
\[
  W_{\mathrm{full}}(L_3)
  =
  W_{\mathrm{enc}}(L_3)
  +
  W_{\mathrm{int}}(L_3)
  +
  W_{\mathrm{gate}}(L_3).
\]
故推论成立。
\end{proof}

\begin{proposition}[下游不可缺依赖度也严格增大]
\label{prop:tex1-ch12-depend}
在完整职责集口径下，
\[
  \boxed{
    D_{\mathrm{full}}(L_3)>D_{\mathrm{min}}(L_3).
  }
\]
更精确地说，极简职责集下依赖 \(L_3\) 的主要只是“首层编码”命题；
而完整职责集下，
\[
  \text{完整性修复},\quad
  \text{高层修复入口},\quad
  \text{删边前可寻址化},\quad
  \text{局部证书槽定位}
\]
都对 \(L_3\) 不可缺。
\end{proposition}

\begin{Certificate}[证书：职责种类增长 + 下游命题扩张]
\label{cert:tex1-ch12-depend}
证书分两条：
\begin{enumerate}
  \item 完整职责集下 \(L_3\) 的不可约职责种类增多；
  \item 每增加一种不可约职责，都会至少新增一类以该职责为前提的下游闭合命题。
\end{enumerate}
\end{Certificate}

\begin{proof}[证明（基于数学符号推演逻辑的谓词因果链）]
由定理~\ref{thm:tex1-ch12-weight} 的证明，
\[
  \mathcal R_{\mathrm{min}}(L_3)\subsetneq \mathcal R_{\mathrm{full}}(L_3).
\]
这意味着极简职责集下只有“首层编码”依赖 \(L_3\)，
而完整职责集下至少还新增了：
\[
  \mathbf{Int}_{1\to2},\qquad
  \mathbf{Adr}_3,\qquad
  \mathbf{Cert}_3,\qquad
  \mathbf{Gate}_{3\to >3}.
\]
每一项新增职责都对应至少一类下游命题：
\begin{enumerate}
  \item \(\mathbf{Int}_{1\to2}\) 对应完整性修复命题；
  \item \(\mathbf{Adr}_3\) 对应删边/可行域重构前的对象定位命题；
  \item \(\mathbf{Cert}_3\) 对应局部可审计证书落点命题；
  \item \(\mathbf{Gate}_{3\to >3}\) 对应高层修复塔入口命题。
\end{enumerate}
故下游不可缺依赖度严格增大：
\[
  D_{\mathrm{full}}(L_3)>D_{\mathrm{min}}(L_3).
\]
命题得证。
\end{proof}

\begin{lemma}[\(L_3\) 的对象定义自带离散--连续--归属三重耦合]
\label{lem:tex1-ch12-coupled}
在 \cite{RepoTex41HomotopyRepairBundle} 的对象口径下，
\[
  L_3
  =
  \{(\sigma_{\le 3},x):x\in K_{\sigma_{\le 3}}\}.
\]
因此 \(L_3\) 从定义开始就同时耦合了
\[
  \text{离散前缀}\quad+\quad \text{连续点}\quad+\quad \text{归属关系 }x\in K_{\sigma_{\le3}},
\]
从而至少满足
\[
  \rho(L_3)\ge 3,\qquad
  \eta(L_3)\ge 1,\qquad
  \chi(L_3)\ge 1,\qquad
  \iota(L_3)\ge 1.
\]
\end{lemma}

\begin{Certificate}[证书：对象定义 + 递归编码 + 证书槽 + 下游闭合]
\label{cert:tex1-ch12-coupled}
证书分四条：
\begin{enumerate}
  \item \cite{RepoTex41HomotopyRepairBundle} 已把 \(L_3\) 写成离散前缀与连续片的耦合对象；
  \item 首层 \(\omega\)-递归编码使 \(L_3\) 成为后续层级生成的首层种子；
  \item 首层前缀同时承载局部证书槽与可审计编码位；
  \item 去掉 \(L_3\) 后，完整性修复、高层入口与证书前缀相关命题都会失去闭合。
\end{enumerate}
\end{Certificate}

\begin{proof}[证明（基于数学符号推演逻辑的谓词因果链）]
由 \cite{RepoTex41HomotopyRepairBundle}，
\[
  L_3
  =
  \{(\sigma_{\le 3},x):x\in K_{\sigma_{\le 3}}\}.
\]
这一定义中已经同时出现：
\[
  \sigma_{\le 3},
  \qquad
  x,
  \qquad
  x\in K_{\sigma_{\le 3}}.
\]
故 \(L_3\) 不是单一类型对象，而是离散前缀、连续点与归属关系的三重耦合对象，
从而
\[
  \rho(L_3)\ge 3.
\]

再由首层 \(\omega\)-递归编码满射接口，
\[
  \operatorname{Enc}:\mathcal C_\omega\to L_2,
\]
\(L_3\) 成为连续对象第一次被可审计地压到富结构链中的落点，
故其递归生成度满足
\[
  \eta(L_3)\ge 1.
\]

又因为第一个可审计编码前缀正是局部证书槽的首个显式落点，
故
\[
  \chi(L_3)\ge 1.
\]

最后，若删去 \(L_3\)，则
\[
  \mathbf{Int}_{1\to2},\qquad
  \mathbf{Adr}_3,\qquad
  \mathbf{Cert}_3,\qquad
  \mathbf{Gate}_{3\to >3}
\]
都无法继续闭合，故
\[
  \iota(L_3)\ge 1.
\]
引理得证。
\end{proof}

\begin{theorem}[\(L_3\) 本来就在富结构宇宙中，且不是低阶平庸截面]
\label{thm:tex1-ch12-gms}
在 \cite{RepoTex41HomotopyRepairBundle,RepoMetaGMSAxiom} 的对象链与元公理口径下，
对任意满足
\[
  0<\theta_0<3\alpha+\beta+\gamma+\delta
\]
的低阶阈值，都有
\[
  \boxed{
    L_3\in \GMS
  }
\]
\[
  \boxed{
    L_3\notin \GMS_{\mathrm{low}}(\theta_0).
  }
\]
\[
  \boxed{
    \begin{aligned}
      &\text{因此关于 }L_3\text{ 的关键不在成员资格，}\\
      &\text{而在其作为富结构的逻辑占位度量已显著高于低阶平庸截面。}
    \end{aligned}
  }
\]
\end{theorem}

\begin{Certificate}[证书：对象三重耦合下界 + 低阶阈值比较 + GMS 元公理]
\label{cert:tex1-ch12-gms}
证书分四段：
\begin{enumerate}
  \item 引理~\ref{lem:tex1-ch12-coupled} 已证 \(L_3\) 自带离散--连续--归属三重耦合；
  \item 定理~\ref{thm:tex1-ch12-axiom-downgraded} 与定理~\ref{thm:tex1-ch12-weight} 已证 \(L_3\) 承接 \((L_1\to L_2)\) 的完整性修复；
  \item \cite{RepoMetaGMSAxiom} 的 A1、A3、A4 说明：可集合化封装、有限异构递归与 GRL 闭合都保留在 \(\GMS\) 内；
  \item 引理~\ref{lem:tex1-ch12-coupled} 给出
  \[
    \Lambda_{\mathrm{log}}(L_3)\ge 3\alpha+\beta+\gamma+\delta>\theta_0.
  \]
\end{enumerate}
\end{Certificate}

\begin{proof}[证明（基于数学符号推演逻辑的谓词因果链）]
由引理~\ref{lem:tex1-ch12-coupled}，
\[
  \rho(L_3)\ge 3,\qquad
  \eta(L_3)\ge 1,\qquad
  \chi(L_3)\ge 1,\qquad
  \iota(L_3)\ge 1.
\]
故
\[
  \Lambda_{\mathrm{log}}(L_3)
  =
  \alpha\,\rho(L_3)
  +
  \beta\,\eta(L_3)
  +
  \gamma\,\chi(L_3)
  +
  \delta\,\iota(L_3)
  \ge 3\alpha+\beta+\gamma+\delta,
\]
而右端又严格大于任意满足
\[
  0<\theta_0<3\alpha+\beta+\gamma+\delta
\]
的低阶阈值。
故
\[
  \Lambda_{\mathrm{log}}(L_3)>\theta_0.
\]

再由 \cite{RepoMetaGMSAxiom} 的公理 A1，
任意广义结构都可封装为广义集合对象；
由公理 A3，
有限层级异构递归构成仍保持结构自嵌套；
由公理 A4，
经 GRL 路径积分后的结果仍闭合在 \(\GMS\) 内；
由公理 A7，
\[
  \GMS
  \equiv
  \GSet_{\mathcal D}
  \equiv
  \GMS_{\mathrm{GRL}}.
\]
因此，凡同时具备结构封装、异构递归与路径积分闭合能力的对象，
都应被定位在 \(\GMS\) 内。

而 \(L_3\) 一方面由 \((\sigma_{\le 3},x,x\in K_{\sigma_{\le3}})\) 给出显式异构耦合，
另一方面又由首层 \(\omega\)-递归编码与完整性修复接口承担
\[
  \mathbf E_3,\qquad
  \mathbf{Int}_{1\to2},\qquad
  \mathbf{Adr}_3,\qquad
  \mathbf{Cert}_3,\qquad
  \mathbf{Gate}_{3\to >3}.
\]
故
\[
  L_3\in \GMS.
\]

再由定义~\ref{def:tex1-ch12-richness}，
\[
  \GMS_{\mathrm{low}}(\theta_0)
  =
  \{X\in \GMS:\Lambda_{\mathrm{log}}(X)\le \theta_0\}.
\]
由于
\[
  L_3\in \GMS
  \quad\text{且}\quad
  \Lambda_{\mathrm{log}}(L_3)>\theta_0,
\]
故
\[
  L_3\notin \GMS_{\mathrm{low}}(\theta_0).
\]
这表明 \(L_3\) 的关键不在“它是不是富结构”，
而在“它作为富结构的逻辑占位度量已显著高于低阶平庸截面”。定理得证。
\end{proof}

\begin{proposition}[\(L_3\) 的稳定落句由证书化显著增重给出]
\label{prop:tex1-ch12-recognize}
在比较基准
\[
  \mathfrak P_2=(L_1,L_2)
\]
下，有
\[
  \boxed{
    \operatorname{Sig}_{\mathcal C}(L_3\mid \mathfrak P_2)
  }
\]
因此，\(L_3\) 的稳定落句应写成：相对于 \(L_1\to L_2\) 止步前缀的、基于对应证书的逻辑占位显著增重层。
\end{proposition}

\begin{Certificate}[证书：富结构成员资格 + 证书差额严格为正]
\label{cert:tex1-ch12-recognize}
证书分四条：
\begin{enumerate}
  \item 定理~\ref{thm:tex1-ch12-plain-rich} 已证“平庸结构也是富结构”，故成员资格已经统一落在 \(\GMS\) 内；
  \item 定理~\ref{thm:tex1-ch12-gms} 已证 \(L_3\in\GMS\) 且 \(L_3\notin \GMS_{\mathrm{low}}(\theta_0)\)；
  \item 定义~\ref{def:tex1-ch12-certlogical} 已把“显著增重”写成证书化逻辑占位差额；
  \item 定理~\ref{thm:tex1-ch12-axiom-downgraded}、引理~\ref{lem:tex1-ch12-entry} 与命题~\ref{prop:tex1-ch12-depend} 给出：
  相对于 \(\mathfrak P_2\)，\(L_3\) 至少新增下列五类正权重证书：
  \[
    \mathbf E_3,\qquad
    \mathbf{Int}_{1\to2},\qquad
    \mathbf{Adr}_3,\qquad
    \mathbf{Cert}_3,\qquad
    \mathbf{Gate}_{3\to >3}.
  \]
\end{enumerate}
\end{Certificate}

\begin{proof}[证明（基于数学符号推演逻辑的谓词因果链）]
由定理~\ref{thm:tex1-ch12-plain-rich}，
\[
  \operatorname{Plain}(X)\Rightarrow X\in \GMS.
\]
因此，平庸结构与显著结构都已经处在富结构宇宙内部。

再由定理~\ref{thm:tex1-ch12-gms}，
\[
  L_3\in \GMS,
  \qquad
  L_3\notin \GMS_{\mathrm{low}}(\theta_0).
\]
这说明 \(L_3\) 的层位已经同时包含成员资格与度量层级两部分。

再由定理~\ref{thm:tex1-ch12-axiom-downgraded}、引理~\ref{lem:tex1-ch12-entry} 与命题~\ref{prop:tex1-ch12-depend}，
\[
  \mathcal C(L_3)\setminus \mathcal C(\mathfrak P_2)
  \supseteq
  \{\mathbf E_3,\mathbf{Int}_{1\to2},\mathbf{Adr}_3,\mathbf{Cert}_3,\mathbf{Gate}_{3\to >3}\}.
\]
由定义~\ref{def:tex1-ch12-certlogical} 中每个证书权重都满足 \(w(c)>0\)，令
\[
  \varepsilon_{\mathcal C}
  :=
  \sum_{c\in \mathcal C(L_3)\setminus \mathcal C(\mathfrak P_2)} w(c)>0,
\]
则有
\[
  \Lambda_{\mathrm{log}}^{\mathcal C}(L_3)-\Lambda_{\mathrm{log}}^{\mathcal C}(\mathfrak P_2)
  \ge
  \varepsilon_{\mathcal C}.
\]
故
\[
  \operatorname{Sig}_{\mathcal C}(L_3\mid \mathfrak P_2)
\]
成立。因此，关于 \(L_3\) 的严格落句可直接写成“基于对应证书的显著增重层”。命题得证。
\end{proof}

\begin{corollary}[\(L_3\) 的新定位]
\label{cor:tex1-ch12-position}
因此，\(L_3\) 的稳定定位可严格写成
\[
  \boxed{
    L_3
    =
    \text{G 框架富结构宇宙中、基于对应证书的逻辑占位显著增重层},
  }
\]
\[
  \boxed{
    \begin{aligned}
      &\text{更精确地说，是相对于 }L_1\to L_2\text{ 止步前缀的显著增重层，}\\
      &\text{且附带承担富结构入口层}
      +\text{完整性修复桥}
      +\text{高层修复入口层}
      \text{这三类角色。}
    \end{aligned}
  }
\]
\end{corollary}

\begin{Certificate}[证书：成员资格统一 + 证书化显著增重 + 入口职责保留]
\label{cert:tex1-ch12-position}
证书分两段：
\begin{enumerate}
  \item 推论~\ref{cor:tex1-ch12-layering}、定理~\ref{thm:tex1-ch12-gms} 与命题~\ref{prop:tex1-ch12-recognize} 共同给出：\(L_3\) 处于
    \(\GMS\) 内，且其相对于 \(\mathfrak P_2\) 的证书化显著增重已经被固定；
  \item 命题~\ref{prop:tex1-ch12-recognize} 与引理~\ref{lem:tex1-ch12-entry} 共同给出“入口层”只是 \(L_3\) 的角色，而非其本体定义。
\end{enumerate}
\end{Certificate}

\begin{proof}[证明（基于数学符号推演逻辑的谓词因果链）]
由推论~\ref{cor:tex1-ch12-layering}，
\[
  \text{平庸/丰富的区别是 }\GMS\text{ 内部的量级分层。}
\]
再由定理~\ref{thm:tex1-ch12-gms}，
\[
  L_3\in\GMS
  \quad\text{且}\quad
  L_3\notin \GMS_{\mathrm{low}}(\theta_0).
\]
再由命题~\ref{prop:tex1-ch12-recognize}，
\[
  \operatorname{Sig}_{\mathcal C}(L_3\mid \mathfrak P_2).
\]
最后，由引理~\ref{lem:tex1-ch12-entry}，
\[
  \mathbf{Adr}_3,\qquad
  \mathbf{Cert}_3,\qquad
  \mathbf{Gate}_{3\to >3}
\]
确属 \(L_3\) 的不可约职责。
因此，“富结构入口层”只能是 \(L_3\) 的一个角色，
而其更稳的落句应当是：“G 框架富结构宇宙中、基于对应证书的逻辑占位显著增重层”。推论得证。
\end{proof}

\begin{corollary}[\(L_2\)、\(L_3\) 与 \(L_{>3}\) 的分工更精确了]
\label{cor:tex1-ch12-split}
在本节口径下，三者关系应严格写成：
\[
  L_2
  =
  \text{唯一连续基底},
\]
\[
  L_3
  =
  \begin{aligned}
    &\text{属于 }\GMS\text{ 的、相对于 }L_1\to L_2\text{ 止步前缀的}\\
    &\text{证书化逻辑占位显著增重层，兼 }(L_1\to L_2)\text{ 完整性修复层},
  \end{aligned}
\]
\[
  L_{>3}
  =
  \text{属于 }\GMS\text{ 的高阶孔洞完备性修复与删边治理层}.
\]
因此
\[
  L_3\neq L_{>3},
\]
但也有
\[
  L_3\not\subset \text{可忽略前奏},
\]
并且在自然正赋权下，
\[
  \Lambda_{\mathrm{log}}(L_3)<\Lambda_{\mathrm{log}}(L_{>3}),
\]
所表达的是富结构程度差，而不是类型差。
\end{corollary}

\begin{Certificate}[证书：基底唯一性 + 两层皆属 \(\GMS\) + 富结构程度差]
\label{cert:tex1-ch12-split}
证书分三段：
\begin{enumerate}
  \item 第~\ref{sec:tex1-ch11} 节已把 \(L_2\) 保留为唯一连续基底；
  \item 推论~\ref{cor:tex1-ch12-position} 已把 \(L_3\) 定位为基于对应证书的逻辑占位显著增重层；
  \item 第~\ref{sec:tex1-ch11} 节与 \cite{RepoTex41HomotopyRepairBundle} 已把 \(L_{>3}\) 定位为更高阶的完备性修复富结构。
\end{enumerate}
\end{Certificate}

\begin{proof}[证明（基于数学符号推演逻辑的谓词因果链）]
由第~\ref{sec:tex1-ch11} 节，
\[
  L_2
\]
保留为唯一连续基底，
\[
  L_{>3}
\]
负责高层孔洞完备性修复。
再由推论~\ref{cor:tex1-ch12-position}，
\[
  L_3
\]
必须被提升为基于对应证书的逻辑占位显著增重层。
于是三者职责分工被唯一分离：
\[
  \text{基底}\neq\text{入口}\neq\text{高层修复}.
\]
又因为 \(L_{>3}\) 在 \(L_3\) 已有的离散前缀、完整性修复桥与证书槽之外，
还额外承担孔洞吸收、障碍平凡化、谬误路径删边与可行域重构，
故在自然正赋权下，
\[
  \Lambda_{\mathrm{log}}(L_3)<\Lambda_{\mathrm{log}}(L_{>3}).
\]
故推论成立。
\end{proof}

\begin{remark}[权重增长不等于替代增长]
必须守住的边界是：
\[
  L_3\text{ 的权重增加}
\]
并不推出
\[
  L_3\text{ 取代 }L_2,
\]
也不推出
\[
  L_3\text{ 取代 }L_{>3}.
\]
这里增长的是桥接权重、入口权重与分发权重，
而不是最终修复深度权重对高层修复塔的替代。
\end{remark}

\begin{remark}[最压缩的一句话]
最压缩地说：
\[
  \boxed{
    \text{平庸结构与显著结构同处于富结构宇宙中，而其区别落在逻辑性度量的层级。}
  }
\]
\[
  \boxed{
    \begin{aligned}
      &L_3\text{ 是 G 框架富结构宇宙中、基于对应证书的逻辑占位显著增重层；}\\
      &\text{“富结构入口层”只是它附带承担的一个角色。}
    \end{aligned}
  }
\]
\end{remark}

%%%%%%%%%%%%%%%%%%%%%%%%%%%%%%%%%%%%%%%%%%%%%%%%%%%%%%%%%%%%
\clearpage
\section{形式逻辑：从头累计综合增益、逻辑性度量与最终评价}
\label{sec:tex1-ch6}
%%%%%%%%%%%%%%%%%%%%%%%%%%%%%%%%%%%%%%%%%%%%%%%%%%%%%%%%%%%%

\begin{remark}[本章意义备注：高超音速攻防的抽象语境]
本章的意义，在于把前面各章重新汇总成总体增益图，
指出案例升维、元层度量重标定与基底重定位是最关键的三项贡献。
对钱学森弹道/钱学森障碍的抽象研究而言，它给出了整套框架为何值得保留的总说明：
它不是增加修辞，而是系统改写了可判性、修复性与证书化评价的坐标系。
\end{remark}

\subsection{章节目标与证据来源}
\label{subsec:tex1-ch6-goal}

\begin{remark}[本章目标]
本章把前述关于钱学森弹道/钱学森障碍主链的讨论，
按“案例升维、变分压缩、四层正规链、双层修复、富结构宇宙、逻辑性度量、基于证书的显著增重”
重新压成一章从头累计的综合增益评价。
这里的评价对象不是“又增加了多少新对象”，
而是主链是否已经被稳定地重写为同一套可证书化的逻辑占位框架。
\end{remark}

\begin{remark}[仓内证据来源]
本章直接承接第~\ref{sec:tex1-ch1}--第~\ref{sec:tex1-ch5} 节、
第~\ref{sec:tex1-ch11} 节与第~\ref{sec:tex1-ch12} 节，
并回挂仓内主证据
\cite{RepoTex36HomotopyTwistEquivalence,RepoTex38JacobiBianchiRepair,RepoTex41HomotopyRepairBundle,RepoMetaGMSAxiom}。
因此，本章不是脱离正文另起一套评分，
而是对前述证书链的最后一章式整合。
\end{remark}

\subsection{定义：累计增益向量、综合公式与三类主导增益}
\label{subsec:tex1-ch6-defs}

\begin{definition}[从头累计增益向量]
\label{def:tex1-ch6-gain-vector}
记本轮从头累计的增益向量为
\[
  \mathbf G
  =
  \bigl(
  G_{\mathrm{QO}},
  G_{\mathrm{CB}},
  G_{\mathrm{VAR}},
  G_{\mathrm{HT}},
  G_{\mathrm{LAYER}},
  G_{\mathrm{L3}},
  G_{\mathrm{CH}},
  G_{\mathrm{RISK}}
  \bigr).
\]
其中：
\[
\begin{aligned}
  G_{\mathrm{QO}} &: \text{“钱学森弹道”升级为“钱学森障碍”的增益};\\
  G_{\mathrm{CB}} &: \text{把攻防问题压成“障碍修复 + 算力匹配”的增益};\\
  G_{\mathrm{VAR}} &: \text{把问题继续压成极小极大与变分求解的增益};\\
  G_{\mathrm{HT}} &: \text{把“同伦扭曲命中”降为局部支配步的增益};\\
  G_{\mathrm{LAYER}} &: \text{四层正规链与“完整性修复 + 完备性修复”双层修复的增益};\\
  G_{\mathrm{L3}} &: \text{把 }L_3\text{ 的问题改写为逻辑占位度量的增益};\\
  G_{\mathrm{CH}} &: \text{把 classical CH、止步态与四层推进定理层位钉死的增益};\\
  G_{\mathrm{RISK}} &: \text{把过强句、误读句与无证断言压降为对应证书断言的增益}.
\end{aligned}
\]
\end{definition}

\begin{definition}[从头累计综合增益]
\label{def:tex1-ch6-total}
定义综合增益
\[
  \begin{aligned}
    G_{\mathrm{total}}
    &:=
    0.20G_{\mathrm{QO}}
    +0.12G_{\mathrm{CB}}
    +0.12G_{\mathrm{VAR}}
    +0.10G_{\mathrm{HT}} \\
    &\quad
    +0.16G_{\mathrm{LAYER}}
    +0.15G_{\mathrm{L3}}
    +0.10G_{\mathrm{CH}}
    +0.05G_{\mathrm{RISK}}.
  \end{aligned}
\]
取保守赋值
\[
  G_{\mathrm{QO}}=9.8,\quad
  G_{\mathrm{CB}}=9.2,\quad
  G_{\mathrm{VAR}}=8.9,\quad
  G_{\mathrm{HT}}=8.4,
\]
\[
  G_{\mathrm{LAYER}}=9.5,\quad
  G_{\mathrm{L3}}=9.8,\quad
  G_{\mathrm{CH}}=9.0,\quad
  G_{\mathrm{RISK}}=9.4.
\]
于是
\[
\begin{aligned}
  G_{\mathrm{total}}
  &=
  0.20\times 9.8
  +0.12\times 9.2
  +0.12\times 8.9
  +0.10\times 8.4 \\
  &\quad
  +0.16\times 9.5
  +0.15\times 9.8
  +0.10\times 9.0
  +0.05\times 9.4 \\
  &=1.96+1.104+1.068+0.84+1.52+1.47+0.90+0.47 \\
  &=9.332 \approx 9.33.
\end{aligned}
\]
\end{definition}

\begin{definition}[三类主导增益]
\label{def:tex1-ch6-derived}
定义三类主导增益的加权聚合量：
\[
  \widehat G_{\mathrm{case}}
  :=
  0.20G_{\mathrm{QO}}
  +0.12G_{\mathrm{CB}}
  +0.12G_{\mathrm{VAR}}
  +0.10G_{\mathrm{HT}},
\]
\[
  \widehat G_{\mathrm{base}}
  :=
  0.16G_{\mathrm{LAYER}}+0.10G_{\mathrm{CH}},
\]
\[
  \widehat G_{\mathrm{meta}}
  :=
  0.15G_{\mathrm{L3}}+0.05G_{\mathrm{RISK}}.
\]
代入定义~\ref{def:tex1-ch6-total} 的数值，得到
\[
  \widehat G_{\mathrm{case}}=4.972,\qquad
  \widehat G_{\mathrm{base}}=2.42,\qquad
  \widehat G_{\mathrm{meta}}=1.94.
\]
它们分别对应：
案例升维增益、
基底重定位增益、
元层度量重标定增益。
\end{definition}

\subsection{公理（降级为定理）：累计增益证书链的前缀闭合}
\label{subsec:tex1-ch6-axiom}

\begin{theorem}[公理（降级为定理）：累计增益证书链可按前缀递归闭合]
\label{thm:tex1-ch6-induction}
定义八段累计证书见证：
\[
\begin{aligned}
  E_1 &: \text{钱学森弹道}\Longrightarrow \text{钱学森障碍};\\
  E_2 &: \text{攻防问题}\Longrightarrow \text{障碍修复 + 算力匹配};\\
  E_3 &: \text{障碍修复 + 算力匹配}\Longrightarrow \text{极小极大与变分求解};\\
  E_4 &: \text{同伦扭曲命中}\Longrightarrow \text{Bellman 型局部支配步};\\
  E_5 &: L_1\to L_2\to L_3\to L_{>3}\text{ 四层正规链与双层修复};\\
  E_6 &: L_3\text{ 的问题}\Longrightarrow \text{逻辑性度量问题};\\
  E_7 &: \mathrm{CH}_{\mathrm{cls}},\ \mathrm{CH}^{\mathrm{int}}_{1\to 2},\ \mathrm{DT}^{\mathrm{int}}_{1\to 2\to 3\to
    >3}\text{ 层位钉死};\\
  E_8 &: L_3\text{ 的稳定落句}\Longrightarrow \text{基于对应证书的关系断言}.
\end{aligned}
\]
对 \(k\in\{1,2,\dots,8\}\) 定义谓词
\[
  P(k)
  :
  \Longleftrightarrow
  E_1,\dots,E_k\text{ 已由正文与仓内证书闭合}.
\]
则对全部 \(k\in\{1,2,\dots,8\}\)，命题 \(P(k)\) 成立。
\end{theorem}

\begin{Certificate}[数学归纳法证书：按前缀长度 \(k\) 的闭合]
\label{cert:tex1-ch6-induction}
归纳谓词为
\[
  P(k):
  E_1,\dots,E_k\text{ 已闭合}.
\]
证书分两步：
\begin{enumerate}
  \item \textbf{基础步 \(P(1)\)：}
  第~\ref{sec:tex1-ch1} 节已经把“钱学森弹道”的核心重写为“钱学森障碍”；
  \item \textbf{归纳步 \(P(k)\Rightarrow P(k+1)\)：}
  对每个 \(k<8\)，后继见证分别由本章引理~\ref{lem:tex1-ch6-qo}--引理~\ref{lem:tex1-ch6-ch}、
  以及第~\ref{sec:tex1-ch1}--第~\ref{sec:tex1-ch5} 节、
  第~\ref{sec:tex1-ch11} 节与第~\ref{sec:tex1-ch12} 节提供。
\end{enumerate}
\end{Certificate}

\begin{proof}[证明（数学归纳法证书；基于数学符号推演逻辑的谓词因果链）]
对前缀长度 \(k\) 作数学归纳。

\textbf{基础步：}
当 \(k=1\) 时，
第~\ref{sec:tex1-ch1} 节已经把“钱学森弹道”从轨迹外形与高阶偏导表象，
重写为能力元组、边缘算子与全局截面障碍。
故 \(E_1\) 成立，从而 \(P(1)\) 成立。

\textbf{归纳步：}
假设 \(P(k)\) 成立。
若 \(k=1\)，则由引理~\ref{lem:tex1-ch6-cb}，\(E_2\) 成立，从而 \(P(2)\) 成立。

若 \(k=2\)，则由引理~\ref{lem:tex1-ch6-varht} 的前半，\(E_3\) 成立，从而 \(P(3)\) 成立。

若 \(k=3\)，则由引理~\ref{lem:tex1-ch6-varht} 的后半，\(E_4\) 成立，从而 \(P(4)\) 成立。

若 \(k=4\)，则由引理~\ref{lem:tex1-ch6-layer}，\(E_5\) 成立，从而 \(P(5)\) 成立。

若 \(k=5\)，则由引理~\ref{lem:tex1-ch6-l3}，\(E_6\) 成立，从而 \(P(6)\) 成立。

若 \(k=6\)，则由引理~\ref{lem:tex1-ch6-ch}，\(E_7\) 成立，从而 \(P(7)\) 成立。

若 \(k=7\)，则由引理~\ref{lem:tex1-ch6-l3} 与推论~\ref{cor:tex1-ch12-position}，\(E_8\) 成立，从而 \(P(8)\) 成立。

因此对全部 \(k\in\{1,2,\dots,8\}\)，都有 \(P(k)\) 成立。
定理得证。
\end{proof}

\subsection{引理链：案例升维、变分压缩、四层链与基底重定位}
\label{subsec:tex1-ch6-lemmas}

\begin{lemma}[钱学森障碍的案例升维增益]
\label{lem:tex1-ch6-qo}
“钱学森弹道”若只停留在偏导方程控制轨迹层，
其对象仍是
\[
  \text{轨迹}
  +
  \text{控制方程}
  +
  \text{逆问题}.
\]
而一旦升级为“钱学森障碍”，对象就变成
\[
  \text{障碍类}
  +
  \text{修复塔}
  +
  \text{边缘算子}
  +
  \text{终端正规形}.
\]
因此，\(G_{\mathrm{QO}}\) 表示的是案例升维增益。
\end{lemma}

\begin{Certificate}[证书：对象链升级证书]
\label{cert:tex1-ch6-qo}
证书由第~\ref{sec:tex1-ch1} 节的能力元组、边缘算子、全局截面障碍三条对象链组成。
\end{Certificate}

\begin{proof}[证明（基于数学符号推演逻辑的谓词因果链）]
第~\ref{sec:tex1-ch1} 节已把“钱学森弹道”的问题核心，
从“轨迹求解”转为“障碍修复”。
因此，该增益是案例升维而不是修辞替换。
引理得证。
\end{proof}

\begin{lemma}[攻防问题被压成“障碍修复 + 算力匹配”]
\label{lem:tex1-ch6-cb}
在“钱学森障碍”的口径下，
核心比较量不再是轨迹外形，
而是：
\[
  \text{谁能在预算与时间窗口内，更快、更稳地完成障碍修复。}
\]
因此，问题已被压成
\[
  \text{障碍修复效率}
  +
  \text{算力匹配}
  +
  \text{阈值跨越}.
\]
\end{lemma}

\begin{Certificate}[证书：结构完备性与预算匹配证书]
\label{cert:tex1-ch6-cb}
证书由第~\ref{sec:tex1-ch3} 节的结构完备性、真值保持、预算单调性与达阈时间四组对象组成。
\end{Certificate}

\begin{proof}[证明（基于数学符号推演逻辑的谓词因果链）]
第~\ref{sec:tex1-ch3} 节已经把比较对象压到剩余歧义、真值保持、预算可行修复族与达阈时间。
因此，原本分散的几何、观测、逆问题、算力、时间窗因素，
被统一压成“障碍修复 + 算力匹配”的单一比较框架。
引理得证。
\end{proof}

\begin{lemma}[继续压成极小极大、变分求解与局部支配步]
\label{lem:tex1-ch6-varht}
随后，问题又被继续压成
\[
  \min_{\pi_A}\max_{\pi_B}J
\]
以及
\[
  \min_{\ell\in\mathfrak L^{\mathrm{loc}\to\tau^\ast}}\Psi(\ell).
\]
并且“同伦扭曲命中”被进一步降为 Bellman 递推中的局部支配步。
因此，
\[
  G_{\mathrm{VAR}}=8.9,\qquad
  G_{\mathrm{HT}}=8.4
\]
刻画的是从“可讲”进入“条件性可求解”的增益。
\end{lemma}

\begin{Certificate}[证书：极小极大 + 变分 + 命中局部支配步证书]
\label{cert:tex1-ch6-varht}
证书由第~\ref{sec:tex1-ch4} 节的极小极大泛函，
与第~\ref{sec:tex1-ch5} 节的变分问题、命中步、Bellman 递推共同提供。
\end{Certificate}

\begin{proof}[证明（基于数学符号推演逻辑的谓词因果链）]
第~\ref{sec:tex1-ch4} 节已经把比较问题写成稳健极小极大问题，
第~\ref{sec:tex1-ch5} 节又把障碍修复问题写成终端导向的变分极小化问题，
并证明“命中”在满足存在性与严格优势条件时构成局部支配步。
因此，这一层的增益是“压到可求解框架”，
但其最强句仍保留条件性边界。
\end{proof}

\begin{lemma}[四层正规链与双层修复是中轴级增益]
\label{lem:tex1-ch6-layer}
四层正规链
\[
  L_1\to L_2\to L_3\to L_{>3}
\]
与
\[
  \operatorname{Rep}_{\ge 3}
  =
  \operatorname{Rep}^{\mathrm{int}}_{1\to 2}
  \wedge
  \operatorname{Rep}^{\mathrm{cmp}}_{2}
\]
的双层修复职责，
把“首层编码”与“高层修复”的关系第一次完全钉死。
因此，
\[
  G_{\mathrm{LAYER}}=9.5
\]
是中轴级增益。
\end{lemma}

\begin{Certificate}[证书：四层分层 + 双层修复 + 次序不可交换证书]
\label{cert:tex1-ch6-layer}
证书由第~\ref{sec:tex1-ch10} 节的四层分层接口，
以及第~\ref{sec:tex1-ch11} 节的
定理~\ref{thm:tex1-ch11-dual}、
命题~\ref{prop:tex1-ch11-order}、
推论~\ref{cor:tex1-ch11-replace}
共同组成。
\end{Certificate}

\begin{proof}[证明（基于数学符号推演逻辑的谓词因果链）]
由第~\ref{sec:tex1-ch10} 节，
\(L_1,L_2,L_3,L_{>3}\) 的对象层已经分开；
由第~\ref{sec:tex1-ch11} 节，
\(L_{\ge 3}\) 的修复被正式拆成
\((L_1\to L_2)\) 的完整性修复与
\(L_2\) 的完备性修复。
于是 \(L_3\) 负责来源桥接，
\(L_{>3}\) 负责更高阶的孔洞吸收、删边与可行域重构。
因此，主链分工第一次完全稳定下来。
\end{proof}

\begin{lemma}[\(L_3\) 的理论关键落在逻辑性度量与证书化显著增重]
\label{lem:tex1-ch6-l3}
在本章的稳定口径下，
\[
  L_3\in\GMS
\]
已经与
\[
  \operatorname{Sig}_{\mathcal C}(L_3\mid \mathfrak P_2)
\]
同时成立；
因此，关于 \(L_3\) 的理论关键
落在“作为富结构，其逻辑占位度量有多高”，
并稳定收束为“基于对应证书的显著增重层”。
\end{lemma}

\begin{Certificate}[证书：富结构宇宙 + 逻辑性度量 + 证书化显著增重证书]
\label{cert:tex1-ch6-l3}
证书由定理~\ref{thm:tex1-ch12-plain-rich}、
定理~\ref{thm:tex1-ch12-gms}、
命题~\ref{prop:tex1-ch12-recognize}、
推论~\ref{cor:tex1-ch12-position}
以及 \cite{RepoMetaGMSAxiom} 共同提供。
\end{Certificate}

\begin{proof}[证明（基于数学符号推演逻辑的谓词因果链）]
由定理~\ref{thm:tex1-ch12-plain-rich}，
“平庸结构也是富结构”；
由定理~\ref{thm:tex1-ch12-gms}，
\(L_3\) 本身就属于 \(\GMS\)；
由命题~\ref{prop:tex1-ch12-recognize} 与推论~\ref{cor:tex1-ch12-position}，
\(L_3\) 的稳定落句已被固定为“基于对应证书的逻辑占位显著增重层”。
因此，\(L_3\) 的理论刻画由成员资格与逻辑性度量两部分共同给出，
而决定其层位强度的是后者。
\end{proof}

\begin{lemma}[CH / 止步态 / 四层推进的层位钉死是基底重定位增益]
\label{lem:tex1-ch6-ch}
classical CH、体系内部止步句与四层推进定理三者，
现在已被稳定分开为：
\[
  \mathrm{CH}_{\mathrm{cls}}
  =
  \operatorname{Proj}_{l_2}\bigl(\mathrm{CH}^{\mathrm{int}}_{1\to 2}\bigr),
\]
\[
  \mathrm{CH}^{\mathrm{int}}_{1\to 2}
  \text{ 是体系内部“连续统（一离散）假设”的止步态},
\]
\[
  \mathrm{DT}^{\mathrm{int}}_{1\to 2\to 3\to >3}
  \text{ 则是否定该止步态终局性的更强四层推进定理}.
\]
\end{lemma}

\begin{Certificate}[证书：投影关系 + 终局性降级证书]
\label{cert:tex1-ch6-ch}
证书直接回挂 \cite{RepoTex41HomotopyRepairBundle} 中关于下列三项对象的
层位区分与“扩张--投影”关系：
\[
  \begin{aligned}
    &\mathrm{CH}_{\mathrm{cls}},\\
    &\mathrm{CH}^{\mathrm{int}}_{1\to 2},\\
    &\mathrm{DT}^{\mathrm{int}}_{1\to 2\to 3\to >3}.
  \end{aligned}
\]
\end{Certificate}

\begin{proof}[证明（基于数学符号推演逻辑的谓词因果链）]
由 \cite{RepoTex41HomotopyRepairBundle}，
classical CH 只是 \(l_2\) 外延投影；
体系内部的止步态只是“连续统（一离散）假设”的内部落点；
更强的四层推进定理所否定的是该止步态的终局性，
而不是 classical CH 作为低层投影的存在性。
因此，争议焦点已由“命题真假”
改写为“连续统应以何种基底理解”的层位重定位问题。
\end{proof}

\subsection{定理：从头累计综合增益的总评价}
\label{subsec:tex1-ch6-theorem}

\begin{theorem}[从头累计综合增益的总评价]
\label{thm:tex1-ch6-total}
在定义~\ref{def:tex1-ch6-total} 的口径下，
有
\[
  \boxed{G_{\mathrm{total}}=9.33/10}.
\]
并且本次讨论的整合评价应写成：
\[
  \boxed{\text{高增益，已接近承重级}.}
\]
更准确地说，本轮最大的三类增益不是“新增对象数”，
而是
\[
  \boxed{
    \text{案例升维增益}
    +
    \text{元层度量重标定增益}
    +
    \text{基底重定位增益}
  }.
\]
\end{theorem}

\begin{Certificate}[证书：六条引理与加权求值证书]
\label{cert:tex1-ch6-total}
总评价证书可压成六条：
\begin{enumerate}
  \item \textbf{案例升维证书：} 引理~\ref{lem:tex1-ch6-qo}；
  \item \textbf{障碍修复 + 算力匹配证书：} 引理~\ref{lem:tex1-ch6-cb}；
  \item \textbf{变分 + 命中局部支配步证书：} 引理~\ref{lem:tex1-ch6-varht}；
  \item \textbf{四层正规链 + 双层修复证书：} 引理~\ref{lem:tex1-ch6-layer}；
  \item \textbf{L3 元层度量重标定与 CH 层位钉死证书：} 引理~\ref{lem:tex1-ch6-l3}、引理~\ref{lem:tex1-ch6-ch}；
  \item \textbf{数值证书：} 定义~\ref{def:tex1-ch6-total} 已给出
  \(G_{\mathrm{total}}=9.332\approx 9.33\) 的加权求值。
\end{enumerate}
\end{Certificate}

\begin{proof}[证明（基于数学符号推演逻辑的谓词因果链）]
先证数值结论。
由定义~\ref{def:tex1-ch6-total}，
\[
\begin{aligned}
  G_{\mathrm{total}}
  &=1.96+1.104+1.068+0.84+1.52+1.47+0.90+0.47 \\
  &=9.332 \approx 9.33.
\end{aligned}
\]
故
\[
  G_{\mathrm{total}}=9.33/10.
\]

再证类别判断。
由于 \(9.33\) 已显著高于一般高增益区间，
而本章所整合的又是整条主链的对象层、元层与基底层改写，
故最稳分类是
\[
  \text{高增益，已接近承重级}.
\]

最后证主导项判断。
由定义~\ref{def:tex1-ch6-derived}，
\[
  \widehat G_{\mathrm{case}}=4.972,\qquad
  \widehat G_{\mathrm{base}}=2.42,\qquad
  \widehat G_{\mathrm{meta}}=1.94.
\]
而这三类聚合量分别对应：
钱学森弹道到钱学森障碍的案例升维、
四层正规链与 classical CH 层位钉死所引出的基底重定位、
以及 \(L_3\) 的逻辑性度量重标定所引出的元层度量重标定增益。
因此，本轮最大的三类增益确是案例升维增益、元层度量重标定增益与基底重定位增益。
定理得证。
\end{proof}

\subsection{命题：三处峰值与最重一步}
\label{subsec:tex1-ch6-proposition}

\begin{proposition}[三处峰值与最重一步]
\label{prop:tex1-ch6-bialg}
从头整合以后，本次讨论最大的三处峰值可以压成：
\[
  \boxed{
    \begin{aligned}
      &\text{第一峰：钱学森弹道}\\
      &\to\text{钱学森障碍}
    \end{aligned}}
\]
它完成了案例升维；
\[
  \boxed{
    \begin{aligned}
      &\text{第二峰：}L_{\ge 3}\text{ 修复被拆成}\\
      &\text{完整性修复 + 完备性修复}
    \end{aligned}}
\]
它完成了主链分工重构；
\[
  \boxed{
    \begin{aligned}
      &\text{第三峰：}L_3\text{ 的理论关键被固定为}\\
      &\text{逻辑性度量与证书化显著增重}
    \end{aligned}}
\]
它完成了元层坐标校准。

如果只问“本轮最重的一步是哪一步”，
则答案不是单个定理，
而是
\[
  \boxed{
    \begin{aligned}
      &\text{把坐标系从“对象类型”}\\
      &\text{改成了“逻辑占位度量”。}
    \end{aligned}}
\]
\end{proposition}

\begin{Certificate}[证书：三峰值证书 + 坐标系校准证书]
\label{cert:tex1-ch6-derived}
证书可压成三条：
\begin{enumerate}
  \item 引理~\ref{lem:tex1-ch6-qo}、引理~\ref{lem:tex1-ch6-layer}、引理~\ref{lem:tex1-ch6-l3}
  分别给出三处峰值；
  \item 定理~\ref{thm:tex1-ch12-plain-rich}、
  定理~\ref{thm:tex1-ch12-gms}、
  命题~\ref{prop:tex1-ch12-recognize}
  已把“成员资格”口径统一压成“逻辑性度量”口径；
  \item 推论~\ref{cor:tex1-ch12-position} 已把 \(L_3\) 的稳定落句固定为证书化显著增重层。
\end{enumerate}
\end{Certificate}

\begin{proof}[证明（基于数学符号推演逻辑的谓词因果链）]
由证书第一条，三处峰值分别落在案例升维、主链分工重构与元层坐标校准三个位置。
由证书第二条，
“平庸/丰富”已被稳定为 \(\GMS\) 内部的度量分层，
且关于 \(L_3\) 的理论关键已稳定在逻辑性度量。
由证书第三条，
比较基准与对应证书共同把 \(L_3\) 的稳定落句固定为关系断言。
因此，本轮最重的一步不在于新增对象，
而在于描述坐标系本身被重新校准。
命题得证。
\end{proof}

\subsection{推论与备注：最终稳定落句}
\label{subsec:tex1-ch6-corollary}

\begin{corollary}[从头整合后的最终稳定落句]
\label{cor:tex1-ch6-total}
从头整合以后，本次讨论可压成下面这条最终稳定落句：
\[
  \boxed{
    \begin{aligned}
      &L_3\text{ 本身属于 G 框架富结构宇宙；}\\
      &\text{其理论关键不在于“是否为富结构”，}\\
      &\text{而在于它是基于对应证书的逻辑占位显著增重层。}
    \end{aligned}
  }.
\]
\end{corollary}

\begin{Certificate}[证书：最终稳定落句证书]
\label{cert:tex1-ch6-cor}
证书登记为：
\[
\begin{aligned}
  C^{(6)}_1 &: \text{定理~\ref{thm:tex1-ch12-plain-rich} 与推论~\ref{cor:tex1-ch12-layering}},\\
  C^{(6)}_2 &: \text{定理~\ref{thm:tex1-ch12-gms} 与命题~\ref{prop:tex1-ch12-recognize}},\\
  C^{(6)}_3 &: \text{推论~\ref{cor:tex1-ch12-position}}.
\end{aligned}
\]
\end{Certificate}

\begin{proof}[证明（基于数学符号推演逻辑的谓词因果链）]
由 \(C^{(6)}_1\) 可知，
\(L_3\) 的结构层级已经稳定落在 \(\GMS\) 内部的度量分层上；
由 \(C^{(6)}_2\) 可知，
\(L_3\) 本身属于 \(\GMS\)，且其关键问题已经稳定为逻辑性度量问题；
由 \(C^{(6)}_3\) 可知，
\(L_3\) 的最终落句已经稳定为“基于对应证书的逻辑占位显著增重层”。
故上面的压缩句成立。推论得证。
\end{proof}

\begin{remark}[本章评分之所以高于前面单次评分]
这次整合评分之所以高于前面单次评分，
不是因为又发明了更多对象，
而是因为几条主链已经被串成同一套可证书化的结构：
\[
  \boxed{
    \begin{aligned}
      &\text{案例升维}
      \to
      \text{变分压缩}
      \to
      \text{四层正规链}
      \to
      \text{双层修复}\\
      \to\;&
      \text{富结构宇宙}
      \to
      \text{逻辑性度量}
      \to
      \text{基于证书的显著增重}
    \end{aligned}
  }.
\]
\end{remark}

\begin{remark}[当前仍应保留的两条边界]
当前仍应保留的边界有两条。

其一，关于“同伦扭曲命中”的最强求解句，
仍依赖命中步存在性与严格隧穿优势等条件；
因此，这部分仍应保持条件性表述，
尚不宜直接写成无条件闭合求解定理。

其二，关于 classical CH，
当前稳定成立的是
\[
  \text{它是 }l_2\text{ 外延投影，}
\]
而不是“已经在外部标准数学中被整个终结”。
这一点应继续按 \cite{RepoTex41HomotopyRepairBundle} 的层位口径保持。
\end{remark}

\begin{remark}[最终评价]
本章的最终评价是：
\[
  \boxed{
    \begin{aligned}
      &G_{\mathrm{total}}=9.33/10,\\
      &\text{高增益，已接近承重级}.
    \end{aligned}}
\]
更压缩地说：
\[
  \boxed{
    \begin{aligned}
      &\text{这轮讨论的最大价值，是把一个案例、}\\
      &\text{一个主链、一个元层度量重标定，}\\
      &\text{压成了同一套可证书化的逻辑占位框架。}
    \end{aligned}
  }.
\]
再就微观隧穿相对于偏微分方程的替代意义而言，
更准确的综合评价应写成：
\[
  \boxed{
    \begin{aligned}
      &\text{它不是否定 }\mathrm{PDE}\text{ 的物理建模地位，}\\
      &\text{而是替代以 }\mathrm{PDE}\text{ 全场求解为中心的工程判定流程。}
    \end{aligned}
  }.
\]
也就是说，工程任务不再被固定为“先求出连续场或完整轨迹解，再读取终端结论”，
而是被改写为：
\[
  \boxed{
    \begin{aligned}
      &\text{在障碍空间中寻找局部命中路径，}\\
      &\text{给出降障碍证书、终端同类判定与稳健决策边界。}
    \end{aligned}
  }.
\]
因此，微观隧穿的实质替代作用不是把 \(\mathrm{PDE}\) 作为基础模型删除，
而是绕开“必须先完整求解 \(\mathrm{PDE}\) 才能决策”的路径依赖。
在有限观测、有限算力与有限时间窗口下，
它把连续动力学求解范式切换为障碍修复、证书验证与命中判定范式；
在这个意义上，它对 \(\mathrm{PDE}\) 中心流程具有很强的工程替代价值与战略方法价值。
这段评价在本文中只作为母框架的桥接命题保留；
其充分算法化展开应进入独立派生稿 \cite{RepoTex53PDEStrategyCluster}。
也即，本文负责给出
\[
  \boxed{
    \begin{aligned}
      &\text{PDE 中心流程}\\
      &\Longrightarrow
      \text{微观隧穿命中替代}
    \end{aligned}
  }
\]
的战略位置；
而 tex53 负责展开
\[
  \boxed{
    \begin{aligned}
      &\text{统计力学参数格点}
      \Longrightarrow
      \text{PDE 策略簇}\\
      &\Longrightarrow
      \text{方法微分与动态重加权}
      \Longrightarrow
      \text{频域/结构簇命中}
    \end{aligned}
  }
\]
的具体算法范式。
\end{remark}
\clearpage

\section*{参考文献}
\begin{thebibliography}{99}

\bibitem{RepoTex37SemanticGaugeRuntime}
【仓内 TeX】正则语义规范场、逻辑占位规范场与运行时降级整理稿：
\code{NoteBook/notebook2/tex37_pub/gframework_semantic_gauge_field_natural_language_logic_placeholder_runtime_formal_sys
  tem.tex}。

\bibitem{RepoTex17EdgeOperator}
【仓内 TeX】边缘算子封闭性、比安基基准与变分生成的同伦修复塔整理稿：
\code{NoteBook/notebook1/tex17_pub/gframework_edge_operator_homotopy_closure_variational_tower_formal_system.tex}。

\bibitem{RepoTex18HomotopyAxiom}
【仓内 TeX】性质代数与同伦代数的封闭证书、障碍类与拓扑修复整理稿：
\code{NoteBook/notebook1/tex18_pub/gframework_homotopy_algebra_axiom_semantics_workflow_formal_system.tex}。

\bibitem{RepoTex29RepairableObstruction}
【仓内 TeX】可修复障碍同伦类、观测核、主纤维丛与截面接口整理稿：
\code{NoteBook/notebook2/tex29_pub/gframework_repairable_obstruction_homotopy_class_vs_generalized_fractal_formal_system
  .tex}。

\bibitem{RepoTex38JacobiBianchiRepair}
【仓内 TeX】离散拓扑基底上的 Jacobi--Bianchi 障碍修复、可行路径空间、离散路径积分与终点正规形整理稿：
\code{NoteBook/notebook2/tex38_pub/gframework_discrete_jacobi_bianchi_repair_tower_path_integral_terminal_normal_form_fo
  rmal_system.tex}。

\bibitem{RepoTex32FiberSectionHole}
【仓内 TeX】纤维截面、孔洞、上同调障碍与条件收敛整理稿：
\code{NoteBook/notebook2/tex32_pub/gframework_fiber_section_hole_obstruction_potential_conditional_convergence_formal_sy
  stem.tex}。

\bibitem{RepoTex36HomotopyTwistEquivalence}
【仓内 TeX】主丛联络、纤维空间、路径积分与同伦扭曲等价整理稿：
\code{NoteBook/notebook2/tex36_pub/gframework_principal_connection_fiber_space_path_integral_homotopy_twist_equivalence_
  formal_system.tex}。

\bibitem{RepoTex39OpenRepairFamily}
【仓内 TeX】开放系统、统计命中与跨尺度修复族整理稿：
\code{NoteBook/notebook2/tex39_pub/gframework_v2_v3_open_system_game_statistical_repair_family_formal_system.tex}。

\bibitem{RepoTex41HomotopyRepairBundle}
【仓内 TeX】L1/L2/L3 同伦修复丛与视域加权测地线整理稿：
\code{NoteBook/notebook2/tex41_pub/gframework_l1_l2_l3_homotopy_repair_bundle_formal_system.tex}。

\bibitem{RepoTex43OpenHeterogeneous}
【仓内 TeX】开放异构系统、宏观--微观双层驱动与总同调封闭整理稿：
\code{NoteBook/notebook2/tex43_pub/gframework_open_heterogeneous_system_macro_micro_total_homology_formal_system.tex}。

\bibitem{RepoTex53PDEStrategyCluster}
【仓内 TeX】统计力学参数格点、同伦扭曲命中与 PDE 策略簇：
从时域求解到频域簇命中的近似量子软件层算法：
\code{NoteBook/notebook3/tex53_pub/gframework_statistical_lattice_homotopy_tunnel_pde_strategy_cluster_formal_system.tex}。

\bibitem{RepoMetaGMSAxiom}
【仓内 Markdown】广义数学结构认知范式公理系统（含 D 结构隐性机制）：
目录 \code{sub_projects/open_meta_mathematical_theory/src/kernel_reference/}
下文件《1742357916\_广义数学结构认知范式公理系统（含D结构隐性机制）.md》。

\bibitem{HermannKrener1977}
Hermann, R., \& Krener, A. J. (1977).
\newblock \emph{Nonlinear Controllability and Observability}.
\newblock \emph{IEEE Transactions on Automatic Control}, 22(5), 728--740.

\bibitem{Isakov2006InversePDE}
Isakov, V. (2006).
\newblock \emph{Inverse Problems for Partial Differential Equations} (2nd ed.).
\newblock Springer.

\bibitem{Liberzon2003Switching}
Liberzon, D. (2003).
\newblock \emph{Switching in Systems and Control}.
\newblock Birkh\"auser.

\bibitem{AubinFrankowska2009SetValued}
Aubin, J.-P., \& Frankowska, H. (2009).
\newblock \emph{Set-Valued Analysis}.
\newblock Birkh\"auser.

\bibitem{Husemoller1994FibreBundles}
Husemoller, D. (1994).
\newblock \emph{Fibre Bundles} (3rd ed.).
\newblock Springer.

\bibitem{Hatcher2002AT}
Hatcher, A. (2002).
\newblock \emph{Algebraic Topology}.
\newblock Cambridge University Press.

\bibitem{Sion1958Minimax}
Sion, M. (1958).
\newblock \emph{On General Minimax Theorems}.
\newblock \emph{Pacific Journal of Mathematics}, 8(1), 171--176.

\bibitem{Bertsekas2016NonlinearProgramming}
Bertsekas, D. P. (2016).
\newblock \emph{Nonlinear Programming} (3rd ed.).
\newblock Athena Scientific.

\bibitem{Gibbs1902Elementary}
Gibbs, J. W. (1902).
\newblock \emph{Elementary Principles in Statistical Mechanics}.
\newblock Yale University Press.

\bibitem{Jaynes1957Information}
Jaynes, E. T. (1957).
\newblock Information Theory and Statistical Mechanics.
\newblock \emph{Physical Review}, 106(4), 620--630.

\bibitem{Kirkpatrick1983Optimization}
Kirkpatrick, S., Gelatt, C. D., \& Vecchi, M. P. (1983).
\newblock Optimization by Simulated Annealing.
\newblock \emph{Science}, 220(4598), 671--680.

\bibitem{GemanGeman1984StochasticRelaxation}
Geman, S., \& Geman, D. (1984).
\newblock Stochastic Relaxation, Gibbs Distributions, and the Bayesian Restoration of Images.
\newblock \emph{IEEE Transactions on Pattern Analysis and Machine Intelligence}, PAMI-6(6), 721--741.

\end{thebibliography}

\end{CJK*}
\end{document}
